% ============================================================================
%  RESULTS -- ONLINE SUPPLEMENT.  project `panel-npiv-discrete`
%
%  WHAT THIS FILE HOLDS.  The 62 theorem-like environments that the article
%  refers to BY NAME but does not print: the supporting lemmas, the technical
%  definitions, the reconciliation and scope remarks, and the results about
%  assumptions the paper retires.
%
%  results_main.tex + results_supp.tex ARE the old results.tex.  The split is a
%  PARTITION: all 85 environments survive, every \label is declared exactly once
%  across the two files, and no statement was rewritten, retitled or renumbered.
%  Each environment carries the comment block that preceded it.  The original
%  results.tex header and the trailing change log are preserved below and at the
%  foot of this file respectively.
%
%  results_main.tex holds the other 23 environments and is \input by
%  paper/qmd/main.qmd.  This file is \input by paper/qmd/supplement.qmd, ahead of
%  proofs.tex, so that the proofs reach the statements they prove.  Putting it
%  first is what repaired the supplement's own cross-reference graph: 110 of the
%  177 targets that were undefined when this document was proofs.tex alone now
%  resolve here.
%
%  LABEL HYGIENE -- READ THIS BEFORE ADDING AN \input ANYWHERE.
%  SUPERSEDED 2026-08-22, AND THIS BLOCK WAS THE LAST PLACE THAT SAID OTHERWISE.
%  results.tex is a RETIRED MIRROR compiled by no document: build.sh:46 carries
%  its own RETIRED marker, driver.tex \inputs the eight split files, and
%  quality.md Sec.3 clause 1 was repointed off it on 2026-08-05.  What remains
%  true of it is only the last clause below -- the symbolic specs cite it by line
%  number -- and that is a follow-up, not a reason to treat it as authority.
%  As written through 2026-08-22 this block read:
%  results.tex SURVIVES UNDIVIDED and is still the source of record: it is what
%  scripts/build-theory/build.sh compiles, what scripts/render-sync/
%  check_rendering_sync.py reads, and what ~100 files under scripts/symbolic/
%  specs/ cite BY LINE NUMBER.  It was therefore not edited.  The consequence is
%  that every \label in this file is declared TWICE on disk, here and in
%  results.tex.  THIS FILE AND results.tex MUST NEVER BE \input INTO THE SAME
%  DOCUMENT, and neither must results_main.tex and results.tex.  That is the same
%  rule results.tex's own header states for paper/sections/theorems/*.tex.
% ============================================================================

% ----------------------------------------------------------------------------
%  PRESERVED VERBATIM: the header of the undivided results.tex.  Its CONTENTS
%  list (Parts A-G) describes the two files JOINTLY; it is retained as
%  provenance and was not updated, on the rule that nothing in the corpus is
%  rewritten by the split.
% ----------------------------------------------------------------------------


% ============================================================================
%  RESULTS -- project `panel-npiv-discrete`
%  Stage 3, dispatch 1: proof_strategy_memo.md Sec.10, Steps 1-2.
%  ROUND 2: revised against proofs_review.md (56/100).  Change log at the foot.
%
%  CONTENTS
%    Part A.  Setup: model, spaces, THE OPERATOR EQUATION (fills the gap at
%             draft `:120--122'), the Frechet derivative, compactness.
%    Part B.  Step 1 of the memo's order of work: lem:shiftinv (Lemma L1) and
%             Remark R1.
%    Part C.  Step 2: lem:bierens (F4 version only, per plan Sec.19.3), with
%             supporting lemmas L-2 (density) and L-3 (quantifier management).
%    Part D.  Reconciliation with the source draft.
%
%    Part E.  Step 3: lem:adjoint and the nu_g algebra; the impossibility
%             result; the kernel characterisation.
%    Part F.  Step 5: def:weight and def:dirclass (the two definitions the
%             round-7 CRITICALs required), L-16, L-6, prp:orthog in sieve form,
%             L-7.
%
%    Part G.  Step 4: L-21 (lem:reprsmooth) with rem:l21route, L-20
%             (lem:formstruct), prp:sharpquad, rem:sharpscope.  Printed AFTER
%             Part F because the file is ordered by dispatch and Step 5 ran
%             first; L-21 is logically PRIOR to Part E's cor:elldiff and
%             retroactively discharges what Parts E and F carried as a
%             hypothesis.  See Part G's header for why that is not a circularity.
%
%  NOT IN THIS FILE, deliberately: ~~prp:sharpquad,~~ prp:phirate, thm:rootn,
%  prp:effbound, prp:plugin, cor:trend.  Those are Steps ~~4 and~~ 6--8.
%  ~~prp:sharpquad in particular is Step 4 and is dispatched separately: it fixes
%  thm:rootn's scope, and stating it here would pre-empt that.~~
%  ROUND 9: prp:sharpquad IS NOW IN THIS FILE, in Part G.  It still fixes
%  thm:rootn's scope; what changes is that the scope statement is now a proved
%  proposition rather than a forward reference.
%
%  Proofs are in proofs.tex.  Assumptions are cited BY TAG and are defined in
%  assumptions.tex.  Symbols follow notation_glossary.md.
%
%  LABEL HYGIENE.  paper/sections/theorems/{lem-shiftinv,lem-bierens}.tex are
%  ALTERNATIVE RENDERINGS of lem:shiftinv and of lem:bierens + cor:opsol,
%  carrying the SAME \label keys so that the paper and this working artifact
%  agree.  They must never be \input into the same document as this file.
%  Within {results,proofs}.tex every \label is declared exactly once.
%
%  REMARK PLACEMENT.  House style puts every remark AFTER the proof of the
%  result it concerns.  This file carries statements only, so a remark here
%  sits immediately after the statement it belongs to; in the paper files it
%  sits after \end{proof}, which is where the ordering is enforced.
%  rem:R1 follows lem:shiftinv; rem:swintercept and rem:Tcond follow
%  lem:bierens and cor:opsol; rem:quadremainder (in proofs.tex) follows
%  prp:frechet; rem:draftrec is expositional and closes the section.
% ============================================================================


\begin{proposition}[$T$ is well defined and bounded]\label{prp:wellposed}
The moment family \eqref{eq:Mdef} defines an element of $L^{2}(\Pi)$ at every trial pair.
Suppose B1(i), (ii), (iv.c) and (iv.d) hold. \emph{\textbf{B1(iv) is cited by clause, and the
factorisation is not among the clauses}}: $\Pi=\Pi_u\otimes\Pi_t\otimes\Pi_{\mathsf w}$ is B1(iv)'s
\emph{declaration}, standing at every statement in this corpus by the assumption block's \textbf{three
registers} paragraph, and it is what lets Step~2's iterated integral be written at all. What the proof
\emph{opens} is \textup{(iv.c)} and \textup{(iv.d)} and nothing else --- Step~2 reads ``by B1(iv.c) and
B1(iv.d)'' in terms --- so a block citation would charge this proposition with \textup{(iv.a)} and
\textup{(iv.b)}, which no step reads. \emph{Through 2026-08-24 this line was the block}, and before that
day it was \textup{(iv.c)} alone. Then for every $(\varphi,\theta)\in\mathcal V\times\IR^{p}$ the function
$\mathcal M(\cdot,\cdot;\varphi,\theta)$ is Borel measurable and
\[
\big\|T(\varphi,\theta)\big\|_{L^{2}(\Pi)}\;\le\;c_\Pi\;<\;\infty ,
\]
with $c_\Pi$ from Assumption~B1(iv.c), so that $T$ maps $\mathcal V\times\IR^{p}$ into a
ball of $L^{2}(\Pi)$ of radius not depending on $(\varphi,\theta)$.
\end{proposition}

\begin{remark}[$\nu$ versus $\nu/\pi$]\label{rem:nupi}
Suppose in addition A1b(iv) holds, so that $\nu=\nu_{\Delta Z}$ is well defined. Then
$\int_{\mathcal X}(\nu/\pi)\,\pi\,dx=\int_{\mathcal X}\nu\,dx=0$, whereas
$\int_{\mathcal X}\nu\,\pi\,dx=E[\Delta Z\{\pi(X_2)-\pi(X_1)\}]$, which does not vanish at a generic
design --- $\pi$ is a feature of $P$ by \cref{def:spaces} and not a shape the analyst may choose to make
this term zero. Since $\mathcal R(\mathcal A^\ast_{\mathrm{cond}})\subseteq\mathcal V$ by
\cref{prp:compact} and $\mathcal V=\{\zeta:\int\zeta\,\pi\,dx=0\}$, the source condition A2 can only be
imposed on $\nu/\pi$, never on $\nu$.
\end{remark}

% THE SPLIT IS DELIBERATE.  Round 2 stated the whole lemma under A1, but Steps
% 1-2 of the proof establish \eqref{eq:shiftinv} from def:admissible ALONE; A1
% enters only the final clause.  Advertising A1 for the equality made cor:opsol
% -- which uses the equality and does not assume A1 -- not strictly citable
% against this lemma, and forced it to spend two sentences explaining which half
% it was using.  The split removes that.

\begin{remark}[R1]\label{rem:R1}
Three things follow from \cref{lem:shiftinv}, and one does not.

\emph{(i) What it establishes.} Restriction (A) carries no first-order information about
$\theta_\dagger$: the set of pairs it admits is a product set, so no functional of the conditional law
of the residual given $(W,\Delta Z)$ can distinguish $\theta_\dagger$ from any other $\theta$.

(ii) What it does not establish. It does not say that $\theta_\dagger$ is unidentified.
$\theta_\dagger$ is identified by (B) alone: under A1 and A1b(iv),
\[
\theta_\dagger\;=\;\big(E[\Delta Z\Delta Z']\big)^{-1}E\big[\Delta Z\,(\Delta Y-\Delta\vdagger)\big],
\]
the inverse existing by A1b(iv) and the expectations being finite by the moment clauses of that same
assumption. This route uses neither A2 nor the retired A3. A3 asserted nonsingularity of
$J=E[\Delta Z\Delta Z']+E[\ell_\star\widetilde W\Delta Z']$ and was itself conditional on A2; routing
identification of $\theta_\dagger$ through it would make identification contingent on a source condition,
which it is not. By \cref{lem:momzero}\textup{(iv)} the second summand is zero under A1, so $J$ is
$E[\Delta Z\Delta Z']$ wherever it is defined, and A3 is retired at A3 as a consequence of
A1 and A1b(iv). \emph{What must not be routed through A2 is the existence of $\ell_\star$}, and the
display above needs none of it. Nor does \cref{lem:shiftinv} say that
(A) is useless. A Neyman correction is mean-zero at the truth \emph{by construction}, that is what
makes it a correction rather than a new restriction, so ``carries no information about $\theta$'' is a
property of every Riesz representer in the debiased-machine-learning construction
\citep{chernozhukov2018dml}, not a defect peculiar to this model.
% ARCHAEOLOGY.  Through dispatch 12 this passage also called J "the different
% matrix"; that half is withdrawn, the separation the paragraph makes being
% unaffected and in fact sharper.  An earlier version of the project inferred from
% lem:shiftinv that the conditional form of A1 forces a two-step plug-in with no
% root-n theorem; the inference does not follow and was withdrawn.

\emph{(iii) What it costs.} Over-identification. Under an \emph{unconditional} independence assumption
one would have $\Delta Y-\Delta\vdagger-\Delta Z'\theta\not\upmodels W$ for $\theta\ne\theta_\dagger$
whenever $\Delta Z\not\upmodels W$, so (A) would carry information about $\theta$ and deliver a testable
restriction. Under A1 as stated, $\theta_\dagger$ is \emph{exactly} identified by (B) and (A) is a pure
debiasing device: \emph{(A) yields no over-identifying restriction on $\theta$}, which is precisely what
\cref{lem:shiftinv} proves. It does not follow that (A) is untestable. Solvability of the operator
equation remains testable in the usual nonparametric-instrumental-variables sense, $\mathcal V_{CI}=\emptyset$
is an observable event, witnessed for instance by $\Delta Y=W\chi+e$ with $\chi\ne0$ and $e\upmodels W$,
where no $\varphi$ makes the residual conditionally independent of $W$. The claim to make in the paper is
about $\theta$, not about (A).

(iv) A population consequence, and a conditional algorithmic one. By \cref{cor:opsol} the
solution set of \eqref{eq:opeq} in $\varphi$ is $\mathcal V_{CI}$ at every $\theta$, so \emph{as a
population statement} the first stage may be posed at any fixed $\theta$, including $\theta=0$, that is,
on raw $\Delta Y$. Two qualifications are required before this says anything about the estimator. First,
``targets $\vdagger$'' needs $\mathcal V_{CI}=\{\vdagger\}$, i.e.\ A1b(iii); without it the first stage
targets a set. Second, consistency of the \emph{regularised sample} solve is not delivered by
\cref{cor:opsol}: it is delivered by \cref{prp:phirate}, whose hypotheses, A7 in its amended form
\eqref{eq:a7amend}, the spectral source condition (P-a.3), and the bandwidth condition (P-f), are conditions on
$\hat T(\cdot,\theta)$ and $\hat{\mathcal A}$, not on $T$, and must hold at the running value of
$\theta$ --- \emph{and it is delivered at the \textbf{infeasible} reading $\hat\varphi^{\pi}$ of
\cref{def:lfstep}, not at the estimator that definition instructs}, which is a third qualification of
the same kind and is \cref{rem:phirscope}\textup{(7)}. \emph{No tangential cone condition is among them}: A6 is retired
(A6, \cref{prp:tccfails}). What is $\theta$-free is the population
solution set.
What is emphatically not $\theta$-free is the conditioning of the inverse problem: $\mathcal A_\varphi$
in \eqref{eq:Adef} evaluates the conditional density at a $\theta$-dependent point, so the singular
values, the degree of ill-posedness $a$ and the source index $\mu$ are all $\theta$-specific.
\end{remark}

% ===========================================================================
%  PART C.  STEP 2 -- THE BIERENS EQUIVALENCE
% ===========================================================================

\begin{lemma}[Bierens core]\label{lem:bierenscore}
A bounded, one-to-one transform and the exponential weight convert a conditional moment restriction into
an unconditional family indexed by \emph{any} set of positive Lebesgue measure.
Suppose B1(i) and (ii) hold, let $\mathcal T'\subseteq\IR^{p}$ be Borel with
$\mathrm{Leb}(\mathcal T')>0$, let $V$ be a random vector in $\IR^{p}$, and let $\varepsilon$ be a random
variable with $E|\varepsilon|<\infty$. Then
\[
E\big[\varepsilon\mid V\big]=0\ \text{a.s.}
\qquad\Longleftrightarrow\qquad
E\big[\varepsilon\,e^{t'\Phi(V)}\big]=0\ \ \text{for every }t\in\mathcal T' .
\]
\emph{\textbf{B1(iii) is not a hypothesis of this lemma, and this is a restatement and not a change of
content}}: the proof reads of the index set only that its Lebesgue measure is positive, and that is now
\emph{quantified in the statement} instead of asserted about the proof. B1(iii) splits under the
assumption block's \textbf{three registers} paragraph: $\mathcal T\subseteq\IR^{p}$ is its
\emph{declaration} and is standing, while $\mathrm{Leb}(\mathcal T)>0$ is its \emph{clause}, and that
clause is exactly the instance $\mathcal T'=\mathcal T$.
\emph{So every result that instantiates this lemma at the maintained index set carries B1(iii) on its own
line} --- \cref{lem:bierens}, \cref{cor:opsol}, \cref{prp:kerchar}, \cref{prp:localise} and
\cref{lem:density} all do, and none of their lists moves ---
\emph{while \cref{lem:adjoint}, which instantiates it at $\mathcal T_1=\mathcal T\cap\operatorname{supp}\Pi_t$
and \textbf{derives} $\mathrm{Leb}(\mathcal T_1)>0$ from B1(iv.b), correctly does not}. That omission was
previously legible only from a comment and from one sentence inside a proof; it is now a consequence of
this statement. \emph{And the two sites that substituted an index set into a \textbf{proof} --- Step~5 of
\cref{lem:bierens}'s and Step~4 of \cref{lem:adjoint}'s --- now instantiate a \textbf{statement}}, which
is the move \cref{lem:sandwich}\textup{(i$'$)} was created to make available.
\emph{Every set this lemma is instantiated at in this corpus is Borel}: $\mathcal T$ by
\cref{def:spaces}, which equips it with $\mathfrak B(\mathcal T)$, and by B1(iv.b), which quantifies over
Borel subsets of it; $\mathcal T_0\setminus N_{\mathcal T}$ because $\mathcal T_0$ is Borel by B1(iv.b)
and $N_{\mathcal T}$ is a finite union of $\Pi_t$-null sets supplied by Fubini on \cref{def:spaces}'s
\emph{uncompleted} product $\sigma$-field, hence Borel; and $\mathcal T_1$ as the intersection of a Borel
set with a closed one. \emph{$\mathcal T'$ is bound by this lemma's own quantifier and is carried to no
other result.}
\end{lemma}

\begin{lemma}[L-3: quantifier management in $u$]\label{lem:quantu}
A family of conditional restrictions indexed by a real number reduces to a countable family, with a
single exceptional null set.
Fix $(\varphi,\theta)\in\mathcal V\times\IR^{p}$. Then a single regular conditional distribution of
$(\Delta U(\theta,\varphi),W)$ given $\Delta Z$ supplies \emph{three} objects at once: a version of the
conditional distribution functions
$F^{(\theta,\varphi)}_j(u\mid z):=P\big(\Delta U(\theta,\varphi)\le u\mid W=j,\Delta Z=z\big)$,
$j\in\{0,\dots,\mathsf L-1\}$, that is a distribution function in $u$ for every $z$; a version
$\varrho_j(\cdot)$ of $P(W=j\mid\Delta Z=\cdot)$ at each of those same $j$; and, for each
$l\in\{1,\dots,\mathsf L-1\}$, a version
of the left-hand side of \eqref{eq:binaryquant}, by which is meant one function of $(u,z)$, jointly
measurable and right-continuous in $u$ at \emph{every} $z$, agreeing at each fixed $u$ almost surely
with the conditional expectation that left-hand side names. \emph{No symbol is introduced for the
third}, the corpus's short capitals being claimed, $\Lambda$ is the integrand of convention
\textup{(P4c)} in \texttt{proofs.tex} and $\mathcal M$ is the unconditional moment of
\cref{def:operator}, so it is called ``the left-hand side, read at the version
fixed here''. Writing
$\bar F^{(\theta,\varphi)}(u\mid z):=\sum_{j=0}^{\mathsf L-1}\varrho_j(z)F^{(\theta,\varphi)}_j(u\mid z)$
for the $W$-free conditional distribution function, with $\varrho_0:=1-\mathbf 1'\varrho$,
\begin{equation}\label{eq:binaryquant}
E\big[\widetilde W_l\,\un\{\Delta U(\theta,\varphi)\le u\}\mid\Delta Z=z\big]
\;=\;\varrho_l(z)\Big\{F^{(\theta,\varphi)}_l(u\mid z)-\bar F^{(\theta,\varphi)}(u\mid z)\Big\}
\end{equation}
for every $u\in\IR$, every $l\in\{1,\dots,\mathsf L-1\}$ and \textbf{every} $z$, both sides being read
at the versions fixed here. \emph{That the second quantifier is ``every $z$'' and not
``a.e.\ $z$'' is the whole content of the lemma}: a null set allowed to move with $u$ says nothing at
uncountably many $u$, the union of uncountably many null sets not being null, and it is the regular
version, one kernel, read at all $u$ simultaneously, that removes the dependence. \emph{The
freedom in $F_j$ at a null cell is untouched}: an arbitrary distribution function there is multiplied by
$\varrho_j(z)=0$ on both sides, so \eqref{eq:binaryquant} holds at every $(u,z)$ for every such choice.
Consequently, if the left-hand side of \eqref{eq:binaryquant}, read at that version, vanishes
for every $u$ in a countable dense subset of $\IR$, every $l$, and a.e.\ $z$, then it vanishes for every
$u\in\IR$, every $l$, and every $z$ outside a single null set not depending on $u$ or on $l$, 
\emph{a finite union of the $\mathsf L-1$ per-cell null sets, each itself a countable union over the
dense set, which is why the index costs nothing here}.
% QUANTIFIERS MADE TO MEET (Phase 6b), STATEMENT AND PROOF IN ONE PASS.  Through
% dispatch 25 this line read "for every u, every l and a.e. z", while the
% "Consequently" clause promised a SINGLE null set not depending on u -- and Step 2
% of the proof delivered only "a.e. z", PER u.  At uncountably many u that is not a
% null set and the two quantifiers did not meet.  What closes it was already in
% Step 1 and unprinted: the regular version is fixed "for EVERY z, not merely for
% almost every z", so the right-hand side is a jointly measurable, right-continuous
% selection valid simultaneously at all u; reading the LEFT-hand side at the version
% the same kernel supplies makes eq:binaryquant an algebraic identity between the
% pieces of one kernel, hence true at every (u,z).  NO HYPOTHESIS MOVES and no
% assumption edge is created: existence of the regular conditional distribution is
% Step 1's Polish-space argument, already printed, and A4(i) remains off the list
% for the reason stated below.  Blast radius nil -- Consumers: none -- but the
% lemma is printed and cited for what the identity delivers.
% PRESERVED DELIBERATELY: the comment under lem:shiftinv in
% results_s2_identification.tex (and its twin in results_main.tex) records that
% eq:binaryquant "holds for ANY version of the conditional distribution functions".
% It still does; the sentence above says so in print.  That comment is STALE in a
% DIFFERENT respect -- it writes the region in the binary form D := {z :
% varrho(1-varrho) = 0}, superseded by eq:degenreg -- and neither file is this
% pass's to edit; reported to the orchestrator.
\emph{At $\mathsf L=2$ the display is the printed one}: $l$ takes the single value $1$, $\varrho_0=1-\varrho_1$,
so $\bar F=(1-\varrho_1)F_0+\varrho_1F_1$ and
$\varrho_1\{F_1-\bar F\}=\varrho_1(1-\varrho_1)\{F_1-F_0\}$.
\emph{\textbf{Consumers: none}, as of 2026-08-09}, on the model of
\cref{lem:nusigma}\textup{(ii)}. Its only proof-step consumer was \cref{lem:bierens} Step~3, which has
been lifted onto the cell-factorisation route of \cref{rem:localise}\textup{(1)}, a route that does
not pass through this lemma, and \cref{cor:opsol} now inherits from that route rather than from here.
The lemma is kept because it is correct, because \eqref{eq:kernelcompareci}'s clause sends a reader here
for what the identity does and does not deliver at $\sigma(\Delta Z)$, and because a reader following a
citation should land on a statement rather than on a gap; it is \textbf{not} available as an inherited
step, since nothing inherits it.
\emph{A4(i) is \textbf{not} among its hypotheses, and that is a consequence of the proof and not an
omission}: the extension from a countable set of $u$ to all of $\IR$ propagates a zero along a
right-continuous function and \emph{divides by no prefactor}, so no lower bound on $\varrho_l$ is needed
and no case split on the degenerate region is made.
\end{lemma}

% NEITHER A1 NOR A4(i) IS A HYPOTHESIS of lem:bierens or cor:opsol.
%
% A1 would make lem:bierens VACUOUS: A1's first clause IS statement (a), so
% under A1 all four of (a)-(d) are true and the printed equivalence asserts
% nothing.  What is proved is that (a) <=> (b) <=> (c) <=> (d) whether or not
% the exogeneity restriction happens to hold.
%
% A4(i) IS IDLE, and the round-2 defence of it was wrong.  It claimed that
% without overlap "the restriction is vacuous wherever varrho in {0,1}, and (c)
% would not recover (a) on that region".  But on that region W is almost surely
% CONSTANT given Delta Z = z, and a conditionally degenerate random variable is
% conditionally independent of everything -- so (a) holds automatically there
% and there is nothing to recover.  Off it, lem:quantu needs only pointwise
% positivity.  And Wtilde = 0 almost surely on the degenerate region, so
% \eqref{eq:binaryquant} is trivially true there.  The old argument conflated
% "the restriction has no bite there" with "the equivalence fails there".
%
% NOTE ALSO WHAT THIS BUYS, and it is worth a sentence in the paper: with A4(i)
% gone, lem:bierens requires NO MOMENT CONDITION on Delta Y, Delta U or
% Delta Z, and NO LEBESGUE DENSITY for Delta U or Delta Z.  That is exactly
% what boundedness of Phi purchases -- see B1(i) and rem:Tcond -- and it is
% currently invisible in the statement.

\begin{lemma}[L-2: density of the closed span]\label{lem:density}
The debiasing weights representable through the Bierens family are dense, so the choice of weighting
measure is immaterial on the $\theta$-side.
Suppose B1(i)--(iii) and (v) hold, and let
\[
\mathcal G\;:=\;\overline{\operatorname{span}}\Big\{(u,z)\mapsto\varsigma(u)\,w\big(t'\Phi(z)\big):\ \varsigma\in C_b(\IR),\ t\in\mathcal T\Big\}
\]
with closure taken in $L^{2}(P_{\Delta U,\Delta Z})$. Then
\begin{enumerate}
\item[(i)] $\mathcal G=L^{2}(P_{\Delta U,\Delta Z})$;
\item[(ii)] the set
$\mathcal G_\gamma:=\big\{L:L(u,z)=\int_{\mathcal T}\lambda(u,t)w(t'\Phi(z))\,\gamma(dt),\ \gamma\text{ finite signed satisfying B1(v)},\ \lambda\text{ bounded}\big\}$
is dense in $L^{2}(P_{\Delta U,\Delta Z})$;
\item[(iii)] $\overline{\{L-E[L\mid\Delta Z]:L\in\mathcal G_\gamma\}}
=\{L\in L^{2}(P_{\Delta U,\Delta Z}):E[L\mid\Delta Z]=0\}$, the subspace of weights \emph{satisfying}
N1.
\end{enumerate}
\end{lemma}

\begin{remark}[Why the weight is the exponential, and what a general analytic weight would cost]\label{rem:swintercept}
B1(ii) fixes $w=\exp$, and the restriction is substantive rather than expositional.

The textbook condition is not enough. Suppose $w$ were required only to be real-analytic and
non-polynomial, as in the usual statement of the device. Take $V=\pm1$ with probability $\tfrac12$ each,
let $\Phi$ be the maintained componentwise $\arctan$, an \emph{odd} function, so
$\Phi(-1)=-\Phi(1)$, and set $\varepsilon:=V$. Then $E|\varepsilon|<\infty$ and
$E[\varepsilon\mid V]=V\neq0$, yet for $w=\cos$ and \emph{every} $t$,
\[
E\big[\varepsilon\,\cos\big(t\Phi(V)\big)\big]
\;=\;\tfrac12\cos\big(t\,\Phi(1)\big)-\tfrac12\cos\big(-t\,\Phi(1)\big)\;=\;0 .
\]
B1(i)--(iii) hold and the converse direction of \cref{lem:bierenscore}, read at
$\mathcal T'=\mathcal T$, fails. The mechanism is that an
odd $\Phi$ against an even $w$ exposes only the even part of the image measure, so Step 4 of the proof
recovers the even moments and no others. The failure is not confined to the abstract lemma: with
$\Delta Z$ scalar and two-valued at $\pm c$,
$\mathcal M(u,t,l)=\cos\big(t\,\Phi(c)\big)\,E[\widetilde W_l\un\{\Delta U\le u\}]$, so (c) of
\cref{lem:bierens} holds
whenever the \emph{unconditional} restriction does, while the $\Delta Z$-conditional restriction can fail
through opposite-signed violations that cancel across the two values of $\Delta Z$.

The correct general statement, and what it would cost. The repair is not to strengthen $\Phi$ but
to add an intercept. Index the family by $(t_0,t)$ over a subset of $\IR^{1+p}$ of positive Lebesgue
measure and use $\{w(t_0+t'\Phi(z))\}$. For $w$ real-analytic and non-polynomial, each derivative
$w^{(k)}$ is analytic and not identically zero, otherwise $w$ is a polynomial, so its zero set is
discrete and hence countable; therefore all but countably many $c\in\IR$ satisfy $w^{(k)}(c)\neq0$ for
every $k\ge0$, and the argument runs at such a $c$. This is the Stinchcombe--White family and the result
is \citet[Thm.~3.1]{stinchcombewhite1998}; it is load-bearing here, not a ``see also''. Adopting it
would change the range space to $L^{2}(\IR\times\mathcal T^{+},\Pi)$ with
$\mathcal T^{+}\subseteq\IR^{1+p}$, and would require the envelope bookkeeping of B1(iv)--(v) to be redone
with $\bar w(t_0,t)=\sup_{|v|\le|t_0|+c_\Phi\|t\|_1}|w(v)|$ throughout \cref{prp:wellposed},
\cref{prp:frechet} and \cref{prp:compact}.

\emph{Why it is not adopted.} $w=\exp$ is the maintained choice of the estimator, it is what
\citet[Thm.~1]{bierens1990} characterises, so that the citation is exact rather than approximate, 
and the additional generality is consumed nowhere in the programme. The general family is recorded here
so that the statement is on the record without the range space paying for it.
\end{remark}

\begin{remark}[On $\mathcal T$, on $\gamma$, and on what compactness is for]\label{rem:Tcond}
It is worth recording what the identification layer costs, because the statement of \cref{lem:bierens}
does not show it. The equivalence requires no moment condition whatsoever on $\Delta Y$, $\Delta U$
or $\Delta Z$, and no Lebesgue density for $\Delta U$ or $\Delta Z$. $\Delta Z$ may be discrete,
$\Delta U$ may have atoms, and neither need have a finite mean. That is exactly what boundedness of
$\Phi$ in B1(i) purchases: it makes $\varepsilon\,w(t'\Phi(\Delta Z))$ integrable for the bounded
$\varepsilon=\widetilde W_l\un\{\Delta U\le u\}$, \emph{scalar at each contrast index, which is the
arity \eqref{eq:kernel} and \cref{lem:bierenscore} are read at}, with no condition on the law of
$\Delta Z$, and it makes
the image measure of \cref{lem:bierenscore} compactly supported, hence determined by its moments. The
smoothness and moment conditions of C1, C2$'$ and A1b(iv) enter only when the \emph{operator} is
differentiated or the \emph{parameter} is identified, never in the equivalence itself.

Compactness of $\mathcal T$ is \emph{sufficient but not necessary} and should not be assumed. What the
proofs use is B1(iv.c) and B1(v), which for $w=\exp$ read
$\Pi_u(\IR)\int e^{2c_\Phi\|t\|_1}\Pi_t(dt)<\infty$ and $\int e^{c_\Phi\|t\|_1}|\gamma|(dt)<\infty$; a
compact $\mathcal T$ carrying a finite $\Pi_t$ and a finite $\gamma$ satisfies both, but so do many
unbounded index sets. Plan \S18.3 introduced compactness in order to license differentiation under the
$t$-integral. That is a specific interchange with a specific dominating function, and B1(v) now supplies
it: differentiating $L(u,z)=\int_{\mathcal T}\lambda(u,t)w(t'\Phi(z))\gamma(dt)$ in $u$ under the integral
sign is licensed by the bound $|\partial_u\lambda(u,t)\,w(t'\Phi(z))|\le\Lambda_1(t)\bar w(t)$, valid for
all $(u,z)$, together with $\int_{\mathcal T}\Lambda_1\bar w\,d|\gamma|<\infty$.

Two clauses of B1(iv) are minimal and should not be strengthened back. Step 5 of the proof of
\cref{lem:bierens} consumes only (iv.a), full topological support of $\Pi_u$, which is what makes the
zero set of $\mathcal M(\cdot,t,l)$ dense in $\IR$, and (iv.b), non-mutual-singularity of
$\mathrm{Leb}|_{\mathcal T_0}$ and $\Pi_t|_{\mathcal T_0}$, which is exactly the statement that no
$\Pi_t$-null set can swallow $\mathcal T_0$. Neither absolute continuity nor equivalence is used in
either factor.

Two things about (iv.b) are worth stating precisely, because it is easy to describe it wrongly. First,
by the Lebesgue decomposition it is equivalent to $\Pi_t$ having a nonzero absolutely continuous part,
so it excludes \emph{every} Lebesgue-singular $\Pi_t$ and not merely the atomic ones: a Cantor measure is
non-atomic, has uncountable support, and is excluded. Second, it is sufficient and not necessary. What
the argument really needs is an analytic-uniqueness condition on $\operatorname{supp}\Pi_t$: by Step 2 of the
proof of \cref{lem:bierenscore}, $t\mapsto\mathcal M(u,t,l)$ is real-analytic on all of $\IR^{p}$ at
each fixed $(u,l)$, so it
suffices that its zero set cannot contain $\operatorname{supp}\Pi_t$ up to a $\Pi_t$-null set unless it
vanishes identically. When $p=1$ that holds for any $\Pi_t$ whose support has an accumulation point, a
Cantor measure included, since a Cantor set is perfect, and \cref{lem:bierens} survives such a
$\Pi_t$. When $p\ge2$ the requirement is that $\operatorname{supp}\Pi_t$ not lie inside a proper
real-analytic variety, which positive Lebesgue measure guarantees and which many Lebesgue-null sets also
satisfy, but which perfectness alone does not: a hyperplane is perfect and is the zero set of a
coordinate. (iv.b) is retained as the clean sufficient version.

What (iv.b) does exclude in every dimension is a finite grid. \emph{The growing-family condition B1(vi)
is \textbf{retired}}, because the clause describes an estimator this programme does not implement:
\cref{def:hatTexp} integrates against $\Pi_t$ over all of $\mathcal T$ and never over a finite subfamily.
Identification of $\vdagger$ here is A1b(iii), assumed outright. The exclusion (iv.b) performs is
recorded here so that it remains visible to a reader who reinstates the growing-grid design.
% ARCHAEOLOGY.  B1(vi) was classed an identification assumption rather than an
% approximation device on exactly the finite-grid ground above; the retirement is
% not a reversal of that classification, and the struck clause is kept in the
% corrected register in assumptions.tex rather than deleted.

The correction measure plays two different roles on the two sides of the construction.
\emph{On the estimation side} the
statement stands as written: by \cref{lem:density} the estimator \emph{solves for} $L$ inside a dense
class rather than integrating a fixed weighting measure against a continuum, so no $t$-integral is ever
computed in forming $\hat L$, and no weighting measure has to be chosen. This is also the surviving
argument against a continuum-GMM formulation: that would require choosing a weighting measure on
$\mathcal T$ and regularising the inverse of its covariance operator. On the population side the
integral representation is not optional. The representer selected by the adjoint equation is an element
of $L^{2}(\Pi)$, and what makes it evaluable at $(\Delta U,\Delta Z)$ is exactly the induction map
$\mathcal W$ of \cref{lem:adjoint}(ii), which is an integral against $\Pi_t$. The correction measure of
B1(v) is therefore not a free object \emph{on that side}: \cref{lem:adjoint}(ii) identifies it with
$\Pi_t$, and \emph{at $\gamma=\Pi_t$, and only there}, the first half of B1(v) is implied by B1(iv.c)
rather than assumed, since
$\int\bar w\,d\Pi_t\le(\Pi_t(\mathcal T)\int\bar w^{2}d\Pi_t)^{1/2}<\infty$ by the Cauchy--Schwarz
inequality. That implication must not be widened.
\Cref{lem:density}(ii) quantifies over \emph{every} finite signed $\gamma$
satisfying B1(v), and its proof constructs an \emph{atomic} $\gamma\ne\Pi_t$ for which the first half is
verified directly; so the first half is live on the estimation side and is retained for that reason,
while the second half is live here, in this remark, and is declined by \cref{cor:elldiff} for the reason
recorded there. What reconciles the two sides is \cref{lem:adjoint}(ii)'s density claim: the sieve and
the Bierens family parameterise the same closure, so the estimator may search the larger class without
losing the population object.

Finally, \cref{lem:bierens} is stated at the equivalence, not at a strict inclusion. The claim that
$\mathcal R(\mathcal A^\ast_{\mathrm{cond}})$ strictly contains the range of the unconditional adjoint is
not proved here and is not claimed: with $W$ taking finitely many values the two moment classes are not
nested, since, under A1,
$\{g:E[g\mid\Delta U,\Delta Z]=0\}=\{\widetilde W'\ell(\Delta U,\Delta Z)\}$ is an equality of sets, 
\emph{the $\mathsf L-1$ contrasts spanning the conditionally centred functions of $W$, which is the
spanning statement accompanying \eqref{eq:contrast}, and at $\mathsf L=2$ the single contrast doing it}
, and
the comparison across the two modelling assumptions is heuristic. The conditional route is defended on
economic grounds and on internal consistency with the maintained exogeneity assumption, not on a weaker
source condition.
\end{remark}

% ===========================================================================
%  PART D.  RECONCILIATION WITH THE SOURCE DRAFT
% ===========================================================================

% ############################################################################
% ##  rem:draftrec.  SHORTENING PASS 2026-08-01: this remark was 94 compiled
% ##  lines of reconciliation with the source draft -- a DRAFT RECORD, in the
% ##  dispatch's own words "named for what it is".  The four defect narratives
% ##  (draft line ranges :93--103, :120--122, :124--126, :172--177, and the
% ##  plan Sec.18.3 assertion they correct) are archived here; what is KEPT
% ##  compiled is the mathematics that later results consume and the two
% ##  caveats a referee must see:
% ##    -- eq:kernelcompare and the p_X = varrho / instrument-exogeneity caveat,
% ##       which is a restriction on the DESIGN and not on a draft;
% ##    -- the surviving pi^{-1} in eq:adjointdisplay, which is why A2 is
% ##       imposed on nu/pi and not on nu.
% ##  THE FULL ARCHAEOLOGY, clause by clause.
% ##  (i)  Draft :93--103 writes T(varphi, theta) as a double integral over
% ##       (x_1, x_2) of a difference of conditional distribution functions and
% ##       then treats it as a function of (w, z, u); the arguments are never
% ##       declared and the (x_1, x_2)-integration is in fact an expectation.
% ##       eq:Mdef makes it the expectation, removes the W-argument by the
% ##       contrast Wtilde (legitimate because W is binary) and the z-argument by
% ##       the Bierens index t, leaving an object indexed by (u, t) with u scalar
% ##       -- which is why the LF machinery carries over with an enlarged index
% ##       set and no new curse of dimensionality.
% ##  (ii) Draft :124--126's derivative is a raw contrast between the two W-cells,
% ##       unconditional in Z, while the draft's own exogeneity assumption
% ##       ass:condind is conditional.  Plan Sec.18.3 asserts that the two
% ##       operators differ only by the scalar weight varrho(1 - varrho) and so
% ##       share a range and a kernel.  THAT IS FALSE; the surviving comparison
% ##       is between kernels and is compiled below.
% ##  (iii) Draft :172--177's adjoint display: read left to right, the outer
% ##       int tilde-varphi'(v) pi(v) dv is the L^2(pi) pairing, everything to its
% ##       right is (A* ell)(v), the inner braces are nu_g(x) -- the xi-integration
% ##       of a conditional expectation against f_{X_1, X_2} marginalises the other
% ##       period out, which is eq:nudef -- and the kernel 1(x >= v) is D^{-1},
% ##       present because the draft parameterises in derivatives.  THREE defects:
% ##       (a) pi(v) appears in numerator and denominator, which is a CATEGORY
% ##       error (the display writes the PAIRING, not the adjoint) and a trap, since
% ##       a reader who cancels them concludes the L^2(pi) weighting is immaterial;
% ##       (b) the subscripts (theta, varphi) are unsupported by the expression
% ##       beneath them, the density being written f_U(u) with neither argument
% ##       shifted, whereas :124--126 evaluates it at varphi(x_2) - varphi(x_1) + u
% ##       -- the defect is SYMMETRIC in the two subscripts and the repair is to
% ##       write the shifted density, not to drop a subscript, both indices being
% ##       substantive (prp:frechet(ii) for varphi, rem:R1(iv) for theta); an
% ##       earlier version of this remark called the theta subscript vacuous while
% ##       defending the varphi one, conflating the operator with its zero set;
% ##       (c) the display has no u-integration -- its integrand carries a free u
% ##       inside psi(u,1) - psi(u,0) and integrates only over (x, xi, v), so the
% ##       left side is a function of u and the right side a number.  The missing
% ##       object is the range measure Pi.  The same omission explains the
% ##       W-argument mismatch between :124--126 and :172--177.
% ##  (iv) Two symbols: the kernel 1(x >= v) is the kernel of D^{-1}, the
% ##       integration operator of the DERIVATIVE parameterisation, and does not
% ##       appear in levels; the draft's psi in the adjoint display is a generic
% ##       range element and is renamed ell, psi being reserved for the moment
% ##       function.
% ############################################################################

% REMARK rem:draftrec MOVED to working_notes.tex on 2026-08-10 -- zero printed
% cross-references anywhere in the compiled corpus, so nothing points at it and
% no stub is needed.  Verbatim in that file, which is input by no document.

% ===========================================================================
%  PART E.  STEP 3 -- THE ADJOINT, THE nu_g ALGEBRA, AND THE IMPOSSIBILITY
%           RESULT
%
%  Contents, in the order of memo Sec.10 Step 3:
%     lem:nusigma   (L-1)  the pair (nu_g, sigma_g); Delta bounded; ker Delta
%     lem:adjfubini (L-4)  the interchange in the ADJOINT direction
%     lem:levels    (L-5)  levels versus derivatives
%     lem:adjoint          the closed form, in three parts.  Part (ii) is what
%                          closes escalation item N-10 BY CONSTRUCTION.
%     prp:nophifree  (R2)  the impossibility result
%     rem:nophifree        its interpretation, and the Feve--Florens link
%     cor:udep      (O-11) the corollary that forces partial_u ell_star != 0
%     prp:kerchar          characterisation of ker A_cond -- ADDITIVE, see the
%     cor:a1cprim          escalation note N-12 in theory_memo.md
%     rem:ffcompare        what is shared with Feve & Florens and what is new
%
%  TYPING NOTE, and it is a CORRECTION (finding F-15).  L, ell = partial_u L,
%  ell_dagger, ell_star and the weight c of cor:udep are R^p-VALUED, one
%  coordinate per coordinate of theta.  The round-2 glossary typed them as
%  scalars while typing psi as R^p-valued, and the two cannot both hold:
%      psi = Delta Z * Delta U + Wtilde * L,     Delta Z is p x 1,
%      J   = E[Delta Z Delta Z'] + E[Wtilde ell_star Delta Z'],  claimed p x p,
%  so a scalar L makes psi ill-formed and a scalar ell_star makes the second
%  term of J a 1 x p row added to a p x p matrix.  A2 already says the right
%  thing -- "EACH COORDINATE of nu/pi lies in R(A*_cond)" -- so the correction
%  is to the type declaration, not to any assumption.
%  CONSEQUENCE FOR PART E.  lem:adjoint, prp:nophifree and cor:udep are stated
%  for a SCALAR ell (resp. a scalar-or-vector g) and are applied COORDINATEWISE:
%  A*_cond acts on each coordinate of ell separately, and \eqref{eq:tilted} is
%  read coordinatewise when g is R^p-valued.  Nothing in the proofs changes;
%  every step is linear in the coordinate index.
%
%  WHY B1(iii) IS NOT A HYPOTHESIS OF lem:adjoint, since a reader will look for
%  it.  The density half of part (ii) invokes lem:density Step 1, which reaches
%  lem:bierenscore, whose only requirement on the index set is positive Lebesgue
%  measure.  The set actually used is T_1 = T intersect supp(Pi_t), and
%  Leb(T_1) > 0 is DERIVED from B1(iv.b) inside the proof rather than inherited
%  from B1(iii).  Listing B1(iii) would therefore be an assumption the result
%  does not consume.
%
%  CLOSED 2026-08-24 BY RESTATEMENT, AND THIS ENTRY IS KEPT AS THE RECORD OF HOW
%  IT WAS OPEN.  Through that date the ground lived HERE, in a comment no document
%  renders, and in one sentence inside lem:adjoint's Step 4 -- i.e. it was a claim
%  about lem:bierenscore's PROOF, which a reader of lem:bierenscore's STATEMENT
%  could not check, and which the register rule could not adjudicate because
%  B1(iii) is half declaration (T subset R^p) and half clause (Leb(T) > 0) and the
%  corpus had no name for the first half.  lem:bierenscore is now STATED at a
%  generic Borel T' with Leb(T') > 0, so B1(iii)'s clause is the instance T' = T
%  and lem:adjoint's omission of it is a consequence of the statement.  This is the
%  lem:sandwich(i') move: STATE the weaker list rather than assert that the proof
%  holds on it.  It also removes Part E's only appeal-to-a-proof-at-substituted-
%  hypotheses, at both of its sites (lem:bierens Step 5, lem:adjoint Step 4).
%  NO CONSUMER'S HYPOTHESIS LIST MOVES: lem:bierens, cor:opsol, prp:kerchar and
%  prp:localise carry B1(i)--(iv) and lem:density carries B1(i)--(iii), so all five
%  instantiate at T' = T with the clause already on their own lines.
% ===========================================================================

\begin{lemma}[L-1: the difference and sum functionals]\label{lem:nusigma}
The period-difference operator is bounded with kernel the constants, and the two functionals it generates
differ in exactly one respect: the difference functional integrates to zero and the sum functional does
not.
Suppose C1(ii) holds and let $g$ be a random variable with $E|g\,h(X_1)|<\infty$ and $E|g\,h(X_2)|<\infty$
for the $h$ named in each clause. Then
\begin{enumerate}
\item[(i)] $\|\Delta h\|_{L^{2}(P)}\le2\,\bar c_\pi^{1/2}\|h\|_{L^{2}(\pi)}$ for every
$h\in L^{2}(\mathcal X,\pi)$, \emph{the whole ambient space and not only $\mathcal V$}, so
$\Delta$ extends from the domain \cref{def:spaces} gives it to a bounded operator
$L^{2}(\mathcal X,\pi)\to L^{2}(P)$ with the same norm bound. \emph{The quantifier is stated at
$L^{2}(\mathcal X,\pi)$ because that is what the proof establishes and because three later results
consume it there}: the proof of \textup{(i)} uses $(a-b)^{2}\le2a^{2}+2b^{2}$ and C1(ii) and \textbf{never}
$\int h\,d\pi=0$, and C1(ii) is also what makes $\Delta h$ well defined for a $\pi$-equivalence class, since
$f_{X_j}\le\bar c_\pi\pi$ makes $P_{X_j}\ll\pi$ for $j=1,2$. \Cref{lem:prelimtheta} Step~4,
\cref{lem:epterm} Step~3 and \cref{lem:sandwich} Steps~3 and~6$'$ all apply this bound at
$h_k=\hat\varphi^{(-k)}-\vdagger$, whose membership of $\mathcal V$ would otherwise need A1b(i);
\emph{the wider quantifier removes that need at all three at once}, which is strictly preferable to
adding A1b(i) to three hypothesis lists.
% ARCHAEOLOGY.  Rounds 1--13 stated clause (i) "for every h in V"; the cost was
% the three unlisted A1b(i) uses named above.  An earlier record proposed adding
% A1b(i) to the three lists instead.
Clause \textup{(ii)} is a different matter and genuinely needs $h\in\mathcal V$;
\item[(ii)] \emph{\textbf{Consumers: none}, as of 2026-08-05.} Its only consumer was the corollary
withdrawn at \cref{rem:a1cwithdrawn}, and with it A1b(ii), which enters no other result, became
a clause no theorem rests on. The clause is kept because it is the statement of a fact about $\Delta$
that the surrounding lemma's other clauses are read against, and because a reader who follows a stale
citation should land on it rather than on a gap; it is not available as an inherited step, since
nothing inherits it. \emph{The content:} $\Delta h=0$ almost surely if and only if $h$ is almost surely
constant, provided A1b(ii)
holds; consequently, under A1b(ii) alone, $\ker\Delta\cap\mathcal V=\{0\}$, \emph{the passage from
``constant'' to ``zero'' is $\int h\,d\pi=0$ for $h\in\mathcal V$, which is \cref{def:spaces} and not an
assumption; A1b(i) is not needed, since it asserts $\vdagger\in\mathcal V$ and $h$ here is generic;}
% ARCHAEOLOGY.  Rounds 1--7 cited A1b(i) at this clause.
\item[(iii)] $E[g\,\Delta h]=\int_{\mathcal X}\nu_g(x)\,h(x)\,dx=\langle\nu_g/\pi,\,h\rangle_{L^{2}(\pi)}$,
for each $h$ satisfying the integrability proviso above, in particular for every $h\in\mathcal V$ and
for $h\equiv1$;
\item[(iv)] $\int_{\mathcal X}\nu_g\,dx=0$ and $\int_{\mathcal X}\sigma_g\,dx=2\,E[g]$, where
\begin{equation}\label{eq:sigmadef}
\sigma_g(x)\;:=\;E[g\mid X_2=x]\,f_{X_2}(x)\;+\;E[g\mid X_1=x]\,f_{X_1}(x);
\end{equation}
\item[(v)] if in addition A1 holds and $g=\Xi(\Delta U,\Delta Z)\,\widetilde W$ for an integrable
$\IR^{p\times(\mathsf L-1)}$-valued $\Xi$, \emph{the direction to the \textbf{left} of the contrast,
which is the only order that multiplies, and the type is \cref{def:dirclass}'s}, then $E[g]=0$, hence
$\int_{\mathcal X}\sigma_g\,dx=0$ as well. \emph{The conclusion is coordinatewise in the $p$ rows}, so it
holds verbatim for an array-valued $\Xi$ whose last index is the one contracted against $\widetilde W$;
that is how \cref{rem:momentread}\textup{(2)} reads it. At $\mathsf L=2$, $\widetilde W$ is a
scalar, $\Xi$ is $\IR^{p}$-valued and $\Xi\widetilde W=\widetilde W\Xi$ is the printed object.
\end{enumerate}
\end{lemma}

\begin{lemma}[L-4: interchange in the adjoint direction]\label{lem:adjfubini}
The $(u,t,l)$-integration defining the adjoint may be taken inside the expectation.
Suppose B1(i), (ii), (iv), C1(ii) and C2$'$(i) hold, and recall from \cref{def:model} that $P$ is a
probability measure, hence $\sigma$-finite. \emph{\textbf{B1(iv) is cited as a block because this proof
opens all four of its clauses}}, which is what the block register requires: the interchange \textbf{is} this lemma's
conclusion and not a constant inside it,
it is Fubini for $\Pi\otimes P$, and that needs $\Pi$ $\sigma$-finite --- which it is because B1(iv)
makes $\Pi$ a product of \emph{three} finite measures, $\Pi_u$ and $\Pi_t$ by \textup{(iv.a)} and
\textup{(iv.b)} and the finitely supported $\Pi_{\mathsf w}$ by \textup{(iv.d)}. Drop any one of the
three and there is no interchange left to state. \emph{\textbf{It is not the only line that carries the
block, and through 2026-08-25 this statement said it was}}: \cref{lem:adjoint} carries it on the same
ground, its own proof opening all four clauses at named steps, and so does \cref{prp:tccfails}; while
\cref{lem:solset} carries it by inheritance, its proof naming no clause of B1 at all and invoking
\cref{lem:adjoint}\textup{(ii)} and \cref{prp:orthog}, which carry the block between them.
\emph{No claim of exhaustiveness is made, and that is deliberate}: the assumption block's own audit of
B1(iv) block-form lines names, besides the three above, \cref{lem:bierens} and \cref{prp:kerchar}, which
consume \textup{(iv.c)} \emph{at the statement} rather than at a step, and \cref{cor:opsol},
\cref{prp:localise}, \cref{lem:biasexp} and \cref{prp:orthog}, which carry it by inheritance as
\cref{lem:solset} does, and the group \textup{(T-3)} of the main theorem. \emph{A sentence whose whole
subject is a miscount is the last place to leave a bare count standing}, which is why the hedge is
attached here and not left to the audit.
\emph{Through 2026-08-23 this line read
``\textup{(iv.c)}'' alone}, so the clause the conclusion rests on was named in the proof and nowhere a
reader of the statement would meet it. Then for every $\ell\in L^{2}(\Pi)$ and every
$h\in\mathcal V$ the function
$(u,t,l,\omega)\mapsto g_{u,t,l}(\omega)\,\Delta h(\omega)\,\ell(u,t,l)$ is $\Pi\otimes P$-integrable, with
\begin{equation}\label{eq:adjenv}
\begin{split}
\sum_{l=1}^{\mathsf L-1}\varpi_{\mathsf w}(l)&\int_{\IR\times\mathcal T}
E\big[\big|g_{u,t,l}\,\Delta h\big|\big]\,\big|\ell(u,t,l)\big|\,\Pi_u(du)\,\Pi_t(dt)\\
&\;\le\;2\,\bar c_\pi^{1/2}\,\|\bar f_0\|_\infty\,c_\Pi\,\|h\|_{L^{2}(\pi)}\,\|\ell\|_{L^{2}(\Pi)}
\;<\;\infty,
\end{split}
\end{equation}
and consequently, with $\int(\cdot)\,d\Pi$ abbreviating the same three-factor integration,
$\int E[g_{u,t,l}\Delta h]\,\ell(u,t,l)\,d\Pi
=E\big[\Delta h\int g_{u,t,l}\,\ell(u,t,l)\,d\Pi\big]$.
\emph{The constant does not move with $\mathsf L$}: the envelope reads the \textbf{per-cell} bound
$|\widetilde W_l|\le1$ inside $g_{u,t,l}$, and the $\Pi$-integration of $\bar w$ against $|\ell|$ is
bounded by $c_\Pi\|\ell\|_{L^{2}(\Pi)}$ by B1(iv.c) and the Cauchy--Schwarz inequality in $\Pi$,
$\Pi_{\mathsf w}$ being a probability measure. At $\mathsf L=2$ the sum has one term with
$\varpi_{\mathsf w}(1)=1$ and every display is the printed one.
\end{lemma}

\begin{lemma}[L-5: levels versus derivatives]\label{lem:levels}
The orthogonality condition and the source condition are the same in the two parameterisations; the
remainder bound and the rate condition are not, and levels is the weaker of the two.
\emph{``The source condition'' in that headline means \textbf{A2}, a plain range condition, and clause
\textup{(iii)} states exactly that. It does \textbf{not} mean the source condition \emph{at index
$\mu_{\text{\textup{H1}}}=s/(2a)$},
which is \textbf{not} parameterisation-invariant}: composing with $\mathcal D^{-1}_0$ raises the degree of
ill-posedness from $a$ to $a+1$ and lowers the index from $s/(2a)$ to $(s-1)/(2(a+1))$, a loss of
$\tfrac83$ units of $s$ uniformly in $a$. \Cref{rem:tccwhy}\textup{(4)} is where that computation is used
and is this lemma's principal consumer.
Suppose C1(ii) and \textbf{A0's compactness clause} hold, and let
$\mathcal D^{-1}_0:L^{2}(\pi)\to\mathcal V$ be
$(\mathcal D^{-1}_0h)(x):=\int_{[\underline x,x]\cap\mathcal X}h(v)\,dv$, with
$\underline x:=\inf\mathcal X$, followed by centring in $L^{2}(\pi)$.
\emph{A0 is cited \textbf{by clause}, and what this lemma reads is the compactness and \textbf{not} the
interval}: $|\mathcal X|<\infty$ in the constant of \textup{(i)}, and $\underline x$ finite. \emph{Through
2026-08-25 this line said the opposite and defended it}, calling the compact interval a
\textbf{declaration} in the sense of the assumption block's register paragraph and arguing that negating
it does not make the lemma false but stops it parsing, ``on a fat Cantor $\mathcal X$ there is no
interval $[\underline x,x]$ to integrate over''. \emph{Clause \textup{(ii)} below refutes that four lines
down in this same statement}: \textup{(ii)} already writes $\int_{\{x\ge v\}\cap\mathcal X}$, the
intersection form, which parses on any measurable $\mathcal X$, and \textup{(i)} is now spelled the same
way --- one intersection symbol, and nothing else moves. \emph{On a fat Cantor $\mathcal X$ the lemma is
then \textbf{true}}: the Cauchy--Schwarz bound of \textup{(i)} is untouched; injectivity holds because
the uncentred integral $F(x):=\int_{[\underline x,x]\cap\mathcal X}h\,dv$ is absolutely continuous with
$F'=h$ almost everywhere on $\mathcal X$, so that at almost every point of a level set of $F$ the
derivative vanishes, and $F$ constant $\pi$-a.e.\ --- hence Lebesgue-a.e., $\pi$ being bounded below by
C1(ii) --- forces $h=0$ a.e.\ on $\mathcal X$; \textup{(ii)} and \textup{(iv)} are unchanged in statement
and in proof; and \textup{(iii)} is unchanged in statement while \emph{its proof is not, and that is
recorded here rather than left to the proof}. \emph{Through 2026-08-26 this sentence read
``\textup{(ii)}--\textup{(iv)} are unchanged'', which is true of the three statements and false of one of
the three proofs}: \textup{(iii)}'s density step asserted that
$\mathcal R(\mathcal D^{-1}_0)$ contains every centred $C^{1}$ function on $\mathcal X$, an interval
argument that fails at $\zeta(x)=x$ on a fat Cantor $\mathcal X$ --- a density-point argument forces
$h'=1$ a.e.\ on $\mathcal X$, so the primitive is $x\mapsto|\mathcal X\cap[\underline x,x]|$, which
differs from $x$ by a function strictly decreasing across every complementary gap and therefore
constant on no set meeting both sides of one. \emph{The clause survives because injectivity of
$(\mathcal D^{-1}_0)^{\ast}$ on $\mathcal V$ is clause \textup{(ii)} itself} --- read at
$\nu_g=\zeta\pi$, after extending the tail function continuously from $\mathcal X$ to $\IR$ --- so
density becomes a conclusion of the proof rather than a step in it, and neither step reads
connectedness. \emph{The other live consequence is the
register}: the interval is a \textbf{clause} and not a declaration, on the split the assumption block now
prints, so a result that reads it carries it --- \cref{lem:intker} does --- and this lemma, which does
not read it, does not.
\emph{A1b(i) is \textbf{not} a hypothesis here}: all
four clauses concern generic $h,h',\nu_g$, the codomain $\mathcal V$ is \cref{def:spaces}, and
$\vdagger$ does not occur in the statement or in any step of the proof. Then
% ARCHAEOLOGY.  A1b(i) was listed on this line through round 7 and is struck.
\begin{enumerate}
\item[(i)] $\mathcal D^{-1}_0$ is bounded, with
$\|\mathcal D^{-1}_0\|\le2\,(\,|\mathcal X|\,\|1/\pi\|_\infty)^{1/2}$, and injective, while
$\mathcal D$ is unbounded on $L^{2}(\pi)$;
\item[(ii)] for integrable $\nu_g$, writing $\mathcal N_g(v):=\int_{\{x\ge v\}\cap\mathcal X}\nu_g(x)\,dx$,
one has $\mathcal N_g\equiv0$ if and only if $\nu_g\equiv0$ almost everywhere on $\mathcal X$;
\item[(iii)] A2 holds for $\mathcal A_{\mathrm{cond}}$ if and only if the corresponding source condition
holds for $\mathcal A_{\mathrm{cond}}\mathcal D^{-1}_0$;
\item[(iv)] $\|h\|_{L^{2}(\pi)}\le\|\mathcal D^{-1}_0\|\,\|h'\|_{L^{2}(\pi)}$ whenever
$h=\mathcal D^{-1}_0h'$, so a rate in the derivative norm implies the same rate in the level norm, while
the converse fails.
\end{enumerate}
\end{lemma}

\begin{lemma}[Closed form of the adjoint]\label{lem:adjoint}
The adjoint of the linearisation is the difference functional of the debiasing weight, divided by the
reference density, once the weight is read through the Bierens family.
Suppose B1(i), (ii), (iv) hold together with C1(ii) and C2$'$(i). Then, with
$\mathsf k$ and $g_{u,t,l}$ as in \eqref{eq:kernel}:
\begin{enumerate}
\item[(i)] \emph{Kernel form.} For every $\ell\in L^{2}(\Pi)$,
\[
\big(\mathcal A^\ast_{\mathrm{cond}}\ell\big)(x)
\;=\;\sum_{l=1}^{\mathsf L-1}\varpi_{\mathsf w}(l)\int_{\IR\times\mathcal T}
\mathsf k(u,t,l;x)\,\ell(u,t,l)\,\Pi_u(du)\,\Pi_t(dt),
\]
and $\mathcal A^\ast_{\mathrm{cond}}\ell\in\mathcal V$. \emph{The integration is over all three factors
of $\Pi$, the third being the weighted sum over the contrast index}; this is \cref{def:spaces}'s pairing
read on the left-hand argument, and it is where the masses
$\varpi_{\mathsf w}(l)$ act that \eqref{eq:Wdef} displays. At $\mathsf L=2$ the third factor is the
one-point set $\{1\}$, $\varpi_{\mathsf w}(1)=1$, $\mathsf k(u,t,1;x)=\mathsf k(u,t;x)$, and the display
is the two-factor integral earlier rounds printed.
\item[(ii)] \emph{The induction map.} The operator
\begin{equation}\label{eq:Wdef}
\big(\mathcal W\ell\big)_l(u,z)\;:=\;\varpi_{\mathsf w}(l)\,\varpi_u(u)\int_{\mathcal T}\ell(u,t,l)\,
w\big(t'\Phi(z)\big)\,\Pi_t(dt),
\qquad (u,z)\in\IR\times\IR^{p},\quad l=1,\dots,\mathsf L-1,
\end{equation}
is a bounded linear map from $L^{2}(\Pi)$ into $L^{2}(P_{\Delta U,\Delta Z};\IR^{\mathsf L-1})$ with
$\|\mathcal W\|\le\big(\|\bar f_0\|_\infty\,\|\varpi_u\|_\infty\int_{\mathcal T}\bar w^{2}\,d\Pi_t\big)^{1/2}$,
its range is dense in $L^{2}(P_{\Delta U,\Delta Z};\IR^{\mathsf L-1})$, and
$\ker\mathcal W\subseteq\ker\mathcal A^\ast_{\mathrm{cond}}$.
Three readings of this clause are fixed here, because each changes with $\mathsf L$ and none is
visible at $\mathsf L=2$. \emph{First, the codomain norm} is
$\|\Xi\|:=(E\|\Xi(\Delta U,\Delta Z)\|^{2})^{1/2}$ with the Euclidean norm on $\IR^{\mathsf L-1}$ inside.
Second, the constant is the printed one at every $\mathsf L$, and R1 is what makes that true.
\eqref{eq:Wdef} acts \emph{diagonally} in $l$, slice $l$ of $\ell$ generates component $l$ of
$\mathcal W\ell$ and no other, so the passage from the slicewise estimate to the vector norm costs
the factor $\big(\max_{l}\varpi_{\mathsf w}(l)\big)^{1/2}\le1$, $\Pi_{\mathsf w}$ being a probability
measure by B1(iv.d); \emph{without the $\varpi_{\mathsf w}(l)$ that R1 puts in \eqref{eq:Wdef} the same
passage would cost $\big(\min_l\varpi_{\mathsf w}(l)\big)^{-1/2}\ge1$ instead}, the domain norm carrying
the masses and the codomain not, so the printed constant would be too small. \emph{Third, density holds
in the $\IR^{\mathsf L-1}$-valued space and not merely componentwise}: the diagonal action makes
$\mathcal R(\mathcal W)$ the algebraic sum of $\mathsf L-1$ subspaces, one dense in each component by
the scalar argument, and a finite sum of dense subspaces of the corresponding factors is dense in the
product. \emph{At $\mathsf L=2$ all three readings are the printed ones}: $\varpi_{\mathsf w}(1)=1$, the
codomain is scalar, and the constant does not move.
% ---------------------------------------------------------------------------
% eq:Wdef CARRIES varpi_w(l) IN COLUMN l.  AUTHOR DECISION, 2026-08-09 (repair
% R1), taken because eq:adjclosed AS PRINTED WAS FALSE AT sans-L >= 3 and true
% at sans-L = 2, which is why no earlier round could see it.  Deriving the
% adjoint from def:spaces' pairing gives
%     A*_cond Xi = nu_{Wtilde' diag(varpi_w)(W Xi)} / pi,
% while eq:adjclosed printed nu_{Wtilde (W ell)}/pi -- SHORT BY diag(varpi_w).
% A uniform Pi_w does not rescue it: the adjoint equation then reads -(L-1) nu.
% It matters rather than moving a constant because psi = DZ DU + L Wtilde
% carries the contrast UNWEIGHTED, so prp:orthog(i)'s condition is unweighted
% while def:weight(c)'s adjoint equation delivers the varpi_w-weighted
% combination; the "consequently" clause identifying the two is what broke.
%
% R1 puts the masses in W rather than printing diag(varpi_w) at every consumer
% (R2).  THE GROUND IS INTERNAL, NOT CONVENIENCE: W already absorbs Pi_u's
% DENSITY varpi_u and already INTEGRATES OUT Pi_t -- it prices two of Pi's
% three factors, and the w-direction was the only one it neither weighted nor
% integrated.  Under R1, eq:adjclosed is verbatim at every sans-L.
%
% THE KERNEL FORMS WERE ALREADY RIGHT and must NOT acquire the masses:
% eq:kernel and lem:adjoint(i) receive them through Pi in the integral itself.
% Only the CLOSED form, which routes through W, must carry them.
%
% AT sans-L = 2: varpi_w(1) = 1 and there is one column, so both displays are
% the pre-dispatch ones character for character.
%
% R1 WAS APPLIED WITHOUT A PROPAGATION CENSUS, WHICH IS HOW eq:Wderiv BELOW
% CAME TO BE FALSE AT sans-L >= 3 AND TRUE AT sans-L = 2 -- the exact failure
% species this phase exists to catch, introduced by the repair for another
% instance of it.  R2's consumers were enumerated when R2 was offered; R1's were
% not, and were not asked for.  Consumers now swept and repaired:
%   * eq:Wderiv (cor:elldiff) -- carries varpi_w(l) and the l argument;
%   * lem:adjoint(ii)'s codomain -- now R^{sans-L - 1}-valued.
% PHASE 4 CLOSED THE ||W|| ITEM, AND THE ANSWER IS THAT THE PRINTED CONSTANT
% SURVIVES.  The passage from the slicewise estimate to the R^{sans-L - 1}-valued
% codomain costs (max_l varpi_w(l))^{1/2} <= 1 -- Pi_w a PROBABILITY measure,
% B1(iv.d) -- so the number does not move and is conservative by that factor.
% Pre-R1 the same passage cost (min_l varpi_w(l))^{-1/2} >= 1 and the printed
% constant was TOO SMALL by exactly that, which is the defect this comment
% recorded.  R1 pays it.  The three readings -- codomain norm, constant, density
% at the VECTOR space -- are now PRINTED in clause (ii), not left to a comment.
%   THE SUP-NORM BOUNDS AFTER eq:Wderiv survive in sans-L for the same reason,
% varpi_w(l) <= 1 -- indeed the l-aggregation costs (sum_l varpi_w(l)^2)^{1/2}
% <= 1, so R1 PAYS FOR the sqrt(sans-L - 1) that dispatch 25 flagged as owed.
%
% *** THE sqrt(p) IS GONE, AND IT WENT UPSTREAM.  AUTHOR DECISION 2026-08-09. ***
% A first pass of this phase found a SECOND and sans-L-FREE shortfall, sqrt(p),
% and printed it here: A2'(i.a)/(i.b), A2b and eq:lambdamu aggregated the budgets
% over the p coordinates of theta by the MAXIMUM (dispatch 25), while every
% consumer of ||ell||_inf reads a EUCLIDEAN/FROBENIUS norm of the assembled
% p x (sans-L - 1) object -- rem:orthoghyp(3), rem:anchor, lem:solset(iii),
% eq:rprodbound.  Pre-dispatch-25 the budget was EUCLIDEAN in j (the retired
% mirror assumptions.tex:717-720 writes |ell^Pi_star(u,t)| for an R^p-valued
% object) and the printed constants were exact, so the max-over-j was a
% regression rather than a design choice.
% THE AUTHOR TOOK THE UPSTREAM REPAIR.  A2'(i.a), A2'(i.b), A2b and eq:lambdamu
% now aggregate EUCLIDEAN OVER j, MAX OVER l.  Consequently the two displays
% above are the PRE-sans-L ones character for character, in the Frobenius norm,
% at every sans-L and every p -- no sqrt(p) and no sqrt(sans-L - 1).  The
% entrywise restatement this file carried for one round is withdrawn: it was a
% way of keeping a true display while a false one was flagged, and it has no
% consumer now that the true display IS the printed one.
% ---------------------------------------------------------------------------
\item[(iii)] \emph{Closed form.} For every $\ell\in L^{2}(\Pi)$,
\begin{equation}\label{eq:adjclosed}
\mathcal A^\ast_{\mathrm{cond}}\ell\;=\;\frac{\nu_{\widetilde W'\,(\mathcal W\ell)}}{\pi} ,
\end{equation}
where $\widetilde W'(\mathcal W\ell)$ abbreviates the Euclidean inner product
$\sum_{l=1}^{\mathsf L-1}\widetilde W_l\,\big(\mathcal W\ell\big)_l(\Delta U,\Delta Z)$;
consequently $\mathcal A^\ast_{\mathrm{cond}}\ell=-\nu/\pi$ if and only if
$\nu_{\Delta Z+\widetilde W'(\mathcal W\ell)}\equiv0$ almost everywhere on $\mathcal X$.
\end{enumerate}
\end{lemma}

\begin{corollary}[to \cref{lem:adjoint}: the representer is differentiable, with explicit bounds]\label{cor:elldiff}
Regularity of a range-space representer transfers to its observable-space image, and the transfer
is what makes the sup-norm requirements of A2$'$ well posed.
Stated at a generic admissible representer, with $\ell^\Pi_\star$ as the specialisation. Suppose
B1(i), (ii), (iv.a) and (iv.d) hold and let $\ell^\Pi\in L^{2}(\Pi)$ be an \emph{admissible representer} in the
sense of A2b, with budgets $\Lambda^{(0)},\Lambda^{(1)}$ and with $\ell:=\mathcal W\ell^\Pi$ defined
pointwise from the version A2b fixes. Then $\ell$ is bounded, is continuously differentiable in $u$ with
\begin{equation}\label{eq:Wderiv}
\partial_u\ell_l(u,z)
=\varpi_{\mathsf w}(l)\Big\{\varpi_u'(u)\!\int_{\mathcal T}\!\ell^\Pi(u,t,l)\,w(t'\Phi(z))\,\Pi_t(dt)
+\varpi_u(u)\!\int_{\mathcal T}\!\partial_u\ell^\Pi(u,t,l)\,w(t'\Phi(z))\,\Pi_t(dt)\Big\},
\qquad l=1,\dots,\mathsf L-1,
\end{equation}
and satisfies
\[
\big\|\ell\big\|_\infty\le\|\varpi_u\|_\infty\,\Lambda^{(0)},
\qquad
\big\|\partial_u\ell\big\|_\infty\le\|\varpi_u'\|_\infty\,\Lambda^{(0)}+\|\varpi_u\|_\infty\,\Lambda^{(1)} ,
\]
where $\|\ell\|_\infty:=\sup_{(u,z)}\|\ell(u,z)\|$ and $\|\partial_u\ell\|_\infty$ likewise, the norm
being \emph{Frobenius} on the $p\times(\mathsf L-1)$ array, \emph{which is the norm every consumer of
these two symbols reads}: \cref{rem:orthoghyp}\textup{(3)}, \cref{rem:anchor},
\cref{lem:solset}\textup{(iii)} and \eqref{eq:rprodbound}.
Neither constant moves with $\mathsf L$ or with $p$, and the two directions are paid for by
different things. The $p$-direction is paid for by A2b, whose budgets are \textbf{Euclidean over the
$p$ coordinates of $\theta$}, so that column $l$ of $\ell$, the $\IR^{p}$-valued $\ell_l$ of
\eqref{eq:Wdef}, obeys
$\|\ell_l(u,z)\|\le\varpi_{\mathsf w}(l)\|\varpi_u\|_\infty\Lambda^{(0)}$ directly, by the
triangle inequality for the $\Pi_t$-integral of a vector-valued integrand. The $\mathsf L$-direction is
paid for by R1: assembling the columns gives
$\|\ell(u,z)\|^{2}\le\|\varpi_u\|_\infty^{2}(\Lambda^{(0)})^{2}\sum_{l}\varpi_{\mathsf w}(l)^{2}$, and
$\sum_{l}\varpi_{\mathsf w}(l)^{2}\le\big(\sum_l\varpi_{\mathsf w}(l)\big)^{2}=1$ because
\eqref{eq:Wdef} carries $\varpi_{\mathsf w}(l)$ in column $l$ and \emph{B1(iv.d), which is on the
hypothesis line above, makes $\Pi_{\mathsf w}$ a probability measure}, so the $\sqrt{\mathsf L-1}$ a
mass-free induction map would have cost is
paid for, and it is paid for exactly.
\emph{\textbf{B1(iv.d) is assumed here; it is no longer also asserted}}. Through 2026-08-23 the clause
stood at this sentence alone, on the \emph{conclusion} side, so a reader who acted on the hypothesis
line had not assumed what the constant needs, while a reader of the whole statement was told it holds
without ever being asked for it. It is cited \emph{by clause} and not through the whole of B1(iv),
because no step of the proof reads \textup{(iv.b)} or \textup{(iv.c)} and a block citation would charge
this corollary --- and everything downstream of it --- for two clauses it never opens. \emph{That is the
assumption block's \textbf{three registers} rule and is no longer this corollary's own defence}: the
factorisation $\Pi=\Pi_u\otimes\Pi_t\otimes\Pi_{\mathsf w}$ is B1(iv)'s \emph{declaration}, standing at
every statement here, so citing the two clauses this proof opens leaves nothing unnamed. \emph{This
corollary's line was the model for the six repaired on 2026-08-24} and is itself unchanged.
\emph{And because it is a hypothesis \textbf{here} it is a hypothesis of every result that runs a bound
through this corollary}, which is why \cref{lem:secondord}, \cref{lem:momzero},
\cref{lem:sandwich}\textup{(i$'$)} and \cref{rem:dirmember} carry \textup{(iv.d)} on their own
lines from the same date. \emph{No conclusion anywhere moves}: \textup{(iv.d)} costs nothing about $P$,
and at $\mathsf L=2$ it is forced.
\emph{Specialisation.} Under A2$'$(i.a)--(i.b) the selected solution $\ell^\Pi_\star$ is an admissible
representer with $\Lambda^{(0)}=\Lambda^{(0)}_\star$, $\Lambda^{(1)}=\Lambda^{(1)}_\star$, so the
corollary applies at $\ell=\ell_\star:=\mathcal W\ell^\Pi_\star$ and yields
$\|\ell_\star\|_\infty\le\|\varpi_u\|_\infty\Lambda^{(0)}_\star$ and
$\|\partial_u\ell_\star\|_\infty\le\|\varpi_u'\|_\infty\Lambda^{(0)}_\star+\|\varpi_u\|_\infty\Lambda^{(1)}_\star$.
\end{corollary}

% =========================================================================
%  WHY cor:elldiff IS NOW STATED AT A GENERIC ADMISSIBLE REPRESENTER --
%  ROUND-8 REPAIR OF A DUAL-STATUS TAG, AND IT GATED STEP 4.
%
%  THROUGH ROUND 7 the hypothesis line read "A2'(i.a)-(i.b), and let
%  ell^Pi_star be the version fixed in A2'(i)".  FOUR SITES APPLIED THE
%  COROLLARY GENERICALLY: def:weight(a) and eq:PsiL in this file, lem:solset(ii)
%  and prp:orthog(i) in proofs.tex.  Each restated the bridge in its own words,
%  so nothing was unsupported -- but the corollary's HYPOTHESIS LINE is what a
%  Step-4 author reads, and L-21's target is A2'(i.a)-(i.b) ON ell^Pi_star.  With
%  the same text serving as a predicate over generic representers, L-21 did not
%  know which object it was discharging.
%
%  THE FIX IS A STRICTLY STRONGER RESULT ON STRICTLY WEAKER HYPOTHESES, which is
%  the same move prp:orthog(i) received in round 7.  The proof does not change by
%  one line: it never used that ell^Pi_star SOLVES the adjoint equation, only that
%  it satisfies the two integrability budgets and the differentiability -- which
%  is exactly A2b.
%
%  ORDERING.  The hypothesis is stated via the TAG A2b (assumptions.tex) and not
%  via def:weight(a), because def:weight appears LATER in this file and a
%  corollary may not cite it.  def:weight(a) now DEFINES the phrase "admissible
%  representer" to mean A2b, which is a naming convention, not a second statement.
% =========================================================================

\begin{corollary}[to \cref{lem:adjoint}: how far A2 can exceed what the moment needs]\label{cor:rangegap}
The source condition and the orthogonality equation the estimator targets can differ only inside the
non-closedness of the range of the adjoint.
Suppose the hypotheses of \cref{lem:adjoint} hold together with A1b(iv), and define the
\emph{tilted difference map}
$\Upsilon:L^{2}(P_{\Delta U,\Delta Z};\IR^{\mathsf L-1})\to\mathcal V$ by
$\Upsilon\Xi:=\nu_{\widetilde W'\Xi}/\pi$, where $\widetilde W'\Xi=\sum_{l}\widetilde W_l\Xi_l$ is the
Euclidean inner product of \eqref{eq:adjclosed}, \emph{$\Upsilon$'s domain is
\cref{lem:adjoint}\textup{(ii)}'s codomain, which is what makes the factorisation below a composition of
maps that meet}. Then $\Upsilon$ is a bounded linear operator with
\[
\big\|\Upsilon\big\|\;\le\;2\,c_{\widetilde W}\,\bar c_\pi^{1/2},
\qquad
c_{\widetilde W}\;:=\;\min\big\{\sqrt2,\ \sqrt{\mathsf L-1}\,\big\},
\]
$\mathcal A^\ast_{\mathrm{cond}}=\Upsilon\circ\mathcal W$, and
\begin{equation}\label{eq:rangegap}
\mathcal R\big(\mathcal A^\ast_{\mathrm{cond}}\big)\;\subseteq\;\mathcal R(\Upsilon)\;\subseteq\;
\overline{\mathcal R\big(\mathcal A^\ast_{\mathrm{cond}}\big)} .
\end{equation}
$c_{\widetilde W}$ is a surely-valid bound on $\|\widetilde W\|$, it is
\textbf{not} a consequence of \textup{(P1)}, and it is the constant every contraction of a
$\widetilde W$ against a vector in this file carries. \emph{It is italic and symbol-subscripted,
joining $c_\Phi$, $c_\Pi$, $c_C$, $c_\mu$ and $c_{\mu''}$}: in the $c$-family every symbol-subscripted
member is italic and every sans-serif member, $\mathsf c_0$ through $\mathsf c_{17}$, is
integer-subscripted, so the font separates the two without a further rule.
By \eqref{eq:contrast}, $\widetilde W=e-\varrho$
is the difference of two points of the simplex $\{q\in\IR^{\mathsf L-1}:q_l\ge0,\ \mathbf 1'q\le1\}$, 
$e$ a vertex and $\varrho$ an interior or boundary point, so $\|\widetilde W\|$ is at most that
simplex's diameter, which is $\sqrt2$ at every $\mathsf L\ge3$ and $1$ at $\mathsf L=2$; and separately
$\|\widetilde W\|^{2}=\sum_l\widetilde W_l^{2}\le\mathsf L-1$, each coordinate lying in $[-1,1]$. The
minimum of the two is $c_{\widetilde W}$, and it equals $1$ at $\mathsf L=2$ and $\sqrt2$ at
every $\mathsf L\ge3$, \emph{so the constant saturates at $\mathsf L=3$ and does not grow with the
number of cells}. \emph{\textup{(P1)} does not deliver it}: \textup{(P1)} is the \emph{per-cell}
statement $|\widetilde W_l|=|\un\{W=l\}-\varrho_l(\Delta Z)|\le1$, one scalar bound at each of the
$\mathsf L-1$ contrast indices, and, at $\mathsf L=2$, the single bound $|W-\varrho(\Delta Z)|\le1$
earlier rounds printed; the Euclidean bound on the vector is a different statement and is not $1$.
\emph{No sharpness is claimed for the bound on $\|\Upsilon\|$}: it is a pointwise Cauchy--Schwarz bound,
and $\|\widetilde W\|$ approaches $\sqrt2$ only where one cell probability approaches $1$, an event whose
probability the bound does not use. At $\mathsf L=2$, $c_{\widetilde W}=1$ and the display
reads $\|\Upsilon\|\le2\bar c_\pi^{1/2}$, the constant earlier rounds printed.
\end{corollary}

% THE OPERATOR IS Upsilon, NOT B -- ROUND-7 RENAME, AND IT IS NOT COSMETIC.
% Through round 6 this corollary's operator and the second-order form operator of
% eq:Bpi both carried the bare letter B, on DIFFERENT SPACES
% (L^2(P_{Delta U, Delta Z}) here, L^2(pi) there), across Parts E and F.  The
% glossary's disambiguation rule -- "B never carries a subscript" -- could not
% separate them, and proofs.tex:lem:balance Step 2 had to write "with B the
% operator of cor:rangegap", which is a file conceding that context does not
% disambiguate.  The eq:Bpi operator keeps the letter, because B_sym is under a
% binding contract with the memo and "sym" presupposes a B.  This one moves.
% Upsilon is a Greek capital, unused anywhere in the project, and takes the
% Greek capital Xi as its argument -- the same type-consistency the glossary
% invoked for F_N and varpi_t.

\begin{remark}[Consequences of \cref{lem:adjoint}, and one independent confirmation]\label{rem:adjointreconcile}
Four things follow, three of which later results consume.

(i) The representer is evaluable at the data, and that is now a theorem. A solution of the adjoint
equation is an element of $L^{2}(\Pi)$, a function of $(u,t)$ with $t$ the Bierens index, while
A2$'$(i), whose closing sentence fixes $\ell_\star:=\mathcal W\ell^\Pi_\star$ pointwise, and
\eqref{eq:JSig} evaluate $\ell_\star(\Delta U,\Delta Z)$ with $z$ a value of $\Delta Z$. Part (ii)
supplies the map between the two and part (iii) shows it commutes with the adjoint equation, so the object
A2$'$(i) and \eqref{eq:JSig} speak about is constructed rather than posited. \emph{And \eqref{eq:Wdef} is
exactly why the evaluation is well posed}: $\ell_\star=\mathcal W\ell^\Pi_\star$ is a Borel map of
$(u,z)$ and of nothing else, in particular not of $W$, which is the measurability
\cref{lem:momzero}\textup{(iv)} consumes. \emph{The construction consumes exactly one thing beyond the
identification block}: absolute continuity of $\Pi_u$ in B1(iv.a), a restriction on a reference
measure the analyst chooses.
% ARCHAEOLOGY.  Through dispatch 13 the first citation on this item read
% "A2'(i)--(ii)"; A2'(ii) is retired and the clause that evaluates ell_star at the
% data was always (i).  Through dispatch 12 the second citation read "A3"; that tag
% is retired at ass:rank and the display it presupposed is eq:JSig.

\emph{(ii) The correction measure is not free \textup{on the population side}}, the map \eqref{eq:Wdef}
integrating against $\Pi_t$; \cref{rem:Tcond} carries the disposition of B1(v)'s two halves.

(iii) What the construction does not give, and what it gives anyway. $\mathcal W$ has dense range
but need not be surjective, so a selection criterion posed over all of
$L^{2}(P_{\Delta U,\Delta Z};\IR^{\mathsf L-1})$ may
be unreachable; and $\mathcal W$ need not be injective, though
$\ker\mathcal W\subseteq\ker\mathcal A^\ast_{\mathrm{cond}}$ means no direction relevant to the selection
is lost. It does not follow that the gap between A2 and the equation the estimator targets is
uncontrolled. By \cref{cor:rangegap}, $\mathcal R(\Upsilon)$, the set of $\nu/\pi$ for which the
primitive orthogonality equation is solvable in $L^{2}(P_{\Delta U,\Delta Z};\IR^{\mathsf L-1})$, is squeezed between
$\mathcal R(\mathcal A^\ast_{\mathrm{cond}})$ and its closure, \emph{with no characterisation of
$\mathcal R(\mathcal W)$ required}. So A2 can exceed what the moment needs only within
$\overline{\mathcal R(\mathcal A^\ast_{\mathrm{cond}})}\setminus\mathcal R(\mathcal A^\ast_{\mathrm{cond}})$.
That set is nonempty under H1: $\mathcal A^\ast_{\mathrm{cond}}$ is compact, being the adjoint of a
compact operator, and a compact operator has closed range if and only if that range is
finite-dimensional, which H1's infinite-rank clause excludes. \emph{Injectivity is neither available nor
needed for this}, $\mathcal A^\ast_{\mathrm{cond}}$ may well have a kernel, since
$\ker\mathcal W\subseteq\ker\mathcal A^\ast_{\mathrm{cond}}$ two sentences above. The
gap is therefore exactly where the ill-posedness lives.
% ARCHAEOLOGY.  An earlier version of this item gave injectivity of A*_cond as the
% reason the gap is nonempty; it was both unavailable and superfluous.

(iv$'$) The derivative of the representer. The sup-norm requirements on $\ell_\star$ remain
assumptions, but they are stated on $\ell^\Pi_\star$, in A2$'$(i.a)--(i.b), and
\cref{cor:elldiff} derives from them both \eqref{eq:Wderiv} and the two bounds. \emph{The ordering
matters}: the $C^{1}$ clause of B1(iv.a) alone controls only the first term of
\eqref{eq:Wderiv}, so $\|\partial_u\ell_\star\|_\infty<\infty$ cannot be assumed directly without
assuming the existence of a derivative nothing else in the block delivers.
\eqref{eq:Wderiv} is what makes the choice of $\Pi_u$ a design lever on $q$ and on
$\sigma_q$; it is a \emph{conditional} identity whose hypothesis is A2$'$(i.b), and A2$'$(i.b) is a
primitive awaiting L-21, not a theorem.

(iv) The $\nu/\pi$ referent is confirmed independently, and in the same form. \emph{Notation:}
\citet{feve2014} write the endogenous regressors as $Z_1,Z_2$ with generic value $z$, the instrument as
$W$ with generic value $w$, and a generic element of the range space as $\psi$; the translation to this
paper is $Z_t\mapsto X_t$, $z\mapsto x$, and their $\psi$ is our $\ell$, the letter $\psi$ being reserved
here for the moment function. For the linear conditional-mean panel operator
$K\varphi=E[\varphi(Z_2)-\varphi(Z_1)\mid W]$, their Proposition~2.1 gives, \emph{verbatim in their
notation},
\[
\big(K^{\ast}\psi\big)(z)\;=\;\int\psi(w)\,\frac{f_{Z_2,W}(z,w)-f_{Z_1,W}(z,w)}{\pi(z)}\,dw ,
\]
followed immediately by the statement that $K^{\ast}\psi$ satisfies the normalising constraint
$\int\big[(K^{\ast}\psi)(z)\big]\pi(z)\,dz=0$. Translated, the numerator is $\nu_{\ell(W)}$ and the
display is $\nu_{\ell(W)}/\pi$. \emph{One difference of setting must be flagged rather than absorbed:}
their $dw$ integral ranges over a general instrument space, whereas $W$ takes the $\mathsf L$ values
$\{0,\dots,\mathsf L-1\}$ here, so the
counterpart in this paper is the finite sum
$\nu_{\ell(W)}(x)=\sum_{j=0}^{\mathsf L-1}\ell(j)\{f_{X_2,W}(x,j)-f_{X_1,W}(x,j)\}$; the difference-of-densities
structure is what carries over, not the integral. At $\mathsf L=2$ this is the two-point sum over
$j\in\{0,1\}$ earlier rounds printed.

\emph{Both features of \eqref{eq:adjclosed} therefore appear in print in the linear case:} the division
by the reference density, with the same letter $\pi$, which the draft takes from them, and the zero
$\pi$-mean, which is membership of $\mathcal V$ and which \cref{prp:compact} Step~2 proves here from
$\int\nu_g\,dx=0$ at $h\equiv1$. This matters because the referent was contested: an earlier version of
this project imposed the source condition on $\nu$ rather than on $\nu/\pi$, and the correction was
forced internally by a Picard membership test. The $\pi^{-1}$ is not an artefact of the
$L^{2}(\pi)$ convention adopted here; it is the shape the adjoint takes in a first-differenced panel.
What transfers is that shape and nothing more: their $K$ is a fixed linear conditional-expectation
operator, while $\mathcal A_{\mathrm{cond}}$ is a Fr\'echet derivative, and the induction map
$\mathcal W$ in \eqref{eq:adjclosed} has no counterpart there because their restriction is already
unconditional in the residual.
\end{remark}

% REMARK rem:nophifree MOVED to working_notes.tex on 2026-08-10 -- zero printed
% cross-references anywhere in the compiled corpus, so nothing points at it and
% no stub is needed.  Verbatim in that file, which is input by no document.

% REMARK rem:ffcompare MOVED to working_notes.tex on 2026-08-10 -- zero printed
% cross-references anywhere in the compiled corpus, so nothing points at it and
% no stub is needed.  Verbatim in that file, which is input by no document.

\begin{remark}[The two anchors of interest, and why the anchor is priced at all]\label{rem:anchor}
Two anchors are canonical and both are admissible with no further assumption. The \emph{zero anchor}
$b\equiv0$ is trivially admissible. The \emph{N1 anchor}
$b(z)=-E\big[\int_0^{\Delta U}\ell(v,z)\,dv\;\big|\;\Delta Z=z\big]$ produces the unique weight generated
by $\ell$ that satisfies N1, and it is admissible whenever A1b(iv) holds, since
$\|b(z)\|\le\|\ell\|_\infty E[|\Delta U|\mid\Delta Z=z]$ and conditional Jensen give
$E\|b(\Delta Z)\|^{2}\le\|\ell\|_\infty^{2}E[\Delta U^{2}]<\infty$. So N1 can always be arranged,
at no cost to $\partial_uL$ and, by \textup{(P3)} and A1, none to $E[\psi]$; it does not hold
automatically, and \cref{prp:orthog}(iii) is the statement of exactly what its failure costs.

\emph{``Arranged'' is $p(\mathsf L-1)$ conditional expectations, and the arithmetic of that count is
where a construction goes wrong.} The N1 anchor is a $p\times(\mathsf L-1)$ matrix computed
\emph{slice by slice}: entry $(j,l)$ is
$-E\big[\int_0^{\Delta U}\ell_{jl}(v,\Delta Z)\,dv\mid\Delta Z\big]$, and the $\mathsf L-1$ columns are
generated by $\mathsf L-1$ separately choosable slices of $\ell^\Pi$ per
\cref{def:weight}. A construction that centres only a \emph{contraction} of the columns against a fixed
$\mathsf v\in\IR^{\mathsf L-1}$ therefore discharges $p$ of the conditions and leaves $p(\mathsf L-2)$
uncentred, which is zero at $\mathsf L=2$ and strictly positive at every $\mathsf L\ge3$;
\cref{cor:centarity} exhibits such a weight \emph{inside} the admissible class and prices what it leaves.
\emph{Nor is targeting this anchor optional once \cref{thm:rootn} is read whole}: \textup{(T-2)} assumes
N1 at $L_\star$ and \textup{(R-a)} forces $\hat b^{(-k)}\to b_\star$, so the two together select this
anchor and no other, and a construction free to choose $b$ is free only before \textup{(T-2)} is
imposed.

The anchor is priced in A2$'$(iii) because it is a \emph{free design object}: the weight is determined by
the pair $(\ell,b)$, the adjoint equation constrains only $\ell$, and an analyst may choose a $b$ with
$E\|b(\Delta Z)\|^{2}=\infty$. \emph{What A2$'$(iii) buys, and nothing more, is that
$E\|L_\star(\Delta U-\tau_1\Delta h,\Delta Z)\|<\infty$ for every $\tau_1$, that
$E[\|L_\star\|\mid\Delta Z]<\infty$ almost surely, and that
$E\big\|L_\star\widetilde W\big\|^{2}<\infty$}, \emph{the third quantity is the squared Euclidean
norm of the $\IR^{p}$-valued product and \textbf{not} $E[\widetilde W^{2}\|L_\star\|^{2}]$, which is not
a well-formed expression once $\widetilde W$ is a vector; at $\mathsf L=2$ the two are the same number}
, the first because $\|L_\star(u,z)\|\le\|b_\star(z)\|+\|\ell_\star\|_\infty|u|$ by \eqref{eq:weightdef}
with all norms Frobenius,
the second by conditional Jensen, the third by
$\|L_\star\widetilde W\|\le c_{\widetilde W}\|L_\star\|$ and
$(\varsigma+\varsigma')^{2}\le2\varsigma^{2}+2\varsigma'^{2}$ applied to
that same bound, together with A1b(iv). \emph{No shift
restriction is needed on the $L_\star$ term for this reason}, which is why the shift proviso of
\cref{def:dirclass} is imposed on the direction and not here.
\end{remark}

\begin{definition}[Admissible perturbation directions]\label{def:dirclass}
Fix $\bar\tau>0$ and $h\in\mathcal V$, and for a Borel
$\Xi:\IR\times\IR^{p}\to\IR^{p\times(\mathsf L-1)}$ write
\begin{equation}\label{eq:shiftnorm}
\big\|\Xi\big\|_{0,h}\;:=\;\sup_{|\tau_1|\le\bar\tau}
\Big(E\big\|\Xi(\Delta U-\tau_1\Delta h,\Delta Z)\big\|^{2}\Big)^{1/2}
\;\in\;[0,\infty] ,
\end{equation}
the norm inside the expectation being Frobenius, as the type declaration of \eqref{eq:contrast} fixes
for every object carrying one column per contrast. A pair $(D,\mathsf r)$ is an \emph{admissible
direction pair at $(h,\bar\tau)$} if
\begin{enumerate}
\item[(a)] $\mathsf r:\IR^{p}\to\IR^{\mathsf L-1}$ is Borel and bounded, and $\varrho+\tau_3\mathsf r$
takes values in the simplex $\{q\in\IR^{\mathsf L-1}:q_l\ge0\ \text{for every }l,\ \mathbf 1'q\le1\}$
for all small $\tau_3$; \emph{the range condition is the \textbf{simplex} and not the box
$[0,1]^{\mathsf L-1}$}, which would permit $\mathbf 1'(\varrho+\tau_3\mathsf r)>1$ and hence a negative
base cell, so that the perturbed vector would not be a propensity at all. It is
\cref{prp:orthog}\textup{(iii)}'s condition verbatim;
\item[(b)] $D:\IR\times\IR^{p}\to\IR^{p\times(\mathsf L-1)}$ is Borel with $\|D\|_{0,h}<\infty$.
\end{enumerate}
It is \emph{differentiable} if in addition $D$ is differentiable in $u$ at every point, $\partial_uD$ is
continuous in $u$, and
\begin{equation}\label{eq:shiftnormone}
\big\|D\big\|_{1,h}\;:=\;\Big(E\Big[\ \sup_{|s|\le\bar\tau}\big\|\partial_uD(\Delta U-s\Delta h,\Delta Z)\big\|^{2}\Big]\Big)^{1/2}\;<\;\infty .
\end{equation}
Taking $\tau_1=0$ in \eqref{eq:shiftnorm} gives $\|D\|_{L^{2}(P)}\le\|D\|_{0,h}$, and taking $s=0$ inside
the supremum in \eqref{eq:shiftnormone} gives $\|\partial_uD\|_{L^{2}(P)}\le\|D\|_{1,h}$; \emph{so the
unshifted $L^{2}(P)$ bounds are implied and are not assumed separately}. Continuity of $\partial_uD$ in
$u$ makes the supremum in \eqref{eq:shiftnormone} equal to the supremum over rational $s$, hence
measurable.

The two value types are forced by what the directions multiply. $D$ is a direction in the debiasing
weight $L$, which is $\IR^{p\times(\mathsf L-1)}$-valued by the type declaration of \eqref{eq:contrast},
and it enters $\psi$ only as $D\widetilde W\in\IR^{p}$; $\mathsf r$ is a direction in the propensity
vector $\varrho\in\IR^{\mathsf L-1}$ and enters only as $\varrho+\tau_3\mathsf r$. \emph{At
$\mathsf L=2$ both are the printed objects}: a
$p\times1$ matrix is a vector in $\IR^{p}$, the Frobenius norm is the Euclidean norm, $\mathsf r$ is
scalar, and the simplex of clause \textup{(a)} is the interval $[0,1]$.
Two pairs a reader will want are admissible. $(0,0)$ is admissible at every $(h,\bar\tau)$ and is
differentiable there, which is what \cref{rem:sharpquadread}\textup{(2)} reads
\cref{lem:secondord} at; and $\mathsf r=\hat\varrho-\varrho$ satisfies clause \textup{(a)}'s range
condition at every $\tau_3\in[0,1]$ by convexity, both $\varrho$ and $\hat\varrho$ taking values in
\eqref{eq:trunset} under A4(i) and \cref{def:hatT}.

Both glyphs suppress a dependence, and the suppression is safe for one reason only, which is worth
printing. $\|\cdot\|_{0,h}$ and $\|\cdot\|_{1,h}$ depend on the \emph{pair} $(h,\bar\tau)$, both are
suprema over $|\tau_1|\le\bar\tau$ or $|s|\le\bar\tau$, while the subscript carries only $h$. The
$\bar\tau$ is suppressed because every statement that uses either seminorm fixes $\bar\tau$ at its head,
exactly as \cref{def:dirclass} does; \emph{and} because the dependence is \textbf{monotone}: enlarging
$\bar\tau$ enlarges the index set of both suprema, so
\[
\bar\tau\le\bar\tau'\quad\Longrightarrow\quad
\|\Xi\|_{0,h}\big|_{\bar\tau}\le\|\Xi\|_{0,h}\big|_{\bar\tau'},
\qquad
\|D\|_{1,h}\big|_{\bar\tau}\le\|D\|_{1,h}\big|_{\bar\tau'} .
\]
\emph{The monotonicity is what forbids one false reading}, and it is recorded here rather than left to
be inferred: in any display that quantifies over $\bar\tau$ with $\bar\tau$ appearing on one side only,
the seminorm side is \emph{also} increasing in $\bar\tau$, so the bound cannot be optimised by sending
$\bar\tau\downarrow0$. See \cref{rem:dirmember}, where the quantifier ``for every $\bar\tau>0$'' stands
in front of exactly such a display.
\end{definition}

% MERGED, shortening pass 2026-08-01: rem:supinside and rem:shiftclass circled the
% same object -- the two seminorms of def:dirclass -- and are now one remark
% carrying both labels, so that every existing \cref to either still resolves.
\begin{remark}[The two seminorms of \cref{def:dirclass}: why the suprema sit on opposite sides, and why the shift-uniform form is not implied]\label{rem:supinside}
The two clauses of \cref{def:dirclass} are not stated in the same form, and the asymmetry is
forced.

\emph{\eqref{eq:shiftnorm} may take the supremum outside}, because the derivatives it licenses are
\emph{exact} rather than limits: $\psi$ is affine in each of $\tau_2,\tau_3$ \emph{separately}, with one
bilinear term $-\tau_2\tau_3\,D(\Delta U-\tau_1\Delta h,\Delta Z)\,\mathsf r(\Delta Z)$, \emph{the
direction to the left of the propensity direction, which is the only order that multiplies}, $D$ being
$p\times(\mathsf L-1)$ and $\mathsf r$ being $(\mathsf L-1)\times1$ by \cref{def:dirclass}, and the
product lying in $\IR^{p}$ as $\psi$ does. \emph{It is
bi-affine, not jointly affine, and the distinction is not pedantic}: read literally, joint affineness
would give $\partial^{2}_{\tau_2\tau_3}\Psi\equiv0$, which is exactly what \cref{lem:secondord}(iii)
denies. So $\partial_{\tau_2}\Psi$, $\partial_{\tau_3}\Psi$ and $\partial^{2}_{\tau_2\tau_3}\Psi$ are
obtained by splitting a finite sum of finite expectations, the bilinear term included, since it too is
one finite expectation with a coefficient, with no interchange of derivative and integral anywhere.
All that is needed is that $\Psi$ be finite at each $\tau$ separately, and
$\sup_{|\tau_1|\le\bar\tau}E[\,\cdot\,]<\infty$ delivers exactly that.

\emph{\eqref{eq:shiftnormone} may not.} The mixed partial $\partial^{2}_{\tau_1\tau_2}\Psi(0)$
differentiates in $\tau_1$, which \emph{is} a limit, and dominated convergence needs an integrable
envelope valid \emph{simultaneously} for all $\tau_1$ in a neighbourhood, that is,
$E[\sup_{|s|\le\bar\tau}\|\partial_uD(\Delta U-s\Delta h,\Delta Z)\|^{2}]<\infty$, with the supremum
\emph{inside}. A sup-outside condition gives $L^{2}(P)$-boundedness of the family
$\{\partial_uD(\cdot-s\Delta h,\cdot)\}_{|s|\le\bar\tau}$ without uniform integrability of its supremum,
and $\sup_s\|\cdot\|_{L^{2}}<\infty$ does not imply $\|\sup_s\cdot\|_{L^{2}}<\infty$. This is
precisely the quantifier that Leibniz's rule consumes, and it is the one A2$'$(i.b) is stated in. Since
$\sup_sE[\,\cdot\,]\le E[\sup_s\,\cdot\,]$, \eqref{eq:shiftnormone} implies the sup-outside version, so
nothing is lost by stating the stronger form where it is needed and the weaker form where it is not.

Why the shift-uniform seminorm at all, and what the alternative would have
cost.\label{rem:shiftclass}
\Cref{def:dirclass}(b) is what makes $\Psi(\tau)$ a real number off $\tau_1=0$, and it is not
implied by anything stated at the true residual. The two conditions are genuinely different, and a
one-line example separates them. Let $\Delta Z$ be degenerate, $\Delta U\sim\mathrm{Unif}[0,1]$, let
$h\in\mathcal V$ be such that $P(\Delta h=-1)=\tfrac12$ and $\Delta h=0$ elsewhere,
and take $D(u,z)=\un\{u>1\}(u-1)^{-1}\,e_{11}$, with $e_{11}$ the $p\times(\mathsf L-1)$ matrix whose
$(1,1)$ entry is $1$ and whose other entries are $0$, one entry carries the example and the
Frobenius norm of $D$ is the modulus of that entry, so nothing below depends on $p$ or on $\mathsf L$;
at $\mathsf L=2$ and $p=1$, $e_{11}=1$ and $D$ is the scalar function earlier rounds printed. Then
$D(\Delta U,\Delta Z)=0$ almost surely, so
$E[D\mid\Delta Z]=0$ and $\|D\|_{L^{2}(P)}=0$; but on $\{\Delta h=-1\}$ the argument at $\tau_1=1$ is
$\Delta U+1$ and
$E|D(\Delta U+1,\Delta Z)|=\tfrac12\int_0^1u^{-1}du=\infty$, so $\|D\|_{0,h}=\infty$ for every
$\bar\tau\ge1$. \emph{With such a $D$ the moment $\Psi(\tau)$ is not merely large at $\tau_1=1$: it is
undefined}, its positive and negative parts both being infinite whenever $\widetilde W$ is
non-degenerate. A statement of the form ``$\partial^{2}_{\tau_2\tau_2}\Psi\equiv0$ identically in
$\tau$'' would then be ill-posed rather than false.

The alternative repair was rejected and it is worth saying why. One could instead require a
uniformly bounded density ratio,
$\sup_{|\tau_1|\le\bar\tau}\big\|\,dP_{(\Delta U-\tau_1\Delta h,\Delta Z)}/dP_{(\Delta U,\Delta Z)}\big\|_\infty<\infty$,
which implies $\|\Xi\|_{0,h}\le C\|\Xi\|_{L^{2}(P)}$ for \emph{every} $\Xi$ at once. That is strictly
stronger and it fails in the most ordinary designs: if $\Delta U$ has compact support and
$\Delta h\not\equiv0$, the shifted law charges points the true law does not. And the failure is
sharper than unboundedness of a ratio: it is failure of absolute continuity, so there is no
Radon--Nikodym derivative to bound. In the display above the law of
$(\Delta U-\tau_1\Delta h,\Delta Z)$ at $\tau_1=1$ assigns mass $P(A)=\tfrac12$ to $\{u>1\}$, which is a
$P_{(\Delta U,\Delta Z)}$-null set; hence
$P_{(\Delta U-\Delta h,\Delta Z)}\not\ll P_{(\Delta U,\Delta Z)}$ and the quotient in the display is not
merely large, it \emph{does not exist}. The rejected condition is therefore not a stronger assumption one
might be willing to pay for, in this design it cannot be stated at all. \Cref{def:dirclass}(b) is a condition on the
\emph{direction class}, where the other direction conditions already live, and it restricts the law of
the data not at all.
\end{remark}

\begin{lemma}[L-16: the solution set and the selection]\label{lem:solset}
Every admissible solution of the adjoint equation yields a valid moment, so the selection is a free
design choice; and the functional it should be chosen to control is the multiplier of the second-order
form, not the sup-norm of the representer.
Suppose A1, A1b(i), A1b(iv), A2, A2$'$(i), A2$'$(iii), B1(i), (ii), (iv) and C1(ii) hold, \emph{A$2'$
is cited by clause and not as a block: A$2'$(ii) is the variance clause, is consumed by no part of this
lemma, and is a Step-8 condition, so listing A$2'$ whole would over-state the hypotheses}, and let
% ARCHAEOLOGY.  The by-clause citations here follow the round-7 narrowings
% C1 -> C1(ii) and B1(iv) -> B1(iv.a).
\emph{admissible representer} (A2b), \emph{admissible weight} and \emph{admissible selection} be as in
\cref{def:weight}. Then
\begin{enumerate}
\item[(i)] the set of $\ell^\Pi$ solving $\mathcal A^\ast_{\mathrm{cond}}\ell^\Pi=-\nu/\pi$ is the affine
set $\ell^\Pi_\dagger+\ker\mathcal A^\ast_{\mathrm{cond}}$, and its image under $\mathcal W$ is the
affine set $\ell_\dagger+\mathcal W(\ker\mathcal A^\ast_{\mathrm{cond}})$ with
$\ell_\dagger:=\mathcal W\ell^\Pi_\dagger$;
\item[(ii)] \emph{every admissible} selection, paired with \emph{any} admissible anchor, gives a moment
satisfying \cref{prp:orthog}(i)--(iii); in particular validity does not select. \emph{The restriction to
admissible members is not decoration:} $\mathcal W(\ell^\Pi_\dagger+\ker\mathcal A^\ast_{\mathrm{cond}})$
contains members with no sup-norm bound, for which \cref{prp:orthog}(i)'s differentiation under the
expectation is unlicensed;
\item[(iii)] fix $\vartheta\in\IR^{p}$ with $\|\vartheta\|=1$ and write
$\mathcal S\ell:=\sigma_{\partial_u\ell\,\widetilde W}$ for the map into the sum functional of
\eqref{eq:sigmadef}, so that $\sigma_q=\mathcal S\ell_\star$, \emph{the contrast standing to the
right of the $p\times(\mathsf L-1)$ derivative, which is the only order that multiplies and is the order
$\psi$ itself carries at \eqref{eq:psidef}}. The second-order form of
\cref{lem:secondord}(i), contracted in the direction $\vartheta$, depends on the selection through the
\emph{pair} $\big(\vartheta'\sigma_q,\ \vartheta'\mathcal K\big)$; \emph{read on $L^{2}(\pi)$, which is
\cref{def:spaces}'s one declared exception}, it is $\langle\mathcal Bh,h\rangle_{L^{2}(\pi)}$ with
\begin{equation}\label{eq:Bpi}
\mathcal B\;=\;\big(\vartheta'\sigma_q/\pi\big)\cdot{}\;-\;2\,\vartheta'\tilde{\mathcal K},
\qquad \tilde{\mathcal K}h:=(\mathcal Kh)/\pi ,
\end{equation}
of which $\vartheta'\sigma_q/\pi$, \emph{not} $\vartheta'\sigma_q$, is the multiplication part;
$\|\vartheta'\sigma_q/\pi\|_\infty\le2\,c_{\widetilde W}\,\bar c_\pi\,\|\partial_u\ell_\star\|_\infty<\infty$, so
$\mathcal B$ is bounded; and $\int_{\mathcal X}\sigma_q\,dx=0$ for every admissible selection.
\emph{The constant carries $c_{\widetilde W}=\min\{\sqrt2,\sqrt{\mathsf L-1}\}$ of
\cref{cor:rangegap} because the bound contracts the $p\times(\mathsf L-1)$ matrix
$\partial_u\ell_\star$ against the vector $\widetilde W$}, and $\|\partial_u\ell_\star\|_\infty$ is
\cref{cor:elldiff}'s Frobenius sup-norm; at $\mathsf L=2$, $c_{\widetilde W}=1$ and the
display is the one earlier rounds printed.
\end{enumerate}
\end{lemma}

\begin{lemma}[L-6: second-order expansion in the three nuisances]\label{lem:secondord}
The moment is affine in two of its three nuisances and quadratic in the third, so of the six second-order
terms three vanish identically and three are products of two estimation errors.
Suppose A1b(i), A1b(iv), B1(i), (ii), (iv.a), (iv.d), C1(ii), A2, A2$'$(i) and A2$'$(iii) hold, \emph{A$2'$
cited by clause; A$2'$(ii) is consumed by no part of this lemma}. Fix $\bar\tau>0$ and $h\in\mathcal V$,
let $(D,\mathsf r)$ be an admissible direction pair at $(h,\bar\tau)$ in the sense of
\cref{def:dirclass}, and put
\[
\zeta_\tau\;:=\;\big(\vdagger+\tau_1h,\;L_\star+\tau_2D,\;\varrho+\tau_3\mathsf r\big),
\qquad
\Psi(\tau)\;:=\;E\big[\psi(\theta_\dagger,\zeta_\tau)\big],
\qquad
\tau\in[-\bar\tau,\bar\tau]\times\IR^{2}.
\]
Then $\Psi$ is finite on $[-\bar\tau,\bar\tau]\times\IR^{2}$ and
\begin{enumerate}
\item[(i)] $\partial^{2}_{\tau_1\tau_1}\Psi(0)=E\big[\partial_u\ell_\star\,\widetilde W\,(\Delta h)^{2}\big]
=\int_{\mathcal X}\sigma_q\,h^{2}\,dx-2\big\langle\mathcal K h,h\big\rangle_{L^{2}(dx)}$, where
$q(x_1,x_2)=E[\partial_u\ell_\star\,\widetilde W\mid X_1=x_1,X_2=x_2]\in\IR^{p}$, in the jointly measurable
version fixed in the proof, and $\mathcal K$ is the integral operator with kernel $q\,f_{X_1,X_2}$;
\emph{the contrast stands to the right of $\partial_u\ell_\star$ in both displays, which is the only
order that multiplies and is the order $\psi$ carries; $(\Delta h)^{2}$ is a scalar and commutes};
\item[(ii)] $\partial^{2}_{\tau_2\tau_2}\Psi\equiv0$ and $\partial^{2}_{\tau_3\tau_3}\Psi\equiv0$
\emph{identically on $[-\bar\tau,\bar\tau]\times\IR^{2}$}, not merely at $\tau=0$;
\item[(iii)] if $D$ satisfies N1$'$ then $\partial^{2}_{\tau_2\tau_3}\Psi(0)=0$; away
from $\tau_1=0$ the cross-partial is $-E[D(\Delta U-\tau_1\Delta h,\Delta Z)\,\mathsf r(\Delta Z)]$, which
need not vanish and is bounded in Euclidean norm by
$\|\mathsf r\|_\infty\,\big\|D(\cdot-\tau_1\Delta h,\cdot)\big\|_{L^{2}(P)}\le\|\mathsf r\|_\infty\|D\|_{0,h}$
, the norm of the \emph{shifted} direction, which is not controlled by $\|D\|_{L^{2}(P)}$, and
$\|\mathsf r\|_\infty:=\sup_z\|\mathsf r(z)\|$ the sup of the Euclidean norm of the
$\IR^{\mathsf L-1}$-valued direction. \emph{No factor $c_{\widetilde W}$ appears here}: the
product $D\mathsf r$ contracts $D$ against $\mathsf r$ and not against $\widetilde W$, and
$\|D\mathsf r\|\le\|D\|\,\|\mathsf r\|$ with the Frobenius norm on $D$;
\item[(iv)] if the pair is differentiable in the sense of \cref{def:dirclass}, then, with the mixed
partials taken in the order $\partial_{\tau_1}\big[\partial_{\tau_2}\Psi\big]$ and
$\partial_{\tau_1}\big[\partial_{\tau_3}\Psi\big]$, \emph{the reverse order returns the same two
values, verified by direct computation in the proof, so no symmetry theorem is invoked}, 
$\partial^{2}_{\tau_1\tau_2}\Psi(0)=-E\big[\partial_uD\,\widetilde W\,\Delta h\big]$ and
$\partial^{2}_{\tau_1\tau_3}\Psi(0)=E\big[\ell_\star\,\mathsf r(\Delta Z)\,\Delta h\big]$, bounded in
Euclidean norm by
\begin{equation}\label{eq:prodbound}
\big\|\partial^{2}_{\tau_1\tau_2}\Psi(0)\big\|
\;\le\;2\,c_{\widetilde W}\,\bar c_\pi^{1/2}\,\big\|\partial_uD\big\|_{L^{2}(P)}\,\|h\|_{L^{2}(\pi)}
\;\le\;2\,c_{\widetilde W}\,\bar c_\pi^{1/2}\,\|D\|_{1,h}\,\|h\|_{L^{2}(\pi)}
\end{equation}
and, respectively, by
\begin{equation}\label{eq:rprodbound}
\big\|\partial^{2}_{\tau_1\tau_3}\Psi(0)\big\|
\;\le\;2\bar c_\pi^{1/2}\,\big\|\ell_\star\big\|_\infty\,\big\|\mathsf r\big\|_{L^{2}(P)}\,\|h\|_{L^{2}(\pi)}
\;\le\;2\bar c_\pi^{1/2}\,\big\|\ell_\star\big\|_\infty\,\big\|\mathsf r\big\|_\infty\,\|h\|_{L^{2}(\pi)},
\end{equation}
where $\|\mathsf r\|_{L^{2}(P)}:=\big(E\big\|\mathsf r(\Delta Z)\big\|^{2}\big)^{1/2}$, \emph{the
second moment of the \textbf{norm} of the $\IR^{\mathsf L-1}$-valued direction, which is what A4(iii)
prices, and not the scalar second moment, which is not a norm on a vector}, in each case a product
of the two direction magnitudes.
\emph{\eqref{eq:prodbound} carries $c_{\widetilde W}$ and \eqref{eq:rprodbound} does not, and
the asymmetry is exactly which vector the matrix is contracted against}: $\partial_uD\,\widetilde W$
contracts against the contrast, costing $\|\widetilde W\|\le c_{\widetilde W}$, while
$\ell_\star\mathsf r$ contracts against $\mathsf r$, whose norm is already the quantity being priced.
At $\mathsf L=2$, $c_{\widetilde W}=1$, $\mathsf r$ is scalar, and both displays are the
ones earlier rounds printed.
\end{enumerate}
In each of \eqref{eq:prodbound} and \eqref{eq:rprodbound} the first inequality is the operative one
and the second is a stated weakening, and in both the operative norm is one the assumption block
prices. $\|\partial_uD\|_{L^{2}(P)}$ is what the proof delivers and what A5(iv) prices, at
$D=\hat L-L_\star$, as $\epsilon_n$; $\|D\|_{1,h}$ is priced by nothing in the assumption block and
appears here only because it is the envelope that licenses the interchange producing the identity.
$\|\mathsf r\|_{L^{2}(P)}$ is what A4(iii) prices, at $\mathsf r=\hat\varrho-\varrho$, as $r_n$;
$\|\mathsf r\|_\infty$ is priced by nothing, no clause of the block bounds
$\|\hat\varrho-\varrho\|_\infty$. The sharper $L^{2}$ form costs one application of the
% ARCHAEOLOGY.  Rounds 1--7 printed only the sup-norm form of eq:rprodbound, which
% is why the instantiation below was mis-cited to A5(iv).
Cauchy--Schwarz inequality against the bounded factor $\ell_\star$, and $\mathsf r$ is bounded by
\cref{def:dirclass}(a), so \eqref{eq:rprodbound} is a strengthening at no cost in hypotheses.
\Cref{rem:dirmember} states what must be established before either bound is instantiated at
$D=\hat L-L_\star$, and gives a sufficient condition under which it is.
\end{lemma}

\begin{remark}[Membership of the estimated direction in \cref{def:dirclass}, and why it is a forward obligation rather than an assumption]\label{rem:dirmember}
\emph{This remark does \textbf{not} import \cref{lem:secondord}'s hypothesis block, and in particular not
its direction-pair hypothesis}, that membership is exactly what \eqref{eq:dirmember} exists to
establish, so importing it would make the remark assume what it delivers. What the displays below use,
exhaustively: A1b(iv); C1(ii), through \cref{lem:nusigma}(i); B1(i), (ii), (iv.a), (iv.d) together with
A2$'$(i.a)--(i.b), through \cref{cor:elldiff}; and A2$'$(iii), through \cref{def:weight}.
% ARCHAEOLOGY.  The round-7 provenance sentence was wrong in both directions --
% too broad, in importing lem:secondord's whole block, and too narrow, in writing
% "and nothing else" after a list omitting B1(i), (ii), (iv.a), which cor:elldiff
% consumes to deliver ||ell_star||_inf < infinity.  Nothing was unsound, those
% clauses being in force through the lemma.

\Cref{lem:secondord}(iii)--(iv) are stated at a \emph{fixed} admissible direction pair $(D,\mathsf r)$ at
$(h,\bar\tau)$. Step 6 or 7 will instantiate them at $D=\hat L-L_\star$, $\mathsf r=\hat\varrho-\varrho$
and $h=\hat\varphi-\vdagger$, all three random, and \emph{that instantiation requires membership, which
is a separate statement from any rate}: clause (iii) needs $\|\hat L-L_\star\|_{0,h}<\infty$, clause (iv)
needs in addition that $\hat L-L_\star$ be differentiable in $u$ with continuous derivative and
$\|\hat L-L_\star\|_{1,h}<\infty$, and both need $\hat\varrho-\varrho$ bounded and satisfying
\cref{def:dirclass}\textup{(a)}'s range condition, \emph{which it does, at every $\tau_3\in[0,1]$, by
convexity of \eqref{eq:trunset}, so the propensity half of the pair costs nothing}. Neither
seminorm is assumed anywhere, and neither is discharged before Step 6.

It is, however, implied by the construction of $\hat L$, and that is why it is recorded here rather
than assumed. Suppose $\hat L=L^{\,\hat b}_{\hat\ell}$ is an admissible weight in the sense of
\cref{def:weight}(b), generated by an $\hat\ell$ that is bounded and continuous in $u$. Then
$D:=\hat L-L_\star$ is differentiable in $u$ at every point with $\partial_uD=\hat\ell-\ell_\star$
continuous in $u$, and for every $h\in\mathcal V$ and every $\bar\tau>0$
\begin{equation}\label{eq:dirmember}
\big\|D\big\|_{1,h}\;\le\;\big\|\hat\ell-\ell_\star\big\|_\infty,
\qquad
\big\|D\big\|^{2}_{0,h}\;\le\;2E\big\|\hat b(\Delta Z)-b_\star(\Delta Z)\big\|^{2}
+4\big\|\hat\ell-\ell_\star\big\|^{2}_\infty
\Big\{E\big[\Delta U^{2}\big]+4\bar\tau^{2}\bar c_\pi\|h\|^{2}_{L^{2}(\pi)}\Big\},
\end{equation}
both finite, $\|\ell_\star\|_\infty<\infty$ by \cref{cor:elldiff} under A2$'$(i), so
$\|\hat\ell-\ell_\star\|_\infty<\infty$. The first holds because the sup-norm bound on $\partial_uD$ is
uniform over \emph{all} arguments, so the supremum inside \eqref{eq:shiftnormone} is dominated by the
constant $\|\hat\ell-\ell_\star\|_\infty$ before any expectation is taken, this is the one place where
a sup-inside condition costs nothing, and it is why the asymmetry of \cref{rem:supinside} is not a
practical burden. The second follows from
$\|D(u,z)\|\le\|\hat b(z)-b_\star(z)\|+\|\hat\ell-\ell_\star\|_\infty|u|$ via \eqref{eq:weightdef},
$(\varsigma+\varsigma')^{2}\le2\varsigma^{2}+2\varsigma'^{2}$ applied twice, A1b(iv), and
$\|\Delta h\|_{L^{2}(P)}\le2\bar c_\pi^{1/2}\|h\|_{L^{2}(\pi)}$ from \cref{lem:nusigma}(i); the first term
is finite because both weights are admissible, $E\|\hat b(\Delta Z)\|^{2}<\infty$ by hypothesis and
$E\|b_\star(\Delta Z)\|^{2}<\infty$ by A2$'$(iii).

One reading of \eqref{eq:dirmember} must be blocked explicitly, because the notation invites it.
The display is quantified over \emph{every} $\bar\tau>0$; its right-hand side carries $\bar\tau^{2}$ and
its left-hand side carries no $\bar\tau$ at all. Read naively one would minimise over $\bar\tau$, send
$\bar\tau\downarrow0$, and conclude a $\bar\tau$-free bound
$\|D\|^{2}_{0,h}\le2E\|\hat b-b_\star\|^{2}+4\|\hat\ell-\ell_\star\|^{2}_\infty E[\Delta U^{2}]$.
\emph{That inference is false.} The seminorm $\|D\|_{0,h}$ suppresses a dependence on $\bar\tau$ in its
glyph but not in its value, it is a supremum over $|\tau_1|\le\bar\tau$, and by the monotonicity
recorded in \cref{def:dirclass} the left-hand side is \emph{also} nondecreasing in $\bar\tau$. The
display is therefore a family of bounds indexed by $\bar\tau$, matched term for term, and neither side
may be optimised in isolation. \emph{Step 6 or 7 carries this display forward at the $\bar\tau$ its own
statement fixes, and at no other.}

Boundedness of $\hat\ell$ is needed \emph{per
realisation only}; no uniformity in $n$ is used, because finiteness at each $n$ is all that clauses
(iii)--(iv) consume and the rate is A5(iv)'s separate job.

\emph{What Step 6 or
7 owes is therefore a verification, not a hypothesis}: that its $\hat L$ is an admissible weight in the
sense of \cref{def:weight}(b) with $\|\hat\ell\|_\infty<\infty$ and $E\|\hat b(\Delta Z)\|^{2}<\infty$.
\emph{It is not a new clause of A5}, for the reasons \cref{rem:obligations} gives, and it is
inventoried there with the other forward obligations.
% ############################################################################
% ##  ARCHAEOLOGY.  "Why this is not a new clause of A5" -- 17 compiled lines
% ##  moved here 2026-08-01, the substance being carried at rem:obligations.
% ##  (a) Every clause of A5 constrains an object the maintained block itself
% ##      introduces -- h, hat-ell, delta-hat_n, hat-varrho -- whereas ||D||_{0,h}
% ##      and ||D||_{1,h} are seminorms of def:dirclass evaluated at a DIRECTION
% ##      PAIR that no maintained clause fixes, so a clause about them would be a
% ##      clause about a definition's argument.  The ground is NOT that A5 prices
% ##      rates only: A5(vi) is ||h||_{H^s} = O_P(1), qualitative, and A5(vi') is a
% ##      w.p.->1 envelope at a fixed radius.  A round-7 draft gave the refuted
% ##      ground.
% ##  (b) Both seminorms are indexed by (h, bar-tau) with h random, so a clean
% ##      assumption would have to quantify over a deterministic class containing h
% ##      with probability approaching one -- the ball cal V_n of L-19.
% ##  (c) eq:dirmember shows the condition is IMPLIED once hat-L is built as an
% ##      admissible weight from a bounded continuous hat-ell.
% ##  It is the same species of forward obligation as A2'(i.b)/L-21 and N1'.
% ############################################################################
\end{remark}

\begin{remark}[The open verification obligations, in one place]\label{rem:obligations}
The block carries four obligations that are \emph{verifications} rather than assumptions.
\textup{(O1)} is discharged, \textup{(O4)} is discharged, and \textup{(O2)} and \textup{(O3)}
\textbf{remain open}. \emph{Neither open obligation is discharged by being carried as a hypothesis},
since a hypothesis of a theorem is not a verification; what each has is a \textbf{named consumer that
carries it explicitly}, which is what this remark's test asks of an open obligation.
\textup{(O3)} is carried by hypothesis (R-a) of \cref{thm:rootn}, which
asserts the construction $\hat L^{(-k)}=L^{\,\hat b^{(-k)}}_{\hat\ell^{(-k)}}$ with a bounded derivative and
so implies \cref{def:dirclass} membership through \eqref{eq:dirmember}. \textup{(O2)} is carried by (R-c),
which is a \emph{disjunction}: (R-c.1) assumes N1$'$ at every $D_k$, and (R-c.2) instead \textbf{prices its
failure} at $r_n\epsilon^{L}_n$, where $\epsilon^{L}_n$ is an assumed $L^{2}(P)$ rate for $\|D_k\|$ and
\eqref{eq:biasb4} is the exact cost. Pricing a failure is the
honest third option and it is neither a discharge nor an evasion.
% ARCHAEOLOGY.  The four are listed together because they were scattered across
% three different KINDS of place -- one a tagged clause, one a tagged property, two
% untagged -- and scattering is what let one of them migrate into the assumption
% block unstated.  (O2) and (O3) acquired their named consumers at Step 7.

\emph{The test used, stated once.} A forward obligation is a verification and not an assumption
\textbf{so long as no stated result's truth depends on it}: every present result carries the condition
as an explicit hypothesis, so that a reader can see what is owed and by whom. \emph{The moment a stated
assumption cannot be written down without the obligation, the obligation has already become an
assumption}, which is why A5(iv)'s first equality is
printed as a conditional consequence rather than as part of the clause.
% ARCHAEOLOGY.  A5(iv) failed that test through round 7.

\begin{enumerate}
\item[(O1)] A$2'$(i.a)--(i.b) in sup-inside form, on $\ell^\Pi_\star$. \textbf{DISCHARGED} at
\cref{lem:reprsmooth}, which derives both clauses from the
source condition A2c on $\ell^\Pi_\star$ at index $\mu''>\tfrac32$, in the sup-inside form and with the
explicit budget \eqref{eq:l21budget}. \emph{A$2'$(i) does not thereby
cease to be an assumption}, an assumption with a proved sufficient condition is still an assumption, 
\emph{but it ceases to be an unverifiable primitive}, A2c being one condition on one object, stated
in the Hilbert-scale language H1 already uses and reportable through \eqref{eq:usobolev}.
% ARCHAEOLOGY.  The discharge is retroactive: eq:Wderiv, cor:elldiff and the
% ||partial_u ell_star||_inf that lem:secondord(i) and lem:solset(iii) consume all
% rest on A2'(i.b) and were written at Step 5, before Step 4 ran.  Nothing was
% assumed silently, each carrying the clause as an explicit hypothesis.  L-21's
% target is A2'(i) and never A2b, which is never asserted.
\item[(O2)] N1$'$ at the estimated direction, $E[\hat L-L_\star\mid\Delta Z]=0$.
\textbf{Where the condition lives now:} the tagged property N1$'$, stated as a hypothesis of
\cref{lem:secondord}(iii), \emph{and, as of Step 7, hypothesis (R-c) of \cref{thm:rootn}}.
\textbf{Discharge site:} Step 6 or 7, by constraining the estimator so that
$\hat L$ is centred; failing that, the leading cross term does not vanish and enters at order
$\|\mathsf r\|_\infty\|D\|_{0,h}$. \textbf{Status: STILL OPEN, and now priced rather than only placed.}
Step 7 does not construct a centred $\hat L$, so the obligation stands; what it adds is (R-c.2), under which
the failure costs exactly $r_n\epsilon^{L}_n$ by \eqref{eq:biasb4}, and (R-c.1), under which
$\Psi^{(3)}_k=0$. \emph{The cheap discharge, re-anchoring $\hat L$ at the N1 anchor of
\cref{rem:anchor}, requires an estimate of $E[\int_0^{\Delta U}\hat\ell\mid\Delta Z]$ and is therefore a
further nuisance}, which is why it is not taken \textbf{for the direction $D_k$}.
\emph{That declining is scoped to $D_k$ and does not extend to $L_\star$}, and the two are not
interchangeable here. N1$'$ at $D_k$ is optional \emph{because it is priced}: (R-c.2) makes its failure
cost exactly $r_n\epsilon^{L}_n$ by \eqref{eq:biasb4}, a product of two estimation errors. N1 at
$L_\star$ is not optional and is not priced: it is hypothesis \textup{(T-2)} of \cref{thm:rootn}, and
\textbf{nothing anywhere prices its failure}, which by \cref{cor:centarity}\textup{(iii)} enters at
$r_n$ \emph{undamped} --- no second small factor multiplies it, so after $\sqrt n$ scaling it diverges.
\emph{A construction that targets a non-centred anchor therefore does not fail \textup{(O2)}; it fails
\textup{(T-2)}}, and reading the decline above as covering both is the error \cref{cor:centarity}
exists to block.
\item[(O3)] Membership of $\hat L-L_\star$ in \cref{def:dirclass}, via $\hat L=L^{\,\hat b}_{\hat\ell}$.
\textbf{Where the condition lives, three homes, and the third is of a different kind:}
\cref{rem:dirmember}, untagged, with the discharge already
written at \eqref{eq:dirmember} under one property of the estimator; hypothesis (R-a) of
\cref{thm:rootn}; and \cref{def:crossfit}\textup{(c)}, which
asserts the construction $\hat L^{(-k)}=L^{\,\hat b^{(-k)}}_{\hat\ell^{(-k)}}$ \emph{together with} the three
properties of $\hat\ell^{(-k)}$, Borel, bounded, $u$-continuous, that make \cref{def:weight}(b) well
posed. \emph{A remark can be declined and a theorem hypothesis can be refused, but a \textbf{definition} is
adopted}: every Part-I result that adopts \cref{def:crossfit} carries (O3) whether or not it lists (R-a),
and \cref{def:onestep}'s Jacobian identity $\partial_\theta\bar\psi_n(\theta)=-\hat J(\theta)$ consumes it
directly through $\partial_u\hat L^{(-k)}=\hat\ell^{(-k)}$.
\textbf{Discharge site:} Step 6 or
7, which must verify that its $\hat L$ is an admissible weight generated by a bounded, $u$-continuous
$\hat\ell$ with $E\|\hat b(\Delta Z)\|^{2}<\infty$. \textbf{Status: STILL OPEN, and carried as
an explicit hypothesis by a stated theorem.} (R-a) of \cref{thm:rootn} asserts exactly that verification, at
each fold and with $\|\partial_u\hat\ell^{(-k)}\|_\infty$ added, so \eqref{eq:dirmember} applies and
\cref{def:dirclass} membership follows; \emph{assuming the construction is not performing it}.
\emph{The obligation is in addition presupposed by A5(iv)}, whose identity
$\|\partial_u(\hat L-L_\star)\|_{L^{2}(P)}=\|\hat\ell-\ell_\star\|_{L^{2}(P)}$ cannot be written without
it; A5(iv) therefore prices $\hat\ell$ alone and states the identity as a consequence conditional
on this obligation.
% ARCHAEOLOGY.  The third home is Step 7's dispatch-11 revision of def:crossfit(c);
% a round-11 draft of this item listed two homes and had itself created the third.
% Through round 7 A5(iv) printed the identity inside the clause, so the obligation
% had migrated into the maintained block unstated.
\item[(O4)] \emph{Construction of $\hat{\mathcal A}_\varphi$.} \textbf{DISCHARGED, with one negative
answer and one amendment.} The operator is
exhibited at \eqref{eq:hatAcon} and characterised by \cref{prp:hatAprop}; $\hat T$ is exhibited at
\eqref{eq:hatT}; A7 is verified in the amended form \eqref{eq:a7amend}
(\cref{prp:stopfeas}); and \emph{A6 is refuted, not verified} (\cref{prp:tccfails}) and has been
retired. \emph{A discharge with a negative answer is a discharge}: the obligation was to verify, and
the verification has an outcome.
% ARCHAEOLOGY.  Through round 7 the three symbols hat-A_{tilde-varphi},
% hat-A_cond and hat-A*_cond appeared in the MAINTAINED assumption A6 and in C3
% with no defining display anywhere, so A6's content was not determined by the
% block that states it -- and once it was determined, it turned out to be false.
\end{enumerate}

\emph{No new clause of A5, and no new maintained
assumption, is created for any of (O1)--(O4). In particular A2c is not maintained}: it is a
property tag with exactly one consumer, and promoting it would convert a discharged
obligation into a new assumption, the trade this remark's test forbids. The ground is
that every clause of A5 constrains an object the maintained block \emph{introduces}, while (O2)--(O4) are
constructive statements about objects the block does not construct.
\end{remark}

% WHY (i) IS STATED AT A GENERIC ADMISSIBLE REPRESENTER -- ROUND-7 REPAIR OF A
% VACUOUS EQUIVALENCE, AND IT REMOVES A HYPOTHESIS RATHER THAN ADDING ONE.
% Through round 6, (i) assumed A2' and then concluded, as the right half of an
% "if and only if", that A*_cond ell^Pi_star = -nu/pi.  BUT A2' DECLARES
% ell^Pi_star TO BE A SOLUTION OF EXACTLY THAT EQUATION.  The right-hand side was
% therefore true by hypothesis and the second "iff" asserted nothing -- the same
% species of defect as M-8 (prp:orthog(iii) read "iff N1, under N1") and as the
% lem:shiftinv split (the equality would have been vacuous under A1), both
% already repaired.
%
% QUANTIFYING OVER lem:solset's "ADMISSIBLE SELECTIONS" WOULD NOT HAVE FIXED IT:
% a selection SOLVES the adjoint equation by definition, so the right-hand side
% would still be true by hypothesis.  What restores content is quantifying over
% admissible REPRESENTERS, which are not required to solve anything.  The
% equivalence then says what it was always meant to say: among weights regular
% enough for the expansion to exist, first-order insensitivity in varphi holds
% EXACTLY at the solutions of the adjoint equation.
% CONSEQUENCE FOR THE HYPOTHESIS LIST: part (i) no longer assumes A2 or A2' at
% all.  Parts (ii) and (iii), and every downstream consumer, still do.
%
% ROUND-7 POSTFIX: Psi_L IS NOW DEFINED IN PLACE (eq:PsiL), NOT BY REFERENCE.
% The round-7 statement read "write Psi_L for the map of lem:secondord with L in
% place of L_star".  That pointed at a lemma whose hypothesis list CONTAINS A2',
% inside the one clause headed "A2 is not assumed", and whose Psi is defined only
% after fixing tau-bar, h and an admissible direction pair (D, r) -- none of
% which part (i) introduces.  The mathematics was right and the STATEMENT was
% under-specified: the hypothesis reduction that is this repair's whole point
% was not readable off it.
%   eq:PsiL fixes that at zero cost, and is STRICTLY LESS demanding than the
% import it replaces: it needs no radius tau-bar (it is finite on all of R), no
% direction pair, and no clause of A2'.  The agreement Psi_{L_star}(tau_1) =
% Psi(tau_1, 0, 0) is stated so that nothing downstream has to guess whether the
% two objects are the same where both exist.  DO NOT re-import: the moment the
% definition is routed through lem:secondord again, part (i) silently reacquires
% A2' and the reduction is lost.

% EXTRACTED FROM THE STATEMENT OF prp:orthog on 2026-08-02.  The finiteness
% argument for eq:PsiL was a PROOF sitting inside a statement, which the house
% style forbids outright; the hypothesis-split paragraphs were commentary.  The
% claim "which is finite for every tau_1" STAYS in the statement -- it is a claim,
% and it is the argument for it that moved.  No \label moved.
\begin{remark}[The hypothesis split of \cref{prp:orthog}, and what makes its displays well posed]\label{rem:orthoghyp}
\begin{enumerate}
\item[(1)] \emph{A1 is \textbf{not} in the global list}: it is consumed by part (ii) alone, where it
is assumed and cited, while parts (i) and (iii) use it nowhere. The split follows the precedent
already in this file at \cref{lem:shiftinv}, and its practical point is the same: \cref{lem:balance},
which inherits ``the hypotheses of \cref{prp:orthog}(i)'', would otherwise inherit A1 as surplus and
consume it in no step.
\item[(2)] \emph{N1 is not assumed in any part}: (i) and (ii) do not use it and (iii) characterises it.
\emph{And part (i) does not assume A2 either}, being stated at an arbitrary admissible representer.
\item[(3)] Why \eqref{eq:PsiL} is finite for every $\tau_1\in\IR$. The second term because
$\|L(u,z)\widetilde W\|\le c_{\widetilde W}\{\|b(z)\|+\|\ell\|_\infty|u|\}$, by
\eqref{eq:weightdef} and $\|L\widetilde W\|\le\|L\|\,\|\widetilde W\|$ with the Frobenius norm on $L$ and
$\|\widetilde W\|\le c_{\widetilde W}$ from \cref{cor:rangegap}, \emph{the vector bound and
not \textup{(P1)}, which is the \textbf{per-cell} statement $|\widetilde W_l|\le1$ and does not bound the
Euclidean norm; at $\mathsf L=2$ the two coincide and $c_{\widetilde W}=1$}, 
$E\|b(\Delta Z)\|^{2}<\infty$ by admissibility, $\|\ell\|_\infty<\infty$ by \cref{cor:elldiff} under A2b
, \emph{the tag is A2b and not A$2'$(i), because $\ell^\Pi$ there is a generic admissible representer
and not $\ell^\Pi_\star$}, and $E|\Delta U|,E|\Delta h|<\infty$ by A1b(iv) and \cref{lem:nusigma}(i);
the first because $E\|\Delta Z\Delta U\|<\infty$ and
$E\|\Delta Z\Delta h\|\le\|\Delta Z\|_{L^{2}(P)}\|\Delta h\|_{L^{2}(P)}<\infty$ by A1b(iv), the
Cauchy--Schwarz inequality and \cref{lem:nusigma}(i).
\item[(4)] \emph{\Cref{prp:orthog}(i) defines $\Psi_L$ and does not import it}: no radius $\bar\tau$, no
direction pair $(D,\mathsf r)$ and no clause of A2$'$ enters \eqref{eq:PsiL}, which is what makes the
hypothesis reduction of that clause readable off its own statement.
\item[(5)] \emph{Part (ii) is the one part that consumes A1}, through the conditional independence that
makes $E\big[D\,\widetilde W\mid\Delta Z\big]=E\big[D\mid\Delta Z\big]\,E\big[\widetilde W\mid\Delta Z\big]=0$
, a $p\times(\mathsf L-1)$ matrix against the zero vector of $\IR^{\mathsf L-1}$, hence $0$ in $\IR^{p}$.
\emph{The product exists only in this order}, $D$ being $p\times(\mathsf L-1)$ and $\widetilde W$ being
$(\mathsf L-1)\times1$ by \cref{def:dirclass} and \eqref{eq:contrast}; \emph{what A1 supplies is
unchanged}, namely conditional independence of $D(\Delta U,\Delta Z)$ and $W$ given $\Delta Z$, which is
what lets the conditional expectation of the product factor at any value type of $D$. At
$\mathsf L=2$, $\widetilde W$ is a scalar and the display is the one earlier rounds printed.
\item[(6)] \emph{Both sides of part (iii)'s equivalence are well defined}, because A2$'$(iii) with
\cref{cor:elldiff} and A1b(iv) give $E[\|L_\star\|\mid\Delta Z]<\infty$ almost surely.
\end{enumerate}
\end{remark}

\begin{lemma}[L-7: the balancing bound on the first-order bias]\label{lem:balance}
The first-order bias in the structural function is negligible under either of two independent
conditions: a product rate for the debiasing weight, or a sieve balancing constraint with an
approximation condition.
Suppose the hypotheses of \cref{prp:orthog}(i) hold together with A2, A2$'$(i) and A2$'$(iii), 
\emph{A$2'$ cited by clause; A$2'$(ii) is consumed by no part of this lemma}, and adopt the following
\emph{sample-splitting convention}: the index set $\{1,\dots,n\}$ is partitioned into an auxiliary fold
$\mathcal I_a$ and a main fold $\mathcal I_m$, independent by A0; $\hat\ell$ is a
$\sigma(\mathcal I_a)$-measurable \textbf{random element} of
$L^{2}(P_{\Delta U,\Delta Z};\IR^{p\times(\mathsf L-1)})$, and one
\textbf{jointly measurable} representative of it is fixed, measurable with respect to
$\sigma(\mathcal I_a)\otimes\mathfrak B(\IR\times\IR^{p})$, which a random element of an $L^{2}$ space
over a separable metric space admits, exactly as \cref{def:crossfit}\textup{(c$'$.1)} fixes one for the
three integrated nuisances of Part~I. \emph{Membership and joint measurability are two statements and
both are needed here}: Steps~0
and~2 freeze $\hat\ell$ inside $\nu_g$ and Step~2's identity
% ARCHAEOLOGY.  Through dispatch 15 this convention stated membership only.
$\nu_g=\nu_{(\hat\ell-\ell_\star)\widetilde W}$ \emph{is} \textup{(P8)}'s freezing lemma at
$\mathcal G=\sigma(\mathcal I_a)$, whose hypothesis is joint measurability of the \emph{composite}, and a
measurable map into $L^{2}$ supplies nothing of the kind on its own. \emph{What this site needs is
strictly \textbf{less} than Part~I needs, and the clause is scoped to it and not imported}: here
$\hat\ell$ is evaluated only at an \textbf{undistorted} $(\Delta U,\Delta Z)$ draw, whose law is the one
the container is built on, so any two representatives agree outside a $P$-null set and the choice is
immaterial; \eqref{eq:prodmeas}, which \cref{def:crossfit}\textup{(c$'$.1)} imposes because Part~I
evaluates $\hat\ell$ at \emph{shifted} first arguments whose law need not be dominated by
$P_{\Delta U,\Delta Z}$, is \textbf{not} required here and is \textbf{not} imposed here
, \emph{and here $\hat\ell$ is a free estimator of $\ell_\star$, not the $u$-derivative of anything:
nothing in this lemma assumes that $\hat\ell=\partial_u\hat L$ for a weight $\hat L$, and A5(iv) prices
$\hat\ell$ in exactly this sense. The relation $\hat L=L^{\,\hat b}_{\hat\ell}$ is a construction
constraint, and it is owed by Step 6 or 7 (\cref{rem:obligations}), not supplied here}, ;
$\hat\varphi$ and $\hat\nu_g$ are computed on $\mathcal I_m$; and every $\nu$-object and every
$O_P$ statement below is conditional on $\sigma(\mathcal I_a)$, so that $\hat\ell$ is frozen wherever it
appears inside a conditional expectation. Put $g:=\Delta Z+\hat\ell\,\widetilde W\in\IR^{p}$, 
\emph{the contrast to the right of the $p\times(\mathsf L-1)$ estimate, the order \eqref{eq:psidef}
carries and the only one that multiplies}, and
$h:=\hat\varphi-\vdagger$, let $\hat\nu_g$ be an estimator of $\nu_g$ computed on $\mathcal I_m$, and let
$P_{L_x}$ denote the $L^{2}(\pi)$-orthogonal projection onto
$\operatorname{span}(b_1,\dots,b_{L_x})$. Then
\begin{equation}\label{eq:balance3}
\big\|\langle\nu_g/\pi,\,h\rangle_{L^{2}(\pi)}\big\|
\;\le\;\underbrace{\big\|(\nu_g-\hat\nu_g)/\pi\big\|\,\|h\|}_{\text{(I)}}
\;+\;\underbrace{\big\|\langle\hat\nu_g/\pi,\,P_{L_x}h\rangle\big\|}_{\text{(II)}}
\;+\;\underbrace{\big\|\hat\nu_g/\pi\big\|\,\big\|(\mathrm{Id}-P_{L_x})h\big\|}_{\text{(III)}},
\end{equation}
% ARCHAEOLOGY.  The left-hand side and term (II) carried a SCALAR MODULUS until the
% present pass -- correct only at p = 1, and contradicted by clause (i) below, which
% prices nu_g / pi in the VECTOR L^2(pi) aggregate.  The proof half moved in the
% SAME pass: proofs.tex, lem:balance Steps 1 and 3.  Same repair at eq:biasb1.
\emph{the bars on the left-hand side and on \textup{(II)} being the \textbf{Euclidean} norm on $\IR^{p}$
and \textbf{not} a modulus}, $g\in\IR^{p}$ here, so \eqref{eq:nudef} makes $\nu_g$ and $\hat\nu_g$
coordinatewise $\IR^{p}$-valued and both pairings are vectors, scalar only at $p=1$, while every norm
on a function is $L^{2}(\pi)$, that of an $\IR^{p}$-valued function being the aggregate
$\big(\int\|\cdot\|^{2}\pi\,dx\big)^{1/2}$ at which clause \textup{(i)} is proved, and
\begin{enumerate}
\item[(i)] $\|\nu_g/\pi\|_{L^{2}(\pi)}\le2\,c_{\widetilde W}\,\bar c_\pi^{1/2}\,\|\hat\ell-\ell_\star\|_{L^{2}(P)}$,
with $c_{\widetilde W}=\min\{\sqrt2,\sqrt{\mathsf L-1}\}$ of \cref{cor:rangegap} and
$\|\cdot\|_{L^{2}(P)}$ the Frobenius $L^{2}$ norm of the $p\times(\mathsf L-1)$ difference.
\emph{The constant is \cref{cor:rangegap}'s and it is $c_{\widetilde W}$ and not $1$, because
Step~2's identity puts $\nu_g$ at $g=(\hat\ell-\ell_\star)\widetilde W$ and the passage
$\|(\hat\ell-\ell_\star)\widetilde W\|\le\|\hat\ell-\ell_\star\|\,\|\widetilde W\|$ reads the Euclidean
bound on the contrast vector and \textbf{not} \textup{(P1)}}, which is the per-cell statement
$|\widetilde W_l|\le1$; \emph{at $\mathsf L=2$ the two agree, $c_{\widetilde W}=1$, and the
display is $2\bar c_\pi^{1/2}\|\hat\ell-\ell_\star\|_{L^{2}(P)}$ as earlier rounds printed}. \emph{The
factor is a constant free of $n$}, so clauses \textup{(ii)} and \textup{(iii)} are unaffected: neither
reads the value of the constant, only that it is finite;
\item[(ii)] \emph{(rate route)} under A5(ii) and A5(iv) alone, the left-hand side is $o_P(n^{-1/2})$
\emph{with no condition on $L_x$, no balancing constraint, no appeal to \eqref{eq:nuhat} below, and no
appeal to A5(vi)};
\item[(iii)] \emph{(balancing route)} suppose the sieve solve sets $\text{\rm(II)}=0$, that A8 holds with
a B-spline basis of degree exceeding $s-\tfrac12$, that A5(vi) holds, and that
\begin{equation}\label{eq:nuhat}
\big\|(\hat\nu_g-\nu_g)/\pi\big\|_{L^{2}(\pi)}=O_P\big(n^{-1/2}\big)
\qquad\text{and}\qquad
\big\|\hat\nu_g/\pi\big\|_{L^{2}(\pi)}=O_P(1).
\end{equation}
Then the left-hand side is $o_P(n^{-1/2})$ \emph{with no rate for $\hat\ell$}.
\end{enumerate}
\end{lemma}

% THREE SCOPE FACTS ABOUT lem:balance, RECORDED BECAUSE EACH WAS WRONG ONCE.
%  1. eq:nuhat IS SCOPED TO PART (iii) ENTIRELY (round-6 MAJ-4, completed here).
%     Round 5 listed its first half among (ii)'s hypotheses and then closed with
%     a sentence contradicting the direct bound two lines above.  The display now
%     sits INSIDE (iii), so no reader can carry it onto (ii).
%  2. N1 IS NOT A HYPOTHESIS (round-7).  N1 governs the varrho-direction; this
%     lemma is entirely about the varphi-direction bias.  It was listed and never
%     consumed, which under the property reading of N1 is a live defect rather
%     than a redundancy.
%  3. THE SPLIT IS A HYPOTHESIS, NOT AN ALGORITHM NOTE (round-7, F-10).  "let
%     ell-hat be computed on an auxiliary fold" is a sentence about the estimator;
%     the proof uses three things it did not state.  (a) eq:nudef defines nu_g for
%     a FIXED integrable g, so with g = Delta Z + Wtilde ell-hat the object must be
%     read with ell-hat FROZEN -- otherwise E[g | X_2 = x] mixes estimation noise
%     into the conditional expectation and Step 2's identity fails.  (b)
%     Independence of the folds, which A0 supplies ONCE THE SPLIT IS DECLARED.
%     (c) ell-hat in L^2 almost surely, so that Upsilon may be applied to
%     ell-hat - ell_star at all; A5(iv) implies it in O_P terms but does not state
%     membership.  All three are now in the hypothesis list.

% ===========================================================================
%  PART G.  STEP 4 -- THE SHARPNESS RESULT
%
%    lem:reprsmooth  (L-21) the source condition A2c delivers A2'(i.a)-(i.b) in
%                    SUP-INSIDE form, on ell^Pi_star and never on A2b
%    rem:l21route          why A2c and not the route the memo registers
%    lem:formstruct  (L-20) structure of the second-order form: decomposition,
%                    essential spectrum, compactness dichotomy, indefiniteness,
%                    and the singular sequence that carries all four
%    prp:sharpquad         the sharp constant at the SUP-OVER-BALL quantifier
%    rem:sharpscope        what prp:sharpquad does and does not claim
%
%  ORDER, AND WHY IT IS NOT A CIRCULARITY.  lem:reprsmooth is logically PRIOR to
%  Part E's cor:elldiff -- it supplies a sufficient condition for the hypothesis
%  cor:elldiff carries -- yet it is printed here, because it is Step 4's
%  deliverable and this file is ordered by dispatch.  There is no circularity, and
%  the reason is worth printing: cor:elldiff, lem:solset, lem:secondord,
%  prp:orthog and lem:balance all take A2'(i) (or A2b) as an EXPLICIT HYPOTHESIS
%  and none of them asserts it.  lem:reprsmooth proves an IMPLICATION into that
%  hypothesis.  So Part E and Part F are unchanged and their proofs are untouched;
%  what changes is that a reader can now see where A2'(i) comes from.
%
%  RETROACTIVE DISCHARGE, STATED PLAINLY BECAUSE A READER SHOULD NOT HAVE TO
%  RECONSTRUCT THE DEPENDENCE DIRECTION.  eq:Wderiv, cor:elldiff, and the
%  ||partial_u ell_star||_inf that lem:secondord(i) and lem:solset(iii) consume
%  ALL REST ON A2'(i.b).  Step 5 was written and closed before Step 4 ran, so
%  those results were carrying A2'(i.b) as an undischarged primitive.
%  lem:reprsmooth now discharges, retroactively, what Step 5 assumed.  That is
%  what escalation N-20 was about and it is what resolution (a) records.
%
%  SYMBOLS INTRODUCED IN THIS PART, all recorded in the glossary Sec.1.5c -- plus
%  four added in the DISPATCH-6 FIX PASS and recorded in Sec.1.5d:
%    ell^Pi_0        the representative A2c supplies, as distinct from the VERSION
%                    lem:reprsmooth Step 2 constructs.  THE SECOND NAME IS THE FIX:
%                    with one name, Step 2 justified its mollification "jointly
%                    measurable by Step 0" -- a property of the PRE-redefinition
%                    representative, i.e. of a different function (C45)
%    cal T_star      the BOREL full-Pi_t-measure set carrying A2c's slice
%                    condition.  Written in round 9 with no construction; now
%                    DEFINED IN A2c (m.1), where the object it qualifies lives, so
%                    that R x cal T_star is a product-measurable set (C46)
%    cal X_+, cal X_-  the two sign sets in lem:formstruct(v).  Clause (ii) is
%                    applied at a near-supremum subset formed INSIDE each, not at
%                    the raw sign sets, which need not meet (ii)'s premise (C47)
%    varkappa_k      the countable pairing family of A2c (m.2), which makes
%                    t -> ||ell^Pi(., t)||_{H^mu''} measurable.  A second LOCAL use
%                    of an existing local letter, integer-subscripted.  Phi_0 and
%                    cal S(R) were both swept and REJECTED -- Phi is the reserved
%                    Bierens transform, cal S the selection functional -- and so
%                    was cal J for the weighted transform, a third font of the
%                    reserved Jacobian J (C48)
%
%  ORIGINAL SEC.1.5c LIST:
%    H^{mu''}(R), F, xi, i   the u-Sobolev scale on R and the Fourier apparatus
%                    that defines it, all confined to A2c and to Step 1 of
%                    lem:reprsmooth's proof, and all declared LOCAL (Sec.1.6)
%    Lambda^{(mu'')} the source budget, joining Lambda^{(0)}, Lambda^{(1)} in the
%                    parenthesised-superscript family: the superscript is the
%                    ORDER, so "the budget at order mu'' dominates the budgets at
%                    orders 0 and 1" is readable off the notation
%    c_{mu''}        the embedding constant, joining c_Phi, c_Pi, c_C
%    h_k             a SEQUENCE of directions, integer-indexed.  Sec.1.7's rule
%                    "a subscripted h is a bandwidth" is AMENDED, not breached:
%                    the bandwidths are h_u, h_x, subscripted by a VARIABLE NAME,
%                    and subscript-type separation is the device Sec.1.7 already
%                    uses three times.  Recorded in Sec.1.5c.
%    cal X_0, c_0    the set on which the multiplier is bounded away from zero and
%                    its local essential sup, both LOCAL (Sec.1.6): cal X_0 joins
%                    cal C_0, cal N_0, cal Z_0 in the subscript-zero calligraphic
%                    set family; sans-serif c_0, c_1 are the generic-constant slot
%                    Sec.1.6 already declares.
%
%  AND TWO ABBREVIATIONS ARE DECLINED, both after sweeping the slot they would
%  have claimed -- the discipline of Sec.1.5b applied to this part's own drafting.
%    NO g_vartheta FOR THE MULTIPLIER.  The multiplier is written out as
%      vartheta' sigma_q / pi at every one of its ~15 sites, which is what EVERY
%      existing site does (lem:solset(iii), def:spaces, glossary Sec.8.4) -- zero
%      sites use an abbreviation.  A one-letter name would also re-collapse the
%      distinction Sec.8.4 item 2 exists to protect (sigma_q is the MULTIPLIER;
%      B_mult is multiplication BY vartheta' sigma_q / pi), and the letter g is
%      the generic MOMENT WEIGHT (C14) and fills the subscript slot of nu_g,
%      sigma_g, hat-nu_g -- so g_vartheta would invite the reading nu_{g_vartheta}.
%    B_{cal V} IS LOCAL TO ONE PROOF ONLY.  Every STATEMENT prints the compression
%      in full.  Inside prp:sharpquad's proof the abbreviation is declared, and the
%      separation from B_sym, B_mult is by subscript TYPE -- word = a part of B,
%      space = the compression to that space -- which is the device Sec.1.7 uses
%      three times.  No new letter and NO FOURTH FONT of B, which is what C33
%      charged against the last enlargement of that letter.
%
%  AND ONE SLOT IS DECLINED RATHER THAN CLAIMED.  The memo writes the second-order
%  remainder Q(h).  ITALIC Q IS OCCUPIED: (P6b)'s regular conditional
%  distributions Q_1, Q_2, Q_{12} live at six sites in proofs.tex, and Q_{12} is
%  what supplies the jointly measurable version of q -- INSIDE the very chain
%  prp:sharpquad runs on.  Per glossary Sec.1.5b the letter being CLAIMED was
%  swept, and the remainder therefore carries NO SYMBOL: it is written as the form
%  <Bh, h>_{L^2(pi)} of eq:quadrem, which is better than a name because the sharp
%  constant is then readable off the notation.  A blocking row is added to
%  scripts/render-sync/check_rendering_sync.py so the slot cannot be claimed later
%  by accident.
% ===========================================================================

\begin{lemma}[L-21: the source condition delivers A$2'$(i), in sup-inside form]\label{lem:reprsmooth}
A source condition of index exceeding the one-dimensional $C^{1}$-embedding threshold on the selected
solution implies both sup-norm clauses of A$2'$(i), with the supremum inside the $t$-integral and with
explicit budgets; neither B1 nor any Picard criterion supplies them.
Suppose $\ell^\Pi_\star$ satisfies \textbf{A2c} at some $\mu''>\tfrac32$, with budget
$\Lambda^{(\mu'')}_\star$ as in \eqref{eq:lambdamu}, and set
\begin{equation}\label{eq:embedconst}
c_{\mu''}\;:=\;\Big(\frac{2}{\pi}\int_{\IR}\big(1+\xi^{2}\big)^{1-\mu''}\,d\xi\Big)^{1/2}\;<\;\infty ,
\end{equation}
together with the two constants the \emph{level} and the \emph{derivative} carry separately,
\begin{equation}\label{eq:embedconstsplit}
c^{(0)}_{\mu''}\;:=\;\Big(\frac{1}{2\pi}\int_{\IR}\big(1+\xi^{2}\big)^{-\mu''}d\xi\Big)^{1/2},
\qquad
c^{(1)}_{\mu''}\;:=\;\Big(\frac{1}{2\pi}\int_{\IR}\xi^{2}\big(1+\xi^{2}\big)^{-\mu''}d\xi\Big)^{1/2}.
\end{equation}
The three constants are one identity apart, and the separation costs nothing, neither in $p$,
nor in $\mathsf L$, nor in the admissible range of $\mu''$. The two integrands of
\eqref{eq:embedconstsplit} add to $(1+\xi^{2})^{1-\mu''}$, so
\begin{equation}\label{eq:embedconstid}
\big(c^{(0)}_{\mu''}\big)^{2}+\big(c^{(1)}_{\mu''}\big)^{2}\;=\;\tfrac14\,c^{2}_{\mu''},
\qquad\text{hence}\qquad
c^{(0)}_{\mu''}+c^{(1)}_{\mu''}\;\le\;\sqrt2\,\Big(\big(c^{(0)}_{\mu''}\big)^{2}+\big(c^{(1)}_{\mu''}\big)^{2}\Big)^{1/2}
\;=\;\frac{c_{\mu''}}{\sqrt2}\;<\;c_{\mu''},
\end{equation}
the inequality being the Cauchy--Schwarz inequality on two terms, with equality exactly when
$c^{(0)}_{\mu''}=c^{(1)}_{\mu''}$. Both split constants are available in closed form against
$c_{\mu''}$, which settles that equality case and makes them computable without a quadrature. One
integration by parts, whose boundary term $\big[\xi(1+\xi^{2})^{1-\mu''}\big]$ vanishes at $\pm\infty$
\emph{precisely} because $\mu''>\tfrac32$, gives
$\int_{\IR}(1+\xi^{2})^{1-\mu''}d\xi=2(\mu''-1)\int_{\IR}\xi^{2}(1+\xi^{2})^{-\mu''}d\xi$, whence
\[
c^{(1)}_{\mu''}\;=\;\frac{c_{\mu''}}{\sqrt{8(\mu''-1)}},
\qquad
c^{(0)}_{\mu''}\;=\;c_{\mu''}\sqrt{\frac{2\mu''-3}{8(\mu''-1)}},
\qquad \mu''>\tfrac32 .
\]
The two therefore agree exactly at $\mu''=2$, where both are $c_{2}/\sqrt8=\tfrac12$, so
\eqref{eq:embedconstid} is tight at that one index and strict at every other.
\emph{And \eqref{eq:embedconstsplit} imposes no new restriction on $\mu''$}: $c^{(0)}_{\mu''}$ is finite
for $\mu''>\tfrac12$ and $c^{(1)}_{\mu''}$ for $\mu''>\tfrac32$, and the second threshold is the one
$c_{\mu''}$ already carries, \eqref{eq:embedconst} diverging at $\mu''=\tfrac32$ for the same reason
\eqref{eq:embedconstsplit}'s second integral does. A2c's index range is untouched.
What else enters clauses \textup{(i)}--\textup{(iii)}, exactly, and it is not nothing. Not C1 and
not A2, A2c does not presuppose that $\ell^\Pi_\star$ solves anything. But \textbf{B1(iv.a)--(iv.b)}
enter \emph{substantively}, at two points and in two different ways: \emph{finiteness} of $\Pi_u$ and
$\Pi_t$ makes the \textbf{two}-factor product $\Pi_u\otimes\Pi_t$ a product of $\sigma$-finite measures,
which is what supplies a product-measurable
representative and licenses Tonelli; and \emph{absolute continuity} of $\Pi_u$ with respect to Lebesgue
measure is what turns a Lebesgue-null modification of a slice into a $\Pi$-null modification, so that the
function clause \textup{(i)} constructs is a \emph{version}.
The third factor of $\Pi$ is reached by \textbf{neither} clause, and the route that leaves it
outside is the more general one. It enters only through \cref{def:spaces}'s declaration that
$\{1,\dots,\mathsf L-1\}$ carries its power set, which makes $\Pi$-measurability equivalent to
measurability of the $\mathsf L-1$ sections; the measure theory therefore runs on
$\Pi_u\otimes\Pi_t$ and the sections are assembled at the end, with the zero function at any $l$ of
zero $\Pi_{\mathsf w}$-mass, where the section is unconstrained, a case B1(iv.d) happens to exclude but
that the argument never needs excluded. \emph{So \textbf{B1(iv.d) does not enter clauses
\textup{(i)}--\textup{(iii)}}}, and that is a
consequence of the route and not an oversight: the three-factor alternative, $\Pi$ itself a product of
finite measures, would consume it, B1(iv.d) being the only clause that makes $\Pi_{\mathsf w}$ a
finite measure. \emph{The scope qualifier is load-bearing and is new on 2026-08-24}: the clause
\textbf{does} enter the concluding sentence, through \cref{cor:elldiff}, whose hypothesis line carries
it from that date. This paragraph printed the claim unqualified while that line was silent, and
unqualified it would now be false of the lemma taken whole --- the reading a reader of a printed
\emph{exact} accounting is entitled to.
In clauses \textup{(i)}--\textup{(iii)} B1(ii) enters only in that it \emph{defines}
$\bar w$, the symbol A2c is written in, its nonnegativity being part of that definition; and B1(i), (ii),
(iv.a), (iv.d) are what the concluding sentence consumes, \emph{but it consumes them in \textbf{two} ways, and
only one of the two is through \cref{cor:elldiff}}. B1(ii)'s envelope $|w(t'\Phi(z))|\le\bar w(t)$
dominates
\eqref{eq:Wdef}'s integrand \emph{directly}, which is what makes $\ell_\star:=\mathcal W\ell^\Pi_\star$
defined pointwise, A2b's last clause, and hence what entitles \cref{cor:elldiff}, stated at an
admissible representer, to be applied at all; the corollary then supplies the two sup-norm bounds under
B1(i), (ii), (iv.a) and (iv.d) jointly. B1(i) and B1(iv.d) enter there and nowhere else in this lemma;
B1(iv.a) enters
there and, for the different reason above, at the two earlier points.
% ARCHAEOLOGY.  An earlier form of the B1(iv.a)--(iv.b) sentence of this paragraph
% claimed "no clause of B1"
% and was too strong: sigma-finiteness is a consequence of B1(iv.a)--(iv.b) and
% not a definition.
% PHASE 4c, 2026-08-10.  THREE CORRECTIONS TO THIS ACCOUNTING, ALL OF THEM CASES
% OF A PRINTED CLAIM THE PROOF CONTRADICTS.
%   (1) "makes Pi a product of sigma-finite measures" OVER-CREDITED B1(iv.a)-(iv.b):
%       Pi_w's sigma-finiteness is not theirs.  Step 0 of the proof runs on the
%       TWO-factor product and reaches the three-fold base by the SECTION CRITERION,
%       which is more general -- it carries zero-mass cells that the three-factor
%       retype would have to assume away, and it is what keeps B1(iv.d) out of this
%       lemma's accounting.  The accounting now names the route the proof takes.
%   (2) "through cor:elldiff" WAS WRONG FOR B1(ii): Step 5 uses the envelope
%       DIRECTLY, to dominate eq:Wdef's integrand; cor:elldiff is invoked separately
%       and after.  In a lemma whose statement prints an EXACT accounting, "through
%       X" where the use is direct is precisely what that accounting exists to rule
%       out.
%   (3) The dangling "Then" that closed this paragraph is the lemma's connector to
%       its conclusion, separated from the clause list by two paragraphs.  It is
%       not deleted -- it is MOVED to the line before the list opens, where it
%       reads as the connector it is.  (The environment name is not spelled here:
%       a begin/end token inside a comment trips balance counters that do not
%       strip comments, and this corpus is checked by one.)
The contrast index is carried \textbf{explicitly} in every clause below, and it is quantified
differently in different ones, which is why the two indices cannot be described by one rule. The $p$
coordinates of $\theta$ are carried \emph{inside} a Euclidean norm, the geometry A2c and
\eqref{eq:lambdamu} are stated in, and one that no display below ever opens. The $\mathsf L-1$ contrasts are
carried as a \emph{free} index: clauses \textup{(i)} and \textup{(ii)}, and the measurability half of
\textup{(iii)}, hold at \emph{every} $l$, and \textup{(i)}'s joint measurability is asserted on the
three-fold base $\IR\times\mathcal T\times\{1,\dots,\mathsf L-1\}$, so it is one statement about all of
them and not $\mathsf L-1$ statements, whereas the three budgets $\Lambda^{(0)}_\star$,
$\Lambda^{(1)}_\star$ and $\Lambda^{(\mu'')}_\star$ that \eqref{eq:l21budget} relates are \emph{maxima}
over $l$ of single-contrast integrals, by A$2'$(i.a)--(i.b) and \eqref{eq:lambdamu}. \emph{One Borel set
$\mathcal T_\ast$ serves every $l$, as A2c requires}, which is what lets the version be constructed by
one argument rather than by $\mathsf L-1$ of them.
\emph{At $\mathsf L=2$ the index is single-valued and at $p=1$ the Euclidean norm is the
modulus, so every display below reduces to the one earlier rounds printed, with \textbf{no} exception}:
under those two identifications, the third argument $l$ dropped, $\|\cdot\|$ read as $|\cdot|$ and
$H^{\mu''}(\IR;\IR^{p})$ as $H^{\mu''}(\IR)$, \eqref{eq:embed}, and \eqref{eq:l21budget} read at its
\emph{outer} members, are those displays character for character.
\emph{And the two sharper inner members of \eqref{eq:l21budget} are sharper at \textbf{every}
$\mathsf L$, $\mathsf L=2$ included}, they are an improvement in the constant and not an
$\mathsf L$-dependence, so nothing here is an $\mathsf L=2$ exception in either direction.
Then, at that $\mu''$:
\begin{enumerate}
\item[(i)] \emph{The version.} $\ell^\Pi_\star$ admits a jointly measurable version, given slicewise by the
\emph{unique continuous representative} of the slice at every $t\in\mathcal T_\ast$, the Borel
full-$\Pi_t$-measure set A2c attaches to its representative, so the version is determined there and not
merely selected, and by $0$ off $\mathcal T_\ast$; for every $t\in\mathcal T_\ast$ and every $l$ the
slice $u\mapsto\ell^\Pi_\star(u,t,l)$ is continuously differentiable on all of $\IR$ with
$\ell^\Pi_\star(\cdot,t,l)$ and $\partial_u\ell^\Pi_\star(\cdot,t,l)$ bounded and vanishing as
$|u|\to\infty$; and both $(u,t,l)\mapsto\ell^\Pi_\star(u,t,l)$ and
$(u,t,l)\mapsto\partial_u\ell^\Pi_\star(u,t,l)$ are jointly measurable on
$\IR\times\mathcal T\times\{1,\dots,\mathsf L-1\}$.
\item[(ii)] The embedding, slicewise, and it holds \textbf{separately} in the level and in the
derivative. For every $t\in\mathcal T_\ast$ and every $l$,
\begin{equation}\label{eq:embedsplit}
\begin{split}
\sup_{u\in\IR}\big\|\ell^\Pi_\star(u,t,l)\big\|\;&\le\;c^{(0)}_{\mu''}\,\big\|\ell^\Pi_\star(\cdot,t,l)\big\|_{H^{\mu''}(\IR;\IR^{p})},\\
\sup_{u\in\IR}\big\|\partial_u\ell^\Pi_\star(u,t,l)\big\|\;&\le\;c^{(1)}_{\mu''}\,\big\|\ell^\Pi_\star(\cdot,t,l)\big\|_{H^{\mu''}(\IR;\IR^{p})} ,
\end{split}
\end{equation}
and therefore, adding the two and applying \eqref{eq:embedconstid},
\begin{equation}\label{eq:embed}
\sup_{u\in\IR}\big\|\ell^\Pi_\star(u,t,l)\big\|\;+\;\sup_{u\in\IR}\big\|\partial_u\ell^\Pi_\star(u,t,l)\big\|
\;\le\;c_{\mu''}\,\big\|\ell^\Pi_\star(\cdot,t,l)\big\|_{H^{\mu''}(\IR;\IR^{p})} .
\end{equation}
\emph{\eqref{eq:embedsplit} is the load-bearing form and \eqref{eq:embed} is its corollary}, not the
reverse: clause \textup{(iii)} takes a maximum over $l$ in each of the two budgets \emph{before} adding
them, and only a bound that has already been separated survives that order of operations. The two are
$\sqrt2$ apart, \eqref{eq:embedsplit} summed gives the constant
$c^{(0)}_{\mu''}+c^{(1)}_{\mu''}\le c_{\mu''}/\sqrt2$, and \eqref{eq:embed} is printed at the weaker
$c_{\mu''}$ because that is the constant every consumer downstream reads.
\emph{Neither constant is inflated by $p$}: Step~1 of the proof is \emph{written} at the Euclidean norm
throughout, at $\IR^{p}$-valued arity and with no coordinate index fixed, the same Fourier
inversion bound, the same Cauchy--Schwarz inequality, now applied once to the level and once to the
derivative rather than once to their sum, so \eqref{eq:embedconst} and \eqref{eq:embedconstsplit}
serve the vector-valued scale verbatim. The sum over the $p$ coordinates sits inside the
definition of $\|\cdot\|_{H^{\mu''}(\IR;\IR^{p})}$, which is where Plancherel deposits it, and
\textbf{no step of that argument ever forms one}.
\item[(iii)] \emph{A$2'$(i.a)--(i.b), in sup-inside form.} Both maps
$t\mapsto\sup_u\|\ell^\Pi_\star(u,t,l)\|$ and $t\mapsto\sup_u\|\partial_u\ell^\Pi_\star(u,t,l)\|$ are
$\Pi_t$-measurable at every $l$, and
\begin{equation}\label{eq:l21budget}
\Lambda^{(0)}_\star+\Lambda^{(1)}_\star
\;\le\;\big(c^{(0)}_{\mu''}+c^{(1)}_{\mu''}\big)\Lambda^{(\mu'')}_\star
\;\le\;\frac{c_{\mu''}}{\sqrt2}\,\Lambda^{(\mu'')}_\star
\;\le\;c_{\mu''}\,\Lambda^{(\mu'')}_\star\;<\;\infty ,
\end{equation}
so A$2'$(i.a) and A$2'$(i.b) hold at that version, with the supremum \emph{inside} the
$\Pi_t$-integral in both.
\eqref{eq:l21budget} is stated at the \textbf{sum}, it holds at every $\mathsf L$ and every $p$,
and read at its outer members, left-hand side against $c_{\mu''}\Lambda^{(\mu'')}_\star<\infty$, 
it is the display earlier rounds printed, character for character. What is new are the two inner
members, and they are a \emph{strengthening} by a factor $\sqrt2$ that is available at
every $\mathsf L$, $\mathsf L=2$ included: nothing in \eqref{eq:l21budget} depends on $\mathsf L$, so
the binary case reduces with no exception.
The one thing $\mathsf L$ does change is the \textbf{route}, and it is worth saying which step
fails, because it is the step a reader will try. All three budgets are maxima over $l$ of
$\bar w$-weighted $\Pi_t$-integrals, write $\Lambda^{(0)}_l$, $\Lambda^{(1)}_l$ and
$\Lambda^{(\mu'')}_l$ for the integral at a single contrast $l$, so that
$\Lambda^{(0)}_\star=\max_l\Lambda^{(0)}_l$ and likewise for the other two, by A$2'$(i.a)--(i.b) and
\eqref{eq:lambdamu}, and the maximum does not distribute over a sum: integrating
\eqref{eq:embed} at a fixed $l$ and then maximising delivers
$\max_l\{\Lambda^{(0)}_l+\Lambda^{(1)}_l\}\le c_{\mu''}\Lambda^{(\mu'')}_\star$, which at $\mathsf L\ge3$
is \emph{weaker} than \eqref{eq:l21budget}, because $\Lambda^{(0)}_\star$ and $\Lambda^{(1)}_\star$ may
be attained at different contrasts and $\max_l a_l+\max_l b_l\ge\max_l\{a_l+b_l\}$. Chaining that route
over the two maximising contrasts recovers only the factor $2$. \emph{That is a defect of the combined
bound \eqref{eq:embed} and not of the statement}: the separated bound \eqref{eq:embedsplit} maximises
over $l$ in each budget \emph{first}, giving $\Lambda^{(0)}_\star\le c^{(0)}_{\mu''}\Lambda^{(\mu'')}_\star$
and $\Lambda^{(1)}_\star\le c^{(1)}_{\mu''}\Lambda^{(\mu'')}_\star$ separately, and the sum of those two
is \eqref{eq:l21budget}. At $\mathsf L=2$ the two routes coincide and the naive one is not wrong, merely
$\sqrt2$ short.
And the representer that would attain the factor $2$ does not exist, so that route was not merely
wasteful but empty. Attaining it asks for one contrast saturating \eqref{eq:embed} in the level and
another saturating it in the derivative; saturation at the first leaves nothing on the right of
\eqref{eq:embed} for its derivative, so it needs $\Lambda^{(1)}_l=0$ at a contrast with
$\Lambda^{(0)}_l>0$. That is unrealisable. $\Lambda^{(1)}_l=0$ integrates a nonnegative function to zero,
so $\sup_u\|\partial_u\ell^\Pi_\star(u,t,l)\|=0$ for $\Pi_t$-almost every $t$ with $\bar w(t)>0$; each
such slice is then constant in $u$, hence identically zero by the vanishing at $|u|\to\infty$ that clause
\textup{(i)} supplies; and the $t$ with $\bar w(t)=0$ contribute nothing to either budget. Hence
$\Lambda^{(0)}_l=0$ too. \emph{The factor $2$ is therefore not attained, and it is not the supremum
either}, \eqref{eq:l21budget} bounding the sum by $c_{\mu''}/\sqrt2$ throughout.
\emph{Every consumer downstream reads the outer form, at $c_{\mu''}$}, the concluding sentence below,
and \cref{cor:elldiff} through it, so the strengthening propagates
nowhere and nothing outside this lemma retypes.
\emph{The $p$-direction costs nothing at all}, \eqref{eq:embed} and \eqref{eq:embedsplit} being stated at
the Euclidean norm that \eqref{eq:lambdamu} is also stated at.
\end{enumerate}
\emph{Consequently \cref{cor:elldiff} applies at $\ell^\Pi=\ell^\Pi_\star$ under B1(i), (ii), (iv.a), (iv.d)},
delivering $\|\ell_\star\|_\infty\le\|\varpi_u\|_\infty\,c_{\mu''}\Lambda^{(\mu'')}_\star$ and
$\|\partial_u\ell_\star\|_\infty\le\big(\|\varpi_u'\|_\infty+\|\varpi_u\|_\infty\big)c_{\mu''}\Lambda^{(\mu'')}_\star$.
\emph{The target is A$2'$(i) and never A2b}: A2b is a predicate over generic representers, is never
asserted, and therefore has nothing to discharge.
\end{lemma}

\begin{remark}[Why A2c, and what the registered route would have cost]\label{rem:l21route}
Four things about \cref{lem:reprsmooth} that a reader will ask, in the order they arise.

\textup{(a)} Why the sup-inside form is the target, and why it is met rather than approximated.
A sup-\emph{outside} form does
not dominate: $\sup_u\int(\cdot)\,d\Pi_t<\infty$ gives $L^{1}(\Pi_t)$-boundedness in $u$ without uniform
integrability, and Leibniz's rule, which is what \eqref{eq:Wderiv} consumes, needs
$\int\sup_u(\cdot)\,d\Pi_t<\infty$. A2c meets the sup-inside target because it is stated in the
\emph{same} $\bar w$-weighted $L^{1}(\Pi_t)$ geometry: each half of \eqref{eq:embedsplit} is a bound
% ARCHAEOLOGY.  A2'(i) carried the supremum outside the Pi_t-integral through
% round 4; the repair moved it inside and raised lem:reprsmooth's target with it.
between two functions
of $t$, holding pointwise in $t$, so \eqref{eq:l21budget} is \emph{two} integrations of \emph{two}
pointwise inequalities, followed by a maximum over the finite contrast index in each and one addition,
and no supremum is ever moved across an integral at any of those steps. The asymmetry of
\cref{rem:supinside} is not reproduced here, because there is no interchange to license.

\textup{(b)} The obvious alternative route, and why it is three steps short. One might instead
require
$\ell_\star\in\mathcal R\big((\mathcal A_{\mathrm{cond}}\mathcal A^\ast_{\mathrm{cond}})^{\mu''}\big)$
followed by an embedding into $C^{1}$. \emph{First}, that names $\ell_\star$, the
observable-space object, whereas A$2'$(i) is stated on the range-space object
$\ell^\Pi_\star$; the bridge between them is \cref{cor:elldiff}, which \emph{consumes} A$2'$(i) and so
cannot be used to reach it. \emph{Second}, and decisively, membership in
$\mathcal R((\mathcal A_{\mathrm{cond}}\mathcal A^\ast_{\mathrm{cond}})^{\mu''})$ is a statement about the
singular system of $\mathcal A_{\mathrm{cond}}$ and implies \emph{nothing} about $\sup_u$ or about
$u$-differentiability unless that scale is \emph{matched} to a $u$-smoothness scale. H1 supplies exactly
such a matching on the domain side; there is no analogue on the range side, and supplying one would be a
\emph{second} Hilbert-scale link, an assumption about how the left singular functions $\tilde e_j$
oscillate in $u$, on an object the analyst does not choose. That is strictly larger than A2c and buys
nothing A2c does not. \emph{Third}, such a range condition says nothing about the $\Pi_t$-integration,
which is half of what (i.a) and (i.b) assert.

\textup{(c)} What A2c costs beyond A2, and it is not only smoothness. Membership
$\ell^\Pi_\star\in L^{2}(\Pi)$ gives, at every $l$ and for $\Pi_t$-almost every $t$, that
$\ell^\Pi_\star(\cdot,t,l)\in L^{2}(\Pi_u;\IR^{p})=L^{2}(\varpi_u\,du;\IR^{p})$, \emph{the $\mathsf L-1$ slices
being separately in $L^{2}(\Pi_u)$ because $\varpi_{\mathsf w}(l)>0$ by B1(iv.d), so no slice is free}.
B1(iv.a) makes $\varpi_u$ bounded and
strictly positive \emph{pointwise}, not bounded below, so $L^{2}(\IR,du)\subseteq L^{2}(\Pi_u)$ and
\emph{not} conversely: a slice may lie in $L^{2}(\Pi_u)$ with tails too fat to lie in $L^{2}(\IR,du)$.
A2c therefore imposes a \emph{tail-decay} condition in $|u|$ on top of smoothness in $u$. \emph{This is
stated rather than hidden}, it is the same species of restriction as C4$'$'s residual truncation, and it
is the mirror image of the asymmetry recorded at A2 for $\pi$, a one-sided density bound gives one
containment and not the other. Two sufficient conditions, in increasing strength: an
$L^{2}(\Pi_t)$-form,
\[
\Lambda^{(\mu'')}_\star\;\le\;\Big(\int_{\mathcal T}\bar w^{2}\,d\Pi_t\Big)^{1/2}
\max_{l\le\mathsf L-1}
\Big(\int_{\mathcal T_\ast}\big\|\ell^\Pi_\star(\cdot,t,l)\big\|^{2}_{H^{\mu''}(\IR;\IR^{p})}\Pi_t(dt)\Big)^{1/2}
\]
by the Cauchy--Schwarz inequality, the first factor being finite by B1(iv.c), \emph{it is
$\int\bar w^{2}d\Pi_t$ and not $c_\Pi$, which carries an extra $\Pi_u(\IR)$ that this inequality does not
produce}; and the uniform form
$\max_{l}\sup_{t\in\mathcal T_\ast}\|\ell^\Pi_\star(\cdot,t,l)\|_{H^{\mu''}(\IR;\IR^{p})}<\infty$, which with B1(v)'s
$\int\bar w\,d\Pi_t<\infty$ implies both. The weakest of the three is the one A2c states.
\emph{The maximum is over a finite index set and introduces no $\mathsf L$-dependent factor}, the
aggregation being the same maximum \eqref{eq:lambdamu} takes; and the $p$ coordinates sit inside
the norm at every one of the three, so none of them carries a $\sqrt p$ either.

\textup{(d)} The integer case needs no transform, and the two constants differ. \emph{Both routes
are written at an $\IR^{p}$-valued $g$, at the Euclidean norm throughout}, and neither constant acquires
a factor in $p$: the identity
$\|g\|^{2}_{H^{\mu''}(\IR;\IR^{p})}=\sum_j\|g_j\|^{2}_{H^{\mu''}(\IR)}$ is A2c's \emph{definition} of the
vector scale and not a step of either argument, the transform route forms no coordinatewise sum at
all, and the elementary route forms only the one that defines $\|g\|^{2}$. At $\mu''=2$,
$H^{2}(\IR)$ is $\{g:g,g',g''\in L^{2}(\IR,du)\}$ with
\begin{equation}\label{eq:h2norm}
\big\|g\big\|^{2}_{H^{2}(\IR)}\;=\;\big\|g\big\|^{2}_{L^{2}}+2\big\|g'\big\|^{2}_{L^{2}}+\big\|g''\big\|^{2}_{L^{2}} ,
\end{equation}
\emph{which is read at the vector scale by reading each $L^{2}$ norm as
$\big(\int\|\cdot\|^{2}du\big)^{1/2}$}, that being A2c's coordinate identity at $\mu''=2$;
and \eqref{eq:embed} holds
with $c_{\mu''}$ replaced by $2$ by the elementary argument of Step 1$'$ of the proof, no Fourier
transform, no complex numbers. \emph{The factor $2$ on the middle term is not decoration and is the one
thing an applied reader will get wrong}: \eqref{eq:usobolev}'s weight at $\mu''=2$ is
$(1+\xi^{2})^{2}=1+2\xi^{2}+\xi^{4}$, and Plancherel in that same normalisation sends $1$, $\xi^{2}$,
$\xi^{4}$ to $\|g\|^{2}_{L^{2}}$, $\|g'\|^{2}_{L^{2}}$, $\|g''\|^{2}_{L^{2}}$, so the cross term of the
square survives as a $2$. \emph{The set $\{g:g,g',g''\in L^{2}\}$ is right either way}, the two
quadratic forms are equivalent, with ratio in $[1,2]$, so printing the sum without the $2$
names the right space and the wrong norm. \emph{And the elementary route's $2$ is a
weaker constant than \eqref{eq:embedconst} gives at the
% ARCHAEOLOGY.  An earlier form of the sentence above printed the H^2 norm without
% the factor 2 on the middle term.  Both constants below are computed from
% eq:h2norm and neither changes.
same index}: $\int_{\IR}(1+\xi^{2})^{-1}d\xi=\pi$, so $c_{2}=\sqrt2<2$. Both are correct and the
difference is the price of the transform-free route; the sharper one is the one \eqref{eq:embedconst}
defines, and \eqref{eq:l21budget}'s last member is stated with it. The separated constants sharpen
the comparison once more, and the arithmetic is worth seeing at one index. At $\mu''=2$,
$\int_{\IR}(1+\xi^{2})^{-2}d\xi=\int_{\IR}\xi^{2}(1+\xi^{2})^{-2}d\xi=\tfrac\pi2$, so
$c^{(0)}_{2}=c^{(1)}_{2}=\tfrac12$ and $c^{(0)}_{2}+c^{(1)}_{2}=1$; \eqref{eq:embedconstid} is tight at
that index, the two integrals being equal there. Step 1$'$ gives the same two constants elementarily at
the value $1$ each, hence $2$ for the sum. A reader who prefers to check A2c the elementary way therefore
loses the range $\mu''\in(\tfrac32,2)$ and a constant: a factor $\sqrt2$ if the comparison is against
\eqref{eq:embedconst}'s combined bound, and a factor $2$ in each budget separately if it is against
\eqref{eq:embedconstsplit}.
\emph{And the threshold is binding, not a margin}: at $\mu''=\tfrac32$ the integral in
\eqref{eq:embedconst} diverges, the integrand behaving like $|\xi|^{-1}$, so $c_{\mu''}$ is undefined
there and Step 1 has no dominating function. The extra $1$ over the $L^{\infty}$ threshold $\mu''>\tfrac12$
is exactly what clause (i.b), the \emph{derivative} clause, costs over clause (i.a),
\emph{and \eqref{eq:embedconstsplit} is where that statement becomes an arithmetic one rather than an
interpretation}: $c^{(0)}_{\mu''}$ is finite for $\mu''>\tfrac12$ and $c^{(1)}_{\mu''}$ only for
$\mu''>\tfrac32$, and it is the second that binds in $c_{\mu''}$, whose integrand is the sum of the
other two. \emph{So separating the constants exhibits a threshold \eqref{eq:embedconst} already carried
and imposes no new one}, and A2c's index range is exactly what it was.
\emph{And $\mu''$ is independent of $\mu$}: $\mu=s/(2a)$ indexes
smoothness of $\vdagger$ in $x$; $\mu''$ indexes smoothness of $\ell^\Pi_\star$ in $u$. Nothing links
them, which is why the smoothness-budget table lists them on separate lines and why the representer line
competes for no part of the $s$ budget.

\emph{What remains an assumption.} A2c is a condition on $\ell^\Pi_\star$, an object the model and the
selection determine jointly; \cref{lem:reprsmooth} does not derive it from anything more primitive and
does not claim to. What it buys is that \emph{one} condition, on \emph{one} object, in a language the
block already uses, replaces two sup-norm primitives with no route to them, and that the condition is
\emph{reportable}: by \eqref{eq:usobolev} it is a weighted summability condition on the $u$-frequency
content of the slices, the same kind of object as A2's Picard sum.
\end{remark}

\begin{lemma}[L-20: structure of the second-order form]\label{lem:formstruct}
The symmetrised second-order form is a bounded multiplication operator perturbed by a compact one; its
essential spectrum is the essential range of the multiplier, it is compact exactly when that multiplier
vanishes, and it is indefinite whenever it does not.
Suppose the hypotheses of \cref{lem:solset} hold \emph{together with} C2$'$(i), \emph{which
\cref{lem:solset} does not carry, and which is consumed by clause \textup{(ii)} alone}, transitively
through \cref{prp:compact}, for the compactness of $\mathcal A_{\mathrm{cond}}$ that delivers
$\|\mathcal A_{\mathrm{cond}}h_k\|\to0$; clauses \textup{(i)} and \textup{(iii)}--\textup{(v)} need it
not. Fix $\vartheta\in\IR^{p}$ with $\|\vartheta\|=1$, and
let $\mathcal B$, $\mathcal B_{\mathrm{mult}}$ and
$\mathcal B_{\mathrm{sym}}:=\tfrac12(\mathcal B+\mathcal B^{\ast})$ be the operators on $L^{2}(\pi)$ of
\eqref{eq:Bpi}, \eqref{eq:Bmult} and \cref{lem:solset}(iii). \emph{The multiplier is written out as
$\vartheta'\sigma_q/\pi$ throughout and is given no abbreviation}: $\sigma_q$ is the multiplier,
$\mathcal B_{\mathrm{mult}}$ is multiplication \emph{by} $\vartheta'\sigma_q/\pi$, and an abbreviation
would collapse the distinction. Then
\begin{enumerate}
\item[(i)] \emph{Decomposition and boundedness.}
$\mathcal B=\mathcal B_{\mathrm{mult}}-2\,\vartheta'\tilde{\mathcal K}$ and
$\mathcal B_{\mathrm{sym}}=\mathcal B_{\mathrm{mult}}-\vartheta'(\tilde{\mathcal K}+\tilde{\mathcal K}^{\ast})$;
$\mathcal B_{\mathrm{mult}}$ is self-adjoint with
$\|\mathcal B_{\mathrm{mult}}\|=\|\vartheta'\sigma_q/\pi\|_\infty\le2\,c_{\widetilde W}\,\bar c_\pi\,\|\partial_u\ell_\star\|_\infty<\infty$
, \emph{\cref{lem:solset}\textup{(iii)}'s bound, carrying $c_{\widetilde W}$ for the contraction
of $\partial_u\ell_\star$ against $\widetilde W$ and equal to the printed constant at $\mathsf L=2$}, 
the essential supremum being with respect to $\pi\,dx$ and hence, by C1(ii), with respect to Lebesgue
measure on $\mathcal X$; $\vartheta'(\tilde{\mathcal K}+\tilde{\mathcal K}^{\ast})$ is compact and
self-adjoint; and $\mathcal B$ is bounded but \emph{not} self-adjoint.
\item[(ii)] The singular sequence, and it carries \textup{(iii)}--\textup{(v)} and
\cref{prp:sharpquad}\textup{(ii)},\textup{(iv)}. For every $\varepsilon>0$ and every Borel
$\mathcal X_0\subseteq\mathcal X$ of positive Lebesgue measure on which $\vartheta'\sigma_q/\pi$ is of one
sign with $|\vartheta'\sigma_q/\pi|\ge\mathsf c_0-\varepsilon$, where
$\mathsf c_0:=\operatorname*{ess\,sup}_{\mathcal X_0}|\vartheta'\sigma_q/\pi|$, there is a sequence
$(h_k)_{k\ge1}\subset\mathcal V$ with $\|h_k\|_{L^{2}(\pi)}=1$, $h_k\rightharpoonup0$ weakly in
$L^{2}(\pi)$, $\|\mathcal A_{\mathrm{cond}}h_k\|_{L^{2}(\Pi)}\to0$, and
\begin{equation}\label{eq:singseq}
\lim_{k\to\infty}\big\langle\mathcal B_{\mathrm{sym}}h_k,h_k\big\rangle_{L^{2}(\pi)}
\;=\;\lim_{k\to\infty}\int_{\mathcal X_0}\big(\vartheta'\sigma_q/\pi\big)\,h_k^{2}\,\pi\,dx ,
\end{equation}
a limit whose modulus is at least $\mathsf c_0-\varepsilon$ and whose sign is that of
$\vartheta'\sigma_q$ on $\mathcal X_0$, along a subsequence. \emph{Such an $\mathcal X_0$ exists whenever
$\vartheta'\sigma_q\not\equiv0$, and the proof constructs one with
$\mathsf c_0\ge\|\vartheta'\sigma_q/\pi\|_\infty-\varepsilon$}, \emph{by taking the near-supremum set
first and splitting it by sign second, an order that cannot be reversed}.
\item[(iii)] \emph{Essential spectrum.}
$\sigma_{\mathrm{ess}}(\mathcal B_{\mathrm{sym}})
=\sigma_{\mathrm{ess}}(\mathcal B_{\mathrm{mult}})=\operatorname{ess\,ran}(\vartheta'\sigma_q/\pi)$, and
the same holds for the compression
$P_{\mathcal V}\mathcal B_{\mathrm{sym}}P_{\mathcal V}|_{\mathcal V}$; consequently
\begin{equation}\label{eq:complower}
\big\|P_{\mathcal V}\mathcal B_{\mathrm{sym}}P_{\mathcal V}|_{\mathcal V}\big\|
\;\ge\;\big\|\vartheta'\sigma_q/\pi\big\|_\infty .
\end{equation}
\item[(iv)] \emph{Compactness dichotomy.} $\mathcal B_{\mathrm{sym}}$ is compact if and only if
$\vartheta'\sigma_q=0$ almost everywhere; and the same dichotomy holds for the compression. In particular
$\mathcal B_{\mathrm{sym}}$ is non-compact whenever $\vartheta'\sigma_q\not\equiv0$.
\item[(v)] \emph{Indefiniteness.} $\int_{\mathcal X}\sigma_q\,dx=0$, so if $\vartheta'\sigma_q\not\equiv0$
then $\vartheta'\sigma_q/\pi$ is strictly positive on a set of positive measure and strictly negative on
another; hence $\sigma_{\mathrm{ess}}(\mathcal B_{\mathrm{sym}})$ contains points of both signs, and
neither $\mathcal B_{\mathrm{sym}}$ nor its compression is positive or negative semidefinite.
\end{enumerate}
\end{lemma}

% EXTRACTED FROM THE STATEMENT OF prp:sharpquad on 2026-08-02.  No \label moved:
% eq:quadrem, eq:sharpconst, eq:bracket and eq:noweak all stayed in the statement.
\begin{remark}[How the hypothesis list of \cref{prp:sharpquad} is read]\label{rem:sharpquadread}
\begin{enumerate}
\item[(1)] \Cref{lem:solset} is cited by clause and not as a block. Clause \textup{(iv)} in
addition requires \textbf{C2$'$(i)}, which it consumes transitively through
\cref{lem:formstruct}\textup{(ii)} and \cref{prp:compact}, for the compactness of
$\mathcal A_{\mathrm{cond}}$ that delivers $\|\mathcal A_{\mathrm{cond}}h_k\|\to0$; \emph{clauses
\textup{(i)}--\textup{(iii)} need it not}, and Step 5 of the proof runs the implication
O-16\textup{(b)}$\Rightarrow$O-16\textup{(a)} through \eqref{eq:complower} and the numerical radius
precisely so that they need not.
\item[(2)] \emph{\Cref{lem:secondord} is read at the admissible direction pair $(D,\mathsf r)=(0,0)$},
which is admissible at every $(h,\bar\tau)$ by \cref{def:dirclass} and is differentiable there, so that
$\Psi(\cdot,0,0)=\Psi_{L_\star}$ by \eqref{eq:PsiL} and \eqref{eq:quadrem} is half the second
$\tau_1$-derivative of the moment along the \emph{structural direction alone}; no condition on a
direction in $L$ or in $\varrho$ is inherited.
\item[(3)] \emph{The remainder carries no symbol of its own}; italic $Q$ names the regular conditional
distributions of \textup{(P6b)}.
\item[(4)] \emph{Clause \textup{(i)}'s sharp constant carries no hypothesis on $q$ beyond
\cref{lem:solset}'s, and clause \textup{(ii)} brackets that constant without pinning it}: the left-hand
member of \eqref{eq:bracket} is a lower bound and is not the sharp constant; the right-hand members are
upper bounds and are not the sharp constant either.
\end{enumerate}
\end{remark}

\begin{remark}[What \cref{prp:sharpquad} claims, and the four things it does not]\label{rem:sharpscope}
The quantifier is a supremum over a ball, and the whole result lives there.
\eqref{eq:sharpconst} is a statement about $\sup_{\|h\|\le\delta}$, hence about the \emph{best bound
valid uniformly on the ball}. Four readings must be blocked, and each has a version that is true.

\textup{(1)} It is not a claim about the estimator. \Cref{prp:sharpquad} does \emph{not} say that
the realised remainder at $h=\hat\varphi-\vdagger$ fails to be $o_P(\|h\|^{2}_{L^{2}(\pi)})$. A
regularisation path confined to directions on which the form happens to be small could do better, and
nothing here excludes that. \emph{What is true:} the direct bound
$\tfrac12|\langle\mathcal Bh,h\rangle|\le\tfrac12\|P_{\mathcal V}\mathcal B_{\mathrm{sym}}P_{\mathcal V}|_{\mathcal V}\|\,\|h\|^{2}$
is \emph{order-sharp}, so A5(ii)'s $\delta_n=o_P(n^{-1/4})$ cannot be relaxed \emph{by finding a better
bound on the remainder}. Relaxing it requires either a different moment, a selection with a smaller
$\sigma_q$ and $\mathcal K$, which is lever~1, or a restriction on the \emph{direction} of the
first-stage error, which is an extra hypothesis and not a sharper argument. The two are different claims
and the paper should print both.

\textup{(2)} The lower bound is not the sharp constant, and the identity
$\sup_{\|h\|\le\delta}\tfrac12|\langle\mathcal Bh,h\rangle|=\tfrac12\delta^{2}\|\mathcal B_{\mathrm{sym}}\|$
holds on the wrong ball. It is exact when the
supremum runs over the ball of $L^{2}(\pi)$ and is an \emph{upper bound} when it runs over the ball of
$\mathcal V$, which is the operative one, since $h=\hat\varphi-\vdagger$ lies in $\mathcal V$ by
\cref{def:spaces}. The exact constant on $\mathcal V$ is the compression's norm, \cref{def:spaces}'s
one declared exception is what forces the distinction, and \eqref{eq:bracket} brackets it. \emph{No
downstream consequence changes}, because what \cref{thm:rootn} consumes is the lower bound.
\emph{And $\tfrac12\|\vartheta'\sigma_q/\pi\|_\infty$ must not be printed as the sharp constant}: it is a
lower bound, delivered by the multiplication part through \cref{lem:formstruct}(iii).

\textup{(3)} $\sigma_q\equiv0$ does not make the remainder vanish, and this is what separates
O-16\textup{(a)} from O-16\textup{(b)}. \Cref{lem:solset}(iii)'s witness exhibits a $q\not\equiv0$ with
$\sigma_q\equiv0$ whose form is nonzero: $\sigma_q$ is a marginalisation and the form sees the
\emph{symmetric part} of $q\,f_{X_1,X_2}$, which a marginalisation does not determine. Consequently, when
$\sigma_q\equiv0$ but O-16(a) holds, clauses (i) and (iii) survive \emph{with the same $\delta^{2}$
order}, while three things change and they change in three different ways: \eqref{eq:complower} becomes
\emph{vacuous}, its right-hand side being zero; the non-compactness of \cref{lem:formstruct}(iv)
\emph{fails}, the form then being compact; and the weak-norm refutation (iv) becomes
\emph{unavailable}, \emph{which is not the same as a weak-norm bound existing}. Whether one exists is
then open, and it is what the Douglas route of the strategy memo \S3.6 would have to settle.
\emph{So lever~1's target $\sigma_q\equiv0$ removes the essential spectrum and does not remove the
$\delta^{2}$ order}, and SG-15, which sweeps $\|\vartheta'\sigma_q/\pi\|_\infty$, minimises a
certified \emph{lower bound} on the sharp constant and not the sharp constant. \emph{The objective is
not even monotone in the quantity of interest, and that is stronger than saying it is a lower bound}:
the compact part $\vartheta'(\tilde{\mathcal K}+\tilde{\mathcal K}^{\ast})$ is \emph{unconstrained} by
$\sigma_q$, so a selection that lowers $\|\vartheta'\sigma_q/\pi\|_\infty$ may \emph{raise}
$\|P_{\mathcal V}\mathcal B_{\mathrm{sym}}P_{\mathcal V}|_{\mathcal V}\|$, and driving the sweep to zero
makes \eqref{eq:complower} vacuous while leaving the sharp constant free. \emph{Any claim that SG-15
improves the constant therefore requires the compression's norm computed directly}, which
\eqref{eq:bracket} brackets but does not deliver. The target that removes the remainder
outright is $q\equiv0$, or a cancellation in the symmetric part of $q\,f_{X_1,X_2}$; probe CP-12 is aimed
at exactly this, and it is the only probe that could reopen the scope. \emph{And $q\not\equiv0$ alone is
not enough for O-16(a) either}: if $q\,f_{X_1,X_2}$ is antisymmetric the form vanishes identically, since
$(\Delta h)^{2}$ is symmetric in $(x_1,x_2)$, which is why (a) is stated as ``the form does not
vanish'' and not as ``$q\not\equiv0$''.

\textup{(4)} O-16 is a hypothesis and not a consequence of O-11. O-11 delivers
$\partial_u\ell_\star\not\equiv0$: a debiasing weight must vary with the residual. It does \emph{not}
deliver $q\not\equiv0$, because $q=E[\partial_u\ell_\star\,\widetilde W\mid X_1,X_2]$ can vanish by
cancellation inside the conditional expectation, still less $\sigma_q\not\equiv0$, which is a second
marginalisation. Both O-16 clauses are therefore carried as explicit hypotheses of
\cref{prp:sharpquad}\textup{(iii)}--\textup{(iv)} and are asserted nowhere. The chain reads: O-11 says
the quadratic term must \emph{exist}; O-16 says it does not degenerate; \cref{prp:sharpquad} says that if
it does not degenerate then its order cannot be improved.
\end{remark}



% ############################################################################
% ##  PART H -- memo Sec.10 STEP 6.  Merged by the orchestrator 2026-07-29.
% ##  DISPATCH-7 FIX PASS APPLIED, 2026-07-29, against proofs_review_dispatch7.md
% ##  (58/100, strike 1 of 3).  archive/results_partH.tex is retained as the record
% ##  ONLY; its header's requests C(1)-C(4) and C(6) are now APPLIED in
% ##  assumptions.tex, and C(5) is HELD one round by the critic's ruling.
% ##  WHAT MOVED OUT OF THIS FILE: def:intker and the definition of T-hat
% ##  (eq:hatT) now live in assumptions.tex Part V, beside A7, because A7 is a
% ##  statement about T-hat and a maintained assumption must not name an
% ##  undefined map.  lem:intker stays here, where its four properties are
% ##  derived from C3.
% ############################################################################

% ===========================================================================

% ---------------------------------------------------------------------------
%  H.1  (O4), EXHIBITION HALF
% ---------------------------------------------------------------------------

% ############################################################################
% ##  lem:bdryfam IS NEW IN DISPATCH 26 AND IT IS THE DISCHARGE OF C3(ii).
% ##  This file is position 7 of 8, so every theorem-like number in the rest of
% ##  this file and in proofs.tex ADVANCES BY ONE.  Nothing breaks -- every
% ##  reference is a \cref or a \suppref -- but printed numbers move.
% ##
% ##  WHY IT IS OWED, AND OWED UNDER EITHER KERNEL ROUTE.  C3(ii) ASSERTS a
% ##  boundary family with a uniform envelope and uniform p_K-moments.  Through
% ##  dispatch 25 NO FILE EXHIBITED ONE.  That is not a defect the envelope route
% ##  introduced: it is equally missing for a compactly supported kernel.  So the
% ##  lemma has to be written whichever kernel the code ships, and writing it
% ##  under the envelope has the bonus that it covers the kernel the code DOES
% ##  ship, k_order(., 8).
% ##
% ##  THE WEIGHT IS phi AND NOT Theta_1.  See the dispatch-26 block at C3 in
% ##  assumptions_main.tex for the full reason; in one line, the Theta_1-weighted
% ##  Gram matrix is a SIGNED Hankel matrix whose form is indefinite, so clause
% ##  (a) is unavailable for it and no floor on |I| repairs that, whereas with the
% ##  Gaussian weight (a) is two lines and the interior member comes out EXACTLY
% ##  the shipped kernel by uniqueness.
% ##
% ##  THE GAUSSIAN WEIGHT CARRIES NO SYMBOL and is written out as
% ##  e^{-varsigma^2/2}, on the C12/C24/C31/C35/C38/C44/C64/C78 write-it-out
% ##  precedent.  phi is NOT used: it is live in this corpus in two readings
% ##  already -- the standard normal density in the Sec.2 counterexample and a
% ##  generic convex/test function in lem:condjensen and prp:bierensid -- both of
% ##  them SEC.1.6 LOCALS bound by their own quantifiers, and promoting a local to
% ##  a binding symbol is the C25 defect.  Dropping the (2 pi)^{-1/2} costs
% ##  nothing: positive-definiteness and the envelope do not see it, and (e)
% ##  carries it explicitly.
% ############################################################################

\begin{lemma}[The boundary family exists, and its interior member is the shipped kernel]\label{lem:bdryfam}
Let $p_K\ge1$ and let $\mathfrak I$ be as in C3. For $I\in\mathfrak I$ put
% LABEL RENAMED 2026-08-15: this was eq:gram, which COLLIDED with prp:tccrepair's
% sample Gram matrix bold-G_n at results_supp.tex:2983 -- a live label with five
% consumers (proofs.tex :6686/:6700, assumptions_supp.tex:2248, results_supp.tex:4067,
% paper/qmd/supplement.qmd:296).  The THEORY BUILD DID NOT CATCH IT and the
% SUPPLEMENT RENDER DID, which is why both are run: theory-corpus "multiply
% defined: none" is not a substitute for the shipped document's own check.
\begin{equation}\label{eq:gramfam}
\mathsf G(I)\;:=\;\Big[\int_I\varsigma^{\,j+k}e^{-\varsigma^{2}/2}\,d\varsigma\Big]_{j,k=0}^{p_K-1},
\qquad
\Theta_I(\varsigma)\;:=\;\Big(\sum_{k=0}^{p_K-1}\big(\mathsf G(I)^{-1}e_0\big)_k\,\varsigma^{k}\Big)
e^{-\varsigma^{2}/2}\,\un\{\varsigma\in I\},
\end{equation}
$e_0$ the first standard basis vector of $\IR^{p_K}$. \emph{Fix the topology on $\mathfrak I$ here,
because \textup{(d)}, \textup{(f)} and \textup{(g)} are cited from other files and a continuity claim
without a topology is not a claim}: $\mathfrak I$ carries \textbf{endpointwise convergence in the
two-point compactification}, the topology induced by
\begin{equation}\label{eq:ifraktop}
\mathfrak I\;\ni\;I\;\longmapsto\;\big(\tanh\inf I,\ \tanh\sup I\big)\;\in\;[-1,0]\times[0,1] .
\end{equation}
\textup{(c)} below shows this map to be injective with closed image, so that $\mathfrak I$ is compact
and metrisable; ``continuous on $\mathfrak I$'' below means continuous for \eqref{eq:ifraktop}.
\emph{$\tanh$ carries no content}: it is one homeomorphism $[-\infty,+\infty]\to[-1,1]$ and any other
returns the same topology, which is why the clauses below are stated in $I$ and not in coordinates.
Then:
\begin{enumerate}
\item[(a)] $\mathsf G(I)\succ0$ for every $I\in\mathfrak I$, so \eqref{eq:gramfam} is well posed.
\item[(b)] $\Theta_I$ vanishes off $I$ and $\int_I\varsigma^{\,j}\Theta_I(\varsigma)\,d\varsigma=\delta_{j0}$
for $0\le j\le p_K-1$; that is, \eqref{eq:truncmom} holds.
\item[(c)] The map \eqref{eq:ifraktop} is injective with closed image, so $\mathfrak I$ is compact and
metrisable for it; and $\underline g:=\inf_{I\in\mathfrak I}\lambda_{\min}\big(\mathsf G(I)\big)>0$.
\item[(d)] Put
\begin{equation}\label{eq:envfam}
\bar\Theta(\varsigma)\;:=\;\sup_{I\in\mathfrak I}\big|\Theta_I(\varsigma)\big| ,
\qquad\text{for which}\qquad
\bar\Theta(\varsigma)\;\le\;\underline g^{-1}\Big(\sum_{k=0}^{p_K-1}\varsigma^{2k}\Big)^{1/2}e^{-\varsigma^{2}/2} .
\end{equation}
Then $|\Theta_I|\le\bar\Theta$ for every $I\in\mathfrak I$ by construction, and $\bar\Theta$ is
continuous --- hence measurable --- bounded and integrable, with $\int_\IR|\varsigma|^{m}\bar\Theta<\infty$
for every $m\ge0$. Hence C3(ii) holds, with $\mathsf c_2$ and $c_C$ the constants \eqref{eq:envconst}
derives from this $\bar\Theta$. Moreover $\bar\Theta$ is \textbf{minimal}: every $\bar\Theta'$ meeting
C3(ii) satisfies $\bar\Theta'\ge\bar\Theta$ pointwise, hence $\|\bar\Theta'\|_\infty\ge\mathsf c_2$ and
$\|\bar\Theta'\|_{L^{1}}\ge2c_C\mathsf c_2$ --- \emph{no admissible envelope returns a smaller pair}, so
\eqref{eq:normconst} and \eqref{eq:tccrad} are at their sharpest here. \emph{The second half of \eqref{eq:envfam} is the
only role the explicit majorant plays \textbf{in any proof}}: a fixed polynomial times a Gaussian, it is
bounded, integrable, of finite moment of every order \emph{and vanishing at infinity}, and $\bar\Theta$
inherits all four from it --- the fourth is not decoration, it is what \textup{(g)} turns into
attainment. \emph{So $\underline g$ reaches every \textbf{consumed} constant only through the finiteness
and the decay it certifies}: no bound, no rate and no constant that any result below reads is computed
from it.
\emph{It does have a second role, that role is expository and not inferential, and the numbers it carries
there are printed at three places --- the sentence this one replaces denied that, and was false at all
three.} The majorant is the \textbf{comparison} envelope, the closed form an analyst writes down first,
against which C3's note and the note below weigh the minimal one. \eqref{eq:envconst} evaluated at it
returns the pair $(4.50\times10^{8},\,1.82)$, printed in C3's note; \emph{of that pair the first entry is
computed from $\underline g=6.59\times10^{-8}$ and the second is \textbf{not}}, $\underline g^{-1}$
cancelling in the ratio that defines $c_C$. The note below prints a third number,
$1.97\times10^{6}$: the same sup-norm formed with the realised family's infimum $1.51\times10^{-5}$ in
place of $\underline g$. \emph{All three are quoted and consumed nowhere}, and the comparison is the whole
of their work --- it is what makes ``take the envelope minimal'' a ruling with a stated price rather than
a preference, and \textup{(g)} now supplies the third term of that comparison, a \emph{proved} pair
between the two.
\item[(e)] At $p_K=8$, $\Theta_{\IR}$ is \emph{exactly} the order-$8$ Gaussian-based kernel
\[
\Theta_1(\varsigma)\;=\;\frac{105-105\varsigma^{2}+21\varsigma^{4}-\varsigma^{6}}{48\sqrt{2\pi}}\;e^{-\varsigma^{2}/2} ,
\]
which is symmetric, $C^\infty$ and Lipschitz; so C3(i) and C3(iii) hold at the same family, and the
interior member is the kernel the application computes.
\item[(f)] $\Theta_{-I}(\varsigma)=\Theta_I(-\varsigma)$, and $I\mapsto\mathsf G(I)^{-1}e_0$ is
continuous on $\mathfrak I$ for the topology \eqref{eq:ifraktop}. \emph{It is the \textbf{coefficient}
map that is continuous and \textbf{not} $I\mapsto\Theta_I(\varsigma)$}, which jumps at every $I$ with an
endpoint at $\varsigma$; the two are one indicator apart, and every consumer of this clause must dispose
of that indicator before invoking it. \emph{Both do, and by different devices}: \textup{(d)}'s proof
restricts the supremum to $\{I:\varsigma\in I\}$, where the indicator is identically $1$; the
measurability paragraph of \cref{def:lfstep} tracks $\varsigma$ and $I$ \emph{together} along
\eqref{eq:ivalx}, where the indicator collapses to $\un\{x\in\mathcal X\}$.
\item[(g)] \emph{A closed-form admissible envelope, two-sided bounds on \eqref{eq:envconst}, and
attainment.} For every $I\in\mathfrak I$ and every $\varsigma\in\IR$,
\begin{equation}\label{eq:envsharp}
\big|\Theta_I(\varsigma)\big|
\;\le\;
2^{1/4}\,p_K^{2}\,e^{1/4}\,
\exp\Big(\frac{\big(|\varsigma|+2^{-1/2}\big)^{2}}{4}-\frac{\varsigma^{2}}{2}\Big) ,
\end{equation}
whose right-hand side is bounded, integrable, of finite moment of every order and hence itself an
admissible envelope for C3(ii), with \eqref{eq:envconst} returning at it the \emph{exactly computed}
pair $\mathsf c_2=2^{1/4}e^{1/2}p_K^{2}$ and $c_C=\sqrt\pi\big(1+\operatorname{erf}(2^{-3/2})\big)$,
the second free of $p_K$. Consequently, at the minimal envelope of \textup{(d)},
\begin{equation}\label{eq:envbnds}
\tfrac12\,p_K^{2}\;\le\;\mathsf c_2\;\le\;2^{1/4}e^{1/2}p_K^{2} ,
\qquad
\|\bar\Theta\|_{L^{1}}=2c_C\mathsf c_2\;\le\;2^{5/4}\sqrt\pi\,e^{1/2}\big(1+\operatorname{erf}(2^{-3/2})\big)\,p_K^{2} ,
\end{equation}
and hence $c_C\le2^{5/4}\sqrt\pi e^{1/2}\big(1+\operatorname{erf}(2^{-3/2})\big)$, also free of $p_K$.
Moreover the supremum defining $\|\bar\Theta\|_\infty$ is \textbf{attained}: there exist
$\varsigma^\ast\in\IR$ and $I^\ast\in\mathfrak I$ with $\varsigma^\ast\in I^\ast$ and
$|\Theta_{I^\ast}(\varsigma^\ast)|=\|\bar\Theta\|_\infty=\mathsf c_2$. Finally, along the slice
$\varsigma=0$ the supremum is carried by the \emph{shortest} admissible intervals:
$\bar\Theta(0)=\max\{\Theta_I(0):I\in\mathfrak I,\ |I|=2\}$.
\end{enumerate}
\end{lemma}

\noindent\emph{Note on \textup{(g)}: which of its bounds to substitute, into what, and in which
direction each consumer reads the pair.} \emph{The two bounds of \eqref{eq:envbnds} are the ones to
substitute and the standalone bound on $c_C$ is \textbf{not}}: that bound is formed with the
\emph{lower} half of \eqref{eq:envbnds} in its denominator, so pairing it with the upper half
double-counts. \emph{Five printed quantities read the pair $(\mathsf c_2,\|\bar\Theta\|_{L^{1}})$, the
list is given in full because the directions are not all the same. \emph{It is not offered as
exhaustive, and no criterion is given for membership}: readers of a single component exist too ---
\cref{lem:intker}\textup{(ii)}'s $\|\partial_v\mathsf C_{h_u}\|_\infty\le\mathsf c_2h_u^{-1}$ is one,
increasing in $\mathsf c_2$ and so safe under the same substitution --- and two drafts of this sentence
were short by a site while claiming not to be. \emph{What is asserted is only the ruling below, and it
holds for any reader whatever, monotone in the pair or invariant in it}: substituting the upper bounds
of \eqref{eq:envbnds} cannot falsify a conclusion.} \textbf{Three read it increasingly}: \eqref{eq:normconst} is
$\|\bar\Theta\|_{L^{1}}^{1/2}\mathsf c_2^{3/2}$ up to $\mathsf L$-free factors, \eqref{eq:cap} is its
square, and $\mathsf c_3$ of \cref{lem:intker}\textup{(i)} is $1\vee\|\bar\Theta\|_{L^{1}}$.
\textbf{One reads it reciprocally}: \eqref{eq:tccrad} carries
$\big(\mathsf c_2\|\bar\Theta\|_{L^{1}}\big)^{-1/2}$, since
$(2c_C)^{1/2}\mathsf c_2=\big(\mathsf c_2\|\bar\Theta\|_{L^{1}}\big)^{1/2}$.
\textbf{And one does not read it at all}: $\mathsf c_{18}$, of \cref{prp:phirate}'s proof, is
\emph{invariant} in the pair. It is a Bernstein exponent, a ratio of a mean bound to a range bound ---
$\mathsf K=4\bar c_\pi\|\bar\Theta\|_{L^{1}}^{2}\|\pi\|_\infty\|1/\pi\|_\infty$ over
$(4c_C\mathsf c_2\mathsf c_{11})^{2}=4\|\bar\Theta\|_{L^{1}}^{2}\mathsf c_{11}^{2}$ --- and both bounds
are quadratic in the same $\Delta_{h_x}$ operator norm, so the envelope cancels exactly and
$\mathsf c_{18}=\tfrac38\bar c_\pi\|\pi\|_\infty\|1/\pi\|_\infty\big/\mathsf c_{11}^{2}$. \emph{So
substituting the \textbf{upper} bounds \eqref{eq:envbnds} is safe throughout, and for one reason}: it
overstates the first three, which are read as upper bounds; it understates the fourth, which is a
certified curvature radius, and a smaller radius is a weaker statement of the same form; and it moves
the fifth not at all. Substituting \eqref{eq:envbnds} therefore leaves every conclusion true.

\noindent\emph{Note on the constants, and read every quantifier --- including the epistemic one.} At
$p_K=8$, numerically: $\mathsf c_2=\|\bar\Theta\|_\infty=34.0$ and $\|\bar\Theta\|_{L^{1}}=10.56$, hence
$c_C=\|\bar\Theta\|_{L^{1}}/(2\mathsf c_2)=0.155$, the supremum in \eqref{eq:envfam} being taken over
\emph{all} of $\mathfrak I$. \emph{These three are \textbf{located} values and not proved ones, and the
distinction is not pedantic here, because the direction of a search error and the direction the
constants must err in are opposite.} By \textup{(g)} the supremum \emph{is} attained, so there is a
maximiser to locate; a search reports its value \emph{at the points searched}, which is a \textbf{lower}
bound on a supremum. \emph{Every consumer needs an \textbf{upper} one, and in both directions of use}:
\eqref{eq:normconst}, \eqref{eq:cap} and $\mathsf c_3$ read the pair increasingly, so understating it
understates $\omega_n$ and \cref{def:lfstep}'s $\|\omega_n^{-1}\hat{\mathcal A}_\varphi
\vert_{\mathcal S_n}\|\le1$ --- which is \cref{prp:hatAprop}(i) and (i$'$) read at $\omega_n$, in each
of the two domain pairings, and holds
\emph{surely} --- ceases to be implied; \eqref{eq:tccrad} reads it reciprocally, so understating it
overstates a certified radius. \emph{Overstating is safe in both directions of use; understating is
unsafe in both, and a number below $\|\bar\Theta\|_\infty$ is the sup-norm of no admissible envelope at
all.} \emph{$\mathsf c_{18}$ is \textbf{not} a third case and an earlier draft of this sentence made it
one}: it is a ratio of a mean bound to a range bound, both quadratic in $\|\bar\Theta\|_{L^{1}}$, so the
envelope cancels and it is invariant --- the note to \textup{(g)} above computes it.
\emph{What is proved, and what it costs.} \textup{(g)} brackets the located values:
$32\le\mathsf c_2\le125.5$ and $\|\bar\Theta\|_{L^{1}}\le615.2$, both in closed form, so $34.0$ sits
$6.4\%$ above a proved floor and a factor $3.7$ below a proved ceiling, and the independently confirmed
ratio $1.06373$ to the Legendre value $\sum_{j=0}^{7}(2j+1)/2=32$ is \textbf{that floor}, not a
coincidence. \emph{A reading that requires its inequality \textbf{proved} rather than \textbf{measured}
runs at \eqref{eq:envbnds}}, and pays a factor $54.0$ in \eqref{eq:normconst}, $2.92\times10^{3}$ in
\eqref{eq:cap} and $58.3$ in $\mathsf c_3$ --- against $6.0\times10^{14}$ and $3.6\times10^{29}$ at the
explicit majorant. \emph{Two parts of the location remain unproved and are named rather than assumed}:
that the maximiser sits at an \emph{endpoint} placement, and that $\varsigma^\ast=0$. \textup{(g)}'s
Step 8 proves that, \emph{given} $\varsigma^\ast=0$, the interval has $|I|=2$; the placement is then a
one-parameter maximisation over a compact interval, provable term-by-term under the Lebesgue weight from
$|\mathrm{Leg}_j|\le1$ and \textbf{not} under the Gaussian one, whose orthonormal family is not
uniformly bounded. \emph{And $\varsigma^\ast=0$ cannot be reached by monotonicity}: $I\mapsto
|\Theta_I(\varsigma)|$ is \textbf{not} antitone at fixed $\varsigma\ne0$, since $|\Theta_I|$ has zeros
that a larger $I$ moves. Over the intervals the application realises, for which
$|I(x^\ast)|=|\mathcal X|/h_x\in[10.1,12.4]$, the same two norms are $17.4$ and $6.22$. \emph{The
uniform pair is what is quoted, and the realised pair is not}, because \eqref{eq:normconst} is a
constant of a theorem while $\mathfrak I$ is what the theorem quantifies over: $h_x$ is free in every
statement below, \cref{lem:intker}\textup{(iv)} admits every $h_x\le(\bar x-\underline x)/2$, and a
number measured at one bandwidth is a function of that bandwidth. The uniformity costs a factor of
$3.6$ in \eqref{eq:normconst} and $13$ in \eqref{eq:cap}, and that is the entire price of quoting a
constant rather than a measurement.

\noindent\emph{The Gram spectrum, reported as an infimum because that is what \textup{(c)} defines.}
Every eigenvalue here is of \eqref{eq:gramfam} as printed, with the bare $e^{-\varsigma^{2}/2}$; a
$\sqrt{2\pi}$-normalised Gram divides each by $\sqrt{2\pi}$, and an earlier draft of this note mixed the
two conventions. $\underline g=6.59\times10^{-8}$ over the grid searched, attained at $I=[-2,0]$, the
shortest admissible interval; $\lambda_{\min}(\mathsf G(\IR))=2.718\times10^{-1}$; and over half-lines
$\lambda_{\min}(\mathsf G([-\infty,0]))=1.51\times10^{-5}$. \emph{The application does \textbf{not}
escape the half-line value, and an earlier draft of this note said it did.} A design point at an
endpoint of $\mathcal X$ realises $I(x^\ast)=[0,|\mathcal X|/h_x]$ by \eqref{eq:ivalx}, a truncated
half-line, and $\mathsf G([0,\infty])-\mathsf G([0,b])=\mathsf G([b,\infty])\succeq0$ gives
$\lambda_{\min}(\mathsf G([0,b]))\le\lambda_{\min}(\mathsf G([0,\infty]))=1.51\times10^{-5}$ at every
right endpoint $b$, by \textup{(f)} for the reflection. So $1.51\times10^{-5}$ --- not the $2.57\times10^{-1}$ that
draft quoted --- is the infimum over the realised family, and it is attained \emph{exactly} at every
bandwidth the application runs; $2.57\times10^{-1}$ is a correct value at the \emph{symmetric} realised
interval, i.e.\ at an interior design point, and is the family's \textbf{maximum}. \emph{The
$\mathsf c_2=115.4$ derived from it is withdrawn}: $\|\cdot\|_\infty$ of the majorant divided by
$\lambda_{\min}$ is an envelope bound only at the infimum, which returns $1.97\times10^{6}$ there.
\emph{Nothing above depends on the repair}, $\mathsf c_2$ and $c_C$ being read off $\bar\Theta$ directly
and never through $\underline g$, which is why \textup{(d)} is written the way it is. \emph{These are
searched values and not proved bounds} --- \textup{(c)} proves the infimum positive; the numbers here
locate it.
% THE HILBERT-MATRIX ROUTE IS NOT PRINTED, AND WHY.  A worst-case-over-all-
% directions bound is available in closed form -- underline-g >= exp(-1/2)
% lambda_min(H_{p_K}) with H the Hilbert matrix, from [0,1] or [-1,0] inside I --
% and gives underline-g >= 6.7e-11, which is loose by THREE ORDERS against the
% measured 6.59042e-8.  (lambda_min(H_8) = 1.111539e-10, confirmed numerically, so
% the quoted Hilbert value was right and only the bound built on it was slack.)
% BOTH NUMBERS ARE IN THE BARE CONVENTION, AND THAT IS THE WHOLE POINT OF SAYING
% SO.  The weight is exp(-sigma^2/2), so the factor at sigma = 1 is exp(-1/2) =
% 0.60653 and NOT phi(1) = 0.24197, which carries the 1/sqrt(2pi) this Gram does
% not.  A 2026-08-16 pass repointed the measurement from the normalised 2.63e-8 to
% the bare 6.59042e-8 and left the bound beside it normalised at 2.7e-11 -- the
% two-conventions defect this very block exists to record, reintroduced inside it.
% The ratio is the tell and it is convention-invariant: 978, i.e. three orders, in
% either convention consistently applied; 2441 in the mixed one.
% Printing it would have cost three symbols the corpus does not otherwise need --
% a Hilbert matrix, a sign-flip diagonal D (italic D is Part F's perturbation
% direction, C2) and a polynomial P (P is probability) -- to state a number three
% orders worse than the one measured.  So the proof of (c) takes the SOFT route,
% which is a complete proof of positivity and needs none of them, and the number
% is reported above.  25,101 grid cells over the tanh-compactified parameter set.
%
% DISPATCH-27 MADE (d) THE MAJORANT.  DISPATCH-28 MAKES IT THE MINIMAL ENVELOPE
% AND KEEPS THE MAJORANT AS THE BOUND.  Both dispatches were repairing the same
% real defect -- the printed constants belonged to no single object -- and d27
% closed it in the direction that costs the most.
%   d27's THREE REASONS, ADJUDICATED.
%   (1) "The reported constants belonged to the majorant."  TRUE, and still the
%       reason an edit was owed.  Fixed here by making the SYMBOL name the object
%       whose numbers are reported, in the other direction: c_2 = 34.039389,
%       ||bar-Theta||_L1 = 10.5572, c_C = 0.155074, all off sup_I |Theta_I| over
%       all of frak-I.  The majorant's 4.50e8 / 1.8216 are printed as the bound.
%   (2) "underline-g^{-1} cancels" was an identity for the majorant.  TRUE, and
%       now MOOT: bar-Theta does not contain underline-g at all.  d27 read the
%       magnitude test off c_C alone (1.82 vs the retired half-width 4) and
%       concluded the majorant passed it.  THE TEST IS ON THE PRODUCT: ~28 sites
%       read 2 c_C c_2 = ||bar-Theta||_L1, whose retired magnitude was ~7.
%       Minimal: 10.56, PASSES.  Majorant: 1.64e9, EIGHT ORDERS OUT.
%   (3) "Measurability was justified by (f), which gives continuity of
%       I -> G(I)^{-1} e_0 only, and I -> Theta_I(s) JUMPS."  THE CRITICISM OF
%       THE ROUTE IS CORRECT; THE CONCLUSION ABOUT THE OBJECT IS NOT.  The jump
%       is only the indicator falling to zero, and it disappears once the sup is
%       restricted to {I : s in I} -- which changes nothing, the omitted terms
%       being zero.  On that set the map is jointly continuous, the constraint
%       set is compact and varies continuously, and (d)'s two semicontinuity
%       paragraphs give CONTINUITY, not merely measurability, in four lines and
%       with no citation.  (f) is proved from (a) and (c), so not circular.
%   WHY IT MATTERS, AND IT IS NOT THE PRINTED CAP.  c_2 enters eq:normconst
%   SQUARED and eq:cap to the FOURTH; it enters eq:tccrad RECIPROCALLY, at
%   (c_2 ||bar-Theta||_L1)^{-1/2}.  At the majorant, prp:tccrepair(ii)'s certified
%   radius is smaller by 4.5e7 -- a SUFFICIENT CONDITION driven to vacuity.
%   Sharpness is its whole content.
%   c_18 IS NOT IN THIS LIST AND AN EARLIER DRAFT PUT IT HERE AT 2.4e16.  c_18 is
%   3K/(8 (4 c_C c_2 c_11)^2) with K = 4 bar-c_pi (2 c_C c_2)^2 ||pi|| ||1/pi||;
%   both carry ||bar-Theta||_L1 SQUARED and it cancels exactly, leaving
%   c_18 = (3/8) bar-c_pi ||pi|| ||1/pi|| / c_11^2, free of the envelope.  A
%   Bernstein exponent is a mean bound over a range bound and both are quadratic
%   in the same Delta_{h_x} norm, which IS ||bar-Theta||_L1 by Schur's test at the
%   (N2.1) display.  The 2.4e16 was computed by moving K's numerator and leaving
%   the denominator fixed.
%
% THE L1 WAS 8.4% LOW AND THE REASON IS THE SAME JUMP AS IN (3).  A tanh-grid
% sweep returns 9.74; the converged value is 10.5572.  The maximiser sits where
% an interval ENDPOINT CROSSES sigma, a coincidence a fixed (alpha,beta) grid
% meets only at isolated sigma.  Refinement ladder: 9.578 (grid only) -> 10.062
% (union with a structured (L,a) sweep) -> 10.514 (+ |I|=2 edge, [0,L] family)
% -> 10.5572 (+ exact corners, right-endpoint family, spike sweep), increments
% 4.5% -> 0.34% -> 0.074% -> 0.000%.  c_2 = 34.039389 needed no such refinement
% and was confirmed independently at ratio 1.06373 to the Legendre value 32.
%
% THE lambda_min QUANTIFIER INVERSION, FOUND AND REPAIRED 2026-08-16.  The note
% read "|I| in [10.1, 12.4] at every design point, where lambda_min(G) = 2.57e-1".
% 2.57e-1 is a CORRECT value -- at the SYMMETRIC realised interval, i.e. at an
% INTERIOR design point (solving for it returns |I| = 10.0990) -- and it is the
% family's MAXIMUM standing in an INFIMUM's slot.  A BOUNDARY design point
% realises I = [0, |X|/h_x], a truncated half-line, and G([0,inf]) - G([0,L]) =
% G([L,inf]) >= 0 gives lambda_min(G([0,L])) <= 1.507063e-5 for every L -- the
% note's OWN half-line figure, four orders below, and ATTAINED EXACTLY at all
% four bandwidths the application realises.  So the endpoint case, the only
% configuration the boundary family exists for, was quantified away.  The derived
% c_2 = 115.4 is WITHDRAWN: it happens to dominate (sup over realised I of
% ||Theta_I||_inf is 17.4) but its derivation is invalid, and reproducing it at
% the true infimum gives 1.97e6.
% ALSO REFUTED, and it was already flagged unverified: "over half-lines
% lambda_min(G) >= 3.40e-2" is measured at 1.507e-5 bare, 6.012e-6 phi-weighted.
% NOT A DEFECT, AND RECORDED SO IT IS NOT RE-RAISED: the equilibrated 1.7368e-7 is
% lambda_min of D^{-1/2} G D^{-1/2}, a DIFFERENT matrix -- the solve's conditioning,
% not underline-g, which (c) defines on eq:gramfam as printed.
% THE HILBERT-MATRIX ROUTE IS STILL NOT PRINTED; see the retained note above.

% PROOF RELOCATED TO proofs.tex, HEAD OF PART H, 2026-08-20.  It was the ONLY
% theorem-like result in the corpus proved where it was stated, and
% .claude/rules/quality.md Sec.4 clause 1 gates on "a matching \begin{proof} in
% proofs.tex, keyed by \label" -- so a literal run reported lem:bdryfam UNPROVED.
% That is a FALSE FAIL, and this project has already paid for both signs of this
% error: a false PASS 2026-08-03 (gating on a retired directory) and a false FAIL
% 2026-08-05 (gating on a withdrawn corollary).  The repair is to make the
% invariant TRUE rather than to qualify the rule: a one-file check cannot drift
% from its artifact, and a two-file label-matching check is the shape of check
% that produced the 2026-08-03 pass.  Nothing about the proof changed in the move.

\begin{lemma}[Properties of the integrated kernel]\label{lem:intker}
Suppose C3 and \textbf{A0's compact-interval clause} hold. \emph{A0 moves \textbf{onto} this list on
2026-08-26, and the ground is the assumption block's own split}: the interval is what A0 \emph{restricts}
$\mathcal X$ to and not what it introduces, a fat Cantor $\mathcal X$ leaves every symbol named and the
restriction false, so it is a \textbf{clause}, and the result that reads it carries it.
\emph{Clause \textup{(iv)} is where it is read, and it reads it twice}: once for
\eqref{eq:ivalx} to be an element of $\mathfrak I$ at all, and once for the substitution
$x=x^\ast+h_x\varsigma$ to carry the $\mathcal X$-integral onto a \emph{single} interval.
\emph{Through 2026-08-25 this line said the opposite} --- that the compact interval is a
\textbf{declaration} standing at every statement, named outside the ``Suppose \dots\ hold'' list and not
a hypothesis --- \emph{and cited two sites as precedent, neither of which supports it}:
\cref{lem:levels}, whose defence is retyped in this same pass because that lemma does \textbf{not} read
the interval, and \cref{lem:unifop}, which \textbf{does} carry A0, as a block and first on its own
hypothesis line, so it was never an instance of the idiom. \emph{And through 2026-08-23
this was the one site that used the interval silently}, which is the reading that survives all of it. Then, with $\mathsf c_2=\|\bar\Theta\|_\infty$ and $c_C$ of \eqref{eq:envconst},
and $\Theta_{I}$ of \eqref{eq:rescaled}:
\begin{enumerate}
\item[(i)] \emph{Tails and normalisation.} For every $\epsilon>0$ there is a finite
$c_\Theta(\epsilon)$ with $\int_{|\varsigma|>c_\Theta(\epsilon)}\bar\Theta\le\epsilon$, and then
$|\mathsf C_{h_u}(v)|\le\epsilon$ for $v\le-c_\Theta(\epsilon)h_u$ and
$|\mathsf C_{h_u}(v)-1|\le\epsilon$ for $v\ge c_\Theta(\epsilon)h_u$; and
$|\mathsf C_{h_u}|\le\mathsf c_3$ with
$\mathsf c_3:=1\vee 2c_C\mathsf c_2$, uniformly in $h_u$. \emph{Nonnegativity is \textbf{not} claimed
and is not available}: at $p_K=8$ the kernel $\Theta_1$ must take negative values, so
$\mathsf C_{h_u}$ does too. Every consumer uses only the two-sided bound.
% ARCHAEOLOGY.  Rounds 1--11 printed "0 <= C_{h_u}" in this clause while the proof
% disclaimed it.
% DISPATCH-26.  Through dispatch 25 (i) read "C(v) = 0 for v <= -c_C h_u and
% C(v) = 1 for v >= c_C h_u" -- EXACT zero and EXACT one, available only under
% compact support.  Under the envelope the exact values are replaced by the
% epsilon-approximation above, which is what every consumer actually uses, and
% THE c_3 CLAUSE SURVIVES VERBATIM because it was never a support statement:
% 2 c_C c_2 is ||bar-Theta||_{L^1} by eq:envconst.  The two consumers of the
% exact form are prp:tccfails Step 3 and prp:norangeinv, both repaired in place.
\item[(ii)] Derivative, and a second-order bound without a second derivative. $\mathsf C_{h_u}$
is continuously differentiable with $\partial_v\mathsf C_{h_u}(v)=h_u^{-1}C_{h_u}(v,\IR)$ and
$\|\partial_v\mathsf C_{h_u}\|_\infty\le\mathsf c_2h_u^{-1}$; $\partial_v\mathsf C_{h_u}$ is Lipschitz
with constant $\mathsf c_4h_u^{-2}$, $\mathsf c_4$ being the Lipschitz constant of
$C_{h_u}(h_u\cdot,\IR)$ furnished by C3(iii); and consequently, for all $v,v'\in\IR$,
\[
\big|\mathsf C_{h_u}(v')-\mathsf C_{h_u}(v)-\partial_v\mathsf C_{h_u}(v)\,(v'-v)\big|
\;\le\;\tfrac12\,\mathsf c_4\,h_u^{-2}\,(v'-v)^{2} .
\]
\emph{No second derivative of $C_{h_u}$ is assumed or used}; C3(iii) gives a Lipschitz kernel and the
display above is what a Lipschitz derivative buys. \emph{The increment carries no symbol}: it is written
out as $v'-v$, in the pair this lemma's own proof already uses for the two arguments of
$\partial_v\mathsf C_{h_u}$.
% ARCHAEOLOGY (dispatch-22 change).  Through dispatch 21 the increment was written
% varDelta, a variant capital, in a LEMMA STATEMENT, against the period-difference
% operator Delta -- not merely a symbol with many sites but a PREFIX applied to
% almost everything in this project (Delta Z, Delta U, Delta Y, Delta varphi,
% Delta_{h_x}, and -- until 2026-08-02 -- lem:bdiff's martingale differences
% Delta_k, which the citation of blm2013 Thm 6.2 deleted along with the whole
% Doob/Azuma construction; the crowding this paragraph records was REAL when the
% rename was made and the rename is not being reopened) -- so
% "varDelta in R" invited the reading "a difference of something", and the consumer
% three Parts later applies the display at Delta_{h_x} v_i.  Glossary C80 records
% why this member of the variant-capital family was renamed while the other three
% were covered by a rule.
\item[(iii)] Order, for a function of \textbf{bounded variation}, which is what its
consumers supply. Let $G:\IR\to\IR$ be of bounded variation
with $G(-\infty)=0$, let $dG$ denote the associated finite signed Borel measure, so that
$G(u)=dG\big((-\infty,u]\big)$ and $|dG|(\IR)<\infty$, and suppose in addition that $G$ is $p_K$ times
continuously differentiable with $\|G^{(p_K)}\|_\infty<\infty$. Then, for every $u$,
\[
\Big|\int_{\IR}\mathsf C_{h_u}(u-a)\,dG(a)-G(u)\Big|
\;\le\;\mathsf c_5\,\|G^{(p_K)}\|_\infty\,h_u^{p_K},
\qquad
\mathsf c_5:=\frac1{p_K!}\int_{\IR}|\varsigma|^{p_K}\big|\Theta_1(\varsigma)\big|\,d\varsigma ,
\]
$\mathsf c_5$ finite and uniform in $h_u$ by C3(i). The two hypotheses on $G$ are not decoration.
A version asking only that $G$ map into $[0,1]$ and be smooth is \emph{false}: at $G\equiv\tfrac12$ it
reads $|0-\tfrac12|\le0$, because $dG=0$ while $G(u)=\tfrac12$. \emph{What the identity needs is that
$G$ be recovered from $dG$}, which is
$G(-\infty)=0$ together with bounded variation, and \emph{that is precisely what its consumers have}:
\cref{lem:opnoise} Step~3 applies it to $\mathcal G_{t,l}=\mathcal M(\cdot,t,l)$, a difference of two
nondecreasing functions each vanishing at $-\infty$ and of total variation at most $\bar w(t)$
, \emph{at every $\mathsf L$, the total variation of the signed measure
$E[\widetilde W_l\,\delta_{\Delta U}\,w(t'\Phi(\Delta Z))]$ being at most
$E[|\widetilde W_l|\,w]\le\bar w(t)$ by the \textbf{per-cell} bound $|\widetilde W_l|\le1$, so the
constant does not move}; and
\cref{lem:aerr} Step~2 applies its argument to conditional-mean products that are likewise differences
of finite measures' distribution functions; and \cref{lem:unifop} Step~2 re-runs
\cref{lem:opnoise} Step~3's computation verbatim at
$\mathcal M(\cdot,t,l;\varphi,\theta_\dagger)$ and carries this clause's $\mathsf c_5$ into its own
aggregation, so it is a third. None of them needs $G\ge0$ and none needs $\int dG=1$. \emph{No claim
of exhaustiveness is made, and that is deliberate}: the draft this clause replaces said ``the two
consumers'' twice and was short by that site both times.
% ARCHAEOLOGY.  Rounds 1--11 stated this clause for G mapping into [0,1] and
% smooth, and the proof silently used both G >= 0 and int dG = 1.
\item[(iv)] Exact $x$-normalisation, and the property of the boundary adjustment it rests on is
now \textbf{stated}. Suppose $h_x\le(\bar x-\underline x)/2$. For every $x^\ast\in\mathcal X$ put
\begin{equation}\label{eq:ivalx}
I(x^\ast)\;:=\;\Big[\frac{\underline x-x^\ast}{h_x},\ \frac{\bar x-x^\ast}{h_x}\Big]\;\in\;\mathfrak I ,
\end{equation}
the \emph{realised truncation interval} at $x^\ast$ in scaled coordinates, which lies in $\mathfrak I$
because $\underline x\le x^\ast\le\bar x$ gives $\inf I(x^\ast)\le0\le\sup I(x^\ast)$ and
$h_x\le(\bar x-\underline x)/2$ gives $\sup I(x^\ast)-\inf I(x^\ast)=(\bar x-\underline x)/h_x\ge2$.
Then, \emph{for every $x^\ast\in\mathcal X$, one case},
\[
h_x^{-1}\!\int_{\mathcal X}C_{h_x}\big(x-x^\ast,\,I(x^\ast)\big)\,dx\;=\;1 .
\]
\emph{The clause of C3 this consumes is \eqref{eq:truncmom} at $j=0$}. \emph{There are no longer three
cases and no reflection convention}: the interval index records both how far the nearer endpoint is
and which endpoint it is, so the interior point, the two boundary points and the point close to
\textbf{both} endpoints --- which a scalar index could not name at all, and which the round-12 form of
this clause excluded by hypothesis --- are one clause at four values of $I$.
\emph{Load-bearing}:
% ARCHAEOLOGY.  Rounds 1--11 left the index to a parenthesis reading "the rescaling
% implicit in C3(ii)", and the proof chose the fraction of the interval
% (x* - underline-x)/(bar-x - underline-x).  Round 12 replaced that by the SCALAR
% eq:s0def, s_0(x*) = min{1, dist(x*, endpoints)/(c_C h_x)}, together with a
% reflection CONVENTION for the right endpoint and the hypothesis
% h_x < (bar-x - underline-x)/(2 c_C), whose only job was to guarantee that at
% most ONE endpoint lies within c_C h_x of x*.
%   DISPATCH 26 RETIRES ALL THREE.  eq:s0def is DELETED and replaced by eq:ivalx.
% The one-endpoint restriction is dissolved by the indexing rather than assumed.
% The reflection is lem:bdryfam(f), proved.  And the hypothesis becomes the
% ABSOLUTE h_x <= (bar-x - underline-x)/2 -- which is the floor |I| >= 2 and
% nothing else -- BECAUSE IT HAD TO: c_C is now DERIVED from bar-Theta, bar-Theta
% from underline-g, and underline-g from the floor, so a floor stated as a
% multiple of c_C would be circular.  The live configuration gives
% |I| = |cal-X|/h_x in [10.1, 12.4], so the new hypothesis binds nowhere.
\cref{prp:hatAprop}(iii) rests entirely on this clause, and without it $\int\hat\nu_{u,t,l}\,dx\ne0$ at
\emph{every} design point --- the defect being the tail $\int_{\IR\setminus I(x^\ast)}\Theta_1$, which is
$O(1)$ only where the nearer endpoint is within $O(h_x)$ of $x^\ast$ and tends to $0$ with $h_x$ at every
interior point, by integrability of $\bar\Theta$ --- so that clause's conclusion fails:
$\hat{\mathcal A}^\ast_\varphi\Xi\notin\mathcal V$ for some $\Xi\in L^{2}(\Pi)$, read at
\cref{prp:hatAprop}\textup{(iii)}'s own bound variable, and (L2) of \cref{prp:phirate} Step~2 loses the
projection identity it runs on.
\emph{The width is written $O(h_x)$ and not $c_Ch_x$}: under the envelope $c_C$ is the ratio
\eqref{eq:envconst} of two norms of $\bar\Theta$ and is not a length.
\end{enumerate}
\end{lemma}

\begin{definition}[The $x$-smoothed sample operator]\label{def:hatTexp}
Throughout Part~H, $\hat T$ is \eqref{eq:hatT} of \cref{def:hatT}, under that definition's
sample-splitting and truncation conventions: $\{1,\dots,n\}=\mathcal I_a\cup\mathcal I_m$ with
$n_{\mathrm a}:=\#\mathcal I_a$, $n_{\mathrm m}:=\#\mathcal I_m$, $\hat\varrho$ computed on
$\mathcal I_a$ and taking values in the $\varepsilon_\varrho$-interior simplex \eqref{eq:trunset}, and
$\widehat{\widetilde W}_{i,l}:=\un\{W_i=l\}-\hat\varrho_l(\Delta Z_i)$,
$\widehat{\widetilde W}_i:=(\widehat{\widetilde W}_{i,1},\dots,\widehat{\widetilde W}_{i,\mathsf L-1})'$,
satisfying
\begin{equation}\label{eq:wbounds}
\varepsilon_\varrho\;\le\;\big|\widehat{\widetilde W}_{i,l}\big|\;\le\;1
\qquad\text{surely, for every }i\text{ and every }l\in\{1,\dots,\mathsf L-1\} ,
\end{equation}
the upper bound as in (P1) \emph{read at the $l$-th contrast} and the lower bound by A4(i) and the
truncation. \emph{Both bounds are \textbf{per cell} and neither is a statement about
$\|\widehat{\widetilde W}_i\|$}: \eqref{eq:trunset} puts $\hat\varrho_l\in[\varepsilon_\varrho,1-\varepsilon_\varrho]$
for every $l$, the lower bound by the first constraint, the upper because
$\hat\varrho_l\le\mathbf 1'\hat\varrho\le1-\varepsilon_\varrho$, so
$\un\{W_i=l\}-\hat\varrho_l$ lies in $[\varepsilon_\varrho,1-\varepsilon_\varrho]$ in modulus at both
values of the indicator. At $\mathsf L=2$ there is one contrast, $\un\{W_i=1\}=W_i$, and
\eqref{eq:wbounds} is the display earlier rounds printed.
\emph{\eqref{eq:wbounds} is used at four places and each would fail without it}: \cref{prp:hatAprop}(v),
\cref{prp:tccfails} and \cref{prp:norangeinv}(ii) all need a design point at which the sample operator
responds, and \cref{prp:tccrepair}\textup{(iii)} needs $\mathbf G_n$ to have no zero row. \emph{All four
therefore carry A4(i) on their own hypothesis lines, and each of the four does}; the free \emph{upper}
half of \eqref{eq:wbounds} is a different matter and is read, without A4(i), at \cref{prp:hatAprop}(i)
and at \cref{lem:aerr} Step~1, whose own hypothesis line carries A4(iii) and not A4(i) and whose
$\sum_l\varpi_{\mathsf w}(l)\widehat{\widetilde W}_{1,l}^{2}\le1$ is that upper half read with
B1(iv.d). \emph{Those two are named as readers of the upper half and no claim of exhaustiveness is
made about it}: the draft this sentence replaces named one.
% CORRECTIONS 2026-08-10 (Phase 6a) AND 2026-08-23, AND THE FIRST IS A MISCOUNT
% NO STRING CENSUS FINDS.  THE HEADING IS NOT NUMBERED, DELIBERATELY: it read
% "TWO CORRECTIONS" and a third was appended below, which is the defect this
% block keeps being bitten by.  (1) The register said "three places" and then NAMED FOUR.  The
% count is four -- prp:hatAprop(v), prp:tccfails, prp:norangeinv(ii),
% prp:tccrepair(iii) -- which is exactly the eq:wbounds row of the A4(i) consumer
% analysis.  (2) The clause number was (ii) and the appeal is in (iii): (iii) is
% where "eq:wbounds rules this out surely" is printed, excluding the zero-row
% degeneracy that would make (ii) vacuous.  A register that points one clause off
% is what let prp:tccrepair's own hypothesis line go without A4(i) for a dispatch.
% prp:hatAprop(i) is named here rather than added to the count BECAUSE IT READS
% ONLY THE UPPER HALF, which is (P1) and holds at every varrho: it would NOT fail
% without A4(i), so it is not a consumer of this display in the register's sense.
% (3) ADDED: lem:aerr Step 1 reads that same upper half, at
% sum_l varpi_w(l) Wtilde-hat_{1,l}^2 <= 1, and is likewise not a consumer in the
% register's sense.  The sentence above named prp:hatAprop(i) ALONE and was short
% by it -- a closure operator inside a repair of a miscount.

For $\varphi\in\mathcal V$ and $i\in\mathcal I_m$ let the \emph{smoothed period difference} be
\begin{equation}\label{eq:smdiff}
\Delta_{h_x}\varphi_i\;:=\;\frac1{h_x}\int_{\mathcal X}
\Big[C_{h_x}\big(x-X_{i2},I_{i2}\big)-C_{h_x}\big(x-X_{i1},I_{i1}\big)\Big]\varphi(x)\,dx ,
\end{equation}
with $I_{ij}:=I(X_{ij})$ of \eqref{eq:ivalx}, which lies in $\mathfrak I$ \emph{at every $n$} by C5$'$(ii), and write
$\Delta_{h_x}\varphi:=(\Delta_{h_x}\varphi_i)_{i\in\mathcal I_m}\in\IR^{n_{\mathrm m}}$. The
\emph{$x$-smoothed sample operator} is \eqref{eq:hatT} with $\Delta\varphi_i$ replaced by
$\Delta_{h_x}\varphi_i$:
\begin{equation}\label{eq:hatTx}
\hat T_{h_x}(\varphi,\theta)(u,t,l)\;:=\;\frac1{n_{\mathrm m}}\sum_{i\in\mathcal I_m}
\widehat{\widetilde W}_{i,l}\;\mathsf C_{h_u}\big(u-\Delta Y_{i,\theta}+\Delta_{h_x}\varphi_i\big)\;
w\big(t'\Phi(\Delta Z_i)\big),
\qquad \Delta Y_{i,\theta}:=\Delta Y_i-\Delta Z_i'\theta ,
\end{equation}
at each $(u,t,l)\in\IR\times\mathcal T\times\{1,\dots,\mathsf L-1\}$, \emph{which is \eqref{eq:hatT}'s
arity read argument for argument}, $\hat T_{h_x}$ differing from $\hat T$ in the period difference alone.
\emph{The two agree to the accuracy of the $x$-smoothing and differ in one respect that decides
everything below}: $\varphi\mapsto\hat T(\varphi,\theta)$ depends on $\varphi$ through the
\emph{point evaluations} $\varphi(X_{ij})$, while $\varphi\mapsto\hat T_{h_x}(\varphi,\theta)$ depends
on it through the \emph{bounded linear functionals} \eqref{eq:smdiff}.
\end{definition}

\begin{definition}[The smoothed sample linearisation, exhibited]\label{def:hatAcon}
Fix $\theta$ and $\varphi\in\mathcal V$, and put, for $i\in\mathcal I_m$,
\begin{equation}\label{eq:hatg}
\hat g_{u,t,l,i}(\varphi,\theta)\;:=\;\widehat{\widetilde W}_{i,l}\;w\big(t'\Phi(\Delta Z_i)\big)\;
\frac1{h_u}\,C_{h_u}\big(u-\Delta Y_{i,\theta}+\Delta_{h_x}\varphi_i,\;1\big) ,
\end{equation}
the sample counterpart of $g_{u,t,l}$ of \eqref{eq:kernel} with the conditional density replaced by the
derivative of the integrated kernel. \emph{The contrast index enters through
$\widehat{\widetilde W}_{i,l}$ and through nothing else}, the kernel and the Bierens factor being
$l$-free; that is what makes every bound below a per-cell bound. \emph{The residual inside \eqref{eq:hatg} carries the
\textbf{smoothed} difference $\Delta_{h_x}\varphi_i$ of \eqref{eq:smdiff}, not the raw $\Delta\varphi_i$,
and the choice is not cosmetic}: with the raw difference every clause of \cref{prp:hatAprop} except
(iv) is unchanged, and (iv) acquires an additional \emph{first-order} term, which is exactly the
discrepancy \cref{prp:tccfails} exploits. Define the \emph{kernel-smoothed sample difference
functional}
\begin{equation}\label{eq:hatnu}
\hat\nu_{u,t,l}(x)\;:=\;\frac1{n_{\mathrm m}h_x}\sum_{i\in\mathcal I_m}\hat g_{u,t,l,i}(\varphi,\theta)
\Big[C_{h_x}\big(x-X_{i2},I_{i2}\big)-C_{h_x}\big(x-X_{i1},I_{i1}\big)\Big] ,
\end{equation}
which replaces each of the two products $E[\,\cdot\mid X_j=x]f_{X_j}(x)$ constituting $\nu_{g_{u,t,l}}$ by
its $C_{h_x}$ kernel estimate, exactly as \cref{def:hatA} prescribes. Then
\begin{equation}\label{eq:hatAcon}
\begin{split}
\hat{\mathcal A}_\varphi h(u,t,l)\;&:=\;\int_{\mathcal X}\frac{\hat\nu_{u,t,l}(x)}{\pi(x)}\,h(x)\,\pi(x)\,dx
\;=\;\int_{\mathcal X}\hat\nu_{u,t,l}(x)\,h(x)\,dx\\
&\;=\;\frac1{n_{\mathrm m}}\sum_{i\in\mathcal I_m}\hat g_{u,t,l,i}(\varphi,\theta)\,\Delta_{h_x}h_i ,
\end{split}
\end{equation}
with kernel $\hat{\mathsf k}_\varphi(u,t,l;x):=\hat\nu_{u,t,l}(x)/\pi(x)$ in the $L^2(\pi)$ pairing, and
\begin{equation}\label{eq:hatAadj}
\hat{\mathcal A}^\ast_\varphi\Xi(x)\;:=\;\frac1{\pi(x)}\sum_{l=1}^{\mathsf L-1}\varpi_{\mathsf w}(l)
\int_{\IR\times\mathcal T}\hat\nu_{u,t,l}(x)\,
\Xi(u,t,l)\,\Pi_u(du)\,\Pi_t(dt),
\qquad \Xi\in L^{2}(\Pi) .
\end{equation}
\emph{Where the first three arguments are suppressed}, $\hat g_{\cdot,\cdot,\cdot,i}(\varphi,\theta)$
denotes the element $(u,t,l)\mapsto\hat g_{u,t,l,i}(\varphi,\theta)$ of $L^{2}(\Pi)$; the third dot
is the contrast index, and at $\mathsf L=2$ it stands for the single value $l=1$, so the glyph is the
two-dot one earlier rounds printed.
\emph{The sum over $l$ against the masses $\varpi_{\mathsf w}(l)$ is \cref{def:spaces}'s pairing and is
what makes \eqref{eq:hatAadj} the adjoint of \eqref{eq:hatAcon}}; it is the sample counterpart of
\cref{lem:adjoint}\textup{(i)}, and it is the third factor of $\Pi$ written out rather than folded into
the measure. \emph{At $\mathsf L=2$ there is one term, $\varpi_{\mathsf w}(1)=1$, and \eqref{eq:hatg},
\eqref{eq:hatnu}, \eqref{eq:hatAcon} and \eqref{eq:hatAadj} are the displays earlier rounds printed}
under the canonical identification of $\IR\times\mathcal T\times\{1\}$ with $\IR\times\mathcal T$.
\emph{The generic element of $L^{2}(\Pi)$ is $\Xi$, as fixed in \cref{rem:spacesread}\textup{(1)}, and not $\alpha$},
which is \textbf{binding} as the propensity rate exponent $r/(2r+d_Z)$ in (P-a.2), (P-f), C5$'$(i),
A5(iii), \eqref{eq:frontier}, \eqref{eq:lxwindow} and \eqref{eq:prelimrate}.
% ARCHAEOLOGY.  Xi is used at Upsilon Xi in cor:rangegap, so no letter is
% introduced.  The declaration is in def:spaces -- which is in
% results_s2_identification.tex:71, NOT in this file (this clause said "of this
% file", stale since the results.tex split; corrected 2026-08-10, Phase C) -- and
% NOT in the glossary's Sec.1.6; a dispatch-14 draft of this sentence cited
% Sec.1.6 and was wrong about where.  Glossary C72.
%   *** PHASE 4 CORRECTION TO THIS BLOCK.  It also claimed a STARRED variant
% "Xi_star in prp:tccfails".  SWEPT: \Xi_\star occurs at ZERO sites in the whole
% corpus -- not in prp:tccfails, not in its proof, not in any file.  The claim was
% false and is struck; a sentence recording it in PRINT was drafted in this pass
% and withdrawn before it shipped.  The four LIVE species of Xi, which do exist,
% are now enumerated in rem:spacesread(1). ***
As in \cref{def:hatA}, $\hat{\mathcal A}_{\mathrm{cond}}:=\hat{\mathcal A}_{\vdagger}$ and
$\hat{\mathcal A}^\ast_{\mathrm{cond}}:=\hat{\mathcal A}^\ast_{\vdagger}$.

Two features of \eqref{eq:hatAcon} are worth naming because they are not accidents. First, the
reference density cancels: the operator is $h\mapsto\int\hat\nu_{u,t,l}h\,dx$ and is therefore
\emph{computable without knowing $\pi$}, which the population identity
$\langle\nu_g/\pi,h\rangle_{L^{2}(\pi)}=\int\nu_gh\,dx$ of (P5) already predicted; only its
\emph{kernel} with respect to $\pi\,dx$ carries $\pi^{-1}$, which is the same bookkeeping
\cref{rem:nupi} records for A2. \emph{That is a statement about \eqref{eq:hatAcon} and it does not
extend to \eqref{eq:hatAadj}, which is the operator the estimator actually iterates}: the adjoint
carries $1/\pi(x)$ in front of the whole expression, $\pi$ is an unknown feature of $P$ by
\cref{def:spaces}, and \emph{every} Landweber step of \cref{def:lfstep}\textup{(b)} is taken through
$\hat{\mathcal A}^\ast_\varphi$. So the forward operator is free of $\pi$, the iteration is not, and the
analyst forms \eqref{eq:hatAadj} at the floored estimate $\hat\pi$ that \cref{def:lfstep} fixes.
\emph{The cancellation above is true and, read by itself, misleading}, which is why the qualification is
printed here, four lines from the display it qualifies, rather than left to a reader who has just been
told that $\pi$ need not be known.
% ADDED 2026-08-26, PHASE 4d, AND THE DEFECT WAS NOT FALSITY.  The sentence above
% is correct about eq:hatAcon and it sits exactly where a reader looks for pi's
% status -- so it read as a claim that the ESTIMATOR does not need pi, which is
% false: eq:hatAadj is two displays up and divides by it.  A true sentence about
% the wrong operator, in the one place the question is asked, is worse than
% silence, because it answers.
Second, \eqref{eq:hatAcon} depends on $h$ through
$\Delta_{h_x}h$ and through nothing else, which is what makes it a bounded operator and is the whole
content of \cref{prp:hatAprop}.
\end{definition}

\begin{proposition}[What $\hat{\mathcal A}_\varphi$ is, and what it is not]\label{prp:hatAprop}
Suppose A0, B1(i), (ii), (iv.a), (iv.b), (iv.c), (iv.d), C1(ii) and C3 hold, and let
$\hat{\mathcal A}_\varphi$
be \eqref{eq:hatAcon}; for clause (v) suppose in addition A4(i), so that \eqref{eq:wbounds} holds.
% B1(iv.d) ADDED TO THIS LIST, PHASE 4c, 2026-08-10.  IT IS AN ASSUMPTION-USAGE
% EDGE THAT WAS CONSUMED BUT NOT LISTED, AND IT IS NOT COSMETIC.  Clause (i) below
% cites B1(iv.d) BY NAME -- "sum_l varpi_w(l) = 1, Pi_w being a probability measure
% by B1(iv.d)" -- and the proof consumes the same device at eight sites.  c_Pi of
% B1(iv.c) is Pi_u(R) int bar-w^2 dPi_t and carries NO Pi_w factor, so the l-sum
% must contribute exactly 1; if Pi_w were merely finite of total mass M, clause
% (i)'s constant would acquire M^{1/2} and the proposition as stated would NOT be
% implied by its listed hypotheses.
% PROPAGATION, AND ONE RESULT THAT MUST NOT BE TOUCHED.  prp:tccrepair and
% prp:norangeinv inherit "the hypotheses of prp:hatAprop" and are repaired by this
% one edit.  prp:tccfails does NOT inherit: it states its own list carrying the
% whole of B1(iv), which already contains (iv.d).  It is left alone deliberately.
Then, conditionally on the sample, and for every $\varphi\in\mathcal V$ and $\theta\in\IR^{p}$:
\begin{enumerate}
\item[(i)] \emph{Boundedness.} $\hat{\mathcal A}_\varphi:\mathcal V\to L^{2}(\Pi)$ is a bounded linear
operator with
$\|\hat{\mathcal A}_\varphi\|\le 2\,(2c_C)^{1/2}\,\mathsf c_2^{2}\,c_\Pi\,\|1/\pi\|_\infty^{1/2}\,h_u^{-1}h_x^{-1/2}$.
This constant does not move with $\mathsf L$, and the derivation is per cell, which is what
\eqref{eq:normconst} asserts of it. At each fixed $l$ the estimate of
$|\hat{\mathcal A}_\varphi h(u,t,l)|$ reads the contrast only through
$|\widehat{\widetilde W}_{i,l}|\le1$ of \eqref{eq:wbounds}, \emph{the per-cell bound, not the
Euclidean bound $\|\widehat{\widetilde W}_i\|\le c_{\widetilde W}$, which would not reproduce
this constant; and the \textbf{free upper half} of that display and not its lower half, the upper bound
being \textup{(P1)} read at the $l$-th contrast, which holds at every $\varrho$. \textbf{So this clause
reads no overlap condition}, and A4(i), which the hypothesis line above scopes to clause \textup{(v)}
alone, is correctly absent from it}, and the resulting bound on $\int_{\IR\times\mathcal T}|\hat{\mathcal A}_\varphi
h(u,t,l)|^{2}\Pi_u(du)\Pi_t(dt)$ is the same at every $l$, the kernel and the Bierens factor in
\eqref{eq:hatg} being $l$-free. Summing over $l$ against the masses therefore contributes
$\sum_l\varpi_{\mathsf w}(l)=1$, $\Pi_{\mathsf w}$ being a probability measure by B1(iv.d), and
$c_\Pi$ of B1(iv.c) is unchanged for the same reason. At $\mathsf L=2$ the sum has one term and
the derivation is the printed one.
\item[(i$'$)] \emph{The same estimate, in the two pairings the estimator puts in play, and neither
implies the other.} The bound behind \textup{(i)} is primitively a bound out of $L^{2}(\mathcal X,dx)$,
and \emph{as such it reads no condition on $\pi$}:
\begin{equation}\label{eq:hatAdx}
\big\|\hat{\mathcal A}_\varphi h\big\|_{L^{2}(\Pi)}\;\le\;
2\,(2c_C)^{1/2}\,\mathsf c_2^{2}\,c_\Pi\,h_u^{-1}h_x^{-1/2}\,\|h\|_{L^{2}(\mathcal X,dx)}
\qquad\text{for every }h\in L^{2}(\mathcal X,dx).
\end{equation}
Clause \textup{(i)} is \eqref{eq:hatAdx} read in $L^{2}(\pi)$, at the conversion cost
$\|1/\pi\|_\infty^{1/2}$ that C1(ii) makes finite. \emph{For any measurable weight $\hat\pi$ on
$\mathcal X$ with $\hat\pi\ge\varepsilon_\pi>0$ Lebesgue-almost everywhere} --- in particular the floored
estimate \cref{def:lfstep} fixes, which is \textbf{not} required to be a probability density, only the
pointwise floor being used --- the same estimate read in $L^{2}(\hat\pi)$ costs $\varepsilon_\pi^{-1/2}$
in place of $\|1/\pi\|_\infty^{1/2}$:
\begin{equation}\label{eq:hatAhatpi}
\big\|\hat{\mathcal A}_\varphi\big\|_{L^{2}(\hat\pi)\to L^{2}(\Pi)}\;\le\;
2\,(2c_C)^{1/2}\,\mathsf c_2^{2}\,c_\Pi\,\varepsilon_\pi^{-1/2}\,h_u^{-1}h_x^{-1/2} .
\end{equation}
\emph{The two readings do not imply one another, and the discrepancy is unbounded in $\varepsilon_\pi$}:
at the constant weight $\hat\pi\equiv\varepsilon_\pi$, which is the very function \cref{def:lfstep}'s
non-emptiness argument exhibits and which the construction convention there admits,
$\|h\|^{2}_{L^{2}(\hat\pi)}=\varepsilon_\pi\|h\|^{2}_{L^{2}(\mathcal X,dx)}$ exactly, whence
\begin{equation}\label{eq:pairgap}
\big\|\hat{\mathcal A}_\varphi\big\|_{L^{2}(\hat\pi)\to L^{2}(\Pi)}
\;\ge\;\big(\varepsilon_\pi\,\|1/\pi\|_\infty\big)^{-1/2}\,
\big\|\hat{\mathcal A}_\varphi\big\|_{L^{2}(\pi)\to L^{2}(\Pi)} ,
\end{equation}
and the right-hand side diverges as $\varepsilon_\pi\downarrow0$ at a fixed sample on which
$\hat{\mathcal A}_\varphi\neq0$. \emph{So a scale calibrated to one of the two pairings is not a scale
for the other}, which is why \eqref{eq:normconst} carries the \textbf{maximum} of the two conversion
costs and why \cref{def:lfstep} can then assert $\|\omega_n^{-1}\hat{\mathcal A}_\varphi\|\le1$ surely in
both. \emph{This clause is stated here and not inside \cref{def:lfstep} on \cref{rem:obligations}'s
test}: a definition may \emph{fix} a constant, but the inequality that constant has to satisfy is a
claim and belongs to a result. \emph{And it consumes \textbf{less} than \textup{(i)} does}: C1(ii) is
read by neither \eqref{eq:hatAdx} nor \eqref{eq:hatAhatpi}, and only by \eqref{eq:pairgap}, which
compares the two.
\item[(ii)] \emph{Finite rank.} $\operatorname{rank}\hat{\mathcal A}_\varphi\le n_{\mathrm m}$, so
$\hat{\mathcal A}_\varphi$ is Hilbert--Schmidt, and $\ker\hat{\mathcal A}_\varphi$ has infinite
dimension. \emph{The rank bound is $n_{\mathrm m}$ and not $n_{\mathrm m}(\mathsf L-1)$}: by
\eqref{eq:hatAcon} the operator factors through the $n_{\mathrm m}$ functionals
$h\mapsto\Delta_{h_x}h_i$, which do not depend on $l$, so the contrast index enlarges the range space
and not the rank.
\item[(iii)] The adjoint maps into $\mathcal V$. $\int_{\mathcal X}\hat\nu_{u,t,l}(x)\,dx=0$ for
every $(u,t,l)$, hence $\hat{\mathcal A}^\ast_\varphi\Xi\in\mathcal V$ for every $\Xi\in L^{2}(\Pi)$,
and \eqref{eq:hatAadj} is the adjoint of \eqref{eq:hatAcon} in the $L^{2}(\pi)$--$L^{2}(\Pi)$ pairing.
\emph{This is what C3(ii)'s boundary index is for}, through \cref{lem:intker}(iv); without it the
identity fails at \emph{every} design point, by an amount that is $O(1)$ only within $O(h_x)$ of an
endpoint of $\mathcal X$ and vanishes with $h_x$ elsewhere, and
$\hat{\mathcal A}^\ast_\varphi$ leaves the parameter space.
\item[(iv)] It is the exact Fr\'echet derivative of $\hat T_{h_x}$. The map
$\varphi\mapsto\hat T_{h_x}(\varphi,\theta)$ is twice Fr\'echet differentiable on $\mathcal V$ with
$\partial_\varphi\hat T_{h_x}(\varphi,\theta)=\hat{\mathcal A}_\varphi$ and, for all
$\varphi,\varphi+v\in\mathcal V$,
\begin{equation}\label{eq:hatTxrem}
\big\|\hat T_{h_x}(\varphi+v,\theta)-\hat T_{h_x}(\varphi,\theta)-\hat{\mathcal A}_\varphi v\big\|_{L^{2}(\Pi)}
\;\le\;\frac{\mathsf c_4\,c_\Pi}{2\,h_u^{2}}\cdot\frac{\|\Delta_{h_x}v\|_2^{2}}{n_{\mathrm m}} .
\end{equation}
\eqref{eq:hatTxrem} is sharper than the closest antecedent, and the sharpening is exactly what
generates (P-f). \citet[Prop.~4.2(a)]{centorrinofeveflorens2019} establish the same quadratic-remainder
property for their sample operator with a constant written only as ``$\hat M<\infty$'';
\eqref{eq:hatTxrem} exhibits that constant, and its bandwidth factor is $h_u^{-2}$. \emph{An unnamed finite
constant and a constant diverging like $h_u^{-2}$ support different theorems}: the second is what
\eqref{eq:effnoise} carries into the stagewise noise level and hence what forces (P-f) and, through it,
\eqref{eq:muhalf}. A treatment that leaves the constant unnamed cannot see that condition.
\item[(v)] It is not, and cannot be, a derivative of $\hat T$, the reading
\cref{def:hatA} rejects is unavailable, not merely different. The Gateaux derivative of
$\hat T(\cdot,\theta)$ at $\varphi$ in the direction $v$ exists for every $v\in\mathcal V$ and equals
\begin{equation}\label{eq:gateaux}
\mathrm d\hat T_\varphi[v](u,t,l)\;=\;\frac1{n_{\mathrm m}}\sum_{i\in\mathcal I_m}
\widehat{\widetilde W}_{i,l}\,w\big(t'\Phi(\Delta Z_i)\big)\,\frac1{h_u}
C_{h_u}\big(u-\Delta U_i(\theta,\varphi),\IR\big)\,\Delta v_i,
\qquad \Delta v_i=v(X_{i2})-v(X_{i1})
\end{equation}
, the \emph{raw} residual and the \emph{raw} difference, which is what distinguishes it from
\eqref{eq:hatAcon}, 
but $v\mapsto\mathrm d\hat T_\varphi[v]$ is \emph{unbounded on $\mathcal V$}: almost surely, the only
exception event being a tie among the $2n$ points $X_{ij}$, which is null because $X$ is continuously
distributed by \cref{def:model}, there is a sequence $v_k\in\mathcal V$ with
$\|v_k\|_{L^{2}(\pi)}\to0$ and $\|\mathrm d\hat T_\varphi[v_k]\|_{L^{2}(\Pi)}\ge\mathsf c>0$.
Consequently $\hat T(\cdot,\theta)$ is \emph{nowhere Fr\'echet differentiable} on $\mathcal V$, and no
bounded operator is its derivative. \emph{No event involving $\hat\varrho$ is excluded here}, because
\eqref{eq:wbounds} holds surely.
\end{enumerate}
\end{proposition}

% ---------------------------------------------------------------------------
%  H.2  (O4), VERIFICATION HALF.  THE HEADLINE IS NEGATIVE.
% ---------------------------------------------------------------------------

\begin{proposition}[A6 is false, and not only for the operator it names]\label{prp:tccfails}
Suppose A0, A4(i), A4(iii), B1(i), (ii), (iv), C1(ii), C3 and C5$'$(i) hold, and, for clause (iii)
only, C2$'$(i), and adopt
\cref{def:hatT}, in particular its fold split and its truncation, so that \eqref{eq:wbounds} holds
surely. Let $\hat T$ be \eqref{eq:hatT} and let $\eta_{\mathrm{TCC}}<1$. Then, on the event that the
$2n$ points $X_{11},X_{12},\dots,X_{n2}$ are pairwise distinct, an event of probability one, because
$X$ is continuously distributed by \cref{def:model}, and for all $n$ large enough that
$h_u<1/\big(2c_\Theta(\epsilon)\big)$ for some $\epsilon<\tfrac14$, with $c_\Theta$ the envelope-tail
radius of \cref{lem:intker}(i):
% DISPATCH-26.  The hypothesis read h_u < 1/(2 c_C) through dispatch 25, c_C then
% being the kernel's support half-width, and Step 3 used lem:intker(i)'s EXACT
% C(w+1) - C(w) = 1 - 0 = 1.  Under the envelope the exact values are gone, so the
% hypothesis is restated at the epsilon-tail radius and Step 3 gets
% C(w+1) - C(w) >= 1 - 2 epsilon > 1/2, which is all it needs -- Step 3 asks for
% STRICT POSITIVITY and nothing more.  C5'(i) delivers h_u -> 0, so the hypothesis
% still holds for all large n at any fixed epsilon.
%   THIS PROPOSITION IS DELIBERATELY *NOT* CONDITIONED ON COMPACT SUPPORT, and
% that is the whole point of repairing it in place rather than adding a clause.
% prp:tccfails is the NEGATIVE result that A6 is false for T-hat; it is what
% licenses the entire regularisation architecture.  Conditioning it on a
% compactly supported kernel would make it say nothing whatever about the
% estimator the paper actually runs, which uses k_order(., 8).  A negative result
% about a kernel nobody uses is not a negative result.
\begin{enumerate}
\item[(i)] \emph{No bounded operator serves.} For \emph{every} $\tilde\varphi\in\mathcal V$, every
$\bar\rho>0$, and \emph{every bounded linear operator} $\mathcal Q:\mathcal V\to L^{2}(\Pi)$ there
exists $\varphi\in\mathcal V$ with $\|\varphi-\tilde\varphi\|_{L^{2}(\pi)}\le\bar\rho$ and
\begin{equation}\label{eq:tccviol}
\big\|\hat T(\varphi,\theta)-\hat T(\tilde\varphi,\theta)-\mathcal Q(\varphi-\tilde\varphi)\big\|_{L^{2}(\Pi)}
\;>\;\eta_{\mathrm{TCC}}\,\big\|\hat T(\varphi,\theta)-\hat T(\tilde\varphi,\theta)\big\|_{L^{2}(\Pi)} .
\end{equation}
\emph{The candidate linearisation is $\mathcal Q$ and not $\mathcal L$}, which in Part~H names the
\emph{one specific} normalised compression
$\omega_n^{-1}\hat{\mathcal A}_{\varphi_\circ}\vert_{\mathcal S_n}$ of \cref{def:lfstep} and
\cref{lem:scaleinv}, an operator that is bounded by construction, whereas the whole point of clause
(iii) below is a candidate that is \emph{not}. $\mathcal Q$ carries no properties beyond those the
statement gives it.
% ARCHAEOLOGY.  cal Q joins lem:srctransfer(ii)'s cal Q_1, cal Q_2 as this file's
% letter for a generic operator in a quantified statement.  Glossary C70.
\item[(ii)] \emph{In particular A6 is false.} Taking $\mathcal Q=\hat{\mathcal A}_{\tilde\varphi}$ of
\eqref{eq:hatAcon}, which is bounded by \cref{prp:hatAprop}(i), the event in A6 has probability $0$, 
for every $n$ large enough, every $\eta_{\mathrm{TCC}}<1$, and every ball of positive
$L^{2}(\pi)$-radius.
\item[(iii)] And the one unbounded candidate does not serve either. Let $\mathcal Q$ be the
Gateaux derivative $\mathrm d\hat T_{\tilde\varphi}[\cdot]$ of \eqref{eq:gateaux}, the reading
\cref{def:hatA} rejects, which by \cref{prp:hatAprop}(v) is unbounded and so is \emph{not} covered by
(i). Fix $i^\ast:=\min\mathcal I_m$. Then \eqref{eq:tccviol} holds at the same $\varphi$, for all $n$
large enough, at every $\tilde\varphi$ for which
\begin{equation}\label{eq:lebpt}
\liminf_{h\downarrow0}\ \frac1h\int_{\IR}C_{h}\big(u-\Delta U_{i^\ast}(\theta,\tilde\varphi),1\big)^{2}\,
\Pi_u(du)\;>\;0 .
\end{equation}
\emph{At each fixed $\tilde\varphi$, \eqref{eq:lebpt} holds with probability one}: under C2$'$(i)
the law of $\Delta U_{i^\ast}(\theta,\tilde\varphi)$ is absolutely continuous, Lebesgue-almost every
point of $\IR$ is a Lebesgue point of $\varpi_u$, and at such a point the displayed limit equals
$\varpi_u(\Delta U_{i^\ast})\int\Theta_1^{2}>0$ by B1(iv.a), $\Theta_1$ being the interior rescaled
kernel \eqref{eq:rescaled} of C3, whose \textbf{envelope} bound C3(ii) makes free of $h$ --- \emph{C3(ii)
posits an envelope and states no support clause}.
\end{enumerate}
A6 is not a strong assumption; it is a false one, and the falsity is a property of $\hat T$
alone.
\emph{Three qualifiers scope this proposition and all three are load-bearing}: it is about $\hat T$
\textbf{on $\mathcal V$}, \textbf{in the levels parameterisation}, and \textbf{off the sieve}. The sieve
qualifier is \cref{rem:tccwhy}\textup{(3)}; \emph{the parameterisation qualifier is
\cref{rem:tccwhy}\textup{(4)} and is the sharper of the two}, because in a space of \emph{derivatives} point
evaluation is bounded, \cref{prp:hatAprop}(v)'s mechanism fails, and the proposition dissolves, at a cost
in the source index that \eqref{eq:muhalf} will not pay.
\end{proposition}

\begin{remark}[Why it fails, what the failure is not, and where it stops]\label{rem:tccwhy}
(1) The mechanism, in one sentence. A tangential cone condition asks a \emph{second-order}
quantity to be dominated by a \emph{first-order} one. Pairing $\hat T$, which sees $\varphi$ through
point evaluations, with $\hat{\mathcal A}_{\tilde\varphi}$, which sees it through the smoothed
functionals \eqref{eq:smdiff}, leaves a \emph{first-order} discrepancy
$\mathrm d\hat T_{\tilde\varphi}[v]-\hat{\mathcal A}_{\tilde\varphi}v$ on the left-hand side, and
by \cref{prp:hatAprop}(v) that discrepancy is unbounded relative to $\|v\|_{L^{2}(\pi)}$ while the
right-hand side is not bounded below by anything. The two readings \cref{def:hatA} separates are
not interchangeable, and C3's bandwidth forces the separately-smoothed one; under that reading A6 pairs
two objects whose difference is first order. The defect is therefore \emph{in the pairing}, not in the
constant $\eta_{\mathrm{TCC}}$ and not in the radius of $\mathcal V_n$, which is why
\cref{prp:tccfails} is stated at every $\eta_{\mathrm{TCC}}<1$ and every $\bar\rho>0$.

(2) The refutation does not depend on the construction, and that matters. Steps~1--4 of the
proof use nothing about $\hat{\mathcal A}_{\tilde\varphi}$ beyond boundedness: a spike of vanishing
$L^{2}(\pi)$-norm is annihilated by \emph{any} bounded operator, while $\hat T$ moves by a fixed
nonzero amount. \emph{So the obstruction is not that \cref{def:hatAcon} chose the wrong linearisation};
no choice of bounded linearisation whatsoever satisfies A6, and \cref{prp:tccfails}(iii) closes the one
unbounded candidate. This is why the disposition is retirement rather than re-specification, and it is
what immunises the refutation against the deviation \cref{def:hatAcon} itself flags, that the
residual inside \eqref{eq:hatg} carries $\Delta_{h_x}\varphi_i$ and not $\Delta\varphi_i$. That
deviation changes which bounded operator is used; \cref{prp:tccfails}(i) quantifies over all of them.

(3) Where the refutation stops: the sieve, and it stops there completely. \Cref{prp:tccfails} is
a statement about $L^{2}(\pi)$-balls in $\mathcal V$, and \cref{def:lfstep} does not run there, it
runs on $\mathcal V_n\cap\mathcal S_n$. On the sieve the spike is not damped; it is not
representable. On A8's quasi-uniform mesh the narrowest feature of $\mathcal S_n$ has width
$\asymp L_x^{-1}$ while the design spacing is $\asymp n^{-1}$, so every nonzero element of
$\mathcal S_n$ has, in the support of its narrowest feature, of order $n/L_x\gg n^{1/2}$ design points;
A8's upper end $L_x\ll n^{1/2}$ is what forces the margin. The inverse inequality
$\|\zeta\|_\infty\le\mathsf c_{11}L_x^{1/2}\|\zeta\|_{L^{2}(\pi)}$ of \cref{lem:unifop}(3d) is the same
fact read as a norm bound. So a cone condition restricted to the sieve is untouched by
\cref{prp:tccfails}, and it is also unproved. The first-order limit below is taken in A6, the
condition on the raw $\hat T$, and \textbf{not} in \eqref{eq:a6b}. Put $\varphi=\tilde\varphi+cv$ with
$v\in\mathcal S_n$
and let $c\downarrow0$. Since $\hat T(\cdot,\theta)$ is Gateaux differentiable in every direction
% ARCHAEOLOGY.  A round-9 draft cited eq:a6b rather than A6 at this display.
(\cref{prp:hatAprop}(v)) and $\hat{\mathcal A}_{\tilde\varphi}$ is linear, A6's inequality divided by
$c$ passes to the limit and becomes
$\|\mathrm d\hat T_{\tilde\varphi}[v]-\hat{\mathcal A}_{\tilde\varphi}v\|_{L^{2}(\Pi)}
\le\eta_{\mathrm{TCC}}\|\mathrm d\hat T_{\tilde\varphi}[v]\|_{L^{2}(\Pi)}$, which at the base point
$\tilde\varphi=0$, where the raw and the smoothed period differences of the base point coincide, both
being zero, so that \eqref{eq:gateaux} and \eqref{eq:hatAcon} carry the \emph{same}
$\hat g_{\cdot,\cdot,\cdot,i}$, is
\begin{equation}\label{eq:a6sieve}
\sup_{v\in\mathcal S_n\setminus\{0\}}
\frac{\big\|\sum_{i\in\mathcal I_m}\hat g_{\cdot,\cdot,\cdot,i}(0,\theta)\,\big(\Delta v_i-\Delta_{h_x}v_i\big)\big\|_{L^{2}(\Pi)}}
{\big\|\sum_{i\in\mathcal I_m}\hat g_{\cdot,\cdot,\cdot,i}(0,\theta)\,\Delta v_i\big\|_{L^{2}(\Pi)}}
\;\le\;\eta_{\mathrm{TCC}} .
\end{equation}
Why not \eqref{eq:a6b}, and the answer is that there the same limit is empty. By
\cref{prp:hatAprop}(iv) the operator $\hat{\mathcal A}_{\tilde\varphi}$ is the \textbf{exact} Fr\'echet
derivative of $\hat T_{h_x}$, with remainder $O(c^{2})$ by \eqref{eq:hatTxrem}; \eqref{eq:a6b}'s
right-hand side is $O(c)$; so dividing by $c$ and letting $c\downarrow0$ gives
$0\le\tfrac{\eta}{1+\eta}\|\hat{\mathcal A}_{\tilde\varphi}v\|$, which is true of every pair and says
nothing. \emph{The pairing defect A6 has and A6b does not is precisely what survives the limit}, and it
is why the display must be read off A6.
\emph{Two scope statements.} \textbf{(a) The base point.}
\eqref{eq:a6sieve} is written at $\tilde\varphi=0$. At a general base point the two displays run on
\emph{different} kernels, \eqref{eq:gateaux} evaluates $C_{h_u}$ at the raw residual
$u-\Delta Y_{i,\theta}+\Delta\tilde\varphi_i$ and \eqref{eq:hatg} at the smoothed one
$u-\Delta Y_{i,\theta}+\Delta_{h_x}\tilde\varphi_i$, and the two arguments differ by
$|\Delta\tilde\varphi_i-\Delta_{h_x}\tilde\varphi_i|\le\mathsf c_{12}\mathsf R\,h_x^{s-1/2}$, which is
\eqref{eq:xsmoothbias}, \emph{whose justification is the $\lfloor s-\tfrac12\rfloor$-th order Taylor
expansion of \cref{lem:unifop} Step~4 and not the vacuous embedding
$H^{s}\hookrightarrow C^{0,s-1/2}$}, and which therefore also inherits the hypothesis
$p_K\ge s-\tfrac12$, and which is
$O\big(h_u\hat\delta_n\big)$ under (P-g). Since $h_u^{-1}\Theta_1(\cdot/h_u)$ is Lipschitz at scale
$h_u$, the resulting relative displacement of the two families is $O(\hat\delta_n)$, so
\eqref{eq:a6sieve} is the exact limit at $\tilde\varphi=0$ and the limit up to that factor elsewhere.
% QUOTATION MOVED WITH ITS SOURCE.  Through dispatch 25 (P-g) read
% R h_x^{s-1/2} = O(hat-delta_n) and this passage read the displacement off it as
% O(h_u^{-1} hat-delta_n).  THE h_u^{-1} PRINTED HERE IS WHAT REFUTED THE OLD (P-g):
% the same factor appears in lem:unifop Step 4's operator-norm bound, so the
% hypothesis had to carry h_u on the other side.  With (P-g) corrected the two
% factors cancel and the displacement is O(hat-delta_n) -- a STRENGTHENING of this
% clause's scope statement, obtained at no cost, since the qualitative claim (the
% limit is exact at the base point and holds up to a vanishing factor elsewhere)
% never depended on which power of h_u the factor carried.
\textbf{(b) ``Only if'', not ``if and only if''.} \eqref{eq:a6sieve} is \emph{necessary} for a cone
condition on $\mathcal V_n\cap\mathcal S_n$, being its first-order shadow. It is not shown sufficient:
recovering A6 at finite $c$ from its first-order limit needs the second-order remainder
\eqref{eq:hatTxrem} to be dominated by $\|\sum_i\hat g_{\cdot,\cdot,\cdot,i}\Delta v_i\|$ uniformly on the
sieve, which is a \emph{lower} bound on that quantity in terms of $\|v\|_{L^{2}(\pi)}$, i.e.\ a bound
on $\lambda_{\min}(\mathbf G_n(\cdot))$, the same object \cref{prp:tccrepair}(iii) declines to supply
and Q-3 records as open. The paragraph below declines it too, so the display is stated as an ``only
if''.
Two things should be read off \eqref{eq:a6sieve}, and the second is why it is not pursued.
First, it is a \emph{first-order} condition and not a curvature condition, the quadratic remainder
scales away, so ``A6 on the sieve'' is not the tangential cone condition restricted but the
statement that \emph{$x$-smoothing acts as a near-identity on $\mathcal S_n$}. Second,
\emph{\eqref{eq:a6sieve} turns on the product $h_xL_x$, which no assumption in this block constrains}:
C5$'$ places no upper bound on $h_x$ (see (P-g)) and A8 no joint restriction with it, and the block's
admissible window at the frontier-clearing configuration $(\mu,s,a,\alpha)=(2,3,\tfrac34,0.32)$ leaves
$h_xL_x$ free over a range containing both $n^{-0.21}\to0$ and $n^{0.19}\to\infty$, \emph{both
verdicts}. Establishing it would cost \emph{a new joint budget
line} $h_xL_x=O(\eta_{\mathrm{TCC}})$ \emph{plus} a lower bound on $\lambda_{\min}(\mathbf G_n)$, which
is the same object \cref{prp:tccrepair}(iii) declines to bound and Q-3 records as open. Neither is
charged here, because \cref{prp:phirate} needs no cone condition at all.
% ARCHAEOLOGY.  Clause (1) survives verbatim on the sieve; what changes there is
% that the pairing defect becomes finite rather than infinite.  Shrinking bar-rho
% does not help; what helps is h_x L_x, since Delta v_i - Delta_{h_x} v_i is the
% x-smoothing bias of v and is O(h_x L_x ||v||_inf) on a sieve of mesh L_x^{-1}.
% The window collects A8, eq:entropysieve, (P-c), C5'(ii) and the hypothesis of
% lem:unifop Step 4.  The new budget line reads "the x-bandwidth may not be coarser
% than the knot spacing".

(4) Where \cref{prp:tccfails} stops, second reading: the \textbf{parameterisation}. And this is the
one that matters for the programme, because the parameterisation that dissolves the proposition is exactly
the one that cannot feed \cref{thm:rootn}. Clause (3) records that the refutation is a statement about
$\hat T$ on $\mathcal V$ and says nothing about $\mathcal V_n\cap\mathcal S_n$. There is a second,
sharper limit of the same kind, and it is a limit on the \emph{space} rather than on the sieve.
\Cref{prp:hatAprop}(v) refutes Fr\'echet differentiability of $\hat T$ because $\hat T$ sees $\varphi$
through point evaluations and point evaluation is unbounded on $L^{2}(\pi)$.
\citet{centorrinofeveflorens2019} face the same construction and \emph{do} assume differentiability of
their sample operator, legitimately, because their parameter space is a space of \textbf{derivatives}, in
which point evaluation \emph{is} bounded: for $\varphi$ with $\varphi(\underline x)=0$,
$|\varphi(x)|=|\int_{\underline x}^{x}\varphi'|\le|\mathcal X|^{1/2}\|\varphi'\|_{L^{2}}$ by the
Cauchy--Schwarz inequality. So the mechanism of \cref{prp:hatAprop}(v), and with it
\cref{prp:tccfails}, dissolves under the derivative parameterisation. \emph{And it cannot be adopted}: by
\cref{lem:levels} the range condition A2 is the same in the two parameterisations, but the \emph{source
index} is not, composing with $\mathcal D^{-1}$ raises the degree of ill-posedness from $a$ to $a+1$ and
lowers the index from $s/(2a)$ to $(s-1)/(2(a+1))$, which at the frontier-clearing configuration
$(s,a)=(3,\tfrac34)$ is $\mu=0.571$ and \textbf{fails} \eqref{eq:muhalf}. \emph{The scope statement is
therefore:} \cref{prp:tccfails} is a
statement about $\hat T$ \emph{on $\mathcal V$, in the levels parameterisation, off the sieve}; each of the
three qualifiers can be removed, and removing any one of them costs something the programme needs, the
sieve costs $h_xL_x$ and $\lambda_{\min}(\mathbf G_n)$ (clause (3)), and the parameterisation costs
$\tfrac83$ units of $s$, uniformly in $a$, which \eqref{eq:muhalf} will not pay.

(5) One sentence about \cref{prp:norangeinv}'s reach, because it is narrower than its title
suggests. Range invariance \emph{implies} the relative-ill-posedness condition along the path that
\citet[Assumption (18)]{centorrinofeveflorens2019} run their nonlinear rate on, and the implication does
\textbf{not} reverse. \emph{So \cref{prp:norangeinv} spends the range-invariance route and leaves the
weaker condition open}: a tangential-cone-free treatment built on (18) rather than on A6 is not closed by
anything in this file. It would still need a lower bound on $\lambda_{\min}(\mathbf G_n)$ (Q-3), which is
why it is recorded as the most promising unexplored route in the neighbourhood and not as an omission.
\end{remark}

\begin{definition}[A6b: restricted curvature (a \emph{property}, not maintained)]\label{def:a6b}
For $\varphi_\circ\in\mathcal V$ and $\rho>0$ write
\begin{equation}\label{eq:ball}
\mathcal V(\varphi_\circ,\rho)\;:=\;\big\{\varphi\in\mathcal V:\ \|\varphi-\varphi_\circ\|_{L^{2}(\pi)}\le\rho\big\}
\end{equation}
for the closed $L^{2}(\pi)$-ball in $\mathcal V$. \emph{The two arguments are the separation}: bare
$\mathcal V_n$ is reserved, throughout Part~H, for the \emph{source-condition} ball \eqref{eq:srcball},
which is a different object and is not a ball of $L^{2}(\pi)$ at all.

Let $\eta\in(0,\tfrac12)$. The pair $(\hat T_{h_x},\hat{\mathcal A})$ \emph{satisfies A6b at
$(\eta,\varphi_\circ,\rho)$} if
\begin{equation}\label{eq:a6b}
\big\|\hat T_{h_x}(\varphi,\theta)-\hat T_{h_x}(\tilde\varphi,\theta)-\hat{\mathcal A}_{\tilde\varphi}(\varphi-\tilde\varphi)\big\|_{L^{2}(\Pi)}
\;\le\;\frac{\eta}{1+\eta}\,\big\|\hat{\mathcal A}_{\tilde\varphi}(\varphi-\tilde\varphi)\big\|_{L^{2}(\Pi)}
\qquad\text{for all }\varphi,\tilde\varphi\in\mathcal V(\varphi_\circ,\rho) .
\end{equation}
\emph{Note the quantifier}: \eqref{eq:a6b} ranges over \emph{all pairs} in the ball, so the base point
$\tilde\varphi$ at which $\hat{\mathcal A}$ is evaluated varies too. That is what
\cref{prp:tccrepair}(ii) must deliver, and it is why its hypothesis is a bound on
$\lambda_{\min}(\mathbf G_n(\cdot))$ \emph{uniform over base points} and not at one point.

\emph{A6b is NOT MAINTAINED}, on the A2c model: it is a property that a pair may or may not have,
it has exactly one consumer (\cref{prp:tccrepair}), it is never asserted, and \emph{no theorem in this
programme uses it}, \cref{prp:phirate} takes a route that needs no cone condition at all. It is
stated so that the paper can say what A6 should have said, and so that the simulation grid has a
testable object (probe CP-8).
\end{definition}

\begin{proposition}[A6b implies a tangential cone condition, and holds on an explicit ball]\label{prp:tccrepair}
Suppose the hypotheses of \cref{prp:hatAprop} hold, A4(i) so that \eqref{eq:wbounds} holds, fix $\theta$
and $\varphi_\circ\in\mathcal V$, and
% A4(i) ADDED TO THIS LINE 2026-08-10 (Phase 6a), IN prp:norangeinv's EXACT WORDING.
% prp:hatAprop's own line scopes A4(i) to ITS CLAUSE (v) alone, so "the hypotheses
% of prp:hatAprop" does NOT import it; but clause (iii) below argues that
% eq:wbounds "rules this out surely" -- the zero-row degeneracy that would make
% clause (ii) vacuous -- and eq:wbounds' LOWER half rests on A4(i).  As stated the
% proposition was therefore not implied by its listed hypotheses: the prp:hatAprop
% species exactly.  THE SIBLING ALREADY HAD THE RIDER.  prp:norangeinv reads
% "hypotheses of prp:hatAprop, A4(i) so that eq:wbounds holds"; Phase 4c's
% propagation note names BOTH as inheritors of that repair, and the pair had
% diverged since.  The wording is copied verbatim so that the two lines can be
% compared by eye.
for $\varphi\in\mathcal V$ let
\begin{equation}\label{eq:gram}
\mathbf G_n(\varphi)\;:=\;\Big(\big\langle \hat g_{\cdot,\cdot,\cdot,i}(\varphi,\theta),\,
\hat g_{\cdot,\cdot,\cdot,k}(\varphi,\theta)\big\rangle_{L^{2}(\Pi)}\Big)_{i,k\in\mathcal I_m}
\end{equation}
be the Gram matrix of the $n_{\mathrm m}$ functions \eqref{eq:hatg} in $L^{2}(\Pi)$, at base point
$\varphi$. Then:
\begin{enumerate}
\item[(i)] A6b delivers A6 for the consistent pair. If $(\hat T_{h_x},\hat{\mathcal A})$
satisfies A6b at $(\eta,\varphi_\circ,\rho)$ then A6 holds verbatim on
$\mathcal V(\varphi_\circ,\rho)$ with $\hat T$ replaced by $\hat T_{h_x}$ and
$\eta_{\mathrm{TCC}}=\eta$, deterministically on the same event; and, on that event, for all
$\varphi,\tilde\varphi\in\mathcal V(\varphi_\circ,\rho)$,
\[
(1-\eta)\big\|\hat T_{h_x}(\varphi,\theta)-\hat T_{h_x}(\tilde\varphi,\theta)\big\|
\;\le\;\big\|\hat{\mathcal A}_{\tilde\varphi}(\varphi-\tilde\varphi)\big\|
\;\le\;(1+\eta)\big\|\hat T_{h_x}(\varphi,\theta)-\hat T_{h_x}(\tilde\varphi,\theta)\big\| .
\]
\item[(ii)] A sufficient condition, with an explicit radius. Fix $\rho_0>0$ and suppose
\begin{equation}\label{eq:gramunif}
\underline\lambda_n\;:=\;\inf_{\varphi\in\mathcal V(\varphi_\circ,\rho_0)}
\lambda_{\min}\big(\mathbf G_n(\varphi)\big)\;>\;0 ,
\end{equation}
and let
\begin{equation}\label{eq:tccrad}
\bar\rho_n\;:=\;\frac{\eta}{1+\eta}\cdot
\frac{h_u^{2}\,h_x^{1/2}}{(2c_C)^{1/2}\,\mathsf c_2\,\mathsf c_4\,c_\Pi\,\|1/\pi\|_\infty^{1/2}}
\cdot\frac{\underline\lambda_n^{1/2}}{n_{\mathrm m}^{1/2}} .
\end{equation}
Then $(\hat T_{h_x},\hat{\mathcal A})$ satisfies A6b at
$\big(\eta,\varphi_\circ,\tfrac12\min\{\bar\rho_n,\rho_0\}\big)$.
\emph{The halving is not cosmetic}: \eqref{eq:a6b} quantifies over pairs, so two points of a ball of
radius $\tfrac12\bar\rho_n$ are at distance at most $\bar\rho_n$, which is what \eqref{eq:tccrad}
governs; and the base point of $\hat{\mathcal A}_{\tilde\varphi}$ moves with the pair, which is what
\eqref{eq:gramunif}'s infimum pays for. A bound at the single base point $\varphi_\circ$ does
\emph{not} suffice, and rounds before this one proved \eqref{eq:a6b} only there.
\item[(iii)] And the radius is governed by the smallest singular value, which is where the
ill-posedness sits. At each $\varphi$, $n_{\mathrm m}^{-1}\lambda_{\min}(\mathbf G_n(\varphi))^{1/2}$
is exactly the smallest ratio
$\|\hat{\mathcal A}_{\varphi}v\|_{L^{2}(\Pi)}\big/\|\Delta_{h_x}v\|_{2}$, i.e.\ the injectivity
modulus of $\hat{\mathcal A}_{\varphi}$ read in the $\Delta_{h_x}$ coordinates, so
\eqref{eq:tccrad} shrinks with the ill-posedness of the empirical linearisation. \emph{No conclusion of
this file rests on $\bar\rho_n$ being bounded away from zero, and none is claimed}; whether it is, is a
design question about the conditioning of $\mathbf G_n$, the same species of question A8's upper end
answers for the balancing solve, and it is the object probe CP-8 should report, \emph{at pairs not
anchored at the centre}. \emph{One degeneracy is excluded by construction and would otherwise be
common}: if $\widehat{\widetilde W}_{i,l}=0$ for some $i$ and \emph{every} $l$ then
$\hat g_{\cdot,\cdot,\cdot,i}\equiv0$, row and
column $i$ of $\mathbf G_n(\varphi)$ vanish and $\lambda_{\min}=0$ exactly, so \eqref{eq:tccrad} returns
$\bar\rho_n=0$ and (ii) is vacuous. \eqref{eq:wbounds} rules this out \emph{surely}, and \emph{a fortiori
at $\mathsf L\ge3$}: it bounds every component away from zero, while the degeneracy needs all of them to
vanish at once. It does not rule
out near-collinearity of the $\hat g_{\cdot,\cdot,\cdot,i}$, which is the residual content of Q-3.
\end{enumerate}
\end{proposition}

% ############################################################################
% ##  rem:tccverdict.  SHORTENING PASS 2026-08-01: this remark was 49 compiled
% ##  lines RECORDING A DECISION -- what obligation (O4) asked, that it was
% ##  executed in assumptions.tex "rather than requested of it", and which four
% ##  statements about A6 were swept in the same pass.  That is process record.
% ##  What is kept compiled is the disposition itself: A6 refuted and retired,
% ##  A6b is what it should have said, nothing downstream needs either, and A7 is
% ##  verified only in an amended form whose cost is named.
% ##  ARCHAEOLOGY.
% ##   -- (O4) asked for two things: exhibit hat-A_varphi, and verify A6 and A7
% ##      for the object exhibited.  Both done, answers differ.  hat-T had to be
% ##      exhibited too, because A6 and A7 are statements ABOUT hat-T and hat-T
% ##      was defined in no file either; eq:hatT now sits in the maintained block
% ##      beside A7.
% ##   -- The retired A6 is kept, struck, at ass:tcc, so that a reader meeting a
% ##      round-9 citation of A6 finds the reason rather than a gap.
% ##   -- Four statements ABOUT A6 elsewhere in these files were swept in the same
% ##      pass, two of them substantively false once A6 is struck: rem:R1(iv), the
% ##      Discussion paragraph of the assumption block, rem:draftrec(iii) and
% ##      rem:quadremainder.
% ##   -- The strategy memo capped the range-invariance attempt at "for hat-T, not
% ##      T, one round only"; the round is spent at prp:norangeinv.
% ##   -- The A7 amendment was MANDATORY rather than optional, because
% ##      def:lfstep(c) already implements the amended level and leaving A7
% ##      unamended would have put the maintained block in contradiction with the
% ##      estimator.
% ############################################################################

\begin{remark}[The disposition of A6 and A7]\label{rem:tccverdict}
\emph{The operator is exhibited} at \eqref{eq:hatAcon}, with $\hat T$ at \eqref{eq:hatT}.
\Cref{prp:hatAprop} establishes that the exhibited object is bounded, finite rank,
adjointable into $\mathcal V$, and is the exact Fr\'echet derivative of the $x$-smoothed operator
$\hat T_{h_x}$, so \cref{def:hatA}'s reading is \emph{realised}, and its rejected
alternative is proved unavailable rather than merely declined.

\emph{A6 is not verified; it is refuted}, by \cref{prp:tccfails}. Three consequences.
\begin{enumerate}
\item[(1)] A6 has left the maintained block. An assumption whose event has probability zero
cannot be maintained. It is not repairable by lowering $\eta_{\mathrm{TCC}}$, by shrinking the ball, or
by re-specifying $\hat{\mathcal A}$: \cref{prp:tccfails} quantifies over all three. The retired
statement is kept, struck, at A6. \emph{Its scope is recorded with it}: refuted on
$L^{2}(\pi)$-balls, untouched on $\mathcal V_n\cap\mathcal S_n$, where the question is
\eqref{eq:a6sieve} and is open, see \cref{rem:tccwhy}(3).
\item[(2)] \emph{What A6 should have said is A6b}, \cref{def:a6b}, which pairs $\hat{\mathcal A}$ with
the operator it actually differentiates. \Cref{prp:tccrepair} proves A6b non-vacuous and gives an
explicit radius. A6b is a \emph{property}, not maintained, exactly as A2b and A2c are.
\item[(3)] \emph{Nothing downstream needs either.} \Cref{prp:phirate} is proved without A6 and without
A6b, by a route that treats the
nonlinearity as an addition to the noise level rather than controlling it by a cone condition;
\cref{rem:lfnonlinear} explains why that route is available here and why the inherited one is not. The
range-invariance route is spent at \cref{prp:norangeinv}: it fails, and it fails for a reason that is
proved rather than reported.
\end{enumerate}

\emph{A7 is verified in an amended form}, by \cref{prp:stopfeas}: at a fixed $\tau_{\mathrm{DP}}$ for the
component of the noise generated by the smoothed indicator and the empirical average, and \emph{not} at
a fixed $\tau_{\mathrm{DP}}$ for the component generated by $\hat\varrho$, because A4(iii) is an $O_P$
statement and an $O_P$ statement carries no fixed-constant exceedance bound. The amendment
\eqref{eq:a7amend} does not cost nothing. It is free
against \eqref{eq:frontier}, whose strictness absorbs the extra $(\ln n)^{2\mu/(2\mu+1)}$
(\cref{prp:stopfeas}(iv)); it is \emph{not} free against (P-f), which is an exact domination in which
both sides now carry logarithms, and it is what turns (P-f)'s thresholds strict, 
see \cref{rem:h6bandwidth}.
\end{remark}

\begin{proposition}[The range-invariance route, spent]\label{prp:norangeinv}
Range invariance in the sense of \citet[(2.6)--(2.7)]{hanke1995} would require a family of bounded
operators $\mathcal O(\varphi,\tilde\varphi)$ on $L^{2}(\Pi)$, with
$\|\mathcal O(\varphi,\tilde\varphi)-\mathrm{Id}\|\le\mathsf c\|\varphi-\tilde\varphi\|_{L^{2}(\pi)}$, such that
$\hat{\mathcal A}_\varphi=\mathcal O(\varphi,\tilde\varphi)\,\hat{\mathcal A}_{\tilde\varphi}$ for all
$\varphi,\tilde\varphi$ in a ball $\mathcal V(\tilde\varphi,\rho)$ of \eqref{eq:ball}.
The intertwining family is written $\mathcal O$, and neither $\Omega$ nor $\mathcal R$.
Calligraphic $\mathcal R(\cdot)$
is the \emph{range} throughout this project, it appears in the next clause four lines below, and
sans-serif $\mathsf R$ is already the source radius of \eqref{eq:srcball}. \emph{And $\Omega$ is the
conditional contrast covariance \eqref{eq:omegadef}, an $\IR^{(\mathsf L-1)\times(\mathsf L-1)}$
matrix, which \eqref{eq:omegaprefac} carries as a prefactor and \cref{rem:htread} inverts}: that is the
only object the symbol names anywhere in this corpus, so no disambiguation rule attaches to it, and the
one that formerly did, the two families having been separated by type and not by arity, both being
written bare at some sites, is discharged by the renaming performed here.
$\mathcal O$ is free: it occurs at no site in the
corpus outside this proposition and its proof. \emph{It is not Landau's $O$}, which this project writes
italic and never calligraphic. \emph{Wherever it \textbf{denotes} a member of the family it carries its
two arguments}; the bare occurrences in this paragraph name the symbol rather than apply it, and the only
bare \emph{applications} are inside the proof,
where the opening line binds $\mathcal O:=\mathcal O(\varphi,\varphi')$ once and the abbreviation runs to
the end of that proof and no further.
% ARCHAEOLOGY.  Glossary C66.  Through dispatch 25 the family was written bare
% Omega, colliding with eq:omegadef's contrast covariance; assumptions_main.tex's
% type declaration recorded the collision as a "declared case pair" and deferred
% the rename to "whenever its own file is next opened".  This is that pass.
%
% CROSS-FILE INCONSISTENCY CLOSED 2026-08-09, PHASE 4 ROUND 2.  The block that
% stood here declared the half-rename as an accepted, temporary price: proofs.tex
% still wrote bare Omega for the intertwining family, including a printed
% Omega^{-1}, so a shipped supplement carried one object under two letters and the
% string Omega^{-1} under two meanings.  THAT WAS THE STATED GROUND FOR RENAMING
% HERE AND IT WAS NOT ACHIEVED UNTIL NOW.  The five proofs.tex sites are swept in
% this pass under a scope grant covering these tokens only; they had drifted from
% the :5044/:5046/:5048/:5050/:5132 recorded above (a 5-O1 paste of 315 lines),
% which is why the sweep was done by content and not by line number.  Statement
% and proof now use one letter for one object and the "declared case pair" is a
% pair no longer.
%
% ONE CONSEQUENCE IS OUT OF THIS PHASE'S SCOPE AND IS RECORDED RATHER THAN FIXED.
% assumptions_main.tex:341-349 is PRINTED text declaring Omega "a declared case
% pair" and promising the rename "when its own file is next opened", with a
% reading rule for deciding which Omega is which.  With the rename complete that
% paragraph describes a collision that no longer exists and a rule with no live
% referent.  It is that file's to retype -- one paragraph, no mathematics -- and
% the notation-warden should see it in Phase 6 if nobody gets there first.
Suppose the hypotheses of \cref{prp:hatAprop}, A4(i) so that \eqref{eq:wbounds} holds, and that
\begin{equation}\label{eq:rirad}
\rho\;>\;\mathsf c_9\,\mathsf D_n\,h_x^{1/2},
\qquad
\mathsf c_9:=2\big(\|\pi\|_\infty/\mathsf c_{13}\big)^{1/2} ,
\end{equation}
with $\mathsf c_{13}$ the uniform lower bound \suppeq{thetasqlb} on $\int_I\Theta_I^{2}$,
where $\mathsf D_n:=\max_{i,k\in\mathcal I_m}\big|(\Delta Y_{i,\theta}-\Delta_{h_x}\tilde\varphi_i)
-(\Delta Y_{k,\theta}-\Delta_{h_x}\tilde\varphi_k)\big|$ is the spread of the $n_{\mathrm m}$
kernel centres, an $O_P(1)$ quantity.
% DISPATCH-26.  Two changes, both consequences of the envelope, both recorded.
%  (1) c_9 was (2 c_C)^{1/2} ||pi||^{1/2}, read off "int Theta^2 >= (2 c_C)^{-1}",
%      which needed a FINITE SUPPORT LENGTH.  It is now read off eq:thetasqlb's
%      c_13 = 1/(8 c_Theta(1/2)).  Same shape; the constant degrades by a factor
%      of 4 inside the square root's reciprocal, i.e. by 2 in c_9.
%      DISPATCH-27 CARRIES A FURTHER FACTOR OF 2, and it is the arithmetic that
%      change (2) below failed to carry through.  With the +4 c_C h_u margin gone,
%      D_n IS max_{i,k}|z_i - z_k| exactly, so a displacement of D_n does not
%      SEPARATE -- the proof's "varsigma* ||b||^2 = D_n exceeds max_{i,k}|...|"
%      read D_n > D_n and the displaced centre coincided with another centre on
%      the event that z_{i*} attains the maximum.  The proof now displaces by
%      2 D_n; this inequality is STRICT, so the factor lands in c_9 and nowhere
%      else.  No downstream result reads c_9.
%  (2) THE +4 c_C h_u TERM IS DROPPED FROM D_n.  It was there so that the
%      displacement would clear the WHOLE SPREAD OF THE SUPPORTS -- half-width
%      c_C h_u on each side of each of two centres.  prp:norangeinv's proof no
%      longer argues by disjoint supports: it argues by linear independence of the
%      translates, via strict positivity of F-Theta_1 at every real frequency, and
%      that needs D_n only to SEPARATE the displaced centre from the others.  So
%      the demand on eq:rirad is strictly WEAKER than before, which is the right
%      direction for a hypothesis, and the dropped term is dropped because it is
%      no longer used and not because it became inconvenient.
% THE NEXT LINE IS PRINTED TEXT AND MUST STAY OUT OF THIS COMMENT BLOCK.  It was
% absorbed into the line above on 2026-08-09 and the proposition then asserted
% (i)-(ii) with NO almost-sure and NO large-n qualifier for a full dispatch.
Then, almost surely for all $n$ large enough:
\begin{enumerate}
\item[(i)] range invariance forces
$\mathcal R(\hat{\mathcal A}_\varphi)=\mathcal R(\hat{\mathcal A}_{\varphi'})$
for every $\varphi,\varphi'\in\mathcal V(\tilde\varphi,\rho)$, first for pairs with
$\mathsf c\|\varphi-\varphi'\|_{L^{2}(\pi)}<1$, and then for \emph{all} pairs, by chaining along the
segment joining them, which \eqref{eq:ball} contains;
\item[(ii)] and that equality fails, \emph{provided} the $n_{\mathrm m}$ design weights
$\mathsf d_i(x):=h_x^{-1}[C_{h_x}(x-X_{i2},I_{i2})-C_{h_x}(x-X_{i1},I_{i1})]$, $i\in\mathcal I_m$,
are linearly independent in $L^{2}(\mathcal X,dx)$: there are then
$\varphi,\varphi'\in\mathcal V(\tilde\varphi,\rho)$ with
$\mathcal R(\hat{\mathcal A}_\varphi)\ne\mathcal R(\hat{\mathcal A}_{\varphi'})$.
\end{enumerate}
Hence no such family exists, for any $\mathsf c<\infty$, on any ball whose radius does not shrink faster
than $h_x^{1/2}$. The mechanism is that $\hat{\mathcal A}_\varphi$'s range is spanned by kernels
\emph{translated} in $u$ by $\Delta_{h_x}\varphi_i$, and a sufficiently displaced translate cannot lie
in their span. \emph{Displacement, and not disjointness of supports, is what the proof uses}: the
translates are linearly independent at distinct centres because the interior kernel's Fourier
transform is strictly positive at every real frequency, so the displacement $\mathsf D_n$ has only to
\textbf{separate} the displaced centre from the others and not to clear their whole spread.

The independence proviso in (ii), and why it is stated rather than read off distinctness of the
design. What (ii) needs is that $\hat g_{\cdot,\cdot,\cdot,i^\ast}$ lie in the \emph{range} of
$\hat{\mathcal A}_{\varphi}$ and not merely in the span of the $\hat g_{\cdot,\cdot,\cdot,i}$; by
\eqref{eq:hatAcon} that is joint surjectivity of the $n_{\mathrm m}$ functionals
$v\mapsto\Delta_{h_x}v_i=\int_{\mathcal X}\mathsf d_iv\,dx$ onto $\IR^{n_{\mathrm m}}$, which is
equivalent to linear independence of $\{\mathsf d_i\}$ in $L^{2}(\mathcal X,dx)$ and is unaffected by
the constraint $\int v\,d\pi=0$ defining $\mathcal V$, since $\int\mathsf d_i\,dx=0$ makes every
functional constant-free. \emph{Distinctness of the $2n$ points $X_{ij}$ is \textbf{not} the reason}: the
minimum design gap is $O_P(n^{-2})$, far below
$h_x$, so the bumps overlap heavily and the family is nowhere near a disjointly supported one.
% ARCHAEOLOGY.  Through round 9 the independence proviso was asserted from
% distinctness of the design.
Where distinctness genuinely enters is one step further in, and \emph{the route through a common
kernel shape is no longer available at all} --- a structural fact, and worth stating as one.
\emph{Exactness of the $j=0$ truncated moment at \textbf{every} design point is incompatible with a
design-point-invariant kernel shape whenever the kernel is not compactly supported.} Under a
compactly supported kernel the incompatibility is confined to a boundary layer: at design points
further than the support half-width from $\partial\mathcal X$ every $\Theta_{I_{ij}}$ is the one
interior kernel, each $\mathsf d_i$ is a difference of two translates of the single function
$x\mapsto h_x^{-1}C_{h_x}(x,\IR)$, its image under the Fourier transform $\mathcal F$ of
\eqref{eq:usobolev} is entire of exponential type and not identically
zero, so its zero set is discrete, and $\sum_ic_i\mathsf d_i=0$ forces
$\sum_ic_i\big(e^{-\mathrm i\xi X_{i2}}-e^{-\mathrm i\xi X_{i1}}\big)=0$ for every $\xi$, whence
$c\equiv0$ by linear independence of the characters $\xi\mapsto e^{-\mathrm i\xi x}$ at the $2n$
\emph{distinct} frequencies $X_{ij}$. Under C3 as it now stands, $I(x^\ast)$ moves with $x^\ast$ at
\emph{every} design point --- there is no interior region on which it is constant --- so the shape
$\Theta_{I_{ij}}$ varies everywhere, the convolution structure is lost everywhere, and this sufficient
condition is unavailable. \emph{That is the one genuine loss of the envelope route, and it is
recorded rather than absorbed.}
% ARCHAEOLOGY.  This paragraph derived the proviso from a convolution structure
% available only where every design point shares ONE kernel shape.  Under the
% envelope no such region exists, so the derivation is retained as a statement
% ABOUT THE COMPACT CASE and the loss is printed.  IT COSTS NOTHING, for the three
% reasons the next sentences already gave -- and note that prp:norangeinv's own
% proof no longer routes through support disjointness either: it uses linear
% independence of the translates via positivity of the interior kernel's Fourier
% transform, which is STRONGER than what this paragraph supplies and does not need
% a common shape.  So the paragraph's loss is a loss of a SUFFICIENT CONDITION for
% a PROVISO, not of any step in a proof.
Three things keep the cost at zero. It is a \emph{qualitative} nonsingularity and
not a lower bound on $\lambda_{\min}$, so it does not reintroduce the conditioning input Q-3 records as
unavailable and \cref{prp:tccrepair}(iii) declines to supply. It is a condition on clause (ii) only;
clause (i) and the population paragraph below use nothing of the kind. And \emph{no result in this
programme consumes \cref{prp:norangeinv}}, \cref{prp:phirate} needs no cone condition, so it needs no
route to one, which is why the proviso is printed as a proviso rather than promoted to a hypothesis
of the block. Probe CP-8's grid can report it directly.
\emph{The step this replaces, and why}: an earlier draft produced a direction annihilating
$n_{\mathrm m}-1$ of the functionals \eqref{eq:smdiff} and read \eqref{eq:hstep1b}, an \emph{upper}
bound on $|\Delta_{h_x}v_i|$ given $\|v\|$, backwards to bound that direction's norm. The norm of
such a direction is governed by the conditioning of $\{\mathsf d_i\}$, i.e.\ by the $\lambda_{\min}$
information Q-3 records as unavailable. The construction below uses \emph{one} explicit direction,
$\mathsf d_{i^\ast}/\pi$, whose norm and whose functional value are both computed directly, and needs no
conditioning input at all. The same argument applies verbatim to the population pair
$(T,\mathcal A_\varphi)$, where by \eqref{eq:Adef} and the shift identity of (P6) the base point enters
$\mathcal A_\varphi$ as an argument shift of the conditional density.
\end{proposition}

% ---------------------------------------------------------------------------
%  H.3  THE NOISE LEVEL.  L-18, THE OPERATOR-ESTIMATION RATE, AND A7.
% ---------------------------------------------------------------------------

\begin{lemma}[L-18: the operator error at the fixed truth]\label{lem:opnoise}
Suppose A0, A1, A4(iii), B1(i), (ii), (iv.c), (iv.d), C2$'$(ii), C3 and C5$'$(i) hold, and adopt
\cref{def:hatTexp}.
\emph{\textbf{B1(iv) is cited by clause, and this line has now moved twice}}: it read
``\textup{(iv.c)}'' alone, which was short by \textup{(iv.d)} --- cited at the
$\sum_l\varpi_{\mathsf w}(l)=1$ aggregations of Steps~2 and~4 --- and the first 2026-08-24 repair
promoted it to the whole block, which was long by \textup{(iv.a)} and \textup{(iv.b)}, neither of which
any step opens. \emph{The two clauses named here are exactly the two the proof opens}, at Steps~0, 2, 3
and~4 for \textup{(iv.c)} and Steps~2 and~4 for \textup{(iv.d)}; and the factorisation
$\Pi=\Pi_u\otimes\Pi_t\otimes\Pi_{\mathsf w}$ is B1(iv)'s \emph{declaration} and carries nothing, on the
assumption block's \textbf{three registers} paragraph. \emph{\Cref{prp:stopfeas} is repaired by this one
edit and must not be repaired a second time}: its own line is ``the hypotheses of \cref{lem:opnoise}'',
so it inherited both the gap and the over-charge, and inherits this fix.
\emph{\textbf{C2$'$ is cited by clause.}} The single C2$'$ citation in the proof is
at Step~3, to (ii), for $\mathfrak m_{p_K}$; no step reads (i)'s envelopes.
\textbf{A1b(iv) is \textbf{not} a hypothesis}, here, or at \cref{lem:aerr}, \cref{lem:unifop}
and \cref{lem:epterm}. The reason is B1(i)'s own: $\Phi$ is \emph{bounded}, so
$|w(t'\Phi(z))|\le\bar w(t)$ uniformly in
$z$ and every $\Delta Z$-integral in Part~H is an integral of a bounded summand, with no moment
condition whatsoever on $\Delta Z$. What A1b(iv) buys is nonsingularity of
$E[\Delta Z\Delta Z']$ and the $\theta$-side moments, and Part~H inverts no such matrix; its consumers
lie in Parts~F and~I --- \cref{lem:momzero} Steps~1 and 3--6, which is clauses \textup{(ii)}--\textup{(v)}
and the reading \textup{(iv.4)}; \cref{lem:prelimtheta} Steps~1 and~3; \cref{lem:biasexp} Steps~0
and~2; \cref{lem:sandwich} Steps~1, 2, 3, 5 and~6; and \cref{thm:rootn} Step~1 --- all of which retain
it. \emph{The list is illustrative and no claim of exhaustiveness is made}: the draft this sentence
replaces named three of them and was short by at least five. \emph{At \cref{lem:epterm} the ground is different and is recorded there}: its own
hypothesis (R-b) gives $E\|\Delta Z\|^{4}<\infty$ and $E[\Delta U^{4}]<\infty$, hence the second moments a
fortiori. Then $T(\vdagger,\theta_\dagger)=0$ and
% ARCHAEOLOGY.  Through dispatch 20 the hypothesis line read "C2'", which on this
% file's register rule -- a block citation of the parent tag is a citation of every
% clause under it -- charged the lemma with C2'(i)'s envelopes.  This is the one
% C2' block citation in Part H that transitivity does not already justify:
% lem:aerr and lem:unifop carry C2' in full and are RIGHT to, the first reading (i)
% at its Steps 2 and 4 and the second inheriting (ii) with the Step-3 computation
% it applies verbatim, so neither was narrowed.  A1b(iv) was carried by four lists
% and cited at no line of any of them; it is struck from all four.
\begin{equation}\label{eq:opnoise}
\big\|\hat T(\vdagger,\theta_\dagger)\big\|_{L^{2}(\Pi)}
\;=\;\underbrace{O_P\big(n^{-1/2}\big)}_{\text{empirical}}
\;+\;\underbrace{O\big(h_u^{p_K}\big)}_{\text{smoothing bias}}
\;+\;\underbrace{O_P(r_n)}_{\text{propensity}}
\;=\;O_P\big(\hat\delta_n\big),
\end{equation}
with $\hat\delta_n$ as in A5(i). Moreover the three components are separately controlled:
\begin{enumerate}
\item[(i)] Empirical, with no logarithm and with a fixed-constant tail. Writing
$\bar T$ for \eqref{eq:hatT} with $\varrho$ in place of $\hat\varrho$ and
for every $\mathsf s>0$
\[
E\big\|\bar T(\vdagger,\theta_\dagger)-E\bar T(\vdagger,\theta_\dagger)\big\|^{2}_{L^{2}(\Pi)}
\le\frac{\mathsf c_3^{2}c_\Pi^{2}}{n_{\mathrm m}},
\qquad
P\Big(\big\|\bar T-E\bar T\big\|_{L^{2}(\Pi)}\ge \mathsf s\Big)\le\frac{\mathsf c_3^{2}c_\Pi^{2}}{n_{\mathrm m}\,\mathsf s^{2}} .
\]
\emph{The $L^{2}(\Pi)$ norm needs no uniformity in $(u,t)$ and hence no logarithm}, because
$L^{2}(\Pi)$ is a Hilbert space and the summands are bounded in \emph{norm}; the logarithm A5(i)
carries is what turns the second display into a \emph{fixed-constant} statement at
$\mathsf s=\mathsf c\,n^{-1/2}\sqrt{\ln n}$, \emph{for every fixed $\mathsf c>0$}, where the right-hand
side is then $O(1/\ln n)\to0$, and that is where it earns its keep. No exponential inequality and no
Hilbert-space concentration result is used; Chebyshev suffices, which matters because no such result is
pinnable in this project's bibliography.
\emph{Sans-serif $\mathsf c$ is a generic finite positive constant} and is \textbf{not} A7's discrepancy
constant, which never appears without its subscript $\tau_{\mathrm{DP}}$; \cref{prp:stopfeas}(i)'s proof
runs the identical Chebyshev step at $\mathsf s=\tfrac12\mathsf c\,n^{-1/2}\sqrt{\ln n}$.
\emph{Bare sans-serif $\mathsf s$ is a \emph{tail level}}: a positive number bound by the quantifier of
the clause in which it appears, ``for every $\mathsf s>0$'' here, and likewise at
\cref{lem:bdiff}\textup{(i)}--(ii), free to take a different value in every such clause and never
carried from one result to another. It is separated by \textbf{font} from the two binding italic
$s$'s: the Sobolev order $s$ of H1 and the generic order $s'$ of A8.
% DISPATCH-26.  This sentence read "the THREE binding italic s's" and listed the
% boundary index s_0 as the third.  s_0 IS RETIRED -- C3's boundary index is the
% interval I in frak-I -- so the count is now two.  The letter s is correspondingly
% less crowded, which is worth noting rather than silently enjoying: it was the
% crowding on s that made sans-serif s a separation worth printing, and two of the
% three occupants remain.
% ARCHAEOLOGY.  Through dispatch 22 the constant above was bare italic tau,
% UNBOUND, in this clause's own statement, so the display was read at an undeclared
% level; the convention is C88's and the row is glossary C91.  Dispatch 22 moved
% prp:stopfeas's twelve bare-tau lines and declared Sec.1.7's rule true, but this
% site is in a different result and a different file and was not swept.
% The font separation on sans-serif s is doing real work and the corpus does not
% pretend otherwise: sans-s and H^s both occur inside lem:unifop's proof, eight
% lines apart -- the r / sans-r situation of C22, accepted there on the same ground,
% that the italic member is binding and cannot move and that the two are of
% different TYPES.  Through dispatch 22 the letter had no glossary row at all, at
% twenty-three lines across two files, two of them in STATEMENTS.  Glossary C97.
\item[(ii)] \emph{Smoothing bias.} $\|E\bar T(\vdagger,\theta_\dagger)-T(\vdagger,\theta_\dagger)\|_{L^{2}(\Pi)}
\le\mathsf c_5\,c_\Pi\,\mathfrak m_{p_K}\,h_u^{p_K}$, where
$\mathfrak m_{p_K}:=\max_{k\le p_K}\max_{l\le\mathsf L-1}\sup_t\bar w(t)^{-1}\sup_u|\partial_u^{k}\mathcal M(u,t,l)|<\infty$ by
C2$'$(ii); this is $O(\hat\delta_n)$ under C5$'$(i) as amended, and $O(n^{-1/2})$ under the unamended
bandwidth $h_u\asymp n^{-1/16}$. \emph{Only the first is needed}: \eqref{eq:opnoise} is a statement at
$\hat\delta_n$, and matching the bias to $n^{-1/2}$ rather than to $\hat\delta_n$ is what C5$'$(i)'s
amendment corrects.
\item[(iii)] \emph{Propensity.}
$\|\hat T(\vdagger,\theta_\dagger)-\bar T(\vdagger,\theta_\dagger)\|_{L^{2}(\Pi)}
\le\mathsf c_3c_\Pi\,n_{\mathrm m}^{-1}\sum_{i\in\mathcal I_m}\big\|\hat\varrho(\Delta Z_i)-\varrho(\Delta Z_i)\big\|
=O_P(r_n)$, \emph{the norm being Euclidean on $\IR^{\mathsf L-1}$ and the passage from the slicewise
bound to it costing nothing, since $\max_l|\cdot|\le\|\cdot\|$ and $\sum_l\varpi_{\mathsf w}(l)=1$; at
$\mathsf L=2$ it is the modulus earlier rounds printed}.
\end{enumerate}
\emph{Scope, per the memo's own correction to L-18}: this is a statement at the \emph{fixed} truth. It
is not sufficient for \cref{prp:phirate}, which evaluates the residual along a data-dependent path;
that is \cref{lem:unifop}.
\end{lemma}

\begin{lemma}[The operator-estimation rate, and the re-derivation C5$'$(ii) asks for]\label{lem:aerr}
\textbf{Consumers: none.} This lemma, together with \eqref{eq:gamman} and \eqref{eq:entropyA},
is consumed by no result in this file. It is
retained for two reasons, both recorded at reading (iv): it is the re-derivation C5$'$(ii) asks for,
and it is the price-tag on the route Step~0 of \cref{prp:phirate} declines.
\textbf{And it is the sole consumer of C1(i) in the corpus}, Step~2 below, at two lines and nowhere
else, so C1(i) is a hypothesis of this lemma and, transitively, of no theorem. C1 is
cited here \emph{in full}: (i) at Step~2 for the $x$-smoothing bias, (ii) at the
Hilbert--Schmidt reduction that opens the proof, for $\|1/\pi\|_\infty$.
Suppose A0, A4(iii), B1(i), (ii), (iv.a), (iv.c), (iv.d), C1, C2$'$, C3 and C5$'$(i) hold.
\emph{\textbf{Three clauses of B1(iv) and not four, and not the block}}: \textup{(iv.b)} is opened
nowhere below, and the factorisation $\Pi=\Pi_u\otimes\Pi_t\otimes\Pi_{\mathsf w}$ is B1(iv)'s
\emph{declaration}, which carries nothing, on the assumption block's \textbf{three registers} paragraph.
\emph{\textup{(iv.c)} is on this line although a draft of the 2026-08-24 pass counted only
\textup{(iv.a)} and \textup{(iv.d)}}: Step~1's variance bound and Step~2's bias bound both integrate
$\bar w^{2}$ against $\Pi_t$, and the stated order $O\big((nh_xh_u)^{-1}\big)$ needs that integral
finite, which is \textup{(iv.c)} and only \textup{(iv.c)}. The clause was \emph{consumed} at two steps
and \emph{named} at neither until that day; \emph{Step~1 was repaired then and \textbf{Step~2 was not},
and named nothing until 2026-08-26}. \emph{Both steps now name it}, and the sentence this replaces said
the repair was complete when it covered one of the two consuming steps --- the same shape of defect,
one level down, as the under-charging it was written to record.
\emph{\textbf{B1(iv.a) is on that line for a reason, and neither promotion retires the reason}}: Step~1 of the
proof writes $\Pi_u(du)=\varpi_u(u)\,du$ and reads $\|\varpi_u\|_\infty<\infty$, which is B1(iv.a) and only
B1(iv.a), it is what turns the $\Pi_u$-integral of $C_{h_u}^{2}$ into a $du$-integral and produces the
$h_u^{-1}$ of \eqref{eq:gamman}'s variance term. \emph{The same clause reaches the step a second time
through \textup{(P7)}}, whose own hypothesis line is B1(iv.a) and C2$'$(i).
\emph{A1b(iv) is \textbf{not} a hypothesis}, on the ground recorded at
\cref{lem:opnoise}: no step below reads a moment of $\Delta Z$, B1(i)'s boundedness of $\Phi$ making every
$\Delta Z$-integral here an integral of a bounded summand.
\emph{C5$'$ is cited \textbf{by clause}, and the narrowing is not cosmetic}: what this lemma consumes in
the $x$-direction
is the proviso it states for itself below, $(nh_xh_u)^{-1}\ln n\to0$, while C5$'$(ii)'s stronger form is
a strictly higher floor for $h_x$, and \emph{a floor stated higher than any proof needs stops being free the
moment it must meet a ceiling}, which is \eqref{eq:hxwindow} at C5$'$. Clauses (iii) and (iv) are a
\emph{cross-reference} and a \emph{statement of scope}, and nothing here reads them.
% RENAMED 2026-08-26.  This sentence read "a pointer and a scope declaration",
% reproducing both technical terms of the assumption block's REGISTER paragraph with
% different meanings -- two words, four meanings.  The register pair is reserved; the
% C5' pair moved.  Companion sites: assumptions_supp.tex, the five-clause Part-VI
% list, items C5'(iii) and C5'(iv).
Define
% ARCHAEOLOGY.  Through round 9 two sentences elsewhere still named this lemma as
% consumed; both were corrected.  The assumption block asserted through dispatch 17
% that C1(i) was consumed by prp:phirate instead -- the one result that declines the
% route C1(i) is used on.  B1(iv.a) was added to the hypothesis line after
% dispatch 20, which carried B1(i), (ii) and (iv.c) alone, and no audit round had
% named the gap.  A1b(iv) stood on the line through dispatch 20 and is struck.
% C5' read as a block through dispatch 20, which on this file's register rule
% charged the lemma with C5'(ii).
\begin{equation}\label{eq:gamman}
\gamma_n\;:=\;\underbrace{\big(n\,h_x\,h_u\big)^{-1/2}\sqrt{\ln n}}_{\text{variance}}
\;+\;\underbrace{h_x^{p_K}+h_u^{p_K}}_{\text{smoothing bias}}
\;+\;\underbrace{r_n\,\big(h_x\,h_u\big)^{-1/2}}_{\text{propensity}} .
\end{equation}
\emph{The propensity term carries the bandwidth factor and the first term is therefore never binding},
since $r_n\ge n^{-1/2}$ by A4(iii) and \eqref{eq:frontier}. \emph{The factor $(h_xh_u)^{-1/2}$ is a
triangle-inequality bound and is very likely improvable}: the mean of that term is a $\nu$-functional
of an integrable function and is $O(r_n)$ with no bandwidth at all, and only its fluctuation carries
the factor, which a sup-norm rate for $\hat\varrho$ would remove, but \emph{nothing in the block
prices $\|\hat\varrho-\varrho\|_\infty$}, as the glossary's $r_n$ row records, so the crude bound is
carried. This is recorded, not claimed away.
Then, with $\mathcal V_n$ the source-condition ball \eqref{eq:srcball} below and
$\mathcal S_n$ the spline sieve of A8,
\begin{equation}\label{eq:aerr}
\sup_{\varphi\in\mathcal V_n\cap\mathcal S_n}\big\|\hat{\mathcal A}_\varphi-\mathcal A_\varphi\big\|_{\mathcal V\to L^{2}(\Pi)}
\;=\;O_P(\gamma_n),
\end{equation}
provided $\big(n\,h_x\,h_u\big)^{-1}\ln n\to0$, $\mathsf R\,h_x^{s-1/2}=o\big(h_u\,\gamma_n\big)$, 
which aligns the base points, $\hat{\mathcal A}_\varphi$ evaluating the residual at
$\Delta_{h_x}\varphi$ and $\mathcal A_\varphi$ at $\Delta\varphi$, \emph{and the factor $h_u$ is there for
the reason \cref{lem:unifop}'s own proviso carries one}: the alignment is charged in \textbf{operator}
norm and therefore at $h_u^{-1}\mathsf R\,h_x^{s-1/2}$, which is $o(\gamma_n)$ exactly when the displayed
condition holds, and
% ARCHAEOLOGY.  This proviso read R h_x^{s-1/2} = o(1) through dispatch 25, which is
% a factor h_u^{-1} too weak for what Step 2's alignment costs and does not deliver
% eq:aerr's O_P(gamma_n) at all: o(1) leaves the alignment term at o(h_u^{-1}), which
% need not be o(gamma_n).  It is lem:unifop's proviso defect at a second site.  THE
% BLAST RADIUS IS NIL -- this lemma declares "Consumers: none" -- but leaving it
% would have been the same defect twice, and the strengthening changes no conclusion
% stated anywhere in the corpus.
\begin{equation}\label{eq:entropyA}
L_x\,\ln n\;=\;o\big(n\,h_x\,h_u\,\gamma_n^{2}\big) .
\end{equation}
\emph{That the supremum in \eqref{eq:aerr} is a random variable at all is \cref{lem:unifop}\textup{(3c)}'s
separability paragraph, read at the subset $\mathcal V_n\cap\mathcal S_n$}; Step~4 of the proof below does the
reading, and the claim is stated once, there, rather than repeated here.
\emph{\textbf{The supremum in \eqref{eq:aerr} is nested, and that paragraph covers the outer one only}}:
$\|\hat{\mathcal A}_\varphi-\mathcal A_\varphi\|_{\mathcal V\to L^{2}(\Pi)}$ is itself a supremum over the
\emph{uncountable} unit ball of $\mathcal V$, so at each fixed $\varphi$ it is not a random variable on the
strength of the definition alone. \emph{The ingredient is present and was uncited}: $\mathcal V$ is a
subspace of $L^{2}(\mathcal X,\pi)$ with $\mathcal X$ compact, hence separable, and
$\zeta\mapsto\|(\hat{\mathcal A}_\varphi-\mathcal A_\varphi)\zeta\|_{L^{2}(\Pi)}$ is continuous in
$\|\cdot\|_{L^{2}(\pi)}$ at every sample point, both operators being bounded there, 
\cref{prp:hatAprop}(i) for the first, \cref{prp:compact} for the second, so the inner supremum equals
the supremum over a fixed countable dense subset of the unit ball, and each of those is a random variable
because $(\hat{\mathcal A}_\varphi-\mathcal A_\varphi)\zeta$ is jointly measurable in $(\omega,u,t)$ by
\cref{def:lfstep}'s measurability verification, first step and $L^{2}(\Pi)$ is separable. The two suprema are then countable
one after the other, and \eqref{eq:aerr} is a statement about a random variable. \emph{This is a
formality with a bounded blast radius, and it is stated rather than left}: this lemma declares
\textbf{Consumers: none}, so no result inherits the gap, but a display that is not defined is not a
display.
The supremum is taken over $\mathcal V_n\cap\mathcal S_n$ and not over $\mathcal V_n$, and the
restriction removes a budget line. On the sieve the covering numbers are elementary, the volumetric
bound of \cref{lem:unifop}(3d), so \eqref{eq:entropyA} is a condition on $L_x$ and rests on no
classical entropy theorem; over all of $\mathcal V_n$ the chaining step would consume the
\emph{imported} sup-norm entropy of a Sobolev ball, item (F-e) of the citation note, and would in
addition require $\mathsf R\,h_x^{s-1/2}=o\big(h_u\,\gamma_n\big)$ to align the two operators' base
points. \emph{The
restriction costs nothing even prospectively}: every iterate of \cref{def:lfstep}(b) lies in
$\mathcal S_n$ by construction, so a later result that did consume this lemma would evaluate it only
there. \emph{No result consumes it now}, and reading (iv) below says why the lemma is nonetheless
retained.
Three readings, and the third is the deliverable.
\begin{enumerate}
\item[(i)] The exponents in \eqref{eq:gamman} are \emph{one} on $h_x$ and \emph{one} on $h_u$, not two
and three. C5$'$(ii)'s $(n h_x^{2}h_u^{3})^{-1}\ln n\to0$ is inherited from a density-\emph{derivative}
estimator in the cross-sectional companion; $\hat{\mathcal A}_\varphi$ estimates a density-weighted
conditional mean and carries one bandwidth in each direction. \emph{C5$'$(ii) \textbf{is} the condition
used here as of this dispatch}: it previously stated the $h_x^{2}h_u^{3}$ pairing and now states the
$h_xh_u$ one, retaining the stronger form only where \eqref{eq:hxwindow} makes it compatible with (P-g).
\emph{The ground formerly given for declining the relaxation, that the stronger clause implies the
weaker one, ``so nothing proved is lost'', was an invalid inference and is struck}: C5$'$(ii) is a
\emph{lower} bound on $h_x$, and a lower bound stated above what any proof needs is free only until it has
to meet an upper bound.
\item[(ii)] $\gamma_n$ is \emph{not} a component of $\hat\delta_n$. It would multiply
$\|\vdagger-\varphi\|_{L^{2}(\pi)}$ wherever it appeared, entering an effective noise level at
$\gamma_n\delta_n$ and not at $\gamma_n$.
\item[(iii)] The $\sqrt{\ln n}$ in \eqref{eq:gamman} buys uniformity over
$\varphi\in\mathcal V_n\cap\mathcal S_n$ and nothing else; at a fixed $\varphi$ it may be deleted, by
the argument of \cref{lem:opnoise}(i).
\end{enumerate}

And a fourth reading, which is a scoping statement rather than a property of the lemma.
\textbf{\Cref{prp:phirate} does not consume this lemma.} Its proof applies the linear regularisation
theory \emph{conditionally on the sample}, at the realised operator
$\hat{\mathcal A}_{\hat\varphi^{(k-1)}}\vert_{\mathcal S_n}$, and the spectral source hypothesis (P-a.3)
is stated for the compression \eqref{eq:compress} of \emph{that} operator and not for
$\mathcal A_{\hat\varphi^{(k-1)}}$. Under that reading
no transfer of a source condition from $\mathcal A$ to $\hat{\mathcal A}$ is needed, which is
fortunate, since \cref{lem:srctransfer}(ii) shows the only general device for such a transfer is
unavailable at the source indices this programme lives at. \emph{So \eqref{eq:aerr} is retained for two
other reasons and neither is a rate}: it is the re-derivation C5$'$(ii) asks for, per reading (i), and
it is what a reader needs in order to see that the conditional-on-sample reading is a choice and not an
evasion, the operator error is $O_P(\gamma_n)$ and is priced, it is simply not the channel through
which the rate is proved.
\end{lemma}

\begin{proposition}[A7, verified in the amended form, and the one place the unamended form does not close]\label{prp:stopfeas}
Suppose the hypotheses of \cref{lem:opnoise}. Then:
\begin{enumerate}
\item[(i)] The $T$-component closes at a fixed constant. There is a finite $\mathsf c$, not
depending on $n$, with
\[
P\Big(\big\|\bar T(\vdagger,\theta_\dagger)\big\|_{L^{2}(\Pi)}>\mathsf c\,\hat\delta_n\Big)\longrightarrow0 ;
\]
$\mathsf c\ge4\mathsf c_3c_\Pi+\mathsf c_5c_\Pi\mathfrak m_{p_K}$ serves. \emph{The statement is at
$\hat\delta_n$ and not at $\hat\delta_n^{(T)}=n^{-1/2}\sqrt{\ln n}$, and the reason is C5$'$(i)}: under
the amended bandwidth the smoothing bias is $\asymp\hat\delta_n$, which is larger than
$\hat\delta_n^{(T)}$ whenever the propensity component binds. Under the unamended bandwidth
$h_u\asymp n^{-1/16}$ the sharper conclusion, with $\hat\delta_n^{(T)}$ in place of $\hat\delta_n$,
holds for every fixed $\mathsf c\ge4\mathsf c_3c_\Pi$. Only the displayed version is consumed.
% ARCHAEOLOGY.  The constant above was written bare italic tau through dispatch 21,
% in this proposition's own statement, at five lines -- against Sec.1.7's rule that
% the discrepancy constant tau_DP never appears without its word subscript, a rule
% clause (iii) kept and clauses (i)--(ii) broke.  Glossary C88.
\item[(ii)] The $\varrho$-component does not close at a fixed constant. Under A4(iii) alone,
$\|\hat T(\vdagger,\theta_\dagger)-\bar T(\vdagger,\theta_\dagger)\|=O_P(r_n)$ and no more; in
particular there is no finite $\mathsf c$ for which
$P(\|\hat T-\bar T\|>\mathsf c\,r_n)\to0$ follows. An $O_P$ hypothesis cannot produce a
fixed-constant exceedance conclusion, and A7 unamended was a fixed-constant statement.
\item[(iii)] \emph{Two repairs, either sufficient.} The unamended A7, the display of (ii) at
threshold $\tau_{\mathrm{DP}}\hat\delta_n$, holds if A4(iii) is strengthened to: there is a finite
$\mathsf c$ with $P(\|\hat\varrho-\varrho\|_{L^{2}}>\mathsf c\,r_n)\to0$. Alternatively, and this is
the route the block takes, \eqref{eq:a7amend} holds under A4(iii) exactly as stated, for every
$\tau_{\mathrm{DP}}\ge1$; and any divergent factor serves in place of $\ln n$.
\item[(iv)] What the amendment costs, and it is not nothing. \emph{The substitution is at the level
of $\hat\delta_n$, and A5(ii) is the record of its effect}: replacing $\hat\delta_n$ by
$\hat\delta_n\ln n$ in the discrepancy threshold is exactly why \eqref{eq:deltalog} sets
$\delta_n:=(\hat\delta_n\ln n)^{2\mu/(2\mu+1)}$ rather than $\hat\delta_n^{2\mu/(2\mu+1)}$, so
$\delta_n$ below already carries the logarithm and is not inflated a second time here.
Against \eqref{eq:frontier} the substitution is free. Since \eqref{eq:frontier} is a \emph{strict}
inequality, $\delta_n=o_P(n^{-1/4})$ survives it at every $(\mu,r,d_Z)$ satisfying it:
if $\tfrac{2\mu}{2\mu+1}\cdot\tfrac{r}{2r+d_Z}=\tfrac14+\varepsilon$ with $\varepsilon>0$ then
$\delta_n=(\hat\delta_n\ln n)^{2\mu/(2\mu+1)}=n^{-1/4-\varepsilon}(\ln n)^{2\mu/(2\mu+1)}=o(n^{-1/4})$.
The strictness of the frontier is what pays for it.
Against \textup{(P-f)} it is not free, and \textup{(P-f)} is the other condition $\delta_n$
enters. \textup{(P-f)} is an \emph{exact domination}, $\delta_n^{2}=O(h_u^{2}\hat\delta_n\ln n)$, in
which the left side now carries $(\ln n)^{4\mu/(2\mu+1)}$ and the right side $(\ln n)^{1}$; and
$\tfrac{4\mu}{2\mu+1}>1$ precisely when $\mu>\tfrac12$, which is everywhere in this block. \emph{So the
logarithms never cancel in \textup{(P-f)} and its thresholds must be strict}: at a configuration where
\textup{(P-f)} holds with equality in powers of $n$ it fails by a power of $\ln n$.
\Cref{rem:h6bandwidth} computes both thresholds and exhibits the \emph{boundary} configuration at which
this is not hypothetical; clause (5) there names the block's two running configurations and says which
is a design and which is an illustration.
\end{enumerate}
\end{proposition}

% ---------------------------------------------------------------------------
%  H.4  L-19.  UNIFORMITY OVER A SOURCE-CONDITION BALL, BY BRACKETING.
% ---------------------------------------------------------------------------

\begin{definition}[The source-condition ball]\label{def:srcball}
Let $s$ and $\mathsf R$ be as in H1 and A5(vi$'$) and set
\begin{equation}\label{eq:srcball}
\mathcal V_n\;:=\;\Big\{\varphi\in\mathcal V:\ \|\varphi-\vdagger\|_{L^{2}(\pi)}\le\bar\rho,\ \
\|\varphi-\vdagger\|_{H^{s}}\le\mathsf R\Big\},
\end{equation}
with $\bar\rho>0$ and $\mathsf R<\infty$ fixed and not depending on $n$, and $H^{s}$ the \emph{classical}
Sobolev space $H^{s}(\mathcal X)$ of order $s$, carrying no boundary condition, as in H1.
\emph{It is \textbf{not} a $\mathcal D$-generated scale}: the two spaces differ at every $s>\tfrac32$,
hence at every $s$ this block admits, and the norm
$\|\varphi-\vdagger\|_{H^{s}}$ in \eqref{eq:srcball} is the classical one.
% ARCHAEOLOGY.  This line read "generated by D" through dispatch 20; the phrase is
% struck here and in H1.
Take $\mathsf R$ large enough that
$\|\vdagger\|_{H^{s}}\le\mathsf R$ as well, which H1's source set permits and which
\cref{lem:unifop} Step~4 uses. \emph{$\mathcal V_n$ is
deterministic}, as L-19 requires, and it is a \emph{source-condition} ball rather than an $L^{2}(\pi)$
ball: the second constraint is exactly A5(vi$'$)'s $H^{s}$ envelope made into a set.

$\mathsf R$ is an $H^{s}$ radius and it is not the source radius (F-f) is applied at. The latter
is $\mathsf R_\ast$, the \emph{spectral pre-image} norm of (P-a.3), which enters
$\mathsf c_{17}=2c_\mu\mathsf R_\ast^{1/(2\mu+1)}$ and hence $\mathsf c_{15}$ of \eqref{eq:lfsched}.
The two are not comparable in either direction without the
operator interpolation \cref{rem:pasrc}\textup{(1)} shows is unavailable at $2\mu>1$. \emph{$\mathsf R$
is the radius that is \textbf{not} free}: it is shared by this definition, A5(vi$'$), (P-a.1), (P-c),
(P-g), \eqref{eq:entropy} and \cref{lem:unifop} Step~4, so ``fixed and large'' tightens two conditions
and inflates none, the hypothesis is an existential over \emph{one} $\mathsf R$ meeting seven
demands, and each of the seven is listed rather than left to be assembled.
\emph{\textbf{This is where the list is maintained}}: A5(vi$'$), \cref{lem:pathcont} and
\cref{rem:paenvelope} point here and maintain no copy, and \cref{def:crossfit}'s clause-\textup{(d)}
reading names three of the seven as instances and says in terms that it maintains none, \emph{as does
\textup{(P-a.1)}, which names four and is itself the third of the seven}.
\emph{No stronger claim is made, and that is deliberate}: the draft this sentence read ``the
corpus's only one'' while that partial copy stood unmarked. \emph{The subscript $n$ is conventional}; nothing
in \eqref{eq:srcball} varies with $n$, and the containment statement that does vary is
\cref{lem:pathcont}.

One symbol, one object, from here on. Bare $\mathcal V_n$ means \eqref{eq:srcball} and nothing
else in Part~H. An $L^{2}(\pi)$-ball is written $\mathcal V(\varphi_\circ,\rho)$ per \eqref{eq:ball},
with its centre and radius shown; the two are separated by the presence of arguments.
% ARCHAEOLOGY.  One symbol carried both the H^s radius and the spectral pre-image
% radius through round 11, on the two sides of an inequality def:lfstep has to
% establish; glossary C69.  A5(vi'), lem:pathcont and rem:paenvelope each carried a
% copy of the R-consumer list through dispatch 20, of lengths five, five and four
% against the seven above, all four maintained independently and all four
% differing, so a reader assembling the consumer set from any one of them got a
% different answer.  Nothing in any of the four was wrong -- each named a subset of
% the seven -- and that is precisely why four copies were worse than one, an
% incomplete list being indistinguishable from a complete one to a reader who meets
% it first.  Through round 9 the symbol cal V_n carried three meanings in Part H --
% the retired A6's Landweber-path ball, def:a6b's bar-rho_n ball, and eq:srcball;
% the first is gone with A6 and the second now carries arguments.
\end{definition}

\begin{lemma}[L-19: uniform operator error over the source-condition ball]\label{lem:unifop}
% p_K >= s - 1/2 ADDED TO THIS LIST 2026-08-05.  Step 4 below consumes it -- the
% Taylor expansion kills every term with 1 <= j <= k only when k <= p_K - 1 -- and
% through dispatch 25 it was carried by prp:phirate alone, i.e. by the downstream
% consumer rather than by the result whose own proof uses it.  This lemma is
% printed with a self-contained "Suppose ... Then ..." and can be cited as stated,
% so the condition has to be here.  Carrying it on prp:phirate as well is free.
Suppose A0, A4(iii), B1(i), (ii), (iv.c), (iv.d), C1(ii), C2$'$, C3, C5$'$(i),
\textbf{A8's basis clause} and H1
hold, adopt \cref{def:hatTexp}, suppose that
$s>\tfrac12$, so that on A0's compact interval the one-dimensional Sobolev embedding
$H^{s}(\mathcal X)\hookrightarrow C(\mathcal X)$ is available, suppose that
$p_K\ge s-\tfrac12$, and
that
\begin{equation}\label{eq:entropy}
\Big(\frac{\mathsf R}{h_u\,\hat\delta_n}\Big)^{1/s}\;=\;o\big(n\,\hat\delta_n^{2}\big) .
\end{equation}
Then, with $\mathcal V_n$ as in \eqref{eq:srcball},
\begin{equation}\label{eq:unifop}
\sup_{\varphi\in\mathcal V_n}\big\|\hat T(\varphi,\theta_\dagger)-T(\varphi,\theta_\dagger)\big\|_{L^{2}(\Pi)}
\;=\;O_P\big(\hat\delta_n\big),
\end{equation}
and the same bound holds with $\hat T$ replaced by $\hat T_{h_x}$ provided in addition
$\mathsf R\,h_x^{s-1/2}=O\big(h_u\,\hat\delta_n\big)$, \emph{whose factor $h_u$ is not decoration and
is what Step~4 actually charges}: that step bounds the
displacement of the two operator families in \textbf{operator norm} by
$\mathsf c_2c_\Pi h_u^{-1}\max_{i}\big|\Delta_{h_x}\varphi_i-\Delta\varphi_i\big|$, the $h_u^{-1}$ being
\cref{lem:intker}\textup{(ii)}'s $\|\partial_vC_{h_u}\|_\infty\le\mathsf c_2h_u^{-1}$, the same
$h_u^{-1}$ that forces the net radius at \textup{(3c)} and the $h_u$ inside \eqref{eq:entropy}, so
with \eqref{eq:xsmoothbias} the displacement is $\mathsf c\,h_u^{-1}\mathsf R\,h_x^{s-1/2}$, and it is
\emph{that product} and not $\mathsf R\,h_x^{s-1/2}$ alone that has to be $O(\hat\delta_n)$.
\emph{This is the same $h_u$ (P-c) carries}, and the proviso is \cref{prp:phirate}'s (P-g).
\emph{\textbf{B1(iv) is cited by clause and A8's is cited as a clause, and the two are different
registers.}} A8's basis clause is a \emph{condition} on the sieve --- it can fail, and A8's own note
exhibits $s=1.7$ with linear splines --- so it is a carrier and belongs on the line. B1(iv)'s
factorisation is a \emph{declaration}, standing at every statement here and carried by nobody; what this
proof opens is \textup{(iv.d)} at Steps~2 and~3(3a) and \textup{(iv.c)} at Steps~3(3a) and~4, and
nothing opens \textup{(iv.a)} or \textup{(iv.b)}. \emph{Through 2026-08-24 this line was the block}.
% ARCHAEOLOGY.  Through dispatch 25 this proviso read R h_x^{s-1/2} = O(hat-delta_n)
% -- a factor h_u too weak for the step that consumes it, since the h_u^{-1} above
% is real and the composed displacement is then n^{alpha/8} hat-delta_n at the top
% of the h_x window.  rem:tccwhy(3) printed the correct arithmetic in this same file
% while this line printed the weak form.  Corrected together with (P-g), lem:aerr's
% proviso, rem:app5read(2), eq:hxwindow and the four arithmetic readings that quote
% them.  THE PROVISO IS DELIBERATELY LEFT INLINE AND UNNUMBERED, as it was before:
% giving it a labelled display would advance the equation counter and renumber every
% later display in the corpus, which the build driver's header warns against and
% which ~33 symbolic specs and the render-sync checker read.  The one site that
% would have cited it -- lem:aerr's own proviso, 200 lines EARLIER in this file --
% would also have been a forward reference, so the label bought nothing.
\emph{Furthermore}, if the supremum is restricted to
$\mathcal V_n\cap\mathcal S_n$ with $\mathcal S_n:=\operatorname{span}(b_1,\dots,b_{L_x})\cap\mathcal V$
the spline sieve of A8, then \eqref{eq:unifop} holds under
\begin{equation}\label{eq:entropysieve}
L_x\,\ln n\;=\;o\big(n\,\hat\delta_n^{2}\big)
\end{equation}
in place of \eqref{eq:entropy}, and that is the version \cref{def:lfstep} and
\cref{prp:phirate} consume; \eqref{eq:entropy} is consumed by nothing.
\eqref{eq:entropy} is written in the primitive variables deliberately. Reading it at a bandwidth
requires choosing one, and that choice is C5$'$(i)'s; \cref{rem:l19cost} unpacks the reading there.

What makes \eqref{eq:unifop} a statement about the distribution of a random variable, and where it is
established. The supremum in \eqref{eq:unifop} is over an \emph{uncountable} index set, so it is not a random
variable on the strength of the definition alone. \emph{It is one, and the proof establishes it at
\textup{(3c)}}: $\mathcal V_n$ is $\|\cdot\|_\infty$-\textbf{separable}, A0's $\mathcal X$ is a compact
interval and $s>\tfrac12$, so $\mathcal V_n\subset C(\mathcal X)$, and the process is sup-norm Lipschitz at
every sample point outside one fixed $\varphi$-free null set, so the supremum agrees almost surely with a
countable one. That paragraph consumes three things this hypothesis list carries \textbf{because}
it consumes them.
% ARCHAEOLOGY.  Through dispatch 19 the hypothesis list carried neither C1(ii) nor
% H1 and named the Sobolev embedding nowhere.
\textbf{C1(ii)} is
load-bearing rather than incidental: it makes $\pi$ strictly positive on $\mathcal X$, so the continuous
representative of an $L^{2}(\pi)$ class is \emph{unique} and $\|\cdot\|_\infty$ is well defined on
$\mathcal V_n$, without it one class can carry two continuous representatives of different sup-norms and
the norm the entire $\epsilon_0$-net runs in is not a function of the class. \textbf{H1} enters through
\cref{def:srcball}, which chooses $\mathsf R$ with $\|\vdagger\|_{H^{s}}\le\mathsf R$ ``which H1's source
set permits'', and hence gives $\|\varphi\|_{H^{s}}\le2\mathsf R$ on $\mathcal V_n$. And the passage from
that bound to continuity is the \textbf{Sobolev embedding} $H^{s}(\mathcal X)\hookrightarrow C(\mathcal X)$
at $s>\tfrac12$, named in the hypothesis list above.
\emph{\textbf{\Cref{def:hatTexp} is adopted rather than presumed}}: the conclusion \eqref{eq:unifop} is
a statement about $\hat T$, which is that definition's object under that definition's sample-splitting and
truncation conventions, and (3c) reads its record that $\varphi\mapsto\hat T_{h_x}(\varphi)$ depends on
$\varphi$ only through the bounded linear functionals \eqref{eq:smdiff}.
\emph{\textbf{A8's basis clause is on the line}} because clause (b) is proved
at (3d) from the B-spline inverse inequality on a quasi-uniform mesh, which carries the entire sieve
entropy bound and hence \eqref{eq:entropysieve}, the version \cref{prp:phirate} consumes.
\emph{\textbf{A1b(iv) is \textbf{not} a hypothesis}}, on the ground recorded at \cref{lem:opnoise}:
nothing in this proof reads a moment of $\Delta Z$.
\emph{And \eqref{eq:wbounds}, cited twice at (3c), is \textbf{not} on this line}: that display is
% ARCHAEOLOGY.  lem:aerr carries C1 in full and was never short; only this lemma
% was.  Three further items were settled in the same dispatch and none changed a
% conclusion: def:hatTexp was presumed rather than adopted (lem:opnoise adopted it
% expressly and this lemma did not); A8's basis clause appeared only in the
% parenthetical "the spline sieve of A8" in the statement, where a reader building
% the matrix from "Suppose" lines does not meet it; A1b(iv) stood on this list
% through dispatch 20 and is struck, a strict weakening with no change of
% conclusion.
two-sided and its lower half rests on A4(i), which neither appeal reads; both are narrowed at the site to
free upper bounds, \textup{(P1)} for the population weight and $\varrho,\hat\varrho\in[0,1]$ for the
propensity difference, which is why \cref{def:hatTexp} is adopted above and A4(i) is not
assumed.
\emph{The claim is made \textbf{once}, there, and referred to at the three other suprema over
$\mathcal V_n$ in this block}: Step~1's and Step~4's in the same proof, and the one in \eqref{eq:aerr}.
\emph{Neither an outer-probability convention nor a Suslin hypothesis is invoked, and none is needed}: the
separability of the index set and the pathwise sup-norm continuity of the process are both available, the
second from clause (3a) of the proof.

What is used, and what is not. The proof of clause (a) uses a \textbf{bracketing} entropy bound
for the source-condition ball \eqref{eq:srcball} in the sup-norm, 
$\log N_{[\,]}(\epsilon,\mathcal V_n,\|\cdot\|_\infty)\le\mathsf c_6(\mathsf R/\epsilon)^{1/s}$, valid
for $s>\tfrac12$, item (F-e) of the citation note and \emph{not proved here}, together with a
bounded-difference inequality at each bracket centre, which \emph{is} proved here, as
\cref{lem:bdiff}(i). The pin for (F-e) is \citet[Ch.~2]{birman1980}; the obvious alternative
is not merely weaker but unavailable. Routing through a
H\"older-ball entropy bound at exponent $s-\tfrac12$ would yield $\epsilon^{-1/(s-1/2)}$ rather than
$\epsilon^{-1/s}$ and would therefore move \eqref{eq:entropy}'s threshold in $s$ up by \emph{exactly}
$\tfrac12$, at every configuration and at every bandwidth, the right-hand side $n\hat\delta_n^{2}$ being
untouched: at \cref{rem:h6bandwidth}\textup{(5)}'s \textbf{boundary} configuration
$(\mu,d_Z,r)=(1,3,4.5)$, equivalently $\alpha=\tfrac38$, from $s>\tfrac{27}{16}$, 
\cref{rem:l19cost}\textup{(1)}'s reading at the amended $h_u\asymp\hat\delta_n^{1/8}$, to
$s>\tfrac{35}{16}$, \emph{but that route needs the embedding
% CONFIGURATION NAMED AND NUMBERS AMENDED 2026-08-10 (Phase C).  This read "at the
% reference configuration from s > 1.75 to s > 2.25".  Two defects.  (1) "The
% reference configuration" names NOTHING: rem:h6bandwidth(5) declares the block's
% two running configurations "named once so that they are not confused" and neither
% is called this.  (2) 1.75 and 2.25 are the RETIRED h_u ~ n^{-1/16} tuning's
% numbers -- rem:l19cost(1)'s own archaeology note says so in terms, forty lines
% below -- while the amended C5'(i) gives s > 27/16 = 1.6875, which the same
% clause prints and which eq:lxwindow matches.  Re-derived here: the condition is
% 9*alpha/(8s) < 1 - 2*alpha, so s > 9*alpha/(8(1-2*alpha)); replacing the entropy
% exponent 1/s by 1/(s-1/2) shifts the threshold by exactly +1/2 whatever alpha and
% whatever h_u, since the right-hand side does not move.  At alpha = 3/8 that is
% 27/16 -> 35/16 = 2.1875; at the retired tuning it was 7/4 -> 9/4, the same +1/2,
% which is the check that the two readings are the same computation.
$H^{s}\hookrightarrow C^{0,s-1/2}$, which is vacuous for
every $s>\tfrac32$}, since $C^{0,\beta}$ on an interval contains only constants at $\beta>1$, and
\eqref{eq:lxwindow} already requires $s>\tfrac{27}{16}$. \emph{The correct H\"older route runs through
$H^{s}\hookrightarrow C^{k,\beta}$ with $k=\lfloor s-\tfrac12\rfloor$}, whose sup-norm entropy is
governed by $k$ and $\beta$ separately and does \emph{not} reduce to a single H\"older exponent, so
Birman--Solomjak is not merely the pin taken but the only one
available. The proof does
\emph{not} use a VC-subgraph bound, and it does \emph{not} use stochastic equicontinuity of the
$\sqrt n$-scaled empirical process at an estimated point. \emph{The VC route registered in the
strategy memo \S3.4 is not available here}: the class
$\{\widetilde W_l\,\mathsf C_{h_u}(u-\Delta Y_\theta+\Delta\varphi)\,w(t'\Phi(\Delta Z)):\varphi\in\mathcal V_n\}$
is indexed by an infinite-dimensional $\varphi$ entering through a smooth transformation, and it is
\emph{not} a VC-subgraph class of indicator-weighted functions, the indicator has been smoothed, and
even before smoothing the index is a function and not a halfspace. \emph{Smoothness of the index set,
not combinatorial dimension, is what supplies the entropy}, and that is precisely why the ball must be
a source-condition ball and not an $L^{2}(\pi)$ ball: the $L^{2}(\pi)$ ball has infinite sup-norm
entropy at every radius.
\emph{On the sieve the entropy is elementary and needs no classical citation}: $\mathcal V_n\cap\mathcal S_n$
is a bounded subset of a space of dimension \emph{at most} $L_x$, exactly $L_x-1$, by the count
\cref{def:lfstep}'s measurability verification, first step performs, and the bound below is stated at $L_x$ because it is an upper
bound and nothing here is sharp in the dimension, so a volumetric bound gives
$\log N(\epsilon)\le L_x\log(\mathsf c\,\mathsf R\,L_x^{1/2}/\epsilon)$. \emph{The two versions
therefore differ in what they rest on as well as in what they require}, and the one consumed is the
elementary one.
\end{lemma}

\begin{remark}[What \cref{lem:unifop} costs, and which clause charges it]\label{rem:l19cost}
\eqref{eq:entropy} is a restriction linking $s$, $r$ and $d_Z$, and it should be read three ways. Write
$\hat\delta_n=n^{-\alpha}$ with $\alpha=r/(2r+d_Z)\in(0,\tfrac12)$ and ignore the fixed $\mathsf R$.
Since $\epsilon_0\asymp h_u\hat\delta_n$, \eqref{eq:entropy} reads
$\big(h_u\hat\delta_n\big)^{-1/s}=o(n^{1-2\alpha})$, and the bandwidth enters through $h_u$ alone.

(0) Which clause charges it, and the answer is neither \cref{prp:phirate} nor
\cref{def:lfstep}. \eqref{eq:entropy} is clause (a)'s condition; what the estimator and the rate
consume is clause (b) under \eqref{eq:entropysieve}, which (P-d) names. \emph{So \eqref{eq:entropy} is
charged to nothing}, and readings (1)--(2) below are a discussion of a clause no result consumes, kept
because a referee will ask what the Sobolev-ball route would have cost.

(1) It is satisfiable, and it is not the binding line. At
the amended bandwidth $h_u\asymp\hat\delta_n^{1/8}$ it
reads $\tfrac98\alpha/s<1-2\alpha$, i.e.\ $s>\tfrac{27}{16}=1.6875$ at $\alpha=\tfrac38$, 
\emph{the same number as \eqref{eq:lxwindow}'s, and not a coincidence}: both
conditions compare a resolution level against $n\hat\delta_n^{2}$ and both resolutions are
$h_u\hat\delta_n$, the bracket radius $\epsilon_0\asymp h_u\hat\delta_n$ here and the sieve
approximation level $\mathsf RL_x^{-s}=O(h_u\hat\delta_n)$ in (P-c). They are two readings of one
budget and only the second is charged. \emph{It also runs the wrong way in $r$}, the left-hand side
increasing in $\alpha$ while $n^{1-2\alpha}$ decreases, so \emph{a smoother propensity makes
\cref{lem:unifop} harder} and the block's two constraints on $r$ pull in opposite directions.
% ARCHAEOLOGY.  A round-9 draft of prp:phirate listed eq:entropy as hypothesis
% (P-d) and thereby charged s > 1.75 for a condition its proof does not use.  At the
% unamended bandwidth h_u ~ n^{-1/16} reading (1) gives s > 1.75, slightly stronger,
% the smaller bandwidth making the bracket radius smaller.

(2) It excludes the configuration in which the propensity component does not bind. As
$\alpha\uparrow\tfrac12$ the right-hand side $n\hat\delta_n^{2}=n^{1-2\alpha}$ of \eqref{eq:entropy}
ceases to diverge and no finite $s$
suffices. \emph{So \cref{lem:unifop} requires $r_n\gg n^{-1/2}$ strictly}, i.e.\ requires
$\alpha<\tfrac12$ with a \emph{polynomial} margin. \emph{A5(i) does not supply it and is not cited for
it}: that clause's floor gives $n\hat\delta_n^{2}\ge\ln n$, a \emph{logarithmic} divergence, and
\eqref{eq:entropy} needs a polynomial one, the same distinction A4(ii)'s $d_Z\ge1$ note draws at
$\alpha=\tfrac12$, where $L_x\ln n=o(\ln n)$ admits no integer $L_x\ge1$.
That is a scope condition on Step 6 and it is stated rather than assumed away: with a parametric
propensity the uniform bound at rate $\hat\delta_n$ is \emph{not} available by this route, and the
$L^{2}(\Pi)$-norm rate would have to be relaxed to $\hat\delta_n$ times a slowly divergent factor, 
which, by \cref{prp:stopfeas}(iv), the strictness of \eqref{eq:frontier} again absorbs.
\end{remark}

\begin{lemma}[Containment of the path]\label{lem:pathcont}
Let $\hat\varphi^{(0)}\in\mathcal S_n$ satisfy (P-a.2), $P(\|\hat\varphi^{(0)}-\vdagger\|_{L^{2}(\pi)}\le\hat\varepsilon_0)\to1$
with $\hat\varepsilon_0\to0$, let \textbf{A5(vi$'$)} hold \emph{in its envelope form at the fixed radius
$\mathsf R$ of \cref{def:srcball}}, 
$P\big(\max_{k\le\bar k}\|\hat\varphi^{(k)}-\vdagger\|_{H^{s}}\le\mathsf R\big)\to1$, and let
$\hat\varphi^{(k)}$ be the stages of \cref{def:lfstep}, each stopped at an index
$\hat m^{(k)}\le\bar m_n$. Then, with that same $\mathsf R$ and with $\bar\rho$ fixed and large enough,
$P\big(\hat\varphi^{(k)}\in\mathcal V_n\cap\mathcal S_n\ \text{for all }k\le\bar k\big)\to1$.

The radius is \textbf{not} chosen here, and that is the point of stating A5(vi$'$) as an envelope.
\cref{def:srcball} requires $\mathcal V_n$ to be \emph{deterministic}, so the $H^{s}$ constraint must be
a fixed number; but $\mathsf R$ also appears at the other sites \cref{def:srcball} enumerates, in two of
which it is an \emph{upper} bound that ``fixed and large'' tightens.
\emph{So the hypothesis is an existential over one $\mathsf R$ satisfying all of them simultaneously},
which is a statement about the logical structure of the hypothesis and not about its size.
The enumeration is \cref{def:srcball}'s and is not repeated here. \emph{An $O_P(1)$ reading of
A5(vi$'$) does
not serve}: $O_P(1)$ is a family of bounds indexed by tolerance, and a deterministic ball cannot be
defined from a family.
% ARCHAEOLOGY.  This sentence carried a copy of def:srcball's enumeration through
% dispatch 20, one entry shorter than the list it pointed at.

Three things about this lemma, and the third is a finding.
\emph{(a)} The $L^{2}(\pi)$ half is the stagewise bound of \cref{prp:phirate} Step 4 and is genuinely
proved; the within-stage Landweber path satisfies
$\|v_m\|_{L^{2}(\pi)}\le m\,\sup_{\mathcal V_n}\|\hat T_{h_x}\|\le 2m\,\mathsf c_3c_\Pi$ under the
normalisation of \cref{def:lfstep}, which bounds the path but by $\bar m_n$ and not by $\bar\rho$; it
is the DP stopping rule, not the Landweber algebra, that keeps the \emph{stopped} iterate inside the
ball. The memo's registered mechanism, ``$\varphi_m-\varphi_0\in\mathcal R(\hat{\mathcal A}^\ast)$
with norm growing at most polynomially in the cap'', gives the first bound and not the second, and
the distinction matters because $\mathcal V_n$'s radius is fixed while $\bar m_n$ diverges.
\emph{(b)} The sieve constraint $\hat\varphi^{(k)}\in\mathcal S_n$ is by construction.
\emph{(c)} \textbf{The $H^{s}$ half is A5(vi$'$), assumed and not proved, and it is assumed as a
w.p.$\to1$ envelope at the fixed $\mathsf R$.} A5(vi) prices
$\|\hat\varphi-\vdagger\|_{H^{s}}$ at the \emph{final} iterate and is consumed by
\cref{lem:balance}\textup{(iii)} alone, on the balancing route. What is needed here is the bound
\emph{at every stage}, because \cref{lem:unifop}'s ball is a source-condition ball and every
$\hat\varphi^{(k)}$ must lie in it; \emph{that is a different statement, not a different route for the
same one}, and it is A5(vi$'$). Since $\bar k$ is finite and does not grow with
$n$, A5(vi$'$) is a finite union of $O_P(1)$ statements and not a uniformity requirement, which is
the only reason it is affordable. Deriving the $H^{s}$ bound instead of assuming it is possible on the
sieve, Bernstein's inequality for splines gives
$\|\zeta\|_{H^{s}}\le\mathsf c_7L_x^{s}\|\zeta\|_{L^{2}(\pi)}$ on $\mathcal S_n$, but at the
cost of $L_x^{s}\delta_n=O(1)$, which is a genuine extra restriction and is \emph{not} imposed;
A5(vi$'$) is cheaper and is what the block now carries.
\end{lemma}

% ---------------------------------------------------------------------------
%  H.5  THE RATE.
% ---------------------------------------------------------------------------

\begin{lemma}[Source-condition transfer: what is free, what is unavailable, and what is merely unproved]\label{lem:srctransfer}
Let $\mathcal A,\hat{\mathcal A}:\mathcal V\to L^{2}(\Pi)$ be bounded with
$\|\hat{\mathcal A}-\mathcal A\|\le\gamma$, and let $v\in\mathcal V$.
\begin{enumerate}
\item[(i)] \emph{$\mu\le\tfrac12$: transfer is free.} If $v=(\mathcal A^\ast\mathcal A)^{\mu}\varsigma$
with $\mu\in(0,\tfrac12]$ and $\|\varsigma\|\le\mathsf R_\ast$, then there is $\tilde\varsigma$ with
\[
\|\tilde\varsigma\|\;\le\;\mathsf c_8\mathsf R_\ast
\qquad\text{and}\qquad
\big\|v-(\hat{\mathcal A}^\ast\hat{\mathcal A})^{\mu}\tilde\varsigma\big\|_{L^{2}(\pi)}
\;\le\;\mathsf c_8\,\mathsf R_\ast\,\gamma^{2\mu} .
\]
Here $\mathsf c_8:=2^{2\mu}$, the constant clause \textup{(ii.b)} also carries. The proof of
\textup{(i)} in fact delivers the smaller $2^{\mu}$; the symbol is written at the larger value so that
\emph{one letter carries one number} across the two clauses, and \textup{(i)} holds \emph{a fortiori}
because $2^{\mu}<2^{2\mu}$ at every $\mu>0$.
% DEFINE-BEFORE-USE, ADDED 2026-08-10 (Phase C).  c_8 was used here and defined
% only at (ii.b), and proofs.tex's clause-(i) step set it to 2^mu -- one constant,
% two values.  The definition is now fixed once, at first use, at the value both
% (ii.b) and proofs.tex:7690 print.  No conclusion moves.
\emph{The radius here is the \textbf{spectral pre-image} norm $\mathsf R_\ast$ throughout this lemma,
never the $H^{s}$ radius $\mathsf R$ of \cref{def:srcball}}, clause (iii) is the one place an
$H^{s}$ radius appears, and it says $\mathsf R$ for that reason.
% ARCHAEOLOGY.  Glossary C69.
\emph{The proof is Heinz's inequality}, operator monotonicity of $z\mapsto z^{2\mu}$ on
$[0,\infty)$, valid for $2\mu\le1$, and $\mu\le\tfrac12$ is exactly its range of validity;
\citet[Prop.~8.21]{engl2000}, which \citeauthor{engl2000} name the inequality of Heinz, state it for a
Hilbert-scale generator and hence under a \emph{strict} positivity that $\mathcal A^\ast\mathcal A$
does not have, so the passage to a merely positive bounded operator is the standard shift: apply it to
$\mathcal A^\ast\mathcal A+\epsilon\,\mathrm{Id}$ and let $\epsilon\downarrow0$.
% (M-12) CLOSED 2026-08-03 AS A PIN WITH A STATED GAP.  MIRROR of results.tex;
% the full archaeology, the verbatim statement and the reason the gap is stated
% rather than papered over are recorded there.  Locator: engl2000 Prop. 8.21,
% p.213, sec 8.4 "Hilbert Scales" -- NOT sec 8.5, which this file's other
% engl2000 cites use for the INTERPOLATION range.  EHN name it at Rem. 8.20.
\item[(ii)] $\mu>\tfrac12$: the device of (i) is unavailable, and the failure is one of
non-commutativity, not of size. Two statements, and \textbf{the second is a limitation on the first.}
\begin{enumerate}
\item[(ii.a)] \emph{The tool fails.} For every $2\mu>1$ there are $\mathcal Q_1,\mathcal Q_2$ on a
two-dimensional space with $0\preceq\mathcal Q_1\preceq\mathcal Q_2$ and
$\mathcal Q_1^{2\mu}\not\preceq\mathcal Q_2^{2\mu}$; at $\mu=1$,
$\mathcal Q_1=\big(\begin{smallmatrix}1&1\\1&1\end{smallmatrix}\big)$ and
$\mathcal Q_2=\big(\begin{smallmatrix}2&1\\1&1\end{smallmatrix}\big)$ serve. \emph{The two clauses of
that sentence are proved by two different arguments}: the general $2\mu>1$ case by L\"owner's
divided-difference criterion for matrix monotonicity of order two \citep[Ch.~V]{bhatia1997}, the
displayed pair by inspection. So step (i)'s passage from
$\mathcal A^\ast\mathcal A\preceq2\hat{\mathcal A}^\ast\hat{\mathcal A}+2\gamma^{2}\,\mathrm{Id}$ to the same
inequality raised to the power $2\mu$ is \emph{not available} above $\mu=\tfrac12$, and no bound in
$\|\hat{\mathcal A}-\mathcal A\|$ alone is obtainable by that route.
\item[(ii.b)] But no counterexample to the \emph{conclusion} is claimed, and none can be
diagonal. If $\mathcal A$ and $\hat{\mathcal A}$ share a singular system, transfer holds at
\emph{every} $\mu>0$, with $\mathsf c_8=2^{2\mu}$ and residual $(2\gamma)^{2\mu}\mathsf R_\ast$: split the
index set at $\sigma_j\ge2\gamma$, where $\hat\sigma_j\ge\sigma_j/2$ bounds $\tilde\varsigma$, and
below it, where $|v_j|\le(2\gamma)^{2\mu}|\varsigma_j|$ bounds the residual. \emph{So the obstruction is
that the perturbation may rotate the singular system, not that it is large}, and whether transfer
itself fails above $\tfrac12$ is left open here. \textbf{What is established is the unavailability of
the device, and that is all this programme uses}: (P-a.3) is stated spectrally on the compression of
$\hat{\mathcal A}$ directly, so no transfer and no interpolation is performed anywhere in Part~H.
\end{enumerate}
\item[(iii)] The Hilbert-scale route: what it needs, and why it is \textbf{not} the route taken.
A Hilbert-scale source condition, $\|v\|_{H^{s}}\le\mathsf R$, is a statement about $v$ alone and
transfers trivially. \emph{But it is not the hypothesis of \textup{(F-f)}}, and it delivers one only
through the \emph{link} condition $\|\hat{\mathcal A}\zeta\|_{L^{2}(\Pi)}\asymp\|\zeta\|_{H^{-a}}$
\emph{together with an operator-interpolation step at exponent $2\mu$}. Two facts about that step, and
they are the reason \cref{prp:phirate} does not take this route.
\begin{enumerate}
\item[(iii.a)] \emph{The link condition cannot hold on $\mathcal V$}: the left side vanishes on the
infinite-dimensional $\ker\hat{\mathcal A}$ of \cref{prp:hatAprop}(ii) and the right side vanishes at no
nonzero $\zeta$, so it must be imposed on the sieve $\operatorname{span}(b_1,\dots,b_{L_x})$ of A8.
Restricting removes the obstruction whenever $L_x\le\operatorname{rank}\hat{\mathcal A}_\varphi$, for
which A8's \emph{upper end} $L_x\ll n^{1/2}\le n_{\mathrm m}$ is sufficient in generic position, 
\emph{and ``generic position'' is an $L_x$-dimensional injectivity statement, hence \textbf{Q-4}'s
object and not Q-3's}, Q-3 being the $n_{\mathrm m}$-dimensional Gram-matrix question
\cref{prp:tccrepair}(iii) and \cref{prp:norangeinv} decline.
\item[(iii.b)] And the interpolation step is unavailable in the regime this block inhabits. Read
in the Löwner order, the link condition says that the compression \eqref{eq:compress} and the compressed
scale operator \emph{each dominate a fixed positive multiple of the other}. What the interpolation step
needs is the corresponding comparison of their $-2\mu$ powers; inverting the two-sided inequality
delivers that comparison at the power $-1$, and \emph{raising a Löwner inequality to a power preserves
it only for exponents in $[0,1]$}, by operator monotonicity of the power map, available at
$2\mu\le1$ by (i)'s Heinz step and \textbf{refuted at every exponent above $1$ by (ii.a)}.
\eqref{eq:muhalf} forces $2\mu>\tfrac53$.
\emph{The two operators are named rather than symbolised}: the
compression and the compressed scale operator are the objects the clause is about, and the L\"owner
statement is more readable said than written. \emph{So the pair
(Hilbert-scale source, link) does not imply the spectral condition (F-f) consumes, at any configuration
% ARCHAEOLOGY.  Through dispatch 23 this clause printed the link condition as
% m^2 B <= A <= M^2 B and the target as A^{-2mu} <= c^2 B^{-2mu}, declaring bare A,
% B, m, M, c and an exponent p in place.  Bare M is the glyph C7 retires; bare m is
% the Landweber-Fridman counter (C3, C21); bare p is RESERVED for dim(Z); bare c
% matches none of the four readings Sec.1.7 declares closed for that letter; and
% A, B sit one font from cal A, cal A_cond, hat-A and from the B family, which
% Sec.1.7 governs as a CLOSED three-font rule (cal B with word and space
% subscripts, bold B, Fraktur B) while the letter in fact carried FIVE fonts.
% Fifteenth application of the C12/C24/C31/C35/C38/C44/C64/C66/C78/C90/C92
% write-it-out precedent.  Glossary C100.
\cref{prp:phirate} admits}, which is why (P-a.3) states the spectral condition directly and why this
clause is now a statement about a route \emph{declined} rather than about the route taken.
\cref{rem:pasrc} carries the full disposition, including the commutativity proviso under which
(iii.a)$\,+\,$(iii.b) \emph{would} deliver $\mu=s/(2a)$, and its price.
\end{enumerate}
\end{enumerate}
\end{lemma}

\begin{lemma}[The normalisation is a relaxation parameter, and \textup{(F-f)}'s bound is invariant under it]\label{lem:scaleinv}
Let $\mathcal L$ be a bounded linear operator from a Hilbert space into $L^{2}(\Pi)$, let
$\bar v_\dagger$ be in its domain, let $y^\delta\in L^{2}(\Pi)$ satisfy
$\|y^\delta-\mathcal L\bar v_\dagger\|\le\delta$, let $\omega_\ast>0$, and put
$\mathcal L_{\omega_\ast}:=\omega_\ast^{-1}\mathcal L$ and
$y^\delta_{\omega_\ast}:=\omega_\ast^{-1}y^\delta$.
\emph{$\omega_\ast$ is the generic scale}; the single value the estimator uses is $\omega_n$, fixed
at \eqref{eq:normconst} below. \emph{The star is not decoration.} Bare $\omega$ is the \textbf{sample
point} of the probability space throughout this project, it is the argument of the regular versions
fixed at (P6), of (P7a)'s ``for a.e.\ $\omega$'', and of the joint-measurability statements in
\cref{lem:adjfubini} and \cref{prp:compact}, and a generic positive scalar and a sample point cannot
share a bare glyph. The three
live objects are separated by form: bare $\omega$ the sample point, $\omega_n$ the estimator's
normalising factor, $\omega_\ast$ this lemma's generic scale. Then:
% ARCHAEOLOGY.  Glossary C67 asserted that "lowercase omega occurs in no file of
% this project"; that claim was false, at fifty-five sites, and the correction is
% recorded at C71.
\begin{enumerate}
\item[(i)] \emph{Same iterates, different relaxation.} The Landweber iteration
$v_{m+1}=v_m+\mathcal L^{\ast}_{\omega_\ast}(y^\delta_{\omega_\ast}-\mathcal L_{\omega_\ast} v_m)$ from
$v_0=0$ is identical, term by term, to the relaxed iteration
$v_{m+1}=v_m+\omega_\ast^{-2}\mathcal L^{\ast}(y^\delta-\mathcal Lv_m)$ for the unscaled pair. In
particular the normalisation is a choice of relaxation parameter and not a change of estimator.
\item[(ii)] \emph{Same stopping index.} For every $m$,
$\|\mathcal L_{\omega_\ast} v_m-y^\delta_{\omega_\ast}\|=\omega_\ast^{-1}\|\mathcal Lv_m-y^\delta\|$, so
the discrepancy rule at level $\tau_{\mathrm{DP}}\omega_\ast^{-1}\delta$ for the scaled pair stops at
exactly the index at which the rule at level $\tau_{\mathrm{DP}}\delta$ stops for the unscaled pair.
\item[(iii)] The source pair transforms, and in opposite directions. If
$\bar v_\dagger=(\mathcal L^{\ast}\mathcal L)^{\mu}\varsigma$ with $\|\varsigma\|\le\mathsf R_\ast$, 
$\varsigma$ being the \emph{spectral} source pre-image and $\mathsf R_\ast$ its norm bound, as in
\cref{lem:srctransfer} and (P-a.3), and \textbf{not} the $H^{s}$ radius $\mathsf R$ of
\cref{def:srcball}, then
$\bar v_\dagger=(\mathcal L^{\ast}_{\omega_\ast}\mathcal L_{\omega_\ast})^{\mu}\big(\omega_\ast^{2\mu}\varsigma\big)$,
so the source radius transforms as $\mathsf R_\ast\mapsto\omega_\ast^{2\mu}\mathsf R_\ast$ while the
data-error bound transforms as $\delta\mapsto\omega_\ast^{-1}\delta$.
\item[(iv)] \emph{Hence \textup{(F-f)}'s error bound is exactly invariant}: for every $\omega_\ast>0$ and
every $\mu>0$,
\begin{equation}\label{eq:scaleinv}
\big(\omega_\ast^{2\mu}\mathsf R_\ast\big)^{1/(2\mu+1)}\big(\omega_\ast^{-1}\delta\big)^{2\mu/(2\mu+1)}
\;=\;\mathsf R_\ast^{1/(2\mu+1)}\,\delta^{2\mu/(2\mu+1)} .
\end{equation}
\item[(v)] The stopping-index bound is not invariant, and that is the whole cost.
\textup{(F-f)}'s bound on $\hat m$ transforms as
$(\mathsf R_\ast/\delta)^{2/(2\mu+1)}\mapsto\omega_\ast^{2}(\mathsf R_\ast/\delta)^{2/(2\mu+1)}$, so the
iteration cap must absorb a factor $\omega_\ast^{2}$.
\item[(vi)] The \textbf{reference measure}'s normalisation is a second free scale, it is
\textbf{not} covered by \textup{(i)}--\textup{(v)}, and \eqref{eq:cap} is proportional to it. Fix
$c_\ast>0$ and replace $\Pi=\Pi_u\otimes\Pi_t\otimes\Pi_{\mathsf w}$ of \cref{def:spaces} by
$\Pi_u\otimes(c_\ast\,\Pi_t)\otimes\Pi_{\mathsf w}$, \emph{the \textbf{second} factor alone is
scaled, the other two are untouched}, with every other primitive of the model
held fixed;
% ARITY RETYPED 2026-08-10 (Phase C).  This clause printed a TWO-factor
% Pi = Pi_u (x) Pi_t against def:spaces's three, stale since the L-category
% instrument was introduced.  THE CONCLUSION IS TRUE VERBATIM AT EVERY L -- Pi_w
% is untouched, so every L^2(Pi) norm still scales by c_*^{1/2} and eq:cap still
% scales by c_* exactly -- so this is stale NOTATION, not a false statement.
% 5d retyped only the proof half (which already writes the three-factor Pi^*) and
% correctly refused a file it did not own; the seventh stranded item, closed here.
the replacement satisfies B1(iv) whenever $\Pi$ does. Then every $L^{2}(\Pi)$ norm is multiplied by
$c_\ast^{1/2}$, hence so are $c_\Pi$ and $\omega_n$ of \eqref{eq:normconst}, while $\hat\delta_n$ of
A5(i), 
$\max\{n^{-1/2}\sqrt{\ln n},r_n\}$, a functional of $P$ and $n$ with no $\Pi$ in it, is unchanged.
So the right-hand side of \eqref{eq:cap} is multiplied by $c_\ast$ \emph{exactly}.
\emph{The three free scales of this lemma are notated alike}: $\omega_\ast$ scales the operator,
$c_\ast$ scales the reference measure, $a_\ast$ scales the index.
% ARCHAEOLOGY.  The reference-measure scale was bare italic c through dispatch 21;
% the subscripted form is dispatch 22's.  Bare c with an argument is cor:udep's
% affine-weight slope c(z), live BARE in that corollary's own statement, and
% Sec.1.6's sans-c row names c(z) as the object the sans-serif constants must stay
% distinct from, so a bare c here would have been a third reading of one letter in
% one file.  The subscripted form joins the constant family c_Phi, c_Pi, c_C,
% c_{mu''}, bar-c_pi, whose defining property is that the subscript is never
% omitted.  Glossary C86.
\item[(vii)] \emph{A third free scale: the index scale, and it is why $c_\Phi$ carries no number.}
Fix $a_\ast>0$ and replace $(\Phi,\mathcal T,\Pi_t)$ by
$(a_\ast\Phi,\ a_\ast^{-1}\mathcal T,\ \Pi_t\circ(a_\ast\,\cdot))$,
every other primitive held fixed. \emph{The subscript is not decoration}: bare $a$ is the degree of
ill-posedness, live as $\mu=s/(2a)$ in the clause that closes the preceding lemma; the starred form
joins $\omega_\ast$ and $c_\ast$ as this lemma's free scales. The family
$\{w(t'\Phi(z)):t\in\mathcal T\}$ is \textbf{unchanged as a
set of functions}, $\bar w$ is unchanged at corresponding indices, $c_\Phi\mapsto a_\ast c_\Phi$ and $c_\Pi$
is \textbf{invariant}; so $\omega_n$ and \eqref{eq:cap} do not move. \emph{Hence normalising $c_\Phi$ is
a change of units in the index and pins nothing.} What is \textbf{not} covered: holding $\Pi_t$ fixed
while changing $c_\Phi$ is a \emph{different family}, not a change of units, and it moves $c_\Pi$ ---
exponentially in $\dim\mathcal T$ through $\int e^{2c_\Phi\|t\|_1}\Pi_t(dt)$.
\end{enumerate}

\emph{Accordingly, and this is a standing convention rather than a clause of the assumption block},
$\Pi$ is normalised to a \textbf{probability measure}, $\Pi_u(\IR)\,\Pi_t(\mathcal T)=1$, throughout this
project. \emph{B1(iv) leaves the normalisation entirely free}: (iv.a) constrains the shape of $\varpi_u$,
(iv.b) is scale-invariant, and (iv.c) bounds $c_\Pi$ without fixing $\Pi(\IR\times\mathcal T)$ --- and,
separately, without fixing the \textbf{dispersion} of $\Pi_t$ --- so
$c_\ast$
above ranges over $(0,\infty)$ inside the maintained block. \emph{Fixing it is a choice of units and not a
restriction on $P$}: $\Pi$ is an object the analyst chooses, on the B1 reading. \emph{No display in this
project needs rewriting}, each being stated in $c_\Pi$, in $\asymp$ or in $O_P$; \emph{the one display
whose numerical content moves is \eqref{eq:cap}}. \emph{The convention pins the \textbf{mass} scale
$c_\ast$ and no more}: the dispersion of $\Pi_t$ is left free by \textup{(iv.b)} and \textup{(iv.c)},
and it enters $c_\Pi$ exponentially through $\int e^{2c_\Phi\|t\|_1}\Pi_t(dt)$. \emph{So $\bar m_n$
remains a number times a free constant, and any instantiation of \eqref{eq:cap} must state the $\Pi_t$
it is computed at.} Pinning the mass is what removes the $c_\ast$ ambiguity of \textup{(vi)} and
nothing else. The inflation figures computed below, and any
simulation that reads them, inherit that constant. \emph{That is not hypothetical}: a Monte Carlo run
of this estimator at an unnormalised $\Pi_t$ reported a cap inflation about $2.9$ times the correct one,
and the discrepancy is exactly $c_\ast$, not an error in \textup{(v)}.
Why \textup{(i)}--\textup{(v)} do not already cover it, since they look as though they should. They
rescale the \emph{operator} by $\omega_\ast$ at a \emph{fixed} $\Pi$, and clause (ii) preserves the
stopping index precisely \textbf{because} the discrepancy threshold rescales with the residual, both
carrying $\omega_\ast^{-1}$. Under a change of $\Pi$ that symmetry breaks: the residual is an $L^{2}(\Pi)$
norm and moves by $c^{1/2}$, while the threshold $\tau_{\mathrm{DP}}\hat\delta^{\mathrm{eff}}_{n,k}$ of
\eqref{eq:lfsched} is built from $\hat\delta_n$, $h_u$ and $\hat\varepsilon_{k-1}$ and does not move at
all. One rescaling is a relaxation parameter; the other is a change of units in which the noise
level is denominated, and only the first is invariant.
Why this is a lemma and not a sentence. \Cref{def:lfstep} normalises so that
$\|\omega_n^{-1}\hat{\mathcal A}_\varphi\vert_{\mathcal S_n}\|\le1$, which is (L1) of
\cref{prp:phirate} Step~2 and
is a hypothesis of (F-f). \emph{That bound is \textbf{pairing-specific} and this lemma is not}: (F-f)
reads the operator norm out of the domain Hilbert space whose inner product defines
$\mathcal L^{\ast}$, so the constant in \eqref{eq:normconst} has to cover the pairing the iteration
runs in --- $L^{2}(\hat\pi)$ for the feasible recursion, $L^{2}(\pi)$ for the analysed one, and
\eqref{eq:normconst} carries the maximum of the two conversion costs by
\cref{prp:hatAprop}\textup{(i$'$)}. \emph{Clauses \textup{(i)}--\textup{(vii)} below are indifferent to
which is in force}, each being stated for an abstract $\mathcal L$ and an arbitrary $\omega_\ast>0$;
what is \emph{not} indifferent is the single number $\omega_n$ they are instantiated at, and that is
fixed in \cref{def:lfstep} and not here. \emph{The normalisation does not cost nothing, because the
only norm bound available in this file diverges}: \cref{prp:hatAprop}(i) gives
$\|\hat{\mathcal A}_\varphi\|\le\mathsf c\,h_u^{-1}h_x^{-1/2}$, so the required $\omega_n$ diverges,
and a divergent rescaling moves $\|\mathcal L\|$, the data error \emph{and} the source radius at once.
Clauses (iii)--(iv) are what make the movement cancel: the exponents $1/(2\mu+1)$ on the radius and
$2\mu/(2\mu+1)$ on the noise are exactly the pair for which it does, which is why \eqref{eq:linlf} must
be written with $\mathsf R_\ast$ shown and not absorbed into an $O_P$.
And the cancellation does not depend on the direction of the movement. If the unscaled norm were
instead $o(1)$, which is plausible, since near-orthogonal $\hat g_{\cdot,\cdot,\cdot,i}$ would give
$\|\hat{\mathcal A}_\varphi\vert_{\mathcal S_n}\|\asymp(n_{\mathrm m}h_u)^{-1/2}$, then
$\omega_n\to0$ would be admissible and (iv) would still hold verbatim, clause (iv) being stated for
\emph{every} $\omega_\ast>0$. Nothing in Part~H
therefore needs $\|\hat{\mathcal A}_\varphi\vert_{\mathcal S_n}\|\asymp1$, and that is the whole of what
this lemma buys.

What the file does and does not know about the unscaled norm. Only the \emph{upper} bound of
\cref{prp:hatAprop}(i) is
available. A matching \emph{lower} bound on $\mathcal S_n$ is the statement
$\|\hat{\mathcal A}_\varphi\zeta\|_{L^{2}(\Pi)}\gtrsim\|\zeta\|$ for $\zeta\in\mathcal S_n$, which is
\emph{one side of the link condition}, and the link condition on the sieve at a random base point is
\textbf{Q-4}'s object, not Q-3's. \emph{Q-3 is a different question and is classified as not a theory
obligation}: it asks for $\lambda_{\min}(\mathbf G_n(\cdot))$ on the $n_{\mathrm m}$-dimensional Gram
matrix of \eqref{eq:gram}, which governs \cref{prp:tccrepair}(iii)'s radius and
\cref{prp:norangeinv}'s proviso and which those two results correctly decline. \emph{Two objects, two
dimensions, two open items}: the sieve-restricted
injectivity modulus is $L_x$-dimensional and is assumed, at (P-a.3); the Gram matrix is
$n_{\mathrm m}$-dimensional and is declined. \Cref{lem:srctransfer}(iii)'s ``generic position'' is the
same $L_x$-dimensional object and carries the same redirect.

What is not free is (v), and it is priced here rather than left as a symbol.
\cref{def:lfstep} pays for (v) by sizing the cap at $\bar m_n\ge\omega_n^{2}\hat\delta_n^{-2}$
rather than at $\hat\delta_n^{-2}$, and \emph{the inflation is polynomial and not negligible}. With
$\omega_n\asymp h_u^{-1}h_x^{-1/2}$ from \eqref{eq:normconst}, $\omega_n^{2}\asymp h_u^{-2}h_x^{-1}$.
\emph{That $\asymp$ survived the 2026-08-27 amendment of \eqref{eq:normconst} and the figures below with
it, and the reason is which kind of quantity moved}: the display now carries
$\max\{\|1/\pi\|_\infty^{1/2},\varepsilon_\pi^{-1/2}\}$ where it carried $\|1/\pi\|_\infty^{1/2}$, and
both entries are fixed positive numbers free of $n$ --- the first finite by C1(ii), the second a
constant \cref{def:lfstep} fixes before the data are seen --- so every \emph{exponent} in this
paragraph is unchanged and the \emph{constant} multiplying $\bar m_n$ rises by at most
$(\varepsilon_\pi\|1/\pi\|_\infty)^{-1}$, and by nothing at all at $\varepsilon_\pi\ge\|1/\pi\|_\infty^{-1}$.
\emph{An instantiation must therefore now state the $\varepsilon_\pi$ it is computed at as well as the
$\Pi_t$}, on exactly the reading clause \textup{(vi)} imposes for $c_\ast$.
At the frontier-clearing configuration $(\mu,s,a,\alpha)=(2,3,\tfrac34,0.32)$ of
\cref{rem:h6bandwidth}(5), where $h_u\asymp n^{-0.04}$ and C5$'$(ii)$\,\wedge\,$(P-g) leave
$h_x\in(n^{-0.44},n^{-0.144})$:
\[
\omega_n^{2}\;\asymp\;n^{0.08}h_x^{-1}\;\in\;\big(n^{0.224},\,n^{0.52}\big),
\qquad
\bar m_n\;\ge\;\omega_n^{2}\hat\delta_n^{-2}\;\in\;\big(n^{0.864},\,n^{1.16}\big)
\quad\text{against}\quad
\hat\delta_n^{-2}=n^{0.64} .
\]
% WINDOW AND BOTH LOW ENDS RECOMPUTED WITH (P-g)'s h_u.  The ceiling is now
% h_x = O(n^{-9 alpha/(8(s - 1/2))}) = O(n^{-9(0.32)/20}) = O(n^{-0.144}), against
% O(n^{-0.128}) before, so the window NARROWS at its cheap end and the three numbers
% that move are the ceiling, omega_n^2's lower end (0.08 + 0.144 = 0.224, was 0.208)
% and bar-m_n's lower end (0.224 + 0.64 = 0.864, was 0.848).  The upper ends are set
% by C5'(ii)'s floor n^{-0.44}, which (P-g) does not touch, and are unchanged.
\emph{At $n=2000$ that is an inflation of the iteration budget by a factor between about $5$ and about
$50$}, $130$ Landweber steps become between $710$ and $6750$, \emph{and the factor is decreasing in
$h_x$, so the upper end of the $h_x$ window is also the cheap end of the cap}. The cap is a computational
budget and not a rate condition, so nothing in \eqref{eq:phirate} moves with it; but it is a real cost of
the normalisation, a reader tuning the estimator will meet it, and it is stated here rather than left to
be discovered. The alternative to paying it is a lower bound on
$\|\hat{\mathcal A}_\varphi\vert_{\mathcal S_n}\|$, which is Q-4 again.
\end{lemma}

\begin{remark}[Why \textup{(P-a.3)} is spectral and at the compression, what H1 still does, and what a derived $\mu$ would cost]\label{rem:pasrc}
Read this before \cref{rem:phirscope}. Four clauses: why the Hilbert-scale pair does not deliver
(F-f)'s hypothesis, what (P-a.3) asserts instead, why that is a \textbf{weakening}, and the alternative,
priced.

(1) The Hilbert-scale pair does not deliver what (F-f) consumes. The pair is (a) a
\emph{Hilbert-scale} source condition,
$\|\bar v^{(k)}_\dagger\|_{H^{s}}\le\mathsf R$, and (b) the \emph{link} condition of
\cref{lem:srctransfer}(iii), $\|\hat{\mathcal A}_{\hat\varphi^{(k-1)}}\zeta\|_{L^{2}(\Pi)}\asymp\|\zeta\|_{H^{-a}}$
on $\mathcal S_n$; what (F-f) consumes is the \textbf{spectral}
condition $\bar v^{(k)}_\dagger=(\mathcal L^{\ast}\mathcal L)^{\mu}\varsigma$ with
$\|\varsigma\|\le\mathsf R_\ast$. The step from (a) and (b) to that is Hilbert-scale operator
interpolation, and it is available only at $2\mu\le1$. Read in the L\"owner order, (b) says that the
compression \eqref{eq:compress} and the compressed scale operator \emph{each dominate a fixed positive
multiple of the other}; what the step needs is the corresponding comparison of their $-2\mu$ powers.
Inverting the two-sided inequality gives that comparison at the power $-1$, and \emph{raising it to a
power preserves it only for exponents in $[0,1]$}, by operator monotonicity of the power map
, Heinz's range, and \citet[\S8.5]{engl2000}'s interpolation range, which is the pin H1 itself names.
Here the exponent is $2\mu$, and \eqref{eq:muhalf} forces $2\mu>\tfrac53$: \emph{no configuration this
block admits escapes}, and \cref{lem:srctransfer}(ii.a) is a proof \emph{in this file} that the device
fails at every exponent above $1$. Finite-dimensionality of $\mathcal S_n$ does not rescue it. On
$\mathcal S_n$ range
membership is vacuous, every positive-definite operator there is surjective, so all the content of
the step is the \emph{uniform-in-$n$ pre-image norm bound}, and that bound \emph{is} the interpolation
inequality. Finite rank converts the failure from qualitative to quantitative; it does not remove it.
% ARCHAEOLOGY.  Rounds 1--11 gave (P-a) ingredients (a) and (b) and Step 2 (L5)
% then asserted (F-f)'s spectral hypothesis.  This passage printed the same
% condition on the bare letters A, B, m, M, c and p through dispatch 23 and is now
% written out with lem:srctransfer(iii.b), where the reasons are given; nothing it
% states changed.

(2) (P-a.3) states the hypothesis where the theorem consumes it. It asserts the spectral
condition directly, at the operator \cref{def:lfstep} actually regularises, named as the compression
\eqref{eq:compress}. Two consequences. \emph{Step~2 (L5) becomes two lines}, no interpolation, no Heinz
range, no L\"owner, and \cref{lem:srctransfer} is used nowhere on the rate route, its clause (ii.a)
surviving as the reason the \emph{inherited nonlinear} theorem is unavailable (\cref{rem:lfnonlinear}(3))
and its clause (iii) as the reason a Hilbert-scale route \emph{would} need the sieve. And the
operator is named correctly. $\mathcal L^{\ast}\mathcal L=\omega_n^{-2}\hat{\mathcal P}^{\,\mathcal S}_{\varphi_\circ}$
is the compression $P_{\mathcal S_n}\hat{\mathcal A}^{\ast}\hat{\mathcal A}\vert_{\mathcal S_n}$, not
$\mathcal R\big((\hat{\mathcal A}^{\ast}\hat{\mathcal A})^{\mu}\big)\cap\mathcal S_n$; the two are nested
in neither direction, and (L2) is the witness that the compression is what the iteration sees, since
``the projection is the identity on the iterates'' holds \emph{because} $P_{\mathcal S_n}$ was absorbed
into $\mathcal L^{\ast}$. \cref{lem:scaleinv}(iii) cannot bridge the gap, it is stated at one fixed
domain and handles only the scalar rescaling.

(3) (P-a.3) is a \textbf{weakening} of the pair, and this is the decisive point. The alternative
is
to keep (a)+(b) and add a \emph{commutativity} proviso: that $\hat{\mathcal A}^{\ast}\hat{\mathcal A}$
and the compressed scale operator share an eigenbasis on $\mathcal S_n$, under which
\cref{lem:srctransfer}(ii.b)'s device licenses interpolation at every $\mu>0$. Under that proviso
(a)+(b) \textbf{imply} (P-a.3). So the hypothesis actually assumed here is strictly weaker than the
alternative's, and it needs strictly less proved, assuming the conclusion of a step whose only licence
is an unverifiable spectral-alignment condition is the \emph{weaker} assumption, not the stronger one.
And nothing is lost that was ever owned. H1's matching is a statement about the \emph{population}
operator $\mathcal A_{\mathrm{cond}}$ on all of $\mathcal V$; (P-a.3) is a statement about a
\emph{realised, finite-rank, sieve-compressed} operator at a \emph{random} base point, and
\cref{lem:srctransfer}(iii) already records that the link condition \emph{cannot} hold on $\mathcal V$
because $\hat{\mathcal A}$ has finite rank by \cref{prp:hatAprop}(ii). H1 therefore never delivers
$\mu=s/(2a)$ for the operator (P-a) is about: (P-a) assumes its own matching, and stating it spectrally
makes that assumption visible rather than adding one.

(4) What H1 still does, and what a derived $\mu$ would cost. H1 is \textbf{not} retired. It fixes
the scale $H^{s}$ in which \cref{def:srcball}, A5(vi$'$), (P-c) and (P-g) are written; it supplies A2's
Picard form; it carries the whole smoothness-budget table. What is withdrawn is only the clause
``$\mu=s/(2a)$'' from \cref{prp:phirate}'s hypothesis line. \emph{A reader who wants $\mu$ derived rather
than primitive may have it, at a stated price}, and the price is exactly one proviso: assume in addition
that $\hat{\mathcal A}^{\ast}_{\hat\varphi^{(k-1)}}\hat{\mathcal A}_{\hat\varphi^{(k-1)}}$ and the
compressed scale operator commute on $\mathcal S_n$; then (a)+(b) at index $s/(2a)$ give (P-a.3) at
$\mu=s/(2a)$ with $\mathsf R_\ast\le c\,\mathsf R$, $c$ being the interpolation constant, which is a
\emph{third} constant \cref{def:lfstep} would then have to print before the data are seen. \emph{That
proviso is not diagnosable and no sufficient condition for it is known}: it asserts that a random,
rank-$\le n_{\mathrm m}$ operator built from $2n$ design points has its singular system aligned with a
spline--Sobolev scale. It is recorded here as the alternative, with its content stated plainly, and it is
not assumed. \emph{The honest positioning statement, which belongs in the paper and not in a footnote}:
$\mu$ is primitive at (P-a.3); what pins it is H1 \emph{plus} a matching proviso, stated and not proved;
and that is the same structure the closest antecedent in the literature carries, where a spectral source
condition and a link condition are two separate assumptions rather than one derived from the other.

(5) The consequence for the \textbf{symbol}: the glyph $\mu$ names two numbers, and they are not
the same number.
\begin{itemize}
\item[$\mu_{\text{\textup{H1}}}$], \emph{H1's index}, $s/(2a)$. It is the source index of
A2, at the \textbf{population} operator $\mathcal A_{\mathrm{cond}}$ on \textbf{all of
$\mathcal V$}, at $\theta_\dagger$. Its ingredients $a$ (the decay of $\sigma_j(\mathcal A_{\mathrm{cond}})$)
and $s$ (the order of the \emph{classical} Sobolev space in which $\vdagger$'s source ball is written,
\emph{not} of a $\mathcal D$-generated scale) are
likewise objects at $\theta_\dagger$, by H1's own sentence.
\item[$\mu$], \emph{(P-a.3)'s index}, and it is \textbf{primitive}. It is the exponent of a
\textbf{spectral} source condition at the \textbf{finite-rank sieve compression} \eqref{eq:compress} of the
\textbf{realised} operator $\hat{\mathcal A}$, at the \textbf{anchor $\theta_\circ$} of
\cref{def:crossfit}(b), which \cref{def:crossfit}(b) also states makes $a$, $\mathsf R_\ast$ and
$\hat\delta_n$ $\theta_\circ$-specific.
\end{itemize}
\emph{Binding reading.} \textbf{The symbol $\mu$, bare, is
(P-a.3)'s primitive index everywhere in Parts~H and~I}, and in particular in \eqref{eq:muhalf},
\eqref{eq:frontier}, \eqref{eq:deltalog}, \eqref{eq:phirate}, $\kappa=2\mu/(2\mu+1)$ and every rate display
built from them. A sentence that computes $\mu$ from $(s,a)$ is a sentence about
$\mu_{\text{\textup{H1}}}$, and it transfers to those displays only under the declined proviso of clause
\textup{(4)}.
\emph{The distinction cuts both ways}: a description of scope in terms of $s>a+\tfrac12$ \textbf{excludes}
designs the theorem covers, a design with $\mu=2$ at the compression and $s/(2a)=0.4$ at the population
operator satisfies every hypothesis of \cref{thm:rootn} and fails $s>a+\tfrac12$, and it
\textbf{claims} designs the theorem does not, since $s>a+\tfrac12$ says nothing about the compression's
spectral index. \emph{The conclusion of \cref{thm:rootn} is untouched}: no proof step computes $\mu$ from
$(s,a)$.
% ARCHAEOLOGY.  Three sites compute the H1 index and read the result into a
% condition on the primitive one, and each is qualified in place:
% rem:rootnscope(1); the "mild ill-posedness" and mu > 1/2 lines of the
% smoothness-budget table, with the sentence deducing mu > 5/6 from s > a + 1/2 at
% a <= 3/4; and the flagship configuration (mu, s, a, alpha) = (2, 3, 3/4, 0.32),
% in which s/(2a) = 2 EQUALS the assumed mu, an accident of the chosen numbers.
% ARCHAEOLOGY.  Why the bare symbol is not renamed: mu carries (P-a.3)'s meaning at
% every rate display in three Parts and in every line of the budget table, while
% mu_H1 appears in one sentence of H1 and in the interpretive lines qualified
% above; under this file's rule the symbol with fewer sites moves, and the move is
% to a subscripted NAME used only where the two must be told apart.
\end{remark}

\begin{remark}[Why \textup{(P-a)}'s three clauses are w.p.$\to1$ envelopes and not $O_P$ statements]\label{rem:paenvelope}
An $O_P$ reading of (P-a)'s first two clauses is strictly weaker than what
\eqref{eq:fixedpoint}'s base case and \cref{lem:pathcont} consume, and the gap cannot be closed by
carrying a constant.
% ARCHAEOLOGY.  Rounds 1--11 wrote those two clauses as O_P statements.
The base case of Step~3's induction is
$\varepsilon_0\le\hat\varepsilon_0$ \emph{with probability tending to one}; an $O_P$ hypothesis gives
instead $\varepsilon_0\le C_0\hat\varepsilon_0$ w.p.$\to1$ for a $C_0$ that is not a number but a family
$\{C_{0,\epsilon}\}$ indexed by the tolerance. \emph{Carrying $C_0$ is not available}, and the reason is
that $\mathsf c_{14}$ and $\mathsf c_{15}$ are \textbf{part of the estimator}: \cref{def:lfstep} prints
them before the data are seen, so ``choose $\mathsf c_{14}\ge\mathsf c_{16}\max\{1,C_0^{2}\}$'' either
says nothing or makes the estimator, hence \eqref{eq:phirate}, indexed by $\epsilon$, at which
point the conclusion is no longer $O_P(\delta_n)$ for \emph{one} estimator. So the envelope is
declared in the hypothesis, which is how \cref{def:lfstep} already speaks of $\hat\varepsilon_0$.

The conversion device is owned by this block and the cost is a logarithm that is already paid.
\cref{prp:stopfeas}(iii) is the same manoeuvre one page earlier: an $O_P(a_n)$ statement becomes a
w.p.$\to1$ envelope at level $a_n\ln n$, at the cost of a divergent factor. Applied here it replaces
$\hat\varepsilon_0$ by $\hat\varepsilon_0\ln n$, which changes $\varsigma_0$ by nothing at all, because
\eqref{eq:prelimrate} is a strict inequality between \emph{polynomial} exponents and absorbs any
$\ln n$. The same repair closes the second site. Step~3 needs
$\mathsf c_{15}\ge\mathsf c_{17}=2c_\mu\mathsf R_\ast^{1/(2\mu+1)}$ with $\mathsf c_{15}$ printed in
advance; that is an inequality between numbers exactly because (P-a.3) fixes $\mathsf R_\ast$
deterministically and asserts the source condition on an event of probability tending to one, rather than
delivering a pre-image of norm $O_P(\mathsf R_\ast)$. \emph{One repair, two sites, and it is
\cref{prp:stopfeas}(ii)'s point applied to the hypothesis that generates it}: an $O_P$ bound cannot
supply a fixed constant to an object printed before the data.

Where the same discipline was already correct, recorded so that it is not re-litigated.
A5(vi$'$) is now also read as a w.p.$\to1$ envelope at the \emph{fixed} radius $\mathsf R$, 
$P(\max_{k\le\bar k}\|\hat\varphi^{(k)}-\vdagger\|_{H^{s}}\le\mathsf R)\to1$, because
\cref{lem:pathcont} needs membership of a \emph{deterministic} ball \cref{def:srcball} fixes, and
$\mathsf R$ is not a free tuning constant: it appears at every site \cref{def:srcball} enumerates,
including, through \cref{lem:unifop} Step~4, in the $x$-smoothing bias. \emph{The hypothesis is therefore
an existential over \textbf{one} $\mathsf R$ shared by all of them}, and it is stated that way in
A5(vi$'$) rather than left to be assembled by a reader. Step~4(a.1) and (a.2) are untouched by any of
this: they run on the deterministic $\hat\varepsilon_k$.
% ARCHAEOLOGY.  This paragraph carried its own four-entry copy of def:srcball's
% enumeration through dispatch 20 -- the shortest of the four the corpus maintained.
\end{remark}

\begin{remark}[F-6 / O10: the nonlinear transfer, and why it is declined rather than performed]\label{rem:lfnonlinear}
The registered route for \cref{prp:phirate} is nonlinear Landweber--Fridman with a tangential cone
condition and the discrepancy principle, after \citet{hanke1995} and \citet[Ch.~11]{engl2000}, with the
operator-noise model of Dunker, Florens, Hohage, Johannes \& Mammen (2014). That route does not
transfer, and the obstruction is not a technicality. Three things go wrong, of increasing severity,
and each is proved or cited rather than reported.

(1) The cone condition is false on $\mathcal V$, and on the sieve, where the estimator lives,
it is contingent on a product no assumption constrains. \Cref{prp:tccfails} refutes it on every
$L^{2}(\pi)$-ball, for every bounded linearisation. \emph{That is not the same as refuting it for a
sieve method}, and the distinction is now printed rather than elided: on $\mathcal V_n\cap\mathcal S_n$
the spike is not representable at all, the exact first-order \emph{necessary} condition is
\eqref{eq:a6sieve}, an ``only if'', for the reason \cref{rem:tccwhy}(3) prints, and
whether it holds turns on $h_xL_x$, unconstrained by C5$'$ and A8, and free over a range containing
both verdicts inside the block's own admissible window (\cref{rem:tccwhy}(3)). \emph{So reason (1) is
the weakest of the three, not the strongest}, and it is repairable in the ambient space too, by A6b and
\cref{prp:tccrepair}, at the cost of replacing $\hat T$ by $\hat T_{h_x}$.

(2) Range invariance, the only route to a cone condition that does not simply assume it, fails.
\Cref{prp:norangeinv}. The strategy memo caps this attempt at one round; the round is spent.

(3) The decisive one, because it is regime-independent: the inherited rate theorem is not
available above $\mu=\tfrac12$, and \eqref{eq:frontier} needs $\mu>\tfrac12$. The rate theorem for
nonlinear Landweber under a tangential cone condition and the discrepancy principle
\citep{hanke1995,engl2000} delivers $\|x_{\hat m}-x^\dagger\|=O(\hat\delta^{2\mu/(2\mu+1)})$ for source
index $\mu\in(0,\tfrac12]$, and above $\tfrac12$ \emph{the analysis gives no improvement}, the
qualification of the method, as its proof uses operator monotonicity of $z\mapsto z^{2\mu}$, which is
the barrier \cref{lem:srctransfer}(ii.a) exhibits.
Two things must be kept apart. The claim made is
that \emph{the theorem is unavailable} above $\tfrac12$; the claim \emph{not} made is that the
\emph{rate} fails there, that would be a saturation statement, requiring a lower bound, and no such
result is proved or cited in this file. The disposition needs only the first: one cannot inherit
a theorem one does not have. At the qualification index the programme then fails, and fails even with
a perfect propensity: at $\mu=\tfrac12$ the available rate is $\hat\delta_n^{1/2}$, and since
$\hat\delta_n\ge\hat\delta_n^{(T)}=n^{-1/2}\sqrt{\ln n}$ by A5(i)'s floor, \emph{not} by
\cref{lem:opnoise}, which bounds the operator error from above and delivers no floor, the best
attainable is $\delta_n\asymp n^{-1/4}(\ln n)^{1/4}$, which is \emph{not} $o(n^{-1/4})$; at
$\mu<\tfrac12$ it fails polynomially. So the inherited rate and the paper's own feasibility
frontier are incompatible, and the incompatibility is exactly at $\mu=\tfrac12$, and it is so on the
sieve and off it alike, which is why this and not (1) is the reason to lead with.
And the ceiling is not an artefact of the regularisation literature's way of stating things: it is
imposed, in the same place and for the same reason, in the closest econometrics treatment.
\citet{centorrinofeveflorens2019} run a nonlinear ill-posed inverse problem of exactly this type and cap
their own source index at $\beta\le1-\epsilon$; in the normalisation of this file $\beta=2\mu$, so their cap
\emph{is} $\mu<\tfrac12$, and it is the Heinz range of \cref{lem:srctransfer}(ii.a) seen from the applied
side. That the same number is reached from operator interpolation here and from an econometric
construction there is the strongest available evidence that $\tfrac12$ is a property of the problem and not
of the proof.

What is done instead, and at what price. \Cref{def:lfstep} linearises and
runs \emph{linear} Landweber--Fridman at each stage, absorbing the nonlinearity into the noise level
through \cref{prp:hatAprop}(iv)'s exact quadratic remainder. Linear Landweber with the discrepancy
principle does not saturate, so the source index is unrestricted from above; what the route needs
instead is that the absorbed remainder be dominated by the noise, which is (P-f) and is
$\mu>\tfrac56$. The condition that removes the inherited route at $\tfrac12$ is the condition
that switches the replacement route on at $\tfrac12$, and the replacement then charges a further
$(\tfrac12,\tfrac56]$. \emph{And the replacement's own load-bearing input must be named in the same
sentence}: it reaches the same place \textbf{under (P-a.3)}, a random \emph{spectral} source condition at
a data-dependent base point, imposed on the finite-rank compression \eqref{eq:compress} of the realised
sample operator and derived from nothing (Q-4). Sold without (P-a.3) the sentence is not true.
\emph{This is the substantive resolution of F-6 / O10 and it should be stated in the paper as such, not
buried:} the paper does not inherit a nonlinear-regularisation rate, it proves a linearised one under an
assumed source condition, and the reason it may is that its own frontier already lives on the side of
$\mu=\tfrac12$ where linearisation closes. \emph{And the assumption is stated spectrally rather than as a
Hilbert-scale pair, for the reason \cref{rem:pasrc} gives}: the pair does not imply the spectral condition
above $2\mu=1$, so a paper claiming to assume the pair and to apply (F-f) would be claiming an implication
\cref{lem:srctransfer}(ii.a) refutes.

\emph{What is given up.} Three things, and they should be printed.
\begin{enumerate}
\item[(a)] \emph{Global convergence.} Nonlinear Landweber with a cone condition converges from any
starting point in the ball where the condition holds. \Cref{def:lfstep} needs a preliminary
$\hat\varphi^{(0)}$ that is not merely consistent but consistent \emph{at a known polynomial rate}
$\hat\varepsilon_0$, because the stagewise schedule \eqref{eq:lfsched} is built from it. That is L-14's
object and is an input, not an output, of \cref{prp:phirate}. The paper must say where
$\hat\varphi^{(0)}$ comes from; a fixed number of Landweber steps from $0$ on the linearised equation
at $\varphi=0$ will not do without an argument (Q-1).
\item[(b)] A sieve, and it is given up twice over. \Cref{lem:srctransfer}(iii.a) forces any
link-type condition onto $\mathcal S_n$, and (P-a.3) states its source condition on the compression to
$\mathcal S_n$ for the same reason, so the estimator is a \emph{sieve} LF and A8's upper end is
consumed by the $\varphi$-step as well as by the balancing solve. But that is not the only thing
the sieve costs, and the second cost was not booked. Even granting a cone condition on
$\mathcal S_n$, the inherited theorem is stated for \emph{unprojected} nonlinear Landweber, with the
cone condition holding on a ball of the \emph{ambient} space containing the solution;
\cref{def:lfstep}(b) projects onto $\mathcal S_n$, projection does not commute with $\hat T_{h_x}$, and
$\vdagger\notin\mathcal S_n$. So the sieve restriction also moves the target to the sieve projection of
$\vdagger$ and would oblige a re-derivation of the cited theorem, not merely its citation. \emph{This
is a defect of the inherited route, not of \cref{def:lfstep}}, whose analysis is conducted on
$\mathcal S_n$ throughout and pays for the projection at (P-c).
\item[(c)] The Dunker et al.\ noise model is adopted and their regulariser is not. The
iteratively regularised Gauss--Newton method \citep{dunker2021} is the standard answer to exactly the
qualification problem in (3), attaining source indices above $\tfrac12$. \emph{Declining IRGN and then
needing $\mu>\tfrac56$ is a live tension in the project's positioning}, and the honest options are two:
adopt IRGN, or state, as this file does, that \emph{under \textup{(P-a.3)}} the linearised route
reaches the same place, at a threshold that is the same at $\tfrac12$ and stricter at $\tfrac56$. The
second is cheaper and is what is proved; the first should be named in the paper as the alternative and
priced, because a methods referee who knows the regularisation literature will ask, and will ask
\emph{about \textup{(P-a.3)}} first.
\emph{Two citations must be kept apart}: the IRGN / above-$\tfrac12$ claim
belongs to \citet{dunker2021}, while the \emph{operator-noise model} this remark's opening adopts is
Dunker, Florens, Hohage, Johannes \& Mammen (2014). Both are needed and they are different papers.
\emph{And the reason for declining IRGN is not that the linearised route is ``cheaper''}: it is
that IRGN, as the literature states it, treats $\hat T'_\varphi$ and $\hat T'^{\ast}_{\varphi_\dagger}$
% ARCHAEOLOGY.  One bib key was carrying both papers; the key was swapped and both
% now resolve.
as operators with the same range and \emph{does not model the estimation error arising from using
regression residuals in place of the true error}, which is exactly the $h_u^{-2}$ object
\eqref{eq:effnoise} and (P-f) exist to price. \emph{The route that would remove (P-f) is the route that
does not model what (P-f) is about}, and that, not cost, is why it is declined.
\end{enumerate}
\end{remark}

\begin{remark}[What \cref{prp:phirate} does not claim]\label{rem:phirscope}
(1) It is a rate, not a sharpness statement. \eqref{eq:phirate} is an upper bound on
$\|\hat\varphi^{\pi}-\vdagger\|_{L^{2}(\pi)}$, \emph{the norm being taken at the infeasible reading of
\cref{def:lfstep}\textup{(b)} and clause \textup{(7)} below being where that matters}. Nothing here says that $\delta_n$ cannot be improved, and
nothing here says that $\delta_n=o(n^{-1/4})$ is \emph{necessary}. \emph{In particular this file does
not reinstate the strategy memo \S2.5(iii) reading} that ``$\delta_n=o(n^{-1/4})$ is sharp for this
moment function'': that is a necessity claim about the estimator, \cref{prp:sharpquad} deliberately
does not prove it, and \cref{rem:sharpscope}\textup{(1)} blocks it. What \cref{thm:rootn} may inherit
from Part G is the \emph{lower bound} on the second-order form, and what it inherits from here is the
\emph{upper bound} \eqref{eq:phirate}; the two are different objects and neither is a sharpness claim
about $\hat\varphi$.

(2) The logarithm is carried in the symbol, not absorbed into a power of $n$.
$\delta_n$ in \eqref{eq:phirate} is $(\hat\delta_n\ln n)^{2\mu/(2\mu+1)}$, and A5(ii) is
that same \emph{rate} by \eqref{eq:deltalog}; \emph{so the symbol in the clause and the symbol in the
conclusion denote one sequence, and that is all this clause settles}. \emph{It does \textbf{not} follow
that the clause is discharged, and the line printed here through 2026-08-27 --- ``this proposition
therefore discharges the clause at the object the block names'' --- confused the two}: A5(ii) is stated
at $\hat\varphi$ and \eqref{eq:phirate} at $\hat\varphi^{\pi}$, so the \emph{rates} agree and the
\emph{objects} do not, which is clause \textup{(7)}. \emph{A downstream result that read A5(ii) as log-free and substituted a
power of $n$ would be off by $(\ln n)^{2\mu/(2\mu+1)}$}; \cref{rem:rootnlog} carries the disposition.
\emph{The logarithm is not absorbed anywhere}: it is genuinely in the rate,
% ARCHAEOLOGY.  Through dispatch 9 A5(ii) carried no logarithm and this remark
% carried the standing warning as a live one; the amendment is dispatch 10's and
% the warning is discharged.
\cref{prp:stopfeas}\textup{(iv)} shows the \emph{strictness} of \eqref{eq:frontier} pays for it in
$\delta_n=o_P(n^{-1/4})$ and in $\delta_n^{2}=o(n^{-1/2})$, and \cref{rem:rootnlog}\textup{(3)} says which
of the four rate products \cref{thm:rootn} needs are paid for that way and which one, A5(iv)'s
$\delta_n\epsilon_n$, is not.

(3) It says nothing about $\hat\ell$, $\hat L$ or $\hat b$. A5(iv) prices
$\|\hat\ell-\ell_\star\|_{L^{2}(P)}$ and is L-13's business; obligations \textup{(O2)} and
\textup{(O3)} of \cref{rem:obligations} are owed by the construction of $\hat L$ and are \emph{not}
discharged here. \emph{This file discharges \textup{(O4)} and nothing else on that list}, and
discharges it with a negative answer on A6, which is a discharge and not an evasion: the obligation was
to verify, and the verification has an outcome.

(4) It does not identify $\vdagger$. A1b(iii) is assumed, as everywhere; \eqref{eq:phirate} is a
rate for an estimator of an object assumed identified, and the discrepancy principle stops on the
residual, which is informative about $\varphi$ only to first order.

(5) It is a rate for the \emph{infeasible reading} of the iteration \cref{def:lfstep} defines --- clause
\textup{(7)} being where that qualification is discharged into five parts --- and that iteration, under
\emph{either} reading, needs four
inputs and two constants the analyst must supply. The four inputs are $\hat\delta_n$ (hence $r$ and
$d_Z$), the bandwidth $h_u$ (now a function of $\hat\delta_n$), the source index $\mu$ through
$\kappa$, and the preliminary's declared rate $\hat\varepsilon_0$. \emph{The two constants are
$\mathsf c_{14}$ and $\mathsf c_{15}$ of \eqref{eq:lfsched}}: they are not identified by the data, they
must dominate the analytic constants $\mathsf c_{16}$
of \eqref{eq:effnoise} and $\mathsf c_{17}$ of \eqref{eq:linlf}, which depend on the kernel, on the
\emph{spectral pre-image} radius $\mathsf R_\ast$ of (P-a.3) and on (F-f)'s $c_\mu$, and \textbf{not} on
the $H^{s}$ radius $\mathsf R$, which is a different object, and if they are set too small the implemented
level sits below the true data error, (L3) fails and (F-f) is void. \emph{So the schedule is
\textbf{computable}, which is what \cref{def:lfstep} needs, and it is not \textbf{known}}; the
distinction is stated in \cref{def:lfstep} and its practical size is what probe CP-13 measures.
\emph{A rule that used only the data would not do, and the reason has a sharp part and a cited part}:
the realised residual cannot bound the unobservable error from below, because
$\|\hat{\mathcal A}v_\dagger\|\ll\|v_\dagger\|$ is what ill-posedness \emph{is}; and the general
statement, that a purely data-driven choice of regularisation parameter cannot be convergent for an
ill-posed problem, is \citet{bakushinskii1984} (see \citet[\S3.3]{engl2000}), which is a worst-case
deterministic result and does \emph{not} exclude Lepskii-type balancing in the statistical setting.
\Cref{def:lfstep}(2a)--(2b) says why the schedule is preferred here rather than forced here. The one
input with no supplier in this file is $\hat\varepsilon_0$, which is Q-1, and which a balancing stop
would soften; that is CP-13's third arm.

(6) It does not derive the source index and does not claim to. $\mu$ is
\textbf{primitive}, at (P-a.3), together with the spectral pre-image radius $\mathsf R_\ast$. H1
does \textbf{not} collapse four smoothness parameters to three by
delivering $\mu=s/(2a)$, and the paper should not be read as claiming it does.
H1 delivers that identity for A2's Picard form and for the
\emph{population} operator on $\mathcal V$, and the passage from it to the hypothesis (F-f) actually
% ARCHAEOLOGY.  Rounds 1--11 made the stronger positioning claim.
consumes, a \emph{spectral} condition on a finite-rank compression at a random base point, is
operator interpolation at exponent $2\mu$, licensed only at $2\mu\le1$ and refuted at every $2\mu>1$ by
\cref{lem:srctransfer}(ii.a), while \eqref{eq:muhalf} forces $2\mu>\tfrac53$. \emph{So a referee may ask
what pins $\mu$, and the answer the paper can give is: H1 \textbf{plus} a matching proviso, stated and
not proved}, \cref{rem:pasrc}(4) states the proviso exactly and prices it, and the same two-assumption
structure (a spectral source condition and a link condition, assumed separately rather than one derived
from the other) is what the closest antecedents in the nonparametric-IV regularisation literature carry.
\emph{What is \textbf{not} conceded}: that the rate is weaker, that H1 is retired, or that any other line
of the budget table moves. \eqref{eq:phirate} is unchanged, H1 still fixes the scale and supplies A2's
Picard form, and the amendment is a change to the \emph{status} of one identity, not to any conclusion.

(7) \emph{It is a rate for the iteration run at $\pi$, and \cref{def:lfstep} runs it at $\hat\pi$.}
\eqref{eq:phirate} is stated in $\|\cdot\|_{L^{2}(\pi)}$ and its proof runs the recursion through
$\hat{\mathcal A}^\ast_\varphi$ as \eqref{eq:hatAadj} writes it, at the population $\pi$. The estimator
cannot: $\pi$ is an unknown feature of $P$ by \cref{def:spaces}, so the adjoint the analyst forms is
\eqref{eq:hatAadj} at the floored $\hat\pi$ that \cref{def:lfstep} fixes. \emph{That definition
accordingly separates the two objects by name from 2026-08-27 --- $\hat\varphi^{\pi}$ for the reading
\eqref{eq:phirate} is scoped to, $\hat\varphi$ for the one \cref{def:crossfit}\textup{(b)} plugs in ---
and this clause is the account of the passage between them.}
\emph{The gap is in five parts
and the division is the point of the clause.} \emph{Two are settled outright}: \textup{(7b)} and
\textup{(7e)}. \emph{One is settled \textbf{with a rider it does not discharge}}: \textup{(7a)}, which
identifies the feasible adjoint and the feasible sieve and then flags that the sieve is \emph{random}
while the measurability verification is written at a fixed one. \emph{A part carrying an undischarged
rider is not settled outright, and the distinction is not bookkeeping here}: that rider is a hypothesis
of another definition, \cref{def:crossfit}\textup{(c$'$.1)}, whose consumer is
\cref{lem:epterm}'s cross-fitting argument --- which kills the empirical-process term off
$\sigma(\mathcal I_{-k})$-measurability and off no rate, so the rider is reached by nothing in
\textup{(7c)} or \textup{(7d)}. \emph{Two are open}: \textup{(7c)}, open and \textbf{priced};
\textup{(7d)}, open and unpriced. \emph{So the honest count is two settled, one settled with a rider,
two open}, and the line this clause printed through 2026-08-27 --- ``\textup{(7a)}, \textup{(7b)} and
\textup{(7e)} settle three of them'' --- was short by exactly that rider, at the clause whose subject is
an accounting. \emph{The parts are
lettered \textup{(7a)}--\textup{(7e)} and not \textup{(a)}--\textup{(e)}, deliberately}:
\cref{def:lfstep} has clauses of the latter shape and this clause cites them.
% "FOUR PARTS" -> "FIVE PARTS", 2026-08-27, AND THE ARITHMETIC WAS THE DEFECT.
% The four parts settle, or price, the ADJOINT IDENTITY and the feasible sieve
% (7a), the NORM THE RATE IS READ IN (7b), the LEVEL of the iterate (7c) and the
% OPERATOR THE SOURCE CONDITION IS POSED AT (7d).  NONE OF THE FOUR REACHES
% (F-f)'s RELAXATION CONDITION ||L|| <= 1 -- (7b) is explicitly about a different
% question, the norm the CONCLUSION is read in, and its two-sided equivalence
% licenses a transfer of RATES and not of a normalisation.  So the enumeration
% claimed a division that was short by one, and the missing part was the one that
% decides whether the iteration converges at all.  (7e) supplies it and closes it.
%   AND THE *ACCOUNTING* WAS SHORT AGAIN, IN A DIFFERENT PLACE, REPAIRED THE SAME
% DAY.  The line "(7a), (7b) and (7e) settle three of them" counted (7a) as
% settled while (7a)'s own text flags a consequence it does not discharge -- the
% RANDOM feasible sieve against a measurability verification written at a fixed
% one.  That rider is not a note: it is a hypothesis of def:crossfit(c'.1), whose
% consumer is lem:epterm Step 6's conditional-on-sigma(I_{-k}) bound, and NEITHER
% (7c) NOR (7d) reaches it, both being about rates and operators.  The count is
% now TWO settled outright, ONE settled with a rider, TWO open.  A part carrying
% an undischarged rider that feeds another definition's hypothesis is not
% settled, which is the ground (7c) is called "open and priced" on.
%   THE PART COUNT IS UNCHANGED AT FIVE.  Only the disposition of (7a) moved.

\emph{(7a) The substitution is not ad hoc.} Write $\hat{\mathcal A}^\ast_{\hat\pi}$ for
\eqref{eq:hatAadj} with $\hat\pi$ in place of $\pi$. It is the adjoint of
$\hat{\mathcal A}_\varphi$ in the $L^{2}(\hat\pi)$--$L^{2}(\Pi)$ pairing: the interchange and the
zero-mean step of \cref{prp:hatAprop}\textup{(iii)} read the weight only as a bounded function bounded
away from zero, which $\hat\pi\in[\varepsilon_\pi,\varepsilon_\pi^{-1}]$ is at every sample point, and
$\int_{\mathcal X}\hat\nu_{u,t,l}\,dx=0$ is a statement about the numerator alone. So the feasible
recursion is a genuine Landweber iteration for the same forward operator, read on
$L^{2}(\mathcal X,\hat\pi\,dx)$ rather than on $L^{2}(\mathcal X,\pi\,dx)$. \emph{What moves with the
weight is where the adjoint lands}: $\hat{\mathcal A}^\ast_{\hat\pi}\Xi$ has zero $\hat\pi$-mean and not
zero $\pi$-mean. So the space the \emph{feasible} construction projects onto in
\cref{def:lfstep}\textup{(b)} is
$\operatorname{span}(b_1,\dots,b_{L_x})\cap\{h:\int_{\mathcal X} h\,d\hat\pi=0\}$, with the projection
$L^{2}(\hat\pi)$-orthogonal, and \textbf{not} the $\mathcal S_n$ that definition writes --- which the
analyst cannot form, $\mathcal V$ being defined by $\pi$. \emph{One consequence of that is flagged and
not discharged, and it is the \textbf{rider} that keeps this part from being settled outright}: the
feasible sieve is a \textbf{random} finite-dimensional subspace, while the
measurability verification in the supplement's discussion of the estimator is written at a fixed
$\mathcal S_n$. The Carath\'eodory device that verification already runs for the random functionals
\eqref{eq:smdiff} is the natural instrument for a random basis as well, \emph{and it is not run for one
here}. \emph{The rider is not a note}: $\sigma(\mathcal I_{-k})$-measurability of
$\hat\varphi^{(-k)}$ is what \cref{def:crossfit}\textup{(c$'$.1)} requires and what
\cref{lem:epterm} consumes, so this part feeds another definition's hypothesis, and \emph{it is
reached by neither \textup{(7c)} nor \textup{(7d)}} --- both of those are about rates and operators,
and \cref{lem:epterm} kills its empirical-process term off measurability and off no rate. \emph{The
membership half of the same requirement, $\hat\varphi^{(-k)}\in\mathcal V$, is a separate matter and
\textbf{is} \textup{(7c)}}, the feasible iterate having zero $\hat\pi$-mean where $\mathcal V$ asks for
zero $\pi$-mean.

\emph{(7b) The norm in which the rate is read is not the difficulty.} Under \cref{def:lfstep}'s floor and
C1(ii), $\hat\pi\ge\varepsilon_\pi\ge(\varepsilon_\pi/\|\pi\|_\infty)\,\pi$ and
$\hat\pi\le\varepsilon_\pi^{-1}\le\varepsilon_\pi^{-1}\|1/\pi\|_\infty\,\pi$ pointwise on $\mathcal X$,
so
\[
\frac{\varepsilon_\pi}{\|\pi\|_\infty}\,\|h\|^{2}_{L^{2}(\pi)}
\;\le\;\|h\|^{2}_{L^{2}(\hat\pi)}\;\le\;
\frac{\|1/\pi\|_\infty}{\varepsilon_\pi}\,\|h\|^{2}_{L^{2}(\pi)}
\qquad\text{surely, for every }h\in L^{2}(\pi),
\]
with both constants deterministic and finite. A rate in either norm is therefore the same rate in the
other, and \eqref{eq:phirate} would not have to be restated if its conclusion survived the substitution.

\emph{(7c) Open, and quantified: the level of the iterate is pinned by $\hat\pi$.} By \textup{(7a)} the
feasible sieve is $\operatorname{span}(b_1,\dots,b_{L_x})\cap\{h:\int h\,d\hat\pi=0\}$, so
$\int_{\mathcal X}\hat\varphi\,d\hat\pi=0$ while A1b(i) puts
$\int_{\mathcal X}\vdagger\,d\pi=0$. Writing $h:=\hat\varphi-\vdagger$ and using the first identity to
subtract,
\[
\Big|\int_{\mathcal X}h\,d\pi\Big|
\;=\;\Big|\int_{\mathcal X}\hat\varphi\,(\pi-\hat\pi)\,dx\Big|
\;\le\;\|1/\pi\|^{1/2}_\infty\,\big\|\hat\varphi\big\|_{L^{2}(\pi)}\,
\big\|\hat\pi-\pi\big\|_{L^{2}(\mathcal X,dx)} ,
\]
by Cauchy--Schwarz in $L^{2}(\mathcal X,dx)$ and
$\|\hat\varphi\|_{L^{2}(\mathcal X,dx)}\le\|1/\pi\|^{1/2}_\infty\|\hat\varphi\|_{L^{2}(\pi)}$.
Since $\pi$ is a probability density, $\|h\|^{2}_{L^{2}(\pi)}=\|P_{\mathcal V}h\|^{2}_{L^{2}(\pi)}
+\big(\int h\,d\pi\big)^{2}$, so the error carries a \emph{constant} component that no argument run
inside $\mathcal V$ can see, and \eqref{eq:phirate} survives this mechanism as soon as
\[
\big\|\hat\pi-\pi\big\|_{L^{2}(\mathcal X,dx)}\;=\;O_P(\delta_n) .
\]
\emph{That is a rate condition on a first-stage estimator of exactly the form A4(iii) imposes on
$\hat\varrho$, and this paper imposes none on $\hat\pi$.} The mechanism is therefore priced but not
discharged: what it needs is one assumption, and the assumption is not made.
\emph{It has a \textbf{second} consumer, and the two are not closed by the same thing, which is why they
are separated here rather than run together}: $\int\hat\varphi\,d\hat\pi=0$ with
$\int\hat\varphi\,d\pi\ne0$ is exactly the failure of $\hat\varphi^{(-k)}\in\mathcal V$ that
\cref{def:crossfit}\textup{(c$'$.1)} requires, so this part is a term in a rate \textbf{and} a membership
hypothesis of another definition. \emph{The displayed rate condition closes the first and \textbf{not}
the second}: it makes the constant component $O_P(\delta_n)$, and $O_P(\delta_n)$ is not $0$, so exact
membership in $\mathcal V$ is not restored by any rate. \emph{The repair available for the membership
half is a container weakening --- read \textup{(c$'$.1)}'s third entry at $L^{2}(\pi)$ rather than at
$\mathcal V$, the move dispatch 14 already made at \cref{lem:nusigma}\textup{(i)} --- and it is
\textbf{not} made here, because it has to be checked at every consumer of that clause and this pass
checked none of them.} And neither half touches \textup{(7d)} or \textup{(7a)}'s measurability rider, so
the rate condition alone does not restore \eqref{eq:phirate} at $\hat\varphi$ by any route.

\emph{(7d) Open, and not quantified: the operator the source condition is posed at.} Since
$\hat{\mathcal A}^\ast_{\hat\pi}=(\pi/\hat\pi)\,\hat{\mathcal A}^\ast_\varphi$, the recursion is governed
by $\hat{\mathcal A}^\ast_{\hat\pi}\hat{\mathcal A}_\varphi
=(\pi/\hat\pi)\,\hat{\mathcal A}^\ast_\varphi\hat{\mathcal A}_\varphi$, \emph{the leading factor acting
as multiplication by the function $\pi/\hat\pi$ and written out rather than given a letter, italic $M$
being retired}: self-adjoint in $L^{2}(\hat\pi)$, \emph{not} in $L^{2}(\pi)$, and a bounded, boundedly
invertible, in general non-constant multiplicative perturbation of the operator (P-a.3) is stated at.
\emph{Norm
equivalence does not repair this, and that is the trap the clause exists to name}: equivalent renorming
preserves rates, and it does not preserve range conditions, since
$\mathcal R\big((\pi/\hat\pi)\hat{\mathcal A}^\ast_\varphi\big)
=(\pi/\hat\pi)\,\mathcal R\big(\hat{\mathcal A}^\ast_\varphi\big)$ need not be the same subspace unless
$\pi/\hat\pi$ is constant, and the ranges of the fractional powers
$(\mathcal A^\ast\mathcal A)^{\mu}$ that (P-a.3) reads are not in general preserved either. So
(P-a.3)'s spectral pre-image $\varsigma^{(k)}$ and its radius $\mathsf R_\ast$, and with them (F-f) and
Steps~2--3 of the proof, are stated for an operator the feasible iteration does not run.

\emph{For \textup{(7d)} two routes are available, neither is taken, and that is why \textup{(7d)} is an
open item and not a lemma.} \textup{(i)} Pose the source condition directly at the feasible operator.
The cost is that the hypothesis is then a condition on an estimated object rather than on $P$, which is
the standard \cref{rem:obligations} holds every other nuisance to. \textup{(ii)} Treat multiplication by
$\pi/\hat\pi$ as a perturbation of the identity and price it. This needs a \emph{uniform} rate for
$\|\hat\pi/\pi-1\|_\infty$ --- no condition in this paper supplies one, $\hat\pi$ being introduced by
\cref{def:lfstep} and constrained there by nothing beyond the floor --- together with a stability result
for source conditions under multiplicative perturbation, which this paper does not carry. \emph{Route
\textup{(ii)} is the honest one and the work is not done here.} \emph{The two open parts have one
prerequisite in common}: on a bounded $\mathcal X$ a uniform rate for $\|\hat\pi/\pi-1\|_\infty$ implies
the $L^{2}(\mathcal X,dx)$ rate \textup{(7c)} asks for, since
$\|\hat\pi-\pi\|_{L^{2}(\mathcal X,dx)}\le|\mathcal X|^{1/2}\|\pi\|_\infty\|\hat\pi/\pi-1\|_\infty$, so
a single assumption on $\hat\pi$ would discharge \textup{(7c)} outright and open \textup{(7d)} to
route~\textup{(ii)}.

\emph{A third route would dissolve the item and is declined for a stated reason.} The results proved
here read $\pi$ through C1(ii)'s two-sided bounds, through $\bar c_\pi$, and through the normalisation
of $\mathcal V$; no step consumes the identity $\pi=\tfrac12(f_{X_1}+f_{X_2})$ itself. \emph{That last
sentence is a reading and not a census: no systematic audit of the corpus for further consumers has been
made, and one may exist.} On that reading a \emph{known} $\pi_0$ bounded above and below on $\mathcal X$
would satisfy C1(ii), would make \eqref{eq:hatAadj}, \eqref{eq:normconst} and \eqref{eq:cap} computable
outright, and would close \textup{(7c)}, \textup{(7d)} and \textup{(7e)} below at once, the two pairings
of \textup{(7e)} collapsing to one at $\hat\pi=\pi_0=\pi$. It is declined because $\pi$ is the design density the
estimator's $\hat\pi$ is built to estimate, and because the exchange is not free: the singular system of
$\mathcal A_{\mathrm{cond}}$, and hence what A2, H1 and (P-a.3) assert, move with the weight, so the
smoothness a source condition posed at $\pi_0$ asserts is smoothness relative to $\pi_0$ and not
relative to the design.

\emph{(7e) Settled, and it was not among the four parts this clause listed until 2026-08-27: the
iteration must be \textbf{relaxation-admissible} in the pairing it runs in.} \textup{(F-f)}'s standing
hypothesis $\|\mathcal L\|\le1$ is a bound on the operator norm \emph{out of the domain Hilbert space
whose inner product defines $\mathcal L^{\ast}$}, \cref{rem:ff}. By \textup{(7a)} that space is
$L^{2}(\hat\pi)$ for the feasible recursion and $L^{2}(\pi)$ for the one \cref{prp:phirate} analyses, so
the condition has to be met in whichever is in force. \emph{The other four parts do not reach it, and
that is why it is a fifth and not a corollary of one of them}: \textup{(7a)} settles the adjoint
identity and identifies the feasible sieve, \textup{(7b)} settles the norm the \emph{conclusion} is read
in --- explicitly a different question --- \textup{(7c)} prices the level of the iterate, \textup{(7d)}
the operator the source condition is posed at. \emph{Nor does \textup{(7b)}'s equivalence do the work,
and the trap is the same species it names}: two-sided equivalence with deterministic constants
transfers \emph{rates}, and a normalisation is not a rate --- it is an inequality against the constant
$1$, which no equivalence preserves.

\emph{What the omission cost, stated at the witness, and stated \textbf{as a gap} because that is what
it is.} Let $\hat\pi\equiv\varepsilon_\pi$, a constant, which the construction convention of
\cref{def:lfstep} admits --- it is the function that definition's own non-emptiness argument exhibits,
and nothing beyond the two-sided floor constrains $\hat\pi$ there. Then
$\|h\|^{2}_{L^{2}(\hat\pi)}=\varepsilon_\pi\|h\|^{2}_{L^{2}(\mathcal X,dx)}$ exactly, and
\eqref{eq:pairgap} gives
$\|\hat{\mathcal A}_\varphi\|_{L^{2}(\hat\pi)\to L^{2}(\Pi)}
\ge(\varepsilon_\pi\|1/\pi\|_\infty)^{-1/2}\|\hat{\mathcal A}_\varphi\|_{L^{2}(\pi)\to L^{2}(\Pi)}$,
whose right-hand side diverges as $\varepsilon_\pi\downarrow0$ at a fixed $\omega_n$. \emph{Against a
scale carrying $\|1/\pi\|_\infty^{1/2}$ alone --- which is what \eqref{eq:normconst} printed through
2026-08-26 --- $\|\mathcal L\|\le1$ in $L^{2}(\hat\pi)$ is then \textbf{not implied} at any
$\varepsilon_\pi<\|1/\pi\|_\infty^{-1}$}, that is, throughout the range \cref{def:lfstep} itself
recommends on the separate ground that the clipping there costs nothing: the only bound this corpus
carries for that operator norm is \eqref{eq:hatAhatpi}, and its constant exceeds the old $\omega_n$ by
the factor $(\varepsilon_\pi\|1/\pi\|_\infty)^{-1/2}>1$.
\emph{Whether the inequality was actually \textbf{violated} at some sample is not settled here, and the
reason it cannot be is itself the point}: that would need a matching \emph{lower} bound on
$\|\hat{\mathcal A}_\varphi\vert_{\mathcal S_n}\|$, which is \textbf{Q-4}'s object and which
\cref{lem:scaleinv}'s closing note records as unavailable. \emph{So the defect is a \textbf{proof gap}
and not an exhibited divergence} --- \textup{(L1)} was checked against a bound in the wrong pairing, and
the corpus could neither verify it in the right one nor produce a sample at which it failed. \emph{What
makes it worth a part of its own rather than a footnote is the failure mode if the inequality does
break}: once $\|\mathcal L\|^{2}>2$ the map $\mathrm{Id}-\mathcal L^{\ast}\mathcal L$ carries an eigenvalue below
$-1$, the domain being finite-dimensional, and \cref{def:lfstep}\textup{(b)} diverges geometrically
rather than converging slowly --- a mode against which a rate theorem offers no protection at all.
\emph{The repair below removes the question rather than answering it}: once \eqref{eq:normconst} covers
both pairings, how tight \cref{prp:hatAprop}\textup{(i)} is stops bearing on \textup{(L1)}, and Q-4
returns to being about the \emph{rate} and not about convergence.

\emph{The repair, and its scope.} \eqref{eq:normconst} now carries
$\max\{\|1/\pi\|_\infty^{1/2},\varepsilon_\pi^{-1/2}\}$, the maximum of the two conversion costs
\cref{prp:hatAprop}\textup{(i$'$)} computes from the single primitive estimate \eqref{eq:hatAdx}, so
$\|\omega_n^{-1}\hat{\mathcal A}_\varphi\|\le1$ holds \emph{surely in both pairings} and \textup{(L1)}
closes whichever the reader runs. \emph{The repair is a constant and not a rate}: $\varepsilon_\pi$ is
fixed before the data and $\|1/\pi\|_\infty$ is finite by C1(ii), so the maximum is a fixed positive
number free of $n$, $\omega_n\asymp h_u^{-1}h_x^{-1/2}$ is unchanged, and every asymptotic statement
reading $\omega_n$ --- the power of $n$ in \eqref{eq:cap}, clause \textup{(v)} of \cref{lem:scaleinv} and
its instantiation, \cref{prp:phirate} Step~2 \textup{(L4)}'s stopping-index bound, and the exact
cancellation \eqref{eq:scaleinv}, which is stated for \emph{every} $\omega_\ast>0$ --- is untouched.
\emph{What moves is one constant}: the right-hand side of \eqref{eq:cap} is multiplied by at most
$(\varepsilon_\pi\|1/\pi\|_\infty)^{-1}$, a computational budget and not a rate condition.
\emph{And the number was already in the corpus, filed one shelf over}: the discussion of the estimator
in the supplement offered $\varepsilon_\pi^{-1/2}$ as the analyst's \textbf{computable} substitute for an
unknown $\|1/\pi\|_\infty^{1/2}$. It is the same number and it was the \textbf{stability} constant all
along --- and, on the range \cref{def:lfstep} recommends, the maximum \emph{is} $\varepsilon_\pi^{-1/2}$,
so the computable choice and the admissible one coincide rather than compete. \emph{What \textup{(7e)}
does \textbf{not} close}: \textup{(7c)} and \textup{(7d)} are untouched by it and remain open, and
\textup{(7a)}'s random-sieve flag is likewise unaffected.

\emph{One structural observation, which \textup{(7e)} settles nothing about and which is recorded
because it is what kept the omission from being a false theorem.} \textup{(P-a.3)} is posed at the
\emph{realised} iterates, $\bar v^{(k)}_\dagger=P_{\mathcal S_n}(\vdagger-\hat\varphi^{(k-1)})$ for every
$k\le\bar k$, and $\bar k\ge2$ always, \cref{prp:phirate} Step~4(b). So on any configuration at which the
stage-1 iteration fails to contract, the hypothesis at $k=2$ demands a bounded spectral pre-image for a
divergent $\hat\varphi^{(1)}$ and \emph{empties}: the proposition is vacuous there, not false.
\emph{That is a thin protection and it should be read as one.} A hypothesis stated at the realised path
is checkable only after the estimator has been run, so it silently absorbs defects in the estimator's own
tuning --- here, a choice of $\varepsilon_\pi$ --- and converts what would otherwise be a visible
falsehood into an invisible vacuity. \emph{The device is legitimate}, it is what Step~0's conditional
route buys and what \cref{rem:pasrc} defends, \emph{and its cost is that the hypothesis set can be
emptied by a tuning constant rather than by a feature of $P$}. \eqref{eq:normconst} as amended removes the
one such mechanism identified in this corpus; \emph{no claim is made that it is the only one}, and a
systematic audit of which hypotheses can be emptied by an admissible choice of a tuning constant has not
been run.
\end{remark}

% ---------------------------------------------------------------------------
%  EXTRACTED FROM THE STATEMENT OF prp:phirate on 2026-08-02.  That statement ran
%  181 live lines -- the longest environment in the corpus -- of which better than
%  half was commentary sitting inside the hypothesis list.  NOTHING WAS DELETED:
%  every removed sentence is below, verbatim except for the connectives needed to
%  make it stand outside the enumeration it used to sit in.  NO LABELLED DISPLAY
%  MOVED: eq:compress, eq:lxwindow, eq:muhalf and eq:phirate all stayed in the
%  statement, so no \eqref anywhere was orphaned and no equation number changed.
% ---------------------------------------------------------------------------

% TITLE SHORTENED 2026-08-05.  The old one ran 133 characters and this remark's
% head is set in a reduced measure, so the tail "load-bearing one)" fell off the
% right edge and was ABSENT FROM THE PDF -- lost text, not a loose rule.  The
% dropped clause "where eq:muhalf comes from" is not lost: clause (4) below
% carries it.
\begin{remark}[What the hypotheses of \cref{prp:phirate} carry, and why \textup{(P-a.3)} is load-bearing]\label{rem:phirhyp}
\begin{enumerate}
\item[(1)] Three tags are narrowed from the blocks ``B1'', ``C1'' and ``C5$'$'' that stood in the
hypothesis list, and each narrowing removes a clause this proposition cannot have. B1(vi) is
\textbf{retired} and is not available as a hypothesis of any result, so a block citation of B1 charged
this proposition with it; C1(i) is consumed at exactly one site in the corpus, Step~2 of
\cref{lem:aerr}, the route this proposition's Step~0 \textbf{declines}, while C1(ii) is consumed
transitively through \cref{prp:hatAprop} and \cref{lem:nusigma}; and C5$'$(ii) is a \emph{lower} bound on
$h_x$ while this proposition's own (P-g) is an \emph{upper} bound on $h_x$, so a block citation of C5$'$
put a floor and a ceiling in one hypothesis set, \emph{and, at the inherited strengthening of
C5$'$(ii), a floor above the ceiling}. \emph{The two do not meet at every admissible parameter, and the
failure is not a corner}: at $(\mu,s,a,r,d_Z)=(2,1.2,0.3,\tfrac83,3)$, which clears \eqref{eq:frontier}
at $\kappa\alpha=0.256$ and clears \eqref{eq:lxwindow} at $s>1$, the strengthened floor is
$h_x\gg n^{-0.44}$ and (P-g) is $h_x=O(n^{-0.5143})$, \emph{empty by the exponent gap $0.0743$}. The
threshold is \eqref{eq:hxwindow}, which reads $s>1.3182$ at $\alpha=0.32$ and is failed by $s=1.2$, and
the emptiness band is non-empty at every $\alpha\in(\tfrac14,\approx0.359)$, which \eqref{eq:frontier}
does not exclude.
% RECOMPUTED WITH (P-g)'s h_u.  (P-g) is now R h_x^{s-1/2} = O(h_u hat-delta_n) =
% O(n^{-9 alpha/8}), so at alpha = 0.32 and s - 1/2 = 0.7 the ceiling is
% h_x = O(n^{-9(0.32)/(8 x 0.7)}) = O(n^{-2.88/5.6}) = O(n^{-0.51429}), against
% O(n^{-0.45714}) before; the gap against the n^{-0.44} floor widens from 0.01714 to
% 0.07429.  eq:hxwindow at alpha = 0.32 moves from 1.22727 to 1.31818, and s = 1.2
% fails it either way, so THE WITNESS IS UNCHANGED IN KIND AND STRONGER IN DEGREE.
% The band's upper root moves from ~0.3494 to ~0.35914, printed rounded down.
\emph{\textbf{What follows from that is vacuity, not falsity}}: on the block-cited reading
this proposition's hypotheses were unsatisfiable there and it held on an empty hypothesis set, and no
conclusion of this corpus is false at that configuration. The narrowing removes the conflict from the
hypothesis set; C5$'$(ii)'s own restatement removes it from the maintained block.
\item[(2)] The clause $p_K\ge s-\tfrac12$ is \textbf{not} free at every $s$ the budget table
admits. \Cref{lem:unifop} Step~4 consumes it, and $p_K=8$ being \emph{fixed} by C3 and by C5$'$(i) it is
a hard cap $s\le\tfrac{17}2$: $p_K$ is not a parameter the analyst may raise against $s$, so the clause
caps $s$, while \eqref{eq:lxwindow} is a \emph{floor} on $s$ that diverges as $\alpha\uparrow\tfrac12$.
\emph{The two are incompatible at every $\alpha\ge\tfrac{68}{145}\approx0.4690$, at every $s$ and every
$a$}: the floor $9\alpha/\big(8(1-2\alpha)\big)$ reaches $\tfrac{17}2$ there. \emph{Witness, and it is
admissible}: $r=8$, $d_Z=1$, so $\alpha=\tfrac8{17}$; with $\mu=2$ the frontier \eqref{eq:frontier}
clears at $\kappa\alpha=\tfrac{8}{17}\cdot\tfrac45\approx0.3765>\tfrac14$; \eqref{eq:lxwindow} demands
$s>9$ and the cap allows $s\le\tfrac{17}2$; \emph{no $s$ exists}.
% DISPOSITION CHANGED 2026-08-05.  This clause used to end "no further
% restriction on alpha is imposed by prp:phirate", and it read the band as one
% where the proposition is VACUOUS.  Both became false when the restriction was
% printed into prp:phirate's own hypothesis list: the band is now EXCLUDED, and
% an excluded region is not a vacuous one -- the proposition simply does not
% speak there, and (T-5) is inapplicable rather than unsatisfiable.  A4(ii)'s
% d_Z >= 1 repair is still the precedent for the shape of the fix, but the two
% dispositions are no longer the same and saying so was the defect.
\emph{The disposition is now a printed restriction, and this remark is where it is derived}:
\cref{prp:phirate} carries $\alpha<\tfrac{68}{145}$ in its own hypothesis list, and
\cref{thm:rootn}'s group (T-5) inherits it, so on $[\tfrac{68}{145},\tfrac12)$ the proposition is
\textbf{excluded rather than vacuous}, it does not speak there, and (T-5) is inapplicable rather
than unsatisfiable. \emph{In primitives the restriction is $9r<68\,d_Z$}, both of which A4(ii)
declares, so a reader can check it without leaving the hypothesis list.
A4(ii)'s $d_Z\ge1$ repair delivers only $\alpha<\tfrac12$ and does not reach this
range; the cap and the restriction it creates are both recorded as rows of the smoothness-budget
table. \emph{And A8 is cited by clause and not
as a block}: A8's own lower end $L_x\gg n^{1/(2s)}$ is the balancing route's, A8's own comment already
records it as unused on the rate route, and item~(8) below prices what carrying it would have cost.
\item[(3)] \textup{(P-a.1)}: the measurability half is a requirement on the preliminary and is
stated \textbf{there} because some hypothesis list must carry it. \Cref{def:lfstep}'s measurability
verification inducts on $k$ \emph{from} it, so the verification that the estimator is a random element
consumes it. \emph{The two remaining halves are jointly satisfiable only under \eqref{eq:sieveinhs}, and
that is not a technicality}: $\hat\varphi^{(0)}\in\mathcal S_n$ forces
$\|\vdagger-\hat\varphi^{(0)}\|_{H^{s}}$ to be read at a sieve element, so unless
$\mathcal S_n\subset H^{s}$ the norm is $+\infty$ surely and the two halves of that one clause are
\emph{jointly unsatisfiable}, the probability statement would then hold on the empty event. A8's basis
clause supplies \eqref{eq:sieveinhs}, with it the well-posedness of (P-a.1), and H1 supplies
$\vdagger\in H^{s}$.
\item[(4)] \emph{\textup{(P-a.3)} is \textbf{not} the pair, a Hilbert-scale source condition on
$\bar v^{(k)}_\dagger$ together with the link condition of \cref{lem:srctransfer}(iii) on
$\mathcal S_n$, and does not follow from it}: the pair delivers (P-a.3) only through operator
interpolation at exponent $2\mu$, which \cref{lem:srctransfer}(ii.a) refutes at \emph{every} $2\mu>1$
while \eqref{eq:muhalf} forces $2\mu>\tfrac53$, so the pair never implies what (F-f) consumes.
\Cref{rem:pasrc} carries the full disposition. \emph{And the operator is named as the compression, which
is what \eqref{eq:compress} is for}: the operator (F-f) is applied to in Step~2 of the proof is
$\mathcal L^{\ast}\mathcal L$ with
$\mathcal L=\omega_n^{-1}\hat{\mathcal A}_{\hat\varphi^{(k-1)}}\vert_{\mathcal S_n}$, and that \emph{is}
$\omega_n^{-2}\hat{\mathcal P}^{\,\mathcal S}_{\hat\varphi^{(k-1)}}$ and not the restriction of
$\hat{\mathcal A}^{\ast}\hat{\mathcal A}$ to $\mathcal S_n$, two operators whose ranges are nested in
neither direction.
\item[(5)] What the empty letter \textup{(P-b)} reserved. An operator-estimation budget
$\gamma_n=O(\hat\delta_n^{1/(2\mu+1)})$ belonging to the perturbation-of-$\mathcal A$ route, which
Step~0 of the proof does not take; nothing is imposed at that letter and the list runs \textup{(P-a)},
\textup{(P-c)}--\textup{(P-g)}. The label is kept empty rather than renumbered so that an existing
citation of (P-b) lands on the reservation.
\item[(6)] \emph{\textup{(P-a.2)}'s requirement on $\varsigma_0$ is strictly weaker than the conclusion
\eqref{eq:phirate} precisely when $\mu>\tfrac56$, which is (P-f), \textbf{the same inequality} rather
than a second one that agrees with it}, as Step 4(a.1) of the proof shows. \emph{And what (P-f) reads is
bandwidth-dependent}: \cref{rem:h6bandwidth} unpacks it, shows that at the unamended bandwidth
$h_u\asymp n^{-1/16}$ it reads
$\tfrac{2\mu-1}{2\mu+1}\cdot\tfrac{r}{2r+d_Z}>\tfrac18$, and shows that at C5$'$(i) as amended, 
$h_u\asymp\hat\delta_n^{1/p_K}$, the largest bandwidth whose bias is $O(\hat\delta_n)$, it reads
$\mu>\tfrac56$, \emph{strictly, with the strictness forced by the logarithms}.
\item[(7)] \textup{(P-c)}: the factor $h_u$ is not decoration. The source and link conditions are
imposed on $\mathcal S_n$, so the proof works with the sieve projection
$P_{\mathcal S_n}(\vdagger-\hat\varphi^{(k-1)})$ and must carry the residual through
$\hat{\mathcal A}$, whose norm on the complement is bounded only by $O(h_u^{-1})$. \emph{(P-c) is a
\textbf{lower} end for $L_x$, and it is this proposition's own; it \textbf{replaces} A8's lower end
rather than restating it}, being read at $h_u\hat\delta_n$ and not at $n^{-1/2}$.
\item[(8)] What \eqref{eq:lxwindow} is numerically, and what carrying A8's lower end as well would
have cost. \eqref{eq:lxwindow} is $s>\tfrac{27}{16}=1.6875$ at $\alpha=\tfrac38$ and $s>1$ at
$\alpha=0.32$. This is a genuine joint line in the smoothness budget and it is recorded there.
\emph{Had A8's lower end been carried as well}, the window would instead require
$s>\big(2(1-2\alpha)\big)^{-1}$, which is $s>2$ at $\alpha=\tfrac38$; that is strictly the stronger line
at every $\alpha<\tfrac49$, and it is why A8 is cited by clause. A configuration on which \emph{both}
routes are to be available at once must clear the stronger one.
\item[(9)] What the conclusion discharges, and what it consumes, \emph{and the list is shorter from
2026-08-27 than it was, because the conclusion is scoped}. A5(i) is \textbf{discharged, not
assumed}, through \cref{lem:opnoise}. \emph{A5(ii) splits}: its second half
$\delta_n=o_P(n^{-1/4})$ is discharged, being a statement about the deterministic sequence $\delta_n$
alone; its \emph{first} half, the rate itself, is \textbf{not}, because that half is stated at
$\hat\varphi$ --- the \emph{feasible} reading of \cref{def:lfstep}\textup{(b)} --- while
\eqref{eq:phirate} is stated at the infeasible $\hat\varphi^{\pi}$, and the passage between them is
clause \textup{(7)} below, open at \textup{(7c)} and \textup{(7d)}. A5(iii) is not discharged and no
proposition can discharge it, being an inequality among primitives. A5(vi$'$) is \textbf{consumed}, it is a hypothesis of
\cref{prp:phirate}, at \cref{lem:pathcont}(c), and A5(vi) follows from it at $k=\bar k$; A5(iv) and
A5(v) concern other objects and are untouched.
\emph{The line this clause printed through dispatch 27 --- ``A5(i)--(iii) are discharged, not assumed''
--- was wrong about A5(iii) before the scoping and about A5(ii)'s first half after it}, and both halves
of that correction are recorded here rather than only at the site that consumes them,
\cref{thm:rootn} Step~0.
\item[(10)] Where the threshold \eqref{eq:muhalf} comes from, and it is two conditions and not
one. The stagewise recursion of Step~4 of the proof is \eqref{eq:fixedpoint},
$\hat\varepsilon_k=\mathsf c_{15}\big(\mathsf c_{14}(\hat\delta_n\ln n+h_u^{-2}\hat\varepsilon_{k-1}^{2})\big)^{2\mu/(2\mu+1)}$,
a recursion in the \emph{deterministic schedule} of \eqref{eq:lfsched} and not in the unobservable error.
\emph{Contraction} needs $2\kappa>1$ with $\kappa=2\mu/(2\mu+1)$, i.e.\ $\mu>\tfrac12$; \emph{the fixed
point being \eqref{eq:phirate}} needs (P-f), which at the amended bandwidth is $\mu>\tfrac56$. The second
implies the first, so \eqref{eq:muhalf} is stated at $\tfrac56$ and the $\mu>\tfrac12$ line of the budget
table is now slack. \emph{The coincidence worth naming is at $\tfrac12$, not at $\tfrac56$}:
$\mu>\tfrac12$ is also exactly what A5(iii) proves necessary for \eqref{eq:frontier} to be satisfiable at
any $r$, and it is exactly the index above which the inherited nonlinear route's rate theorem is
unavailable (\cref{rem:lfnonlinear}(3)). The two arise independently, one from the propensity
frontier, one from the linearisation recursion, and they coincide. \emph{What \eqref{eq:muhalf}
charges beyond that is $(\tfrac12,\tfrac56]$, and that is a real line in the budget}, listed with the
others in the assumption block's table.
% ARCHAEOLOGY.  The symbol delta_n^{(k)} that a round-9 draft used for the
% unobservable error is withdrawn; glossary C58.
\item[(11)] And what happens at $\tfrac56$ is \textbf{not} a second coincidence. Step 4(a.1) of
the proof exhibits a requirement on the preliminary, \eqref{eq:prelimrate}, and shows it is weaker than
the rate delivered exactly when $\mu>\tfrac56$. That is (P-f) again, \emph{the same inequality},
written in the same primitives, up to the strictly positive factor $\kappa/(2\kappa-1)$. \emph{The
statement worth printing, and it is not a restatement of (P-f) either}: (P-f) is exactly the condition
under which the newly exposed preliminary-rate requirement is \emph{not self-defeating}, so that
\eqref{eq:prelimrate} adds a genuine new demand on $\hat\varphi^{(0)}$, open item Q-1, without
making the procedure ask its input to be as good as its output.
% ARCHAEOLOGY.  A round-9 draft called the agreement at 5/6 a coincidence and said
% the two conditions were not restatements of one another; that sentence is struck.
\item[(12)] And the load-bearing hypothesis is \textup{(P-a.3)}, which should be read in the same
breath as \eqref{eq:muhalf}. (P-a.3) is a \emph{random spectral source condition} on a
\emph{finite-rank, sieve-compressed} operator at a \emph{data-dependent} base point, assumed to hold with
probability tending to one, with no sufficient condition supplied anywhere in this file; it is what makes
the conditional-on-sample route legitimate, and it is what carries the source index above $\tfrac12$.
\emph{Any statement that ``the linearised route reaches the same place under the same threshold'' is true
only with (P-a.3) priced in the same sentence}, see \cref{rem:lfnonlinear}, \cref{rem:pasrc} and open
item Q-4.
\end{enumerate}
\end{remark}

% The optional argument is BRACED because it contains a `]' inside math:
% \begin{lemma}[... $E[\Delta Z\Delta Z']$] closes the optional argument at the
% FIRST unbraced `]', which is the one in E[...], leaving the `$' unmatched.
% Four LaTeX errors, no PDF.  An environment-balance count cannot see this --
% nothing is unbalanced; the argument is simply cut short.
\begin{lemma}[{The moment is valid, its variance is finite, its $\theta$-Jacobian is $-J$, and $J=E[\Delta Z\Delta Z']$}]\label{lem:momzero}
The moment has mean zero at the truth, a finite second moment, and derivative $-J$ in $\theta$; the
finiteness is delivered by the \emph{level} clause of A$2'$ and not by its derivative clause; $J$ is
the second-moment matrix $E[\Delta Z\Delta Z']$, so no rank condition beyond A1b(iv) is needed anywhere;
and, with one further clause of C$2'$, the variance is strictly positive definite, so \cref{thm:rootn}'s
studentised conclusion needs no positivity proviso.
Suppose A1, A1b(iv), A$2'$(i), \textbf{both} clauses (i.a) and (i.b), A$2'$(iii), B1(i), (ii),
(iv.a), (iv.d) hold. \emph{Clause \textup{(i.b)} is in the list and not an ``if in
addition'', for two reasons.} First, every bound below runs through \cref{cor:elldiff}, which
% ARCHAEOLOGY.  A round-10 draft carried (i.b) only as an "if in addition".
is stated at an admissible representer \emph{in the sense of A2b}, and A2b requires
$\Lambda^{(0)},\Lambda^{(1)}<\infty$ and continuous $u$-differentiability, that is (i.a) \textbf{and} (i.b).
Second, A$2'$(i)'s closing clause fixes $\ell_\star$ \emph{pointwise} ``from that version'', the version
satisfying both, so \emph{$\ell_\star$ is not even pointwise well defined from \textup{(i.a)} alone}, and a
sup-norm statement about it is not well posed. \cref{thm:rootn}'s group (T-2) lists both clauses, so no
result downstream is affected. Then
\begin{enumerate}
\item[(i)] $E\|\psi_\star\|<\infty$ and $E\big[\psi_\star\big]=0$;
\item[(ii)] $E\|\psi_\star\|^{2}<\infty$, so $\Sigma$ in \eqref{eq:JSig} is a well-defined finite positive
semidefinite matrix, with
\[
E\|\psi_\star\|^{2}\;\le\;2\,E\big\|\Delta Z\,\Delta U\big\|^{2}
+4\,c_{\widetilde W}^{2}\Big\{E\big\|b_\star(\Delta Z)\big\|^{2}+\big\|\ell_\star\big\|_\infty^{2}\,E\big[\Delta U^{2}\big]\Big\} ,
\]
\emph{the factor $c_{\widetilde W}^{2}$ of \cref{cor:rangegap} entering through
$\|L_\star\widetilde W\|\le c_{\widetilde W}\|L_\star\|$ and equal to $1$ at $\mathsf L=2$};
\item[(iii)] $\theta\mapsto E[\psi(\theta,\zeta_\star)]$ is twice differentiable on $\IR^{p}$ with
$\partial_\theta E[\psi(\theta,\zeta_\star)]\big|_{\theta=\theta_\dagger}=-J$, $\|J\|<\infty$, and second
derivative bounded in norm by $c_{\widetilde W}\|\partial_u\ell_\star\|_\infty E\|\Delta Z\|^{2}$ uniformly in $\theta$
, the first derivative using $\Lambda^{(0)}_\star$ and the second using both budgets;
\item[(iv)] the debiasing summand of $J$ vanishes, and the $\theta$-map is exactly affine. Under
A1's \emph{first} clause,
\[
E\big[\ell_\star(\Delta U,\Delta Z)\,\widetilde W\,\Delta Z'\big]\;=\;0,
\qquad\text{so}\qquad
J\;=\;E\big[\Delta Z\,\Delta Z'\big] ;
\]
consequently $J$ is symmetric positive semidefinite, and $J$ is nonsingular by A1b(iv) alone. Moreover
$E\big[L_\star\big(\Delta U(\theta,\vdagger),\Delta Z\big)\,\widetilde W\big]=0$ at every
$\theta\in\IR^{p}$, so
\begin{equation}\label{eq:psiaffine}
E\big[\psi(\theta,\zeta_\star)\big]\;=\;-\,E\big[\Delta Z\,\Delta Z'\big]\,\big(\theta-\theta_\dagger\big)
\qquad\text{for every }\theta\in\IR^{p} ,
\end{equation}
which is affine with nonsingular gradient: the \emph{population} moment has $\theta_\dagger$ as its unique
root on all of $\IR^{p}$, and its second derivative is the zero matrix;
\item[(v)] the variance is strictly positive definite, and no hypothesis beyond \textup{(ii)}'s list
plus C2$'$(i) is needed for it. If in addition C2$'$(i) holds, then
\begin{equation}\label{eq:sigmapd}
\Sigma\;\succ\;0,
\qquad\text{hence}\qquad
J^{-1}\Sigma J^{-1\prime}\;\succ\;0
\quad\text{and}\quad
\vartheta'J^{-1}\Sigma J^{-1\prime}\vartheta>0\ \text{ for every }\vartheta\ne0 .
\end{equation}
\end{enumerate}
\emph{Clause \textup{(v)} converts a \textbf{proviso} of \cref{thm:rootn} into
a \textbf{consequence} of its hypotheses.} That theorem's studentised conclusion is stated for directions
$\vartheta$ with $\vartheta'J^{-1}\Sigma J^{-1\prime}\vartheta>0$, which reads as a condition the user must
check; \eqref{eq:sigmapd} says it holds at every $\vartheta\ne0$ under hypotheses the theorem already
carries in (T-1) and (T-3), namely A1b(iv) and C2$'$(i). The two ingredients are exactly those and
nothing more. By \eqref{eq:sigsplit} below, $\vartheta'\Sigma\vartheta\ge E[(\vartheta'\Delta Z)^{2}\Delta U^{2}]$,
both summands of $\Sigma$ being positive semidefinite; A1b(iv) makes $E[\Delta Z\Delta Z']$ nonsingular, so
$P(\vartheta'\Delta Z\ne0)>0$ for every $\vartheta\ne0$; and C2$'$(i) gives $\Delta Y$ a conditional density
given $\mathcal F_{XZW}$, hence, since $\Delta U=\Delta Y-\Delta\vdagger-\Delta Z'\theta_\dagger$ differs
from $\Delta Y$ by an $\mathcal F_{XZW}$-measurable shift, gives $\Delta U$ one too, so
$P(\Delta U=0\mid\mathcal F_{XZW})=0$ almost surely and $P(\Delta U=0)=0$. On the positive-probability event
$\{\vartheta'\Delta Z\ne0\}$ the integrand is therefore almost surely strictly positive, and
$E[(\vartheta'\Delta Z)^{2}\Delta U^{2}]>0$; finiteness is A1b(iv)'s $E\|\Delta Z\Delta U\|^{2}<\infty$.
\emph{The second summand of \eqref{eq:sigsplit} is not used}, so no overlap condition enters: A4(i) is what
makes that summand bounded \emph{below}, and clause \textup{(v)} does not need it.
Two things clause \textup{(v)} does not say. It says nothing about the \emph{size} of $\Sigma$, which
is where A4(i) does bind (as $\varrho$ approaches a boundary the second summand diverges); and it is a
statement at a fixed $P$, not a bound uniform over a class of laws, $\lambda_{\min}(\Sigma)$ may be
arbitrarily small at nearby $P$, which is the same pointwise-versus-uniform boundary
\cref{rem:rootnscope}\textup{(6)} draws for the theorem itself.

Three readings of \textup{(iv)} must be blocked, and each of the three is a place a reader will
over-generalise it.

\textup{(iv.1)} Clause \textup{(iii)}'s second-derivative envelope is a correct bound on the zero
matrix, hence carries no information. It is marked here rather than struck, on this file's standing rule
that a proved clause is retired in place. \emph{$\|\partial_u\ell_\star\|_\infty$, and hence A$2'$(i.b)'s
budget $\Lambda^{(1)}_\star$, is \textbf{not} thereby freed}, for two independent reasons: it is consumed
by \cref{prp:sharpquad}'s inheritance in \cref{lem:biasexp}\textup{(2)}, and A$2'$(i.b) is in any case
required to \emph{license} \cref{cor:elldiff} at all, which is what makes $\ell_\star$ pointwise well
defined and $\|\ell_\star\|_\infty$ finite, and hence what makes \eqref{eq:JSig} a matrix. Only the
attribution at \cref{rem:rootnhyp}\textup{(1)} changes.

\textup{(iv.2)} It is a statement about $J$ and not about $\hat J$. A4(i) puts $\varrho$ inside
the $\varepsilon_\varrho$-interior simplex \eqref{eq:trunset}, so \cref{def:crossfit}(a)'s truncation is free at the
population propensity and the population $\widetilde W$ is conditionally centred; but
$\widehat{\widetilde W}^{(-k)}_l=\un\{W=l\}-\hat\varrho^{(-k)}_l(\Delta Z)$ is \textbf{not} conditionally centred, at
any $n$, and the second summand of \eqref{eq:Jhat} does not vanish. \cref{lem:sandwich} Step~3 is where its
cost $\|\hat\ell^{(-k)}\|_\infty\|\mathsf r_k\|_{L^{2}(P)}$ is paid, and that step is unchanged.
\emph{Replacing $\hat J$ by $\IE_n[\Delta Z\Delta Z']$ would also be consistent for $J$; it is not
done}, because it would break the exact identity $\partial_\theta\bar\psi_n(\theta)=-\hat J(\theta)$ of
\cref{def:onestep} on which \cref{thm:rootn} Step~1's algebra rests.

\textup{(iv.3)} It does not propagate to the $\varphi$-direction, and the block would collapse if it
did. $\Delta h$ is $\sigma(X_1,X_2)$-measurable and \textbf{not}
$\sigma(\Delta U,\Delta Z)$-measurable, so $E[\ell_\star\widetilde W\Delta h]$ is not zero, indeed A2 is
precisely the assertion that it equals $-E[\Delta Z\Delta h]$ for every $h\in\mathcal V$
(\cref{prp:orthog}\textup{(i)} and \eqref{eq:orthphi}), and were it zero for every $h$ one would have
$\nu_{\Delta Z}\equiv0$, which fails at any design in which the instrument moves the regressor.
\emph{The boundary between the two sides is exactly the one \textup{(P4b)}'s warning and \textup{(P4c)}
were split to police}, and $q=E[\partial_u\ell_\star\,\widetilde W\mid X_1,X_2]$ sits on the non-vanishing
side of it.

\textup{(iv.4)} It does not touch $\Sigma$, and the one thing it \emph{does} say about $\Sigma$ is a
decomposition and not a vanishing. The same tower step applied to the cross term of
$\Sigma=E[\psi_\star\psi_\star']$, legitimate because $\Delta Z\Delta U\,L_\star'$ is
$\sigma(\Delta U,\Delta Z)$-measurable and integrable by clause (ii) and the Cauchy--Schwarz inequality
, gives $E[\Delta Z\Delta U\,\widetilde W'L_\star']=0$ and hence, exactly,
\begin{equation}\label{eq:sigsplit}
\Sigma\;=\;E\big[\Delta Z\Delta Z'\,\Delta U^{2}\big]\;+\;E\big[L_\star\,\widetilde W\widetilde W'\,L_\star'\big] .
\end{equation}
\emph{Both summands are positive semidefinite, and the second does not vanish, for which one further
tag is needed and is \textbf{not} on this lemma's list. \textbf{And the route to it changes at
$\mathsf L\ge3$, where the binary one is unavailable rather than merely weaker.}} At $\mathsf L=2$,
$\widetilde W^{2}\ge\varepsilon_\varrho^{2}$ surely under A4(i), so
$E[L_\star\widetilde W\widetilde W'L_\star']\succeq\varepsilon_\varrho^{2}E[L_\star L_\star']$, which is
nonzero whenever $L_\star$ is. \emph{At $\mathsf L\ge3$ that argument fails, and not by a constant}:
$\widetilde W\widetilde W'$ is a rank-one matrix, so no sure inequality
$\widetilde W\widetilde W'\succeq cI$ holds at any $c>0$, and $L_\star\widetilde W\widetilde W'L_\star'$
is singular at every sample point once $p>1$. \emph{The correct route is the conditional one, and
\eqref{eq:omegafloor} was written for it}: under A1, $\Delta U\upmodels W\mid\Delta Z$ gives
$E[\widetilde W\widetilde W'\mid\Delta U,\Delta Z]=\Omega(\Delta Z)$ of \eqref{eq:omegadef}, and A4(i)
gives $\Omega(\Delta Z)\succeq\varepsilon_\Omega\,\mathrm{Id}$ with $\varepsilon_\Omega=\varepsilon_\varrho/\mathsf L$,
so that
$E[L_\star\widetilde W\widetilde W'L_\star']=E[L_\star\Omega(\Delta Z)L_\star']\succeq\varepsilon_\Omega
E[L_\star L_\star']$, nonzero whenever $L_\star$ is. The conditional route uses A1, which is on this
lemma's list, and it is valid at every $\mathsf L$ including $2$, the binary display is retained
because its constant $\varepsilon_\varrho^{2}$ is the one earlier rounds printed and, at
$\varepsilon_\varrho<\tfrac12$, is the weaker of the two. \emph{A4(i) is cited
for this sentence and for nothing else here}; \eqref{eq:sigsplit} itself, and every clause
\textup{(i)}--\textup{(iv)}, hold without it. \emph{The reading to
take from \eqref{eq:sigsplit} is the standard debiasing trade and not a defect}: with
$J=E[\Delta Z\Delta Z']$, the sandwich $J^{-1}\Sigma J^{-1\prime}$ is the variance the infeasible plug-in
\eqref{eq:prelimtheta} would have at a \emph{known} $\vdagger$, plus a positive semidefinite term; what the
debiasing weight buys is not variance but the removal of the first-order bias
$\langle\nu/\pi,h_k\rangle$ that \cref{lem:prelimtheta} Step~4$'$ shows the plug-in carries at exact order
$\|h_k\|$. \emph{Neyman orthogonality is untouched throughout}, being a statement about the three
\emph{nuisance} directions of \cref{prp:orthog} and not about $\partial_\theta$.
\emph{One simplification is available and is \textbf{not} taken into the displays}: since $J$ is
symmetric, $J^{-1\prime}=J^{-1}$, so \cref{thm:rootn}'s $J^{-1}\Sigma J^{-1\prime}$ and
\eqref{eq:Sighat}'s $\hat V$ could drop a transpose. They do not, because $\hat J$ is \emph{not} symmetric
by \textup{(iv.2)}, and a variance estimator that dropped the transpose would not be symmetric either.
\emph{The transpose is carried in both places so that the population display and the sample display have
the same shape}, which is the convention \cref{lem:sandwich}\textup{(ii)} states the limit in.

A$2'$\textup{(ii)} is not among the hypotheses and is consumed nowhere in Part~I. The variance object
is $E\|L_\star\widetilde W\|^{2}$, which clause (ii) bounds through A$2'$(iii) and
$\|\ell_\star\|_\infty$; A$2'$(ii) bounds $E\|\ell_\star\widetilde W\|^{2}$, a different quantity, and
that quantity is in any case \emph{implied by A$2'$\textup{(i)} in full together with B1(i), (ii), (iv.a), (iv.d)}:
\cref{cor:elldiff} gives $\|\ell_\star\|_\infty\le\|\varpi_u\|_\infty\Lambda^{(0)}_\star<\infty$ and
$\|\widetilde W\|\le c_{\widetilde W}$ by \cref{cor:rangegap}, so
$E\|\ell_\star\widetilde W\|^{2}\le c_{\widetilde W}^{2}\|\ell_\star\|_\infty^{2}<\infty$ with no
further hypothesis. \emph{The bound on $\|\widetilde W\|$ is the Euclidean one and not \textup{(P1)}},
which is the per-cell statement.
\emph{The attribution matters}: the redundancy is \textbf{not} ``A$2'$(ii)
is implied by A$2'$(i.a)'', which is not a statement this file can make, since the only route to
$\|\ell_\star\|_\infty<\infty$ available here is \cref{cor:elldiff} and that corollary is stated at A2b.
The redundancy is real; its ground is A$2'$(i) in full plus the four B1 clauses.
% CORRECTED 2026-08-24.  This read "the three B1 clauses" while the sentence NINE
% LINES ABOVE names four -- B1(i), (ii), (iv.a), (iv.d).  The count was three
% before (iv.d) was added to cor:elldiff's line earlier the same day and did not
% move with it: a count inside the very paragraph that exists to police
% attributions.
A$2'$(ii) is \textbf{retired in place} at A2$'$, on the
A3 precedent, a proved clause is retired, its number kept so a stale citation lands
somewhere. \emph{The forward
Step-8 need is the only ground for retention that survives, and it is a forward obligation rather than a
present consumer}; it is recorded at the retirement note.
% ARCHAEOLOGY.  A round-10 draft recorded the redundancy as "A2'(ii) is implied by
% A2'(i.a)".  This is the referent slip A4's comment (d) recorded, seen from the
% other side; the retirement was acted on in dispatch 14, and A4's comment (d),
% whose retention ground was "the clause the second-order and bias terms consume",
% was rewritten, that ground being false in this file's own record.
\end{lemma}

% ===========================================================================
%  I.3  THE FOUR LEMMAS
% ===========================================================================

\begin{lemma}[L-14 with a rate: the plug-in preliminary]\label{lem:prelimtheta}
The preliminary inherits the first-stage rate, and it inherits it because $\theta_\dagger$ is identified by
restriction (B) alone.
Suppose A0, A1, A1b(iv), C1(ii) hold, adopt \cref{def:crossfit}, and suppose A5(ii)'s \textbf{first half,
read at each fold}, the first display of \eqref{eq:foldrates}. \emph{The fold index is cited for the
reason \cref{def:crossfit}(d) gives}: A5(ii) prices one $h$, and \eqref{eq:prelimtheta} averages over
$\mathsf F$ of them. \emph{A5\textup{(ii)}'s second half $\delta_n=o_P(n^{-1/4})$ is \textbf{not} in this
list}; it is the hypothesis of clause \textup{(ii)} alone. Then
\begin{enumerate}
\item[(i)] \emph{the rate, on a list with no threshold in $n$;}
\begin{equation}\label{eq:prelimthetarate}
\big\|\hat\theta^{(0)}-\theta_\dagger\big\|\;=\;O_P\big(n^{-1/2}\big)+O_P\big(\delta_n\big) ,
\end{equation}
and if in addition $\delta_n\to0$ then $\hat\theta^{(0)}\xrightarrow{p}\theta_\dagger$;
\item[(ii)] the threshold, which is the only clause A5\textup{(ii)}'s second half enters. If in
addition $\delta_n=o_P(n^{-1/4})$, then
\begin{equation}\label{eq:prelimthetathresh}
\big\|\hat\theta^{(0)}-\theta_\dagger\big\|\;=\;o_P\big(n^{-1/4}\big) .
\end{equation}
\emph{Both summands of \textup{(i)} are needed for this conclusion and neither is redundant to it}: the
first is $o(n^{-1/4})$ trivially and the second is A5\textup{(ii)}'s second half.
\end{enumerate}
Why the split, and why it is \textbf{not} the \textup{(i)}/\textup{(i$'$)} device of
\cref{lem:sandwich}. That device states \emph{one} conclusion on two lists; here the two clauses have
\emph{different} conclusions, so they are numbered and not primed. The split exists for
\cref{cor:ratefloor}, whose hypotheses deliberately omit A5(ii)'s second half, doing without
\eqref{eq:frontier}, hence without that threshold, is what the corollary is \emph{for}, and which
consumes this lemma twice: once for the fourth hypothesis of \cref{lem:sandwich}\textup{(i$'$)},
$\|\hat\theta^{(0)}-\theta_\dagger\|=o_P(1)$, and once for
$\|\Psi^{\theta}_n\|=O_P(\|\hat\theta^{(0)}-\theta_\dagger\|^{2})$. \emph{Clause \textup{(i)} supplies both
from its statement.}
% ARCHAEOLOGY.  A round-11 draft stated only the conjunction of eq:prelimthetarate
% and eq:prelimthetathresh under a list containing A5(ii) in full, so cor:ratefloor
% had to be read against this lemma's PROOF -- the same defect
% lem:sandwich(i') was written to end, one lemma short of the end.
\emph{If in addition A5(i) and
A5(iii) hold}, then $\delta_n\gg n^{-1/2}$, because $\hat\delta_n=n^{-\alpha}$ with
$\alpha=r/(2r+d_Z)<\tfrac12$ and $\kappa=2\mu/(2\mu+1)<1$ give $\kappa\alpha<\tfrac12$, so the two
summands of \eqref{eq:prelimthetarate} collapse to $O_P(\delta_n)$. \emph{That simplification is cosmetic
and is used by neither clause}:
\cref{thm:rootn} Step~2 works at the uncollapsed rate \eqref{eq:prelimthetarate}, since
$\sqrt n(n^{-1/2}+\delta_n)^{2}\le2(n^{-1/2}+\sqrt n\,\delta_n^{2})$ and both summands vanish under A5(ii);
which is why A5(i) and A5(iii) are not among the hypotheses above.
% ARCHAEOLOGY.  A round-10 draft of this sentence said the collapse was "not on the
% path to thm:rootn" while Step 2 opened by using the collapsed form; Step 2 is
% restated at the uncollapsed rate and the two agree.
\emph{No condition on $\hat\ell$, on $\hat L$, on $\hat\varrho$, on $L_x$ or on A2 is used}: the preliminary
does not touch the debiased moment.
\end{lemma}

\begin{lemma}[L-6 and L-7 at the estimated nuisances: the exact bias decomposition]\label{lem:biasexp}
The conditional bias of the debiased moment at the estimated nuisances is the balancing pairing plus exactly
three remainders, one quadratic in the structural error, two products of two estimation errors, and no
Taylor expansion in the nuisances is taken.
Suppose A1, A1b(i), A1b(iv), A2, A$2'$(i), A$2'$(iii), N1, B1(i), (ii), (iv), C1(ii) and C2$'$(i) hold;
adopt \cref{def:crossfit}; and suppose hypothesis (R-a) of \cref{thm:rootn}, so that $\hat\ell^{(-k)}$ is
bounded and continuously differentiable in $u$ with
$\max_{k\le\mathsf F}\|\partial_u\hat\ell^{(-k)}\|_\infty=O_P(1)$, \emph{an $O_P(1)$ envelope and
\textbf{not} a bound at every sample point}: \eqref{eq:biasb2}
carries the envelope on its \emph{right} and is read, as the displays below all are, on the event where
it is finite. Put $g_k:=\Delta Z+\hat\ell^{(-k)}(\Delta U,\Delta Z)\,\widetilde W\in\IR^{p}$, the contrast
standing to the right of the $p\times(\mathsf L-1)$ estimate as in \eqref{eq:psidef}. Then, for every $k$, almost
% ARCHAEOLOGY.  This gloss read "||partial_u hat-ell^{(-k)}||_inf < infinity almost
% surely" through dispatch 20, which is neither what (R-a) states nor what
% eq:biasb2 needs.
surely
\begin{equation}\label{eq:biasdecomp}
E\big[\hat\psi^{(-k)}\,\big|\,\sigma(\mathcal I_{-k})\big]
\;=\;-\big\langle\nu_{g_k}/\pi,\;h_k\big\rangle_{L^{2}(\pi)}
\;-\;\Psi^{(1)}_k\;+\;\Psi^{(2)}_k\;-\;\Psi^{(3)}_k ,
\end{equation}
where, all expectations being conditional on $\sigma(\mathcal I_{-k})$,
\[
\Psi^{(1)}_k:=E\Big[\Delta h_k\!\int_0^1\!\big\{\hat\ell^{(-k)}(\Delta U-\lambda\Delta h_k,\Delta Z)-\hat\ell^{(-k)}(\Delta U,\Delta Z)\big\}d\lambda\;\widetilde W\Big],
\]
\[
\Psi^{(2)}_k:=E\Big[\Delta h_k\!\int_0^1\!\hat\ell^{(-k)}(\Delta U-\lambda\Delta h_k,\Delta Z)\,d\lambda\;\mathsf r_k(\Delta Z)\Big],
\qquad
\Psi^{(3)}_k:=E\big[D_k(\Delta U,\Delta Z)\,\mathsf r_k(\Delta Z)\big] ,
\]
\emph{each an $\IR^{p}$-valued object in which the $p\times(\mathsf L-1)$ factor stands to the left of
the $(\mathsf L-1)$-vector it is contracted against and the scalar $\Delta h_k$ commutes}, and
\begin{align}
\big\|\big\langle\nu_{g_k}/\pi,h_k\big\rangle\big\|
&\;\le\;2\,c_{\widetilde W}\,\bar c_\pi^{1/2}\,\big\|\hat\ell^{(-k)}-\ell_\star\big\|_{L^{2}(P)}\,\big\|h_k\big\|_{L^{2}(\pi)},
\label{eq:biasb1}\\
\big\|\Psi^{(1)}_k\big\|&\;\le\;2\,c_{\widetilde W}\,\bar c_\pi\,\big\|\partial_u\hat\ell^{(-k)}\big\|_\infty\,\big\|h_k\big\|^{2}_{L^{2}(\pi)},
\label{eq:biasb2}\\
\big\|\Psi^{(2)}_k\big\|&\;\le\;2\,\bar c_\pi^{1/2}\,\big\|\hat\ell^{(-k)}\big\|_\infty\,\big\|\mathsf r_k\big\|_{L^{2}(P)}\,\big\|h_k\big\|_{L^{2}(\pi)},
\label{eq:biasb3}\\
\big\|\Psi^{(3)}_k\big\|&\;\le\;\big\|\mathsf r_k\big\|_{L^{2}(P)}\,\big\|D_k\big\|_{L^{2}(P)},
\qquad\text{and}\qquad \Psi^{(3)}_k=0\ \text{ if $D_k$ satisfies N1}',
\label{eq:biasb4}
\end{align}
where $\|\mathsf r_k\|_{L^{2}(P)}:=\big(E\|\mathsf r_k(\Delta Z)\|^{2}\big)^{1/2}$, \emph{the second
moment of the norm of the $\IR^{\mathsf L-1}$-valued propensity error, which is what A4(iii) prices, and
not the scalar second moment}, and $\|D_k\|_{L^{2}(P)}$, $\|\hat\ell^{(-k)}\|_\infty$,
$\|\partial_u\hat\ell^{(-k)}\|_\infty$ are Frobenius.
\emph{The four left-hand sides are \textbf{Euclidean} norms on $\IR^{p}$, and \eqref{eq:biasb1}'s is a
norm for the same reason the other three are}: the objects bounded are the four summands of
\eqref{eq:biasdecomp}, which is an identity between $\IR^{p}$ vectors, $\psi$ being $\IR^{p}$-valued by
\eqref{eq:psidef}. \emph{That is exactly what the triangle inequality in $\IR^{p}$ needs, and it is the
form in which \cref{thm:rootn} Step~4 consumes these displays}: applied to \eqref{eq:biasdecomp} it gives
$\|E[\hat\psi^{(-k)}\mid\sigma(\mathcal I_{-k})]\|$ at most the sum of the four right-hand sides.
% ARCHAEOLOGY.  eq:biasb1 carried a SCALAR MODULUS until the present pass, which
% supplies that step only at p = 1, while eq:biasdecomp above types the very same
% object as an R^p summand.  Statement and proof moved together: proofs.tex,
% lem:biasexp Step 3 and thm:rootn Step 4(a).  Same repair at lem:balance,
% eq:balance3.
\emph{\eqref{eq:biasb1} and \eqref{eq:biasb2} carry
$c_{\widetilde W}=\min\{\sqrt2,\sqrt{\mathsf L-1}\}$ of \cref{cor:rangegap} and
\eqref{eq:biasb3}--\eqref{eq:biasb4} do not, for the reason \cref{lem:secondord} gives}: the first two
contract a matrix against $\widetilde W$, the last two against $\mathsf r_k$ or $D_k$, whose norms are
already the quantities being priced. At $\mathsf L=2$, $c_{\widetilde W}=1$ and all four
displays carry the constants earlier rounds printed.

The three remainders are written $\Psi^{(j)}_k$ and not $\mathcal R^{(j)}_k$, and the choice is
forced. $\mathcal R(\cdot)$ is the \emph{range} of an operator throughout this file, A2, \cref{lem:solset},
\cref{lem:srctransfer}, and calligraphic $\mathcal C$, $\mathcal E$, $\mathcal G$, $\mathcal H$,
$\mathcal U$ are all claimed as well. $\Psi$ is the letter this project already binds to $E[\psi]$ along a
nuisance path (\cref{lem:secondord}), and these three objects \emph{are} components of $E[\psi]$ at the
estimated nuisances, so the letter is semantically right; $\Psi(\tau)$ always carries an argument and
$\Psi^{(j)}_k$ never does, which is the same separation device that keeps
$T(\cdot,\cdot)$ apart from $\mathcal T$. \emph{\Cref{prp:sharpquad}'s precedent, ``the remainder carries
no symbol of its own'', is followed \textbf{where it can be, which is at two objects and not at three}}:
the Hessian envelope of \cref{thm:rootn} Step~1 and the event $\{\hat J\ \text{nonsingular}\}$ are written
out, while the $\theta$-remainder \textbf{does} carry a symbol, $\Psi^{\theta}_n$, because it is referred to
at six separate sites across two proofs and writing it out at each would be unreadable.
% ARCHAEOLOGY.  A round-10 draft of this sentence claimed the theta-remainder was
% written out; it is glossary slot C77 and the glossary row is the accurate one.

Three structural facts are visible in \eqref{eq:biasdecomp} and each is a claim of the paper.
\textup{(1)} \emph{No second-order term in $L$ and none in $\varrho$ appears, exactly}, because $\psi$ is
bi-affine in $(L,\varrho)$, this is \cref{lem:secondord}\textup{(ii)} and it is used here as an
\emph{identity} rather than as a derivative statement, which is why \eqref{eq:biasdecomp} carries no Taylor
remainder in those two directions. \textup{(2)} \emph{The only square is $\Psi^{(1)}_k$}, and by
\cref{prp:sharpquad}\textup{(i)},\textup{(iii)} its order in $\|h_k\|_{L^{2}(\pi)}$ cannot be improved
uniformly over an $L^{2}(\pi)$-ball, what is inherited from Part~G is the \emph{lower} bound on the
leading term, and \cref{rem:sharpscope}\textup{(1)} is what forbids reading it as a statement about the
realised $h_k$. \textup{(3)} \emph{$\Psi^{(3)}_k$ is the entire cost of N1$'$ failing at the estimated
direction}, which is obligation \textup{(O2)} of \cref{rem:obligations}, and \eqref{eq:biasb4} prices it.
\end{lemma}

\begin{lemma}[L-8: cross-fitting removes the empirical-process term]\label{lem:epterm}
Once the estimated moment is $L^{2}(P)$-consistent at a frozen nuisance, the centred empirical average of
the estimated moment minus the true one is $o_P(n^{-1/2})$, no entropy condition, no Donsker property,
and no stochastic equicontinuity at the estimated point. \emph{What is \textbf{not} claimed is that the
consistency is free}: \eqref{eq:psiL2} is established below from the rate and envelope conditions listed,
and the headline is a statement about what \eqref{eq:psiL2} \emph{buys} at scale $n^{-1/2}$, not about what
\eqref{eq:psiL2} costs.
Suppose A0, C1(ii), A4(ii)--(iii) hold, adopt \cref{def:crossfit}, suppose hypotheses (R-a) and (R-b) of
\cref{thm:rootn}, and suppose A5(ii) and A5(vi) with $s>\tfrac12$ \textbf{read at each fold}, the first
two displays of \eqref{eq:foldrates}.
\emph{\textbf{A4\textup{(ii)} is on this line and it is not decoration}}: Step~5 concludes
$O_P(r_n^{1/2})=o_P(1)$, which needs $r_n\to0$, and A4(iii) states only the \emph{form}
$r_n=n^{-r/(2r+d_Z)}$ and nothing about $r$. It is A4(ii)'s $r>0$ together with $1\le d_Z\le p-1$ that
makes the exponent strictly positive; without it $r_n$ need not vanish, Step~5 delivers $O_P(1)$,
\eqref{eq:psiL2} fails and \eqref{eq:eptermbound} with it.
% ADDED 2026-08-10 (Phase 6a).  This list read "A4(iii)" alone through dispatch 25
% while Step 5 printed "O_P(r_n^{1/2}) = o_P(1) by A4(iii)", so the lemma's stated
% conclusion did not follow from its stated hypotheses.  THE ALTERNATIVE REPAIR IS
% NOT TAKEN AND THE CHOICE IS DELIBERATE: cor:ratefloor lists "r_n -> 0" separately
% and that precedent was available here, but the primitive that delivers it sits on
% the same tag one clause away, and this block everywhere else refuses to let a
% clause citation carry its neighbour silently -- so the clause is named rather than
% its consequence assumed.  Only one of the two repairs is made; making both would
% put a hypothesis and its own corollary on one line, which is what the A3 and
% A2'(ii) retirements removed.  lem:sandwich inherits this list wholesale and needs
% no separate edit; its clause (i') is narrowed in the same pass for a different
% reason and is not this repair's mirror.
\emph{Of A5(ii) this lemma uses the \textbf{first half}, the rate itself, together with
$\delta_n\to0$, and nothing else}: Steps~2 and~3 of the proof read $O_P(\delta_n)$ and conclude $o_P(1)$,
and A5\textup{(ii)}'s second half $\delta_n=o_P(n^{-1/4})$ is consumed at no step. \emph{This is the split
\cref{lem:prelimtheta} states on its own hypothesis line}, and it is stated here for the same reason: an
asymmetry between two adjacent lemmas over one clause reads as a gap in the seam rather than as a
difference in what each consumes. \emph{The clause is nonetheless left on the list whole rather than
narrowed to its first half}, because $\delta_n\to0$ is what the second half supplies a fortiori and no
other hypothesis on this list delivers it.
% ASYMMETRY CLOSED 2026-08-10 (Phase C).  lem:prelimtheta is explicit that it takes
% "A5(ii)'s FIRST HALF, read at each fold" and that the second half "is NOT in this
% list"; this lemma said "A5(ii)" plain, over an identical split.  Re-checked at
% the artifact before writing: Step 2 reads (E[(Delta h)^4])^{1/4} = O_P(delta_n^{1/2}),
% Step 3 reads O_P(delta_n) by (R-a) and A5(ii), Steps 4-6 read neither -- so the
% n^{-1/4} threshold is consumed NOWHERE.  THE LIST IS NOT NARROWED: delta_n -> 0
% is needed by Steps 2 and 3 and is supplied here only by the second half, so
% striking it would have been a mathematical change, not a wording one.  No edge
% of the assumption-usage matrix moves.
\emph{The fold index is cited rather than left implicit because A5(ii)
and A5(vi) are stated at a \textbf{single} $h=\hat\varphi-\vdagger$ while this lemma consumes them at
$\mathsf F$ distinct estimators}, one per fold complement; \cref{def:crossfit}(d) is where the reading is
declared and \cref{thm:rootn} Step~0 is where it is discharged.
A1b(iv) is \textbf{not} a hypothesis, on a ground of its own and
not on \cref{lem:opnoise}'s. No step below cites it, and the two moments any step of this proof or of the
clause of \cref{def:crossfit}\textup{(c$'$.1)} it adopts could want from it, $E\|\Delta Z\|^{2}$ at
Step~2 and $E[\Delta U^{2}]$ in that clause's derived $L^{2}$ membership of $\hat L^{(-k)}$, are
implied a fortiori by (R-b), which this list already carries and which asks for the \emph{fourth} moments
of both. What this lemma
does \textbf{not} use is A1b(iv)'s nonsingularity clause, which is \cref{lem:sandwich}'s. Then
% ARCHAEOLOGY.  A1b(iv) stood on this line through dispatch 20; the deletion is a
% strict weakening and the two displays are unchanged.
\begin{equation}\label{eq:psiL2}
\max_{k\le\mathsf F}\big\|\hat\psi^{(-k)}-\psi_\star\big\|_{L^{2}(P)}\;=\;o_P(1),
\end{equation}
where the norm is taken at frozen $\hat\zeta^{(-k)}$, and consequently
\begin{equation}\label{eq:eptermbound}
\sqrt n\,\sum_{k=1}^{\mathsf F}\frac{n_k}{n}\,
\Big(\IE_{n,k}-E\Big)\big[\hat\psi^{(-k)}-\psi_\star\big]\;=\;o_P(1).
\end{equation}
What is and is not conceded, stated because leg~3 of the Chen--Pouzo answer is a positioning
commitment. \eqref{eq:eptermbound} is proved by conditional Chebyshev at a \emph{frozen} nuisance, so the
only input it requires is the $L^{2}(P)$ consistency \eqref{eq:psiL2}. \emph{This is a different
statement from \cref{lem:unifop}, and the difference is the point}: \cref{lem:unifop} concedes a uniform law
of large numbers with a rate over a slowly growing ball, at scale $\hat\delta_n\gg n^{-1/2}$, and that
concession is charged in the $\varphi$-step; the $\theta$-step concedes nothing at scale $n^{-1/2}$, which is
where a penalised sieve minimum-distance treatment of $\theta$ would need stochastic equicontinuity of the
$\sqrt n$-scaled empirical process at the estimated point.
\end{lemma}

\begin{lemma}[L-11: the sandwich, and no additional correction for the propensity]\label{lem:sandwich}
Both factors of the sandwich are consistent, and the sandwich is the \emph{whole} asymptotic variance: no
Neyman correction for the propensity appears in it.
Suppose the hypotheses of \cref{lem:epterm} \textbf{and of \cref{lem:momzero}}, together with A2,
(R-e) of \cref{thm:rootn}, and $\|\hat\theta^{(0)}-\theta_\dagger\|=o_P(1)$. \emph{\cref{lem:momzero}'s list
is inherited in full and not abbreviated}: Steps~3, 6, 6$'$ and~7 of the proof cite
\cref{lem:momzero}\textup{(ii)}--\textup{(iv)}, whose hypotheses run A1, A1b(iv), A$2'$(i), A$2'$(iii),
B1(i), (ii), (iv.a), (iv.d), the last four being the bridge \cref{cor:elldiff} needs to make
$\|\ell_\star\|_\infty$ finite, hence to make $J$ of \eqref{eq:JSig} a finite matrix. \emph{A3 is
\textbf{not} on this list}: Step~6 needs $J$
nonsingular, \cref{lem:momzero}\textup{(iv)} delivers $J=E[\Delta Z\Delta Z']$ and A1b(iv) delivers its
nonsingularity, and both are already on the inherited list. A3 is retired at A3. Then
% ARCHAEOLOGY.  A round-10 draft listed A2'(i) and A2'(iii) and omitted A1 and the
% three B1 clauses, so it was short by exactly the cor:elldiff bridge.  A3 was on
% this list through dispatch 12; the deletion is a strict weakening with no change
% of conclusion.
\begin{enumerate}
\item[(i)] $\hat J(\theta_\dagger)\xrightarrow{p}J$ and $\hat J=\hat J(\hat\theta^{(0)})\xrightarrow{p}J$;
consequently $P(\hat J\ \text{nonsingular})\to1$ and $\hat J^{-1}\xrightarrow{p}J^{-1}$;
\item[(i$'$)] clause \textup{(i)} uses no rate in $n$, and holds on a strictly weaker list. The
conclusions of \textup{(i)} hold under A0, A1, A1b(iv), C1(ii), A2, A$2'$(i), A$2'$(iii), B1(i), (ii),
(iv.a), (iv.d), A4(iii), \cref{def:crossfit}, (R-a), $\|\hat\theta^{(0)}-\theta_\dagger\|=o_P(1)$ and the three
consistency statements
\[
\max_{k\le\mathsf F}\big\|h_k\big\|_{L^{2}(\pi)}=o_P(1),
\qquad
\max_{k\le\mathsf F}\big\|\hat\ell^{(-k)}-\ell_\star\big\|_{L^{2}(P)}=o_P(1),
\qquad
\max_{k\le\mathsf F}\big\|\mathsf r_k\big\|_{L^{2}(P)}=o_P(1),
\]
the last of which A4(ii)--(iii) would supply on their own, \emph{and neither is on the list above,
because all three statements are assumed \textbf{outright} here rather than derived}. \emph{In particular
A4(i), A4(ii), A5(ii), A5(vi), (R-b) and
(R-e) are not used, and no product of two rates and no threshold in $n$ appears}; \eqref{eq:psiL2} is not
used either, clause \textup{(i)} being about $\hat J$ alone.
% NARROWED FROM A4(i)--(iii) TO A4(iii) 2026-08-10 (Phase 6a), AND THE NARROWING IS
% WHAT MAKES THIS CLAUSE'S OWN CLAIM TRUE.  The clause advertises a hypothesis set
% "strictly weaker" than the lemma's main list; that main list carries A4(ii)--(iii)
% and not A4(i), so a list reading A4(i)--(iii) was NOT a subset of it and the
% advertised comparison was false.  Sweeping Steps 1-6': the only A4 citation is
% A4(iii) at Step 3.  A4(ii) is not needed because the three consistency statements
% are hypotheses here, not conclusions -- which is exactly what distinguishes this
% clause from lem:epterm, where A4(ii) IS needed and was added in the same pass.
% With A4(iii) alone the list is now a subset of the main one and drops A5(ii),
% A5(vi), (R-b), (R-e) and A4(ii), so "strictly weaker" is true as printed.
This is the clause \cref{cor:ratefloor} invokes,
and it is stated because that corollary's hypotheses contain neither A5(ii) nor a rate product, so an appeal
to \textup{(i)} as stated would be an appeal to the \emph{proof} rather than to the statement;
\item[(ii)] if in addition $\hat\theta\xrightarrow{p}\theta_\dagger$ then
$\hat\Sigma\xrightarrow{p}\Sigma$ and $\hat V\xrightarrow{p}J^{-1}\Sigma J^{-1\prime}$.
\end{enumerate}
A flag, not a side taken. The influence function of \cref{thm:rootn} is $J^{-1}\psi_\star$, in which
$\varrho$ appears only through $\widetilde W$ at its true value, and no term correcting for the estimation of
$\varrho$ is present. That is a \emph{consequence} of \cref{prp:orthog}\textup{(iii)} under N1: the
$\varrho$-derivative of the moment vanishes \emph{exactly}, so there is nothing for a correction to remove.
\emph{The project's positioning section states the opposite}, that the sandwich carries an additional
correction for $\varrho$; the two statements cannot both be right, and the contradiction is recorded for the
writer at \S9 item~10 of the strategy memo rather than resolved here.
\end{lemma}

\begin{proposition}[L-10: the median over partitions]\label{prp:medianpart}
Averaging or taking medians over a fixed number of independent fold partitions changes neither the influence
function nor the limit.
Let $\mathsf P\ge1$ be fixed and not depending on $n$, let
$\hat\theta_1,\dots,\hat\theta_{\mathsf P}$ be the estimators \eqref{eq:onestep} computed on $\mathsf P$
partitions drawn independently of the data and indexed by $\mathsf p\in\{1,\dots,\mathsf P\}$
, \emph{sans-serif, the lowercase of the count $\mathsf P$, and \textbf{not} $q$, which is bound to
$q=E[\partial_u\ell_\star\,\widetilde W\mid X_1,X_2]$}, 
and suppose the hypotheses of \cref{thm:rootn} hold
\textbf{at each partition}. \emph{The per-partition reading is part of the hypothesis and not a formality}:
the $\mathsf P$ partitions have $\mathsf P\mathsf F$ distinct fold complements, and (T-5), (R-a), (R-c),
(R-e) and \eqref{eq:foldrates} are conditions on the nuisances computed on a fold complement, so what is
required is that they hold at all $\mathsf P\mathsf F$ of them, with the maxima in \eqref{eq:foldrates} and
(R-a) read over that index set. \emph{Since $\mathsf P$ and $\mathsf F$ are both fixed this is finitely many
statements and every order is preserved}; it is not a uniformity requirement, and it is the \emph{first}
of the three uses of $\mathsf P$ fixed that this proposition's argument makes, enumerated below.
% RE-COUNTED 2026-08-10 (Phase C), AND THE MINIMAL DELETION WOULD NOT HAVE CLOSED
% IT.  This read "a THIRD place P fixed is doing work, BEYOND THE TWO THE PROOF
% NAMES".  Both halves fail against the artifact: the proof's Step 3 is headed
% "where P fixed is used, THREE times" and names three, of which THE FIRST IS THIS
% ONE -- the hypothesis, in the same words ("the maxima in eq:foldrates and (R-a)
% may be read over PF fold complements at unchanged order").  So it is not a third
% place beyond two; it is the first of three.  Deleting only the "beyond the two"
% clause would have left "third", which is the false half.  This paragraph's own
% later enumeration lists the same three and needed nothing.
Let
$\hat\theta^{\mathrm{md}}$ be their coordinatewise
median and $\hat\theta^{\mathrm{av}}:=\mathsf P^{-1}\sum_{\mathsf p}\hat\theta_{\mathsf p}$. Then
\[
\sqrt n\big(\hat\theta_{\mathsf p}-\hat\theta_{\mathsf p'}\big)=o_P(1)
\quad\text{for every }\mathsf p,\mathsf p',
\qquad
\sqrt n\big(\hat\theta^{\mathrm{md}}-\theta_\dagger\big)
=\sqrt n\big(\hat\theta^{\mathrm{av}}-\theta_\dagger\big)+o_P(1)
\;\xrightarrow{\;d\;}\;\text{the limit law of \eqref{eq:rootn}} .
\]
\emph{The limit is named by pointing at \eqref{eq:rootn} rather than by printing a symbol}, which is
\cref{thm:rootn}'s disposition carried to its third statement site and is also the more informative
form: what this proposition asserts \emph{is} that the median and the average have
\cref{thm:rootn}'s limit.
% ARCHAEOLOGY.  Glossary C90.
\emph{The argument is that all $\mathsf P$ estimators share one influence function, and it uses
$\mathsf P$ fixed three times}: once in the hypothesis, so that ``the hypotheses of \cref{thm:rootn} at each
partition'' is finitely many statements; once so that finitely many $o_P(1)$ terms may be collected; and once
so that the median is a continuous function of finitely many arguments. \emph{It says nothing about
$\mathsf P\to\infty$}, where the median would have to be handled as a functional of a random measure and the
hypothesis itself would become a uniformity requirement over a growing family of fold complements.
\end{proposition}

% ===========================================================================
%  I.5  THREE REMARKS
% ===========================================================================

% ############################################################################
% ##  rem:rootnlog.  SHORTENING PASS 2026-08-01.  Clause (1), "the defect", was
% ##  17 compiled lines of process record and is archived here.  It read:
% ##  prp:phirate delivers ||hat-varphi - varphi_dagger|| = O_P(delta_n) with
% ##  delta_n = (delta-hat_n ln n)^{2mu/(2mu+1)}, while A5(ii) as written through
% ##  dispatch 9 set delta_n = delta-hat_n^{2mu/(2mu+1)}; the two differ by
% ##  (ln n)^{2mu/(2mu+1)}.  The block was already inconsistent about which object
% ##  delta_n names -- (P-f), rem:h6bandwidth(0) and prp:stopfeas(iv) all read the
% ##  LOG-INFLATED delta_n while A5(ii) read the log-free one -- and
% ##  prp:stopfeas(iv) was inconsistent WITH ITSELF, its opening display reading
% ##  "delta_n |-> delta_n (ln n)^{2mu/(2mu+1)}" (the log-FREE delta_n being
% ##  inflated) while its (P-f) paragraph six lines later computed with the
% ##  log-INFLATED delta_n^2 carrying (ln n)^{4mu/(2mu+1)}.  So route (B) -- leave
% ##  A5(ii) log-free -- would have preserved a clause contradicting itself inside
% ##  one proposition.  That opening display was repaired in dispatch 11.
% ############################################################################
\begin{remark}[The logarithm in \eqref{eq:phirate}, resolved in A5(ii), and what it costs]\label{rem:rootnlog}
\emph{(1) Which object $\delta_n$ names.} \cref{prp:phirate} delivers
$\|\hat\varphi^{\pi}-\vdagger\|_{L^{2}(\pi)}=O_P(\delta_n)$ with
$\delta_n=(\hat\delta_n\ln n)^{2\mu/(2\mu+1)}$ --- \emph{at the infeasible reading $\hat\varphi^{\pi}$
of \cref{def:lfstep}, which is a separate matter from the one this remark settles and does not bear on
it}: this clause is about the \textbf{symbol} $\delta_n$ and the logarithm inside it, and $\delta_n$ is
the same deterministic sequence under either reading --- and (P-f), \cref{rem:h6bandwidth}\textup{(0)} and
\cref{prp:stopfeas}\textup{(iv)} all read that \emph{log-inflated} object. \emph{A log-free reading of
A5(ii) is off by $(\ln n)^{2\mu/(2\mu+1)}$}, and \cref{rem:phirscope}\textup{(2)} is where the warning
is discharged.

\emph{(2) The resolution taken.} A5(ii) is stated so that the symbol is the rate
\cref{prp:phirate} delivers,
\begin{equation}\label{eq:deltalog}
\delta_n\;:=\;\big(\hat\delta_n\ln n\big)^{2\mu/(2\mu+1)},
\qquad\text{and}\qquad \delta_n=o_P\big(n^{-1/4}\big) .
\end{equation}
\emph{The alternative, leave A5(ii) log-free and let \cref{thm:rootn}'s rate conditions carry
the logarithm, is rejected}, because it leaves (P-f) and A5(ii) reading the same symbol
differently and because it would put the logarithm in five places instead of one.

(3) What the resolution costs, term by term. This is the frontier arithmetic and it is not uniform
across the four products \cref{thm:rootn} needs.
\begin{enumerate}
\item[(a)] \emph{$\delta_n^{2}=o(n^{-1/2})$: free.} \eqref{eq:frontier} is a \emph{strict} inequality
between polynomial exponents, so writing $\hat\delta_n=n^{-\alpha}$, $\kappa=2\mu/(2\mu+1)$ and
$\kappa\alpha=\tfrac14+\varepsilon$ with $\varepsilon>0$ gives
$\delta_n^{2}=n^{-1/2-2\varepsilon}(\ln n)^{2\kappa}=o(n^{-1/2})$. This is
\cref{prp:stopfeas}\textup{(iv)} and it is why \eqref{eq:frontier} is \textbf{unchanged}.
\emph{And the strictness is not decoration}: at $\kappa\alpha=\tfrac14$ exactly,
$\delta_n=n^{-1/4}(\ln n)^{\kappa}$ is \emph{not} $o(n^{-1/4})$, so the boundary of \eqref{eq:frontier}
fails, which is what makes A5(iii)'s ``if and only if'' literally correct at the boundary
under \eqref{eq:deltalog}.
\item[(b)] $r_n\delta_n=o(n^{-1/2})$: free, and implied by \eqref{eq:frontier} rather than assumed.
Since $r_n=n^{-\alpha}$, $r_n\delta_n=n^{-\alpha(1+\kappa)}(\ln n)^{\kappa}$, and \eqref{eq:frontier}'s
$\kappa\alpha>\tfrac14$ with $\kappa<1$ gives $\alpha>1/(4\kappa)$, hence
\[
\alpha(1+\kappa)-\tfrac12\;>\;\frac{1+\kappa}{4\kappa}-\frac12\;=\;\frac{1-\kappa}{4\kappa}\;>\;0 ,
\]
a strictly positive polynomial margin that absorbs $(\ln n)^{\kappa}$. \emph{So \eqref{eq:frontier} does
double duty and no separate condition on $r_n\delta_n$ is imposed}; the margin is
\begin{equation}\label{eq:rndeltamargin}
\frac{1-\kappa}{4\kappa}\;=\;\frac{(2\mu+1)^{-1}}{8\mu/(2\mu+1)}\;=\;\frac1{8\mu} ,
\end{equation}
which shrinks as $\mu$ grows and is $\tfrac3{20}$ at $\mu=\tfrac56$ and $\tfrac1{16}$ at $\mu=2$.
\emph{The constant is exact at the frontier boundary and therefore cannot be improved from
\eqref{eq:frontier} alone}: the middle inequality above is an equality at $\alpha=1/(4\kappa)$, so
$1/(8\mu)$ is the infimum of $\alpha(1+\kappa)-\tfrac12$ over the configurations \eqref{eq:frontier}
admits.
% ARCHAEOLOGY.  A round-10 draft of this clause wrote 1/(4 mu), and with it 3/10 and
% 1/8; the error was a factor of two in 4 kappa and it moved no conclusion, the
% implication and its direction in mu being unaffected.  eq:rndeltamargin is
% displayed so that the arithmetic can be checked at sight.
\item[(c)] $\delta_n\epsilon_n=o(n^{-1/2})$: \textbf{not} free, and this is the one place the
logarithm bites. A5(iv) is a stated condition on $\epsilon_n$, and under
\eqref{eq:deltalog} it is a strictly stronger demand than under the log-free reading, by the factor
$(\ln n)^{\kappa}$. \emph{Nothing in the block pays for it}: it is a demand on the estimator of the
representer, is L-13's business, and it is charged here rather than absorbed.
\item[(d)] $r_n\epsilon^{L}_n=o(n^{-1/2})$: not affected by the logarithm, and zero under N1$'$.
(R-c) is where it lives.
\end{enumerate}

\emph{(4) What does \textbf{not} move.} \eqref{eq:muhalf}'s $\mu>\tfrac56$ is unchanged, because
\cref{rem:h6bandwidth}\textup{(2)} computes it at the log-inflated $\delta_n$ already and derives the
\emph{strictness} from precisely the failure of the logarithms to cancel in (P-f). \eqref{eq:frontier} is
unchanged by (a). A5(v) becomes \emph{weaker}, its right-hand side being inflated.
\eqref{eq:lxwindow} is unchanged, being a condition on $L_x$ and $h_u$ and not on $\delta_n$.
\end{remark}

\begin{remark}[The scope: what \cref{thm:rootn} covers, and the six things it does not]\label{rem:rootnscope}
(1) Mild ill-posedness only, and severe ill-posedness is not a gap in the proof but a place the
theorem does not reach. \eqref{eq:frontier} requires $\kappa\alpha>\tfrac14$ with
$\kappa=2\mu/(2\mu+1)<1$ and $\alpha=r/(2r+d_Z)<\tfrac12$, so it requires $\delta_n$ to be a
\emph{polynomial} rate faster than $n^{-1/4}$. Under H1's $\sigma_j\asymp j^{-a}$ that is what the source
condition delivers. \emph{Under severely ill-posed decay, $\sigma_j$ decaying exponentially, the rate
attainable from a source condition is logarithmic in $\hat\delta_n$, hence in $n$, and
$\delta_n=o(n^{-1/4})$ fails at every $r$}; A5(ii) is then not satisfiable and
\cref{thm:rootn} is unavailable. \emph{This is plan \S2.6's scope and it is restated here without
improvement}: mild ill-posedness with $s>a+\tfrac12$, which under H1 reads
$\mu>\tfrac12+1/(4a)$ and which implies \eqref{eq:muhalf} when $a\le\tfrac34$ and is an additional
restriction otherwise (\cref{rem:h6bandwidth}\textup{(2)}).
That last sentence is an \textbf{interpretation} and not a scope statement. The two displays it
links carry two different objects: $s>a+\tfrac12$
and the computed $\mu>\tfrac12+1/(4a)$ are about $\mu_{\text{\textup{H1}}}=s/(2a)$, the index of A2 at the
\emph{population} operator $\mathcal A_{\mathrm{cond}}$ on all of $\mathcal V$ at $\theta_\dagger$, while
\eqref{eq:muhalf}, \eqref{eq:frontier} and $\kappa$ carry (P-a.3)'s \emph{primitive} $\mu$, at the
finite-rank compression of the realised operator at $\theta_\circ$, and $a$ and $s$ are themselves
$\theta_\circ$-specific in (T-5) and $\theta_\dagger$-specific in H1. \emph{The bridge between them is the
operator interpolation \cref{rem:pasrc}\textup{(1)} shows unavailable at every $2\mu>1$, and
\cref{rem:pasrc}\textup{(4)} declines the commutativity proviso that would restore it}, so the implication
``$s>a+\tfrac12\Rightarrow$ \eqref{eq:muhalf}'' holds \textbf{under that declined proviso and not
otherwise}. \emph{Read literally without the qualification, this item both excludes designs the theorem
covers and claims designs it does not}; \cref{rem:pasrc}\textup{(5)} states the binding reading and gives
the two-sided example. \emph{What survives unqualified is the substance of the item}: severe ill-posedness
, exponentially decaying $\sigma_j$, makes the rate attainable from \emph{any} source condition
logarithmic, so A5(ii) fails at every $r$ whichever index the source condition is stated at, and that
argument does not go through $\mu=s/(2a)$ at all.

(2) The rate condition cannot be relaxed by a better remainder bound, and that is a weaker claim
than necessity. By \cref{prp:sharpquad}\textup{(i)},\textup{(iii)} the direct bound on
$\Psi^{(1)}_k$ is order-sharp: $\sup_{\|h\|\le\delta}\tfrac12|\langle\mathcal Bh,h\rangle|\asymp\delta^{2}$
under O-16(a), so no argument that merely bounds the second-order remainder more carefully will weaken
\eqref{eq:frontier}. \emph{Three readings must be blocked and \cref{rem:sharpscope} blocks them}: the result
does \emph{not} say that the realised $\Psi^{(1)}_k$ at $h_k=\hat\varphi^{(-k)}-\vdagger$ fails to be
$o_P(\|h_k\|^{2})$; it does \emph{not} say that $\delta_n=o(n^{-1/4})$ is necessary for
$\sqrt n$-estimability of $\theta_\dagger$ by \emph{some} estimator; and it does not say that the moment class
is exhausted. \emph{What relaxing \eqref{eq:frontier} would require} is a different selection, one with a
smaller $\sigma_q$ and $\mathcal K$, which is lever~1 of the strategy memo and is open, or a restriction on
the \emph{direction} of the first-stage error, which is an extra hypothesis and not a sharper argument.

(3) It does not derive its own inputs. $\mu$ and $\mathsf R_\ast$ are primitive at (P-a.3)
(\cref{rem:pasrc}); the preliminary $\hat\varphi^{(0)}$ that \cref{def:lfstep} needs at the rate
\eqref{eq:prelimrate} has no supplier in this file (open item Q-1); the representer's estimator $\hat\ell$ is
priced by A5(iv) and constructed nowhere; N1 at $L_\star$ is \emph{assumed} in (T-2) and
\emph{arranged} nowhere, the re-anchoring of \cref{rem:anchor} being available but not performed, and by
\cref{cor:centarity} it is $p(\mathsf L-1)$ conditions rather than the $p$ a contracted centring
supplies; and \eqref{eq:nuhat}, if (R-e.2) is taken, is
assumed. \emph{\cref{thm:rootn} is a theorem about an estimator five of whose ingredients are assumed to
exist with stated properties}, and saying so is not a hedge: each such ingredient is named in (T-2), (T-5), (R-a) or
(R-e) and none is hidden in a regularity condition.

(4) It says nothing about efficiency, about the plug-in, or about what happens when A2 fails.
Those are \texttt{prp:effbound}, \texttt{prp:plugin} and \texttt{cor:trend}, Step~8. In particular
A5(v) is \textbf{not} consumed anywhere in Part~I: \cref{def:crossfit}(b) runs the
$\varphi$-step at the fixed anchor $\theta_\circ$, so no Landweber pass on
$\Delta Y-\Delta Z'\hat\theta$ occurs and the plug-in cost $\sqrt{\hat m}\,n^{-1/2}$ never arises. And when
A2 fails, $\ell_\star$ does not exist, $L_\star$ is not defined, $\psi_\star$ is not defined and
\eqref{eq:JSig} is not a well-formed display; the object to report is then the Picard diagnostic and the
plug-in at its slower rate, and pivotal inference is declined. \emph{What fails under $\neg$A2 is the
\emph{existence} of the objects and not a rank condition}:
\cref{lem:momzero}\textup{(iv)} shows the rank condition never had content beyond A1b(iv), and A1b(iv)
holds whether or not A2 does.
% ARCHAEOLOGY.  Through dispatch 12 this was stated as "A3 has no truth value,
% being stated conditional on A2"; the substance survives the retirement of that tag.

(5) It is an asymptotic statement whose asymptotics do not bind at any sample size a user will run,
and \cref{cor:ratefloor} and not \cref{thm:rootn} is what describes a feasible $n$. With
$\delta_n=(\hat\delta_n\ln n)^{\kappa}$, $\hat\delta_n=n^{-\alpha}$ and the notation of item~(1), the
$\theta$-remainder of Step~2 of the proof is $O_P(n^{-1/2}+\sqrt n\,\delta_n^{2})$, and its second term is
\begin{equation}\label{eq:rootnturn}
\sqrt n\,\delta_n^{2}\;=\;n^{-\beta}\,(\ln n)^{2\kappa},
\qquad \beta:=2\kappa\alpha-\tfrac12>0\ \ \text{by \eqref{eq:frontier}} ,
\end{equation}
and $t\mapsto e^{-\beta t}t^{2\kappa}$ is \emph{increasing} for $t<2\kappa/\beta$. So \eqref{eq:rootnturn}
\textbf{rises} in $n$ up to
\[
n^{\ast}\;=\;\exp\Big\{\frac{2\kappa}{2\kappa\alpha-\tfrac12}\Big\}
\]
and only then falls. \emph{The threshold is astronomically large near the frontier and that is not a
pathology of a bad parameterisation but a property of \eqref{eq:frontier}}: $\beta\downarrow0$ as
$\kappa\alpha\downarrow\tfrac14$, so $n^{\ast}\uparrow\infty$ there. At $\mu=2$ (so $\kappa=\tfrac45$) and
$\alpha=0.32$, inside \eqref{eq:frontier}, since $\kappa\alpha=0.256>\tfrac14$, one gets
$\beta=0.012$ and $n^{\ast}=e^{400/3}\approx10^{58}$. Two things follow and neither weakens
\eqref{eq:rootn}. First, \eqref{eq:rootn}'s $o_P(1)$ is a statement about the limit and is proved;
what \eqref{eq:rootnturn} says is that the \emph{approach} to that limit is not monotone and that
simulation evidence at $n\le10^{4}$ should be read against \cref{cor:ratefloor}'s
\eqref{eq:ratefloor}, in which $\delta_n^{2}$ appears as a term and not as a remainder. Second, the
honest reading of \eqref{eq:frontier} is that it is a \emph{qualitative} frontier: crossing it makes the
remainder vanish eventually, and says nothing about where. \emph{Whether the companion term $r_n\delta_n$
turns over as late is a separate computation and is not made here}; \cref{rem:rootnlog}\textup{(3)}(b)'s
\eqref{eq:rndeltamargin} prices its \emph{margin} to $n^{-1/2}$ and is the place to make it.

(6) It is \textbf{pointwise in $P$}, and no uniformity over any class of laws is claimed or proved.
Item~(5) frames the $n^{\ast}$ point as a
finite-sample caveat; the more exact statement is that there is no uniformity claim for it to be a caveat
about. Every hypothesis of (T-1)--(T-5) and (R-a)--(R-e) is imposed at a \emph{single} data-generating
law: $\mu$, $\mathsf R_\ast$, $a$, $s$, $\varepsilon_\varrho$, $r$, $\bar c_\pi$, $\Lambda^{(0)}_\star$,
$\Lambda^{(1)}_\star$ and the singular system of $\mathcal A_{\mathrm{cond}}$ are all $P$-specific
constants, and the conclusion \eqref{eq:rootn} is a convergence in distribution under that $P$.
Three consequences follow and each is a real limitation rather than a formality. First, the
$o_P(1)$ in \eqref{eq:rootn} is not uniform over any neighbourhood of $P$, so the theorem does not deliver
confidence sets whose coverage is \emph{honest} over a class, for that one needs the supremum over
$P\in\mathcal P$ of the absolute coverage error of the studentised interval to vanish, at a fixed nominal
level and over a class $\mathcal P$, and no such class is even defined here.
\emph{The standard reference for that notion in ill-posed models is}
\citet{babii2016}, \emph{whose honest sets are built at a linear operator under
completeness and so do not cover the equation solved here}.
% CORRECTED 2026-08-03.  This read that the standard reference "is not cited
% because its key is absent from this project's bibliography".  IT IS PRESENT:
% babii2016 was entered the same day, copied verbatim from the author's master
% bibliography.  The claim was true when written and stopped being true when the
% key landed, which is why a sentence asserting an absence has to be swept
% whenever the bibliography grows.
Second, \emph{that is exactly the gap item~(5)
displays numerically}: $n^{\ast}\uparrow\infty$ as $\kappa\alpha\downarrow\tfrac14$, so within any class
containing laws arbitrarily close to the frontier the convergence cannot be uniform, and the $10^{58}$ is
the pointwise statement's honest shadow rather than a bad constant. Third, the positive-definiteness
\eqref{eq:sigmapd} is likewise pointwise: $\lambda_{\min}(\Sigma)>0$ at each $P$ and may be arbitrarily
small along a sequence of laws. \emph{What would be needed for a uniform statement}, and it is open work
and not a hedge: a class $\mathcal P$ on which the constants above are bounded uniformly, a uniform version
of \eqref{eq:phirate}, which would require (P-a.3)'s $\mathsf R_\ast$ and $\mu$ to be uniform over
$\mathcal P$ at a \emph{random} compressed operator, and a Lindeberg condition uniform in $P$. None is
attempted. \emph{The convention that a limit theorem is read pointwise unless uniformity is asserted is
standard and is not being restated as a defect}; what is a defect is asserting a scope remark that
enumerates finite-sample caveats and omits the quantifier they all sit inside.
\end{remark}

\begin{remark}[Where each hypothesis of \cref{thm:rootn} is used, by clause]\label{rem:rootnhyp}
\emph{This is the hypothesis-usage audit.} ``Step'' refers to the proof of \cref{thm:rootn} in
\texttt{proofs.tex}.
% ARCHAEOLOGY.  Four inherited lists are merged here.

\begin{enumerate}
\item[(1)] \emph{Consumed, with the site.} A0, fold independence and the i.i.d.\ central limit theorem
(Steps~3, 5), and the compact \emph{interval} $\mathcal X$ for the Sobolev embedding (\cref{lem:epterm}).
A1, first clause, $E[\Xi(\Delta U,\Delta Z)\,\widetilde W]=0$ for the
$\IR^{p\times(\mathsf L-1)}$-valued generic direction $\Xi$ of \cref{def:dirclass}, hence
$\partial_{\tau_2}\Psi\equiv0$ exactly (\cref{lem:momzero}(i), \cref{lem:biasexp} Step~2, through
\cref{prp:orthog}(ii)); second clause, $E[\Delta Z\Delta U]=0$ (\cref{lem:momzero}(i),
\cref{lem:prelimtheta}). A1b(i), $\vdagger\in\mathcal V$, so that $h_k\in\mathcal V$ and the pairing
$\langle\nu_{g_k}/\pi,h_k\rangle$ is the one \cref{lem:nusigma}(iii) computes. A1b(iii), so that the
$\varphi$-step targets $\vdagger$ and not a set; consumed through (T-5). A1b(iv), nonsingularity of
$E[\Delta Z\Delta Z']$ (\cref{lem:prelimtheta}) and the second-moment clauses (\cref{lem:momzero}(ii)).
A2, existence of $\ell^\Pi_\star$, hence of $\nu_{\Delta Z+\ell_\star\widetilde W}\equiv0$
(\cref{lem:biasexp} \eqref{eq:biasb1}). A$2'$(i.a) \textbf{and} (i.b), \emph{jointly}, because both are
required for \cref{cor:elldiff} to apply at all: A2b is what that corollary assumes, and A$2'$(i)'s closing
clause defines $\ell_\star$ pointwise from the version satisfying both. Of the two \emph{bounds} the
corollary returns, $\|\ell_\star\|_\infty\le\|\varpi_u\|_\infty\Lambda^{(0)}_\star$ uses only the (i.a)
budget while $\|\partial_u\ell_\star\|_\infty$ uses both; the latter is consumed by
\cref{prp:sharpquad}'s inheritance in \cref{lem:biasexp}(2). \emph{\cref{lem:momzero}\textup{(iii)} is
\textbf{not} a second consumer of $\|\partial_u\ell_\star\|_\infty$}: by
\cref{lem:momzero}\textup{(iv.1)} that clause's second-derivative envelope bounds the
zero matrix. \emph{A$2'$(i.b) itself is nonetheless consumed}, it
licenses \cref{cor:elldiff} at all, hence $\ell_\star$ pointwise, and it keeps the
\cref{prp:sharpquad} consumer.
A$2'$(iii), $E\|b_\star\|^{2}<\infty$ for $\Sigma<\infty$. N1, $E[L_\star\mathsf r_k]=0$
% ARCHAEOLOGY.  lem:momzero(iii) was listed as a second consumer through dispatch 12.
(\cref{lem:biasexp} Step~2, through \cref{prp:orthog}(iii)).
B1(i), (ii), (iv.a) and (iv.d), \cref{cor:elldiff}; B1(iv) in full, \cref{prp:orthog}'s list, inherited by
\cref{lem:balance}. C1(ii), $\bar c_\pi$, at every display. C2$'$(i), \cref{prp:orthog}(i)'s
differentiation, inherited, \emph{and a \textbf{second and direct} consumer}, namely
\cref{lem:momzero}\textup{(v)}, where the conditional density of $\Delta Y$ given $\mathcal F_{XZW}$ is
what gives $P(\Delta U=0)=0$ and hence $\Sigma\succ0$. A4(i), \emph{consumed \textbf{definitionally} and at \textbf{no step} of this proof, which is the
stronger of the two statements and the one Section~6 needs}. It reaches this theorem through
\cref{def:crossfit}\textup{(a)}, which the theorem adopts: A4(i)'s $\varepsilon_\varrho\in(0,1/\mathsf L)$
is what makes \eqref{eq:trunset} non-empty, its quantifier over \emph{every} $l\in\{0,\dots,\mathsf L-1\}$
, the base cell included, is what puts $\varrho$ itself inside \eqref{eq:trunset}, so that truncating
$\hat\varrho^{(-k)}$ onto that set cannot degrade A4(iii)'s rate, and it is what grounds that clause's
per-cell bound $\|\mathsf r_k\|_\infty\le1-2\varepsilon_\varrho$. \emph{\textbf{It is not the ground of
the sure bound the proof actually uses}}: that bound is \eqref{eq:estcontrnorm}'s \textbf{Euclidean}
$\|\mathsf r_k\|\le c_{\widetilde W}$, whose derivation states in terms that ``no overlap condition is
consumed and A4(i) is not on this lemma's list'', $\varrho$ lies in the simplex for free and
\eqref{eq:trunset} is contained in it at \emph{every} $\varepsilon_\varrho$, with
$c_{\widetilde W}=\min\{\sqrt2,\sqrt{\mathsf L-1}\}$ the simplex diameter of \cref{cor:rangegap}, equal to
$1$ at $\mathsf L=2$. \emph{So no step of the $\theta$-step reads $\varepsilon_\varrho$ at a value}, and
what A4(i) asks of the application here is positivity at every cell and nothing more.
A4(ii)--(iii), $r_n$.
% RETYPED 2026-08-10 (Phase 6a).  THREE THINGS WERE WRONG AND THIS IS THE PER-THEOREM
% USAGE AUDIT, so a reader learns from it where A4(i) binds in the theta-step.
%  (1) THE GROUND.  The entry attributed ||r_k|| <= c_Wtilde to "eq:trunset and
%      A4(i)".  eq:estcontrnorm's derivation says the OPPOSITE IN PRINT, and so does
%      lem:epterm Step 5: "the ground of that sure bound is the truncation alone and
%      NOT A4(i)".  The audit pointed at a site that had been regrounded under it.
%  (2) THE NORM.  It wrote ||r_k||_inf <= c_Wtilde, fusing eq:estcontrnorm's
%      EUCLIDEAN bound with def:crossfit(a)'s SUP-norm constant and standing on the
%      ground of neither.  eq:estcontrnorm is Euclidean; the sup-norm statement is
%      def:crossfit(a)'s 1 - 2 eps_varrho, a different display with a different
%      constant, and it IS the one that needs A4(i).
%  (3) THE CONSUMER.  Once (1) and (2) are fixed the honest answer is that A4(i) has
%      no STEP-level consumer in Part I at all -- only a DEFINITIONAL one, at
%      def:crossfit(a) -- and that is printed above rather than left to be inferred,
%      because it is the answer the empirical section reads: the burden is
%      positivity at every cell, not a declared value of eps_varrho.
A5(i), $\hat\delta_n$, through (T-5).
A5(ii) as \eqref{eq:deltalog}, read at each fold per \eqref{eq:foldrates}, \emph{and the two halves are
consumed at different places, which the split of \cref{lem:prelimtheta} makes visible}: the \emph{first}
half supplies $\delta_n$ at Step~0 and is the whole of what \cref{lem:prelimtheta}\textup{(i)} needs, while
the \emph{second} half, $\delta_n=o_P(n^{-1/4})$, is consumed at Steps~1, 2 and~4 of this proof directly and
\textbf{not} through \cref{lem:prelimtheta}, whose clause \textup{(ii)} this theorem never invokes, Step~2
takes clause \textup{(i)}'s uncollapsed rate and applies the threshold itself. A5(iii), \eqref{eq:frontier},
which delivers (a) and (b) of \cref{rem:rootnlog}(3), \emph{and which is the one clause of (T-4) that no
result in this file delivers}. A5(vi) with $s>\tfrac12$, read at each fold, the $L^{4}$ bound in
\cref{lem:epterm}. (R-a), (R-c), (R-e), as stated. (R-b), $E\|\Delta Z\|^{4}$ at \cref{lem:epterm}
Step~2 and \cref{lem:sandwich} Step~9; $E[\Delta U^{4}]$ at \cref{lem:epterm} Step~5.
\item[(2)] \textbf{Not} listed in \textup{(T-1)}--\textup{(T-5)} and not consumed, the five clauses
a reader will look for, and why each is absent. A$2'$(ii), \textbf{retired}: the
variance needs the \emph{level} and the clause priced the \emph{derivative}, and the clause is in any case
implied by A$2'$(i) in full together with B1(i), (ii), (iv.a), (iv.d), so it is retired in place at
A2$'$ exactly as A3 was; see \cref{lem:momzero}'s closing paragraph. A5(v), see
\cref{rem:rootnscope}\textup{(4)}. A1b(ii), \cref{lem:nusigma}(ii) is not used; only
(i), (iii), (iv) and (v) are, and those need C1(ii) and, for (v), A1. A6, retired
(\cref{prp:tccfails}); A6b and \cref{prp:tccrepair} are not consumed, \cref{def:lfstep} needing no cone
condition. A3, retired (A3), being proved from (T-1)'s A1 and A1b(iv) together
with (T-2)'s A2, A$2'$(i), A$2'$(iii) and (T-3)'s B1(i), (ii), (iv.a), (iv.d), at
\cref{lem:momzero}\textup{(iv)}. \emph{The retirement removes a hypothesis and adds none}; the
nonsingularity \cref{lem:sandwich}(i) Step~6 needs is A1b(iv)'s, listed in (T-1).
% ARCHAEOLOGY.  A round-10 draft headed this item "Listed and not consumed", which
% was wrong: none of the five appears in any group above.  A2'(ii) was retired in
% dispatch 14 and is a SIXTH entry of the A3 species rather than a fifth of this
% item's; A3 was in (T-2) through dispatch 12 and is a FIFTH entry here, not a
% fourth.
\item[(3)] Inherited as a block through (T-5) and not re-derived. H1, A5(vi$'$), A8's basis and
upper end, $p_K\ge s-\tfrac12$, C3, C5$'$(i), (P-a.1)--(P-a.3), (P-c)--(P-g), \eqref{eq:muhalf}, and A7 in the
amended form \eqref{eq:a7amend}. \emph{All are read at $\theta_\circ$}, per \cref{def:crossfit}(b), and each
is $\theta$-specific in the sense of \cref{rem:R1}\textup{(iv)}.
\emph{C5$'$ is inherited \textbf{by clause}, and a block citation would charge the theorem with
C5$'$(ii)}. That clause is a lower bound on $h_x$ and (P-g),
inherited on the same line, is an upper bound on $h_x$; at the strengthened form of C5$'$(ii) the two have no
common solution whenever $s\le\tfrac12+18\alpha/(8-3\alpha)$, \eqref{eq:hxwindow}, a band that
\eqref{eq:frontier} leaves open at every $\alpha\in(\tfrac14,\approx0.359)$. \emph{The consequence would be
that (T-5) is unsatisfiable and this theorem holds \textbf{vacuously}, not that it is false}; the by-clause
citation is what makes (T-5) a satisfiable group at those configurations, and C5$'$(ii) carries the
corresponding repair in the maintained block.
% ARCHAEOLOGY.  This line read "C5'" through dispatch 20, which on this file's
% register rule is a block citation.
\item[(4)] What \eqref{eq:prelimrate} does and does not discharge, and the answer is negative. A
one-step needs a \emph{rate} for its
preliminary and L-14 supplies consistency only. \emph{\eqref{eq:prelimrate} does not close that gap,
% ARCHAEOLOGY.  This item answers the round-3 strategist-critic's m-23.
because it is about a different object}: it is a condition on $\varsigma_0$, the exponent of
$\hat\varepsilon_0$, which is the $L^{2}(\pi)$ rate of the \emph{preliminary structural estimator}
$\hat\varphi^{(0)}$ inside \cref{def:lfstep}, and it is hypothesis (P-a.2) of the $\varphi$-step. The
one-step's preliminary is $\hat\theta^{(0)}$, an estimator of $\theta_\dagger$; nothing in
\eqref{eq:prelimrate} bears on it. \emph{What closes the gap is \cref{lem:prelimtheta}\textup{(i)}}: the
plug-in \eqref{eq:prelimtheta} has rate $O_P(n^{-1/2}\vee\delta_n)$ on a list containing no threshold, and
$\delta_n=o_P(n^{-1/4})$ is A5(ii)'s second half, applied where it is needed and not before. So the gap is real, the registered instrument is the wrong one, and the
discharge required specifying the preliminary, which no file did before Part~I.
\item[(5)] Why \textup{(R-a)}--\textup{(R-e)} are stated on the theorem and \textbf{not} promoted to
A5. Every clause of A5 constrains an object the \emph{maintained block} introduces, at a single
estimator, while (R-a)--(R-e) constrain $\hat\varphi^{(-k)}$, $\hat L^{(-k)}$, $\hat\ell^{(-k)}$,
$\hat b^{(-k)}$ and $D_k$, objects that exist only after \cref{def:crossfit}, a definition Part~I
introduces. A maintained clause cannot quantify over objects the block does not define. In
addition (R-a) is a construction claim governed by \cref{rem:obligations}'s test, (R-b) is a moment
condition on the primitives that would belong beside A1b(iv) rather than in the nuisance-rate block, and
(R-c) and (R-e) are \textbf{disjunctions}. \emph{The ground is \textbf{not} that A5 prices rates only},
A5(vi) and A5(vi$'$) being qualitative.
% ARCHAEOLOGY.  A round-10 draft gave the refuted ground.  def:crossfit(d) is the
% narrow exception to the object-existence rule: it re-reads three existing A5
% clauses at the folds rather than adding a clause.
\end{enumerate}
\end{remark}

% ---------------------------------------------------------------------------
%  THE TWO REMARKS BELOW WERE EXTRACTED FROM THE STATEMENT OF thm:rootn on
%  2026-08-02, when that statement ran 130 live lines against a house style
%  (.claude/references/theorem-environment-style.md) that asks a theorem for one
%  descriptive sentence and a "Suppose ... Then ..." pair.  NOTHING WAS DELETED:
%  every sentence removed from the statement is below, verbatim except for the
%  connectives needed to make it stand outside the enumeration it used to sit in.
%  The statement now carries conditions and conclusions and no commentary, and it
%  points here by name.  These two are placed at the END of Part I on purpose:
%  nothing follows them in this file, so adding them to the shared theorem
%  counter renumbers no existing environment.
% ---------------------------------------------------------------------------

\begin{remark}[What the hypothesis groups of \cref{thm:rootn} carry, and the six things they do not]\label{rem:rootngroups}
Each item below stood as a parenthesis inside the theorem's own hypothesis list. None of it is a
hypothesis, and the theorem's content is unchanged by its removal.
\begin{enumerate}
\item[(1)] \textup{(T-2)}: no local rank condition on $\ell_\star$ is in the group, and neither is
A3. The nonsingularity such a moment would ordinarily require follows from (T-1) and the rest of (T-2) at
\cref{lem:momzero}\textup{(iv)}, and the rank condition the theorem needs is A1b(iv), which is in (T-1).
A3 stood in this group through dispatch 12 and is retired on that same ground.
\emph{Two things about the group's N1 clause are recorded here rather than in the statement.} Its
\textbf{arity}: N1 at $L_\star$ is $p(\mathsf L-1)$ scalar conditions, one per entry, and not a
condition on a contraction of the $\mathsf L-1$ columns, so it is $p(\mathsf L-2)$ conditions
\emph{stronger} than a contracted centring at every $\mathsf L\ge3$ and identical to one at
$\mathsf L=2$ (\cref{cor:centarity}). And its \textbf{anchor-pinning}: A$2'$(iii) leaves the anchor
$b_\star$ free, but (T-2) and (R-a) together do not --- (R-a) forces $\hat b^{(-k)}\to b_\star$ and
(T-2) requires $L_\star$ centred, so $b_\star$ is the N1 anchor of \cref{rem:anchor} and any
construction of $\hat L^{(-k)}$ must target that one.
\item[(2)] \textup{(T-4)} is \textbf{not} homogeneous in status, and that is why the group was
annotated; \emph{it is \textbf{less} homogeneous from 2026-08-27, one clause having moved from
delivered to assumed}. A5(i) is \textbf{delivered} by (T-5)
rather than required, through \cref{lem:opnoise}, and is listed all the same, because Steps~2 and~4 cite
it by name.
\emph{A5(ii)'s first half was delivered through dispatch 27 and is now \textbf{assumed}}: the delivery
ran through \eqref{eq:phirate}, which \cref{prp:phirate} is now scoped to state at the infeasible
$\hat\varphi^{\pi}$, whereas A5(ii) is stated at the $\hat\varphi$ that
\cref{def:crossfit}\textup{(b)} plugs in, and \cref{rem:phirscope}\textup{(7)} is open at
\textup{(7c)} and \textup{(7d)}. \emph{The listing is therefore no longer redundant at that clause; it
is the charge}, and \cref{thm:rootn} is a conditional theorem on it. \emph{The clause is still listed in
the form \eqref{eq:deltalog} for the reason it always was}, which is that the display fixes the object
the symbol $\delta_n$ denotes and Steps~2 and~4 read that symbol by name. A5(ii)'s \emph{second} half,
$\delta_n=o_P(n^{-1/4})$, is still delivered, being about the deterministic $\delta_n$ alone.
A5(iii) sits at the
other end: \emph{no result in this file delivers it and it is therefore assumed}, being a restriction on
the primitives $(\mu,r,d_Z)$ that \cref{prp:phirate} presupposes rather than establishes, that
proposition returns \eqref{eq:phirate} \emph{whenever} \eqref{eq:frontier} holds and says nothing when it
fails. A5(vi) with $s>\tfrac12$ is a third case: (T-5)'s A5(vi$'$) delivers the $H^{s}$ bound at the
\emph{final linearisation stage} $\bar k$ of \cref{def:lfstep}, \emph{a stage index, not a fold index},
the two being one glyph apart, while $s>\tfrac12$ is a restriction on the primitive $s$ and is
delivered by nothing.
\item[(3)] \textup{(R-a)}: the \emph{sure} boundedness of $\partial_u\hat\ell^{(-k)}$ is
\textbf{not} required. No first-step construction the block contemplates delivers a sure bound on an
estimated object, what a series or kernel estimator of $\ell_\star$ supplies being a bound that is finite
with probability tending to one, exactly (R-a)'s $O_P(1)$ display. The sites that read finiteness
of the derivative envelope --- \cref{lem:biasexp} Step~4, \cref{lem:sandwich} Step~5 and
\cref{thm:rootn} Step~1, where it bounds $\partial^{2}_{\theta\theta'}\bar\psi_n$ and the Taylor
remainder $\Psi^{\theta}_n$ --- carry it on the
\textbf{right} of an inequality and are read on the event where it is finite, which is where every $O_P$
conclusion in those results is read in any case. \emph{No claim of exhaustiveness is made, and that
is deliberate}: the draft this sentence replaces said ``the two sites'' and was short by the third,
while \cref{def:onestep}'s own reading named that third and was short by these two --- one three-site
set described twice, in halves, by two claims neither of which held all of it. \emph{Boundedness of $\hat\ell^{(-k)}$ itself is
retained and is not part of this}: it is what \cref{def:crossfit}\textup{(c)} already requires, for
\cref{def:weight}(b) to be well posed. \emph{And \cref{def:crossfit}\textup{(c$'$.2)}'s $[0,+\infty]$
typing is unchanged by (R-a)}: that typing is the \emph{definition}'s, fixing the range of the map
$\omega\mapsto\|\cdot\|_\infty$ before any hypothesis is imposed, and it must stay extended-valued because
\cref{def:crossfit} is adopted by results that do not list (R-a); \emph{tightness} of that
extended-valued map, and hence finiteness of it with probability tending to one, but \textbf{not} at
every sample point, is what (R-a) adds to it, at the results that do.
\item[(4)] \textup{(R-b)}: a fourth moment of $b_\star$ is \textbf{not} needed. $b_\star$ enters
Part~I at two lines only, both at \emph{second} moment through A$2'$(iii) (\cref{lem:momzero} Step~1), and
nothing in Part~I needs a fourth moment of $\psi_\star$, \cref{lem:sandwich} Step~8 applies conditional
Markov to squares, its Step~10 needs only $\IE_n\|\psi_\star\|^{2}=O_P(1)$, and the central limit theorem
used is Lindeberg--L\'evy. The \emph{fourth}-moment need is on $\hat b^{(-k)}$, and (R-a) supplies it.
% ARCHAEOLOGY.  A round-10 draft carried a third clause, E||b_star||^4 < infinity.
\item[(5)] \textup{(R-e.2)}: \eqref{eq:re2window} is stated on that branch because that branch is
the one place both of its ingredients arrive at once. \Cref{lem:balance}\textup{(iii)} carries A8
\emph{in full}, so (e.2) imports A8's lower end $L_x\gg n^{1/(2s)}$; (T-5) simultaneously imports (P-d),
$L_x=o\big(n^{1-2\alpha}/\ln n\big)$; and \eqref{eq:frontier} forces $\alpha>\tfrac14$, so (P-d) binds and
the joint window for $L_x$ is non-empty if and only if \eqref{eq:re2window}. \emph{The arithmetic is
printed at (P-d) already}, ``had A8's lower end been carried as well, the window would instead require
$s>(2(1-2\alpha))^{-1}$'', and \cref{rem:h6bandwidth}\textup{(5)} records that the reference design
clears it. (T-4)'s only $s$-floor is $s>\tfrac12$, which does not reach it.
\emph{Witness, chosen so that nothing else is empty there}: $(\mu,s,a,\alpha)=(2,1.35,0.3375,0.32)$ clears
\eqref{eq:frontier} at $\kappa\alpha=0.256$, clears \eqref{eq:lxwindow} at $s>1$ and clears
\eqref{eq:hxwindow} at $s>1.3182$, while A8's lower end asks $L_x\gg n^{0.3704}$ and (P-d) allows
$L_x=o(n^{0.36}/\ln n)$: \emph{empty}, the threshold \eqref{eq:re2window} being $s>1.3889$.
% WITNESS RE-CHOSEN 2026-08-10 (Phase 6a) BECAUSE THE (P-g) REPAIR FALSIFIED THE OLD
% ONE, AND THIS IS RECORDED RATHER THAN PATCHED SILENTLY.  The old witness was
% (mu, s, a, alpha) = (2, 1.3, 0.325, 0.32) and it read "clears eq:hxwindow at
% s > 1.2273".  With (P-g)'s h_u restored, eq:hxwindow at alpha = 0.32 is s > 1.3182,
% which s = 1.3 FAILS -- so the clause "chosen so that nothing else is empty there"
% became false of its own witness.  The witness needs s strictly between eq:hxwindow
% (1.3182) and eq:re2window (1/(2(1 - 2 alpha)) = 1.3889), which is a real interval,
% and s = 1.35 sits in it.  a = 0.3375 keeps the block's convention that the chosen
% numbers satisfy s/(2a) = mu = 2 -- a coincidence the block keeps deliberately and
% does not invoke as an identity, per the note at the flagship configuration.
% RE-VERIFIED AT THE NEW s: A8's lower end is n^{1/(2s)} = n^{1/2.7} = n^{0.3704} and
% (P-d) allows o(n^{0.36}), so 0.3704 > 0.36 and the joint L_x window is STILL empty,
% which is the whole point of the witness; eq:lxwindow is s > 1 and clears; mild
% ill-posedness s > a + 1/2 = 0.8375 clears; p_K >= s - 1/2 clears at 8 >= 0.85;
% alpha < 68/145 clears at 0.32.  The point the clause makes is unchanged.
\emph{So without \eqref{eq:re2window} the branch is vacuous below that threshold and the theorem is
\textbf{not}}: (e.1) is unaffected and carries the conclusion there. \emph{An either/or that advertises a
choice which does not exist at part of the parameter set is a defect of the enumeration, not of the
theorem}, and \eqref{eq:re2window} is what makes the choice real wherever it is offered.
\item[(6)] \emph{\textup{(R-d)}'s letter is \emph{reserved and not reused}}, on the \textup{(P-b)}
precedent, so that a citation to ``(R-d)'' lands on the reservation and not on the wrong hypothesis.
\item[(7)] \emph{\textup{(R-e.1)}'s fold index is part of the hypothesis and not a reading of it}, for
the reason \cref{def:crossfit}(d) gives, and it is what \cref{thm:rootn} Step~4(a) and
\cref{lem:sandwich} Step~3 consume.
\end{enumerate}
\end{remark}

\begin{remark}[How the conclusion of \cref{thm:rootn} is read]\label{rem:rootnread}
\begin{enumerate}
\item[(1)] The quantifier is ``every'' and not ``every $\vartheta$ with
$\vartheta'J^{-1}\Sigma J^{-1\prime}\vartheta>0$''. By \cref{lem:momzero}\textup{(v)} and
\eqref{eq:sigmapd}, $\Sigma\succ0$ under A1b(iv) and C2$'$(i), both already in (T-1) and (T-3), so
$J^{-1}\Sigma J^{-1\prime}\succ0$ and the positivity guard is \emph{implied by the theorem's own
hypotheses}. A proviso that a theorem's hypotheses entail should be stated as a consequence and not
as a condition the user checks.
% ARCHAEOLOGY.  Through dispatch 13 the quantifier read "every vartheta with
% vartheta' J^{-1} Sigma J^{-1'} vartheta > 0"; the change is a strengthening of
% the statement at no cost in hypotheses, and the entailment was found by probing
% that very display.
\item[(2)] What the guard is \textbf{not} replaced by is a uniform lower bound.
$\lambda_{\min}(\Sigma)>0$ at the fixed $P$ of the theorem and may be arbitrarily small at nearby laws,
which is \cref{rem:rootnscope}\textup{(6)}.
\item[(3)] \emph{The denominator $\vartheta'\hat V\vartheta$ can still be zero or negative in a finite
sample}, $\hat V$ is a sample sandwich and is set to $0$ off $\{\hat J\ \text{nonsingular}\}$ by
\eqref{eq:Sighat}, so the studentised display is read, as such displays always are, on the event of
probability tending to one where the denominator is strictly positive. That is the reading the statement
now records in one clause.
\item[(4)] \emph{The contraction direction is $\vartheta$, which is this file's binding letter for exactly
this object}: $\vartheta'\psi$ is the moment for $\vartheta'\theta_\dagger$. It is \textbf{not}
$\varkappa$, which is a local bound variable in \eqref{eq:jackson}, nor $\mathsf r$, which is the
$\varrho$-direction.
% ARCHAEOLOGY.  Glossary C25.
\item[(5)] \emph{Both limit laws of \cref{thm:rootn}, the one in \eqref{eq:rootn} and the studentised
one, are written out and carry no symbol}: italic $N$ is a count, a null set or a covering number
everywhere else in this corpus, and calligraphic $\mathcal N$ is unavailable, being \cref{lem:density}'s
closed subspace and A2c's null Borel set.
\end{enumerate}
\end{remark}

% EXTRACTED FROM THE STATEMENT OF cor:ratefloor on 2026-08-02, same operation and
% same rule.  Item (c)'s unlabelled display went with it: it is a SUBSTITUTION
% into eq:ratefloor under two of the corollary's own branches, not a further
% conclusion, and the corollary's printed statement already says in words what
% (R-e.2) replaces.  No \label moved.
\begin{remark}[How much of \eqref{eq:ratefloor} is in primitives, by branch]\label{rem:ratefloorprim}
The answer differs by branch, which is why it is stated by branch.
\begin{enumerate}
\item[(a)] \emph{Always in primitives.} $\delta_n=(\hat\delta_n\ln n)^{2\mu/(2\mu+1)}$ with
$\hat\delta_n=n^{-r/(2r+d_Z)}$ and $r_n=n^{-r/(2r+d_Z)}$, so the terms $n^{-1/2}$, $\delta_n^{2}$ and
$r_n\delta_n$ have exponents that are functions of $(\mu,r,d_Z)$ alone, through \cref{prp:phirate} and
A4(iii).
\item[(b)] Never in primitives on the branch \cref{cor:ratefloor} states. $\epsilon_n$ is priced
by A5(iv) and $\epsilon^{L}_n$ by (R-c.2), and \emph{neither has a supplier anywhere in the maintained
block}: $\hat\ell$ is constructed nowhere (\cref{rem:rootnscope}\textup{(3)}) and $\epsilon^{L}_n$ is
priced by nothing at all. So on the hypotheses as listed, (R-e.1) with (R-c.2), the terms
$\delta_n\epsilon_n$ and $r_n\epsilon^{L}_n$ carry two exponents that are \textbf{assumed and not
derived}, and the claim ``the rate is in primitives'' is false for them. \emph{Note also that $s$ enters
no exponent on this branch}, A8's \textbf{lower} end being consumed only through (R-e.2); its basis
and upper-end clauses reach the theorem through (T-5), and neither carries $s$ into an exponent.
\item[(c)] Entirely in primitives on one branch, and that branch is what \S19.1's sieve decision
buys. Under (R-e.2) \textbf{and} (R-c.1), the sieve solve for the first-order bias, N1$'$ at each
$D_k$ for the $\varrho\times L$ term, the two unsupplied symbols both disappear, and
\[
\big\|\hat\theta-\theta_\dagger\big\|
=O_P\Big(n^{-1/2}+n^{-1/2}\delta_n+\mathsf R\,L_x^{-s}+\delta_n^{2}+r_n\delta_n\Big),
\]
every exponent of which is a function of $(\mu,r,d_Z,s)$ through \cref{prp:phirate}, A4(iii), A8's window
and the spline direct theorem. That is the sense in which this is a result and not bookkeeping, and
it holds on that branch only.
\end{enumerate}
\emph{No further index enters.} A2c's $\mu''$ supplies a finite \emph{constant} through
\cref{lem:reprsmooth} and \cref{cor:elldiff} and appears in no exponent, and the project's second
smoothness index is $\mu''$ and no other.
\end{remark}

% ---------------------------------------------------------------------------
%  THE REMARKS BELOW WERE EXTRACTED ON 2026-08-03, in the second pass of the
%  same operation that produced rem:rootngroups and rem:rootnread on 2026-08-02
%  and against the same house style (.claude/references/theorem-environment-style.md):
%  an ASSUMPTION carries conditions only, a DEFINITION carries construction only,
%  and commentary -- motivation, comparisons, worked derivations, "what this rules
%  out" -- goes after the statement.  NOTHING WAS DELETED: every sentence removed
%  is below, verbatim except for the connectives needed to make it stand outside
%  the environment it used to sit in.  They are placed at the END of Part I for
%  the reason given above: nothing numbered follows them in this file, so adding
%  them to the shared theorem counter renumbers no existing environment.
% ---------------------------------------------------------------------------

\begin{remark}[How \cref{def:moment}'s two matrices are read]\label{rem:momentread}
Three clauses, each of which stood inside the definition. None of it is construction, and the
objects \cref{def:moment} defines are unchanged by its removal.
\begin{enumerate}
\item[(1)] \emph{The second summand of \eqref{eq:JSig} carries $\ell_\star=\partial_uL_\star$ and not
$L_\star$}, because $\theta$ enters $\psi$ only through $\Delta U(\theta,\varphi)$ and differentiating
$L_\star$ in that argument returns its $u$-derivative. \emph{And under A1 that summand is the zero matrix,
so $J=E[\Delta Z\Delta Z']$ exactly}, \cref{lem:momzero}\textup{(iv)}, whose whole content is
$E[\widetilde W\mid\Delta U,\Delta Z]=0$, \emph{the zero vector of $\IR^{\mathsf L-1}$}, applied
between the two $\sigma(\Delta U,\Delta Z)$-measurable factors
$\ell_\star(\Delta U,\Delta Z)$ and $\Delta Z'$, which is where the contrast sits in
$E[\ell_\star\widetilde W\Delta Z']$. \eqref{eq:JSig} is nevertheless left in the two-summand form,
because $J$ is \emph{defined} as the $\theta$-derivative of \eqref{eq:psidef}, which has two summands
whatever their values, and because that is the form the estimator \eqref{eq:Jhat} takes, where the
second summand does \textbf{not} vanish.
\item[(2)] \emph{$J$ must \textbf{not} be described as ``a different matrix'' from
$E[\Delta Z\Delta Z']$}: wherever \eqref{eq:JSig} is well formed the two are the same matrix, by
\cref{lem:nusigma}\textup{(v)} at the array $\Xi(u,z)$ with entries
$\Xi_{jj'l}(u,z)=\ell_{\star,jl}(u,z)\,z_{j'}$, whose last index is the one contracted against
$\widetilde W$, \emph{which is the array-valued reading that clause states, and at $\mathsf L=2$ it is
the matrix $\ell_\star(u,z)z'$ earlier rounds printed}. \emph{What \textbf{is} true, and is why
identification of $\theta_\dagger$ must not be routed through \eqref{eq:JSig}, is the conditionality}:
$\ell_\star$ exists only under A2, so \eqref{eq:JSig} is a well-formed object only under A2, while
identification of $\theta_\dagger$ is A1b(iv)'s and holds whether or not A2 does.
\item[(3)] The typing note is about the \emph{form} and not the value. $\ell_\star$ is
$\IR^{p\times(\mathsf L-1)}$-valued by the type declaration of \eqref{eq:contrast}, $\widetilde W$ is
$(\mathsf L-1)\times1$ and $\Delta Z'$ is $1\times p$, so
$E[\ell_\star\widetilde W\Delta Z']$ is $p\times p$, \emph{and the three factors multiply in that
order and in no other}, and its \emph{form} is neither coordinatewise in the
$p$ coordinates of $\theta$ nor symmetric, which is what
makes the sum with $E[\Delta Z\Delta Z']$ well formed rather than a type error. Its \emph{value} is zero,
so $J$ itself is symmetric and positive semidefinite. \emph{The warning is therefore not idle; it
transfers to $\hat J$}, whose second summand in \eqref{eq:Jhat} is a sample average at
$\widehat{\widetilde W}^{(-k)}$ and $\hat\ell^{(-k)}$, does not vanish, and leaves $\hat J$
non-symmetric, which is why \eqref{eq:Sighat} writes $\hat J^{-1}\hat\Sigma\hat J^{-1\prime}$ with the
transpose and not $\hat J^{-1}\hat\Sigma\hat J^{-1}$.
\end{enumerate}
\end{remark}

\begin{remark}[How \cref{def:spaces}'s domain declaration is read]\label{rem:spacesread}
Four clauses. The first three stood inside the definition, and none of that is construction: the spaces,
the operator, the two inner products, the declared exception and \eqref{eq:nudef} are unchanged by their
removal. \emph{The fourth never stood there}: it prices the declaration \cref{def:spaces} makes about
$\pi$, which is construction and stays.
\begin{enumerate}
\item[(1)] The bound variables of the two inner products. They are $h,h'$ on the
domain side and $\Xi,\Xi'$ on the range side, being this paper's letters for a direction in
$L^{2}(\pi)$ and a generic element of a range-type space respectively.
$\Xi$ carries \textbf{five} value types in this corpus, each declared at its own site, and they
must not be conflated, at $\mathsf L=2$ the value types of \textup{(a)} and \textup{(b)} are both
scalar, and at $\mathsf L=2$ \emph{and} $p=1$ together three of the first four collapse, which is why the
separation is recorded here rather than left to context. \textup{(a)} In \eqref{eq:hatAadj} and
\cref{prp:hatAprop}\textup{(iii)}, $\Xi$ is a \emph{scalar} element of $L^{2}(\Pi)$, a function of
$(u,t,l)$; that is the declaration this clause makes and it does \textbf{not} move with $\mathsf L$.
\textup{(b)} In \cref{lem:adjoint}\textup{(ii)} and in \cref{cor:rangegap}, $\Xi$ is
$\IR^{\mathsf L-1}$-valued on
$\IR\times\IR^{p}$: the first binds it as the generic element of the codomain
$L^{2}(P_{\Delta U,\Delta Z};\IR^{\mathsf L-1})$, in which that clause writes its norm, and the second
inherits the type because $\Upsilon$'s domain is \cref{lem:adjoint}\textup{(ii)}'s codomain.
\textup{(c)} In \cref{def:dirclass}, \cref{lem:nusigma}\textup{(v)} and \cref{rem:rootnhyp}\textup{(1)},
$\Xi$ is the generic \emph{direction}, $\IR^{p\times(\mathsf L-1)}$-valued, and is contracted against
$\widetilde W$ from the left. \textup{(d)} In \cref{rem:momentread}\textup{(2)} it is array-valued, which
is \cref{lem:nusigma}\textup{(v)}'s stated extension of \textup{(c)}.
\textup{(e)} In \cref{lem:pgence} and its proof, $\Xi$ and $\Xi'$ are \emph{generic real random
variables} on the underlying probability space, bound in that statement and used nowhere else.
\emph{This one is of a different kind from \textup{(a)}--\textup{(d)}, which are deterministic functions
on a range-type index set, and it is the reason the count is five and not four}: it is a random variable
on the sample space and not a function of $(u,t,l)$ or of $(u,z)$, it carries neither $\mathsf L$ nor
$p$, and it therefore collapses into none of the others at $\mathsf L=2$. \emph{It is listed rather than renamed} because \cref{lem:pgence} is a self-contained
measure-theoretic lemma whose $(\Xi,\Xi';\varsigma,\varsigma')$ pairing is itself the product of an
earlier collision repair, recorded in its own statement.
\item[(2)] Why the restriction to $\mathcal V$ costs nothing. The
orthogonal decomposition $L^{2}(\pi)=\mathcal V\oplus\{\text{constants}\}$ holds because
$\langle c,h\rangle_{L^2(\pi)}=c\int h\,d\pi$, and $\Delta$ annihilates constants, so nothing is lost by
the restriction while a zero singular value is avoided. In particular the singular system
$(\sigma_j;e_j,\tilde e_j)_{j\ge1}$ of \cref{prp:compact}, the completeness condition A1c, the source
condition A2 and the Hilbert-scale link H1 all refer to
$\mathcal A_{\mathrm{cond}}:\mathcal V\to L^{2}(\Pi)$ and to that same singular system.
\item[(3)] Why $\mathcal B$ is the exception, and what governs the form's sharp constant.
$\mathcal B$ is read on all of $L^{2}(\pi)$ because $\mathcal V$ is \emph{not} invariant under
multiplication by a non-constant
function and $\mathcal B$ contains such a multiplication. The quadratic form of
\cref{lem:secondord}(i) is evaluated only at $h\in\mathcal V$, so the form's own sharp constant is
governed by the compression $P_{\mathcal V}\mathcal B_{\mathrm{sym}}P_{\mathcal V}|_{\mathcal V}$, where
$P_{\mathcal V}$ is the $L^{2}(\pi)$-orthogonal projection onto $\mathcal V$; \cref{lem:solset}(iii)
records the relation between the two.
\item[(4)] What \cref{def:spaces}'s declaration of $\pi$ costs, and the one thing it does not settle.

\textup{(4a)} \emph{$\pi$ and $\Pi$ do not have the same status and the shared word ``reference'' is a
trap.} B1's $\Pi=\Pi_u\otimes\Pi_t\otimes\Pi_{\mathsf w}$ is a set of \emph{choices made by the
analyst}, which is why \cref{rem:localise}\textup{(2)} can say that B1 survives the passage to a
sub-population unchanged and why \cref{lem:scaleinv}\textup{(vi)}--\textup{(vii)} can rescale its
\emph{second} factor in two independent ways with nothing else moving. $\pi$ is a \emph{functional of
$P$}. Nothing in this paper selects it, and no result below is licensed to treat it as a design lever.

\textup{(4b)} \emph{Two consequences of the pooled form, both free, and neither of them a weakening of
C1(ii).} First, $f_{X_j}\le f_{X_1}+f_{X_2}=2\pi$ pointwise for $j=1,2$, so
$\bar c_\pi=\max_j\|f_{X_j}/\pi\|_\infty\le2$: the ratio constant that every boundedness statement about
$\mathcal A_{\mathrm{cond}}$ and about the second-order form carries has a numerical ceiling it did not
have while $\pi$ was read as an arbitrary reference measure. \emph{The worked example of
\cref{rem:localise}\textup{(7)} is the check, and it passes}: its marginals are
$f_{X_1}(x)=1+\tfrac{\gamma}{3}(\tfrac12-x)$ and $f_{X_2}(x)=1+\tfrac{\gamma}{3}(x-\tfrac12)$, whose
average is $1$ at every $x$, so that example's declared $\pi\equiv1$ \emph{is} its pooled design
density and was consistent with this reading before the reading was stated; its
$\bar c_\pi=1+\tfrac{\gamma}{6}<\tfrac76$ at $\gamma\in(0,1)$ sits inside the ceiling.
Second, C1(ii)'s bounds on $\pi$ follow from its bounds on $f_{X_1}$ and $f_{X_2}$,
since an average of two functions bounded above and below away from zero is one. \emph{C1(ii) is
retained verbatim and is not narrowed}: it also bounds $f_{X_1,X_2}$, $f_{X_1}$ and $f_{X_2}$
themselves, which no pooled identity delivers, and every citation of it stands where it is.

\textup{(4c)} \emph{Why $\pi$ is fixed at $P$ and not re-read at a sub-population.}
\cref{rem:localise}\textup{(2)} transfers C1(ii)'s \emph{upper} bounds to $Q=P(\cdot\mid W\in\mathsf S)$
with $\bar c_\pi$ inflated to $\bar c_\pi/P(W\in\mathsf S)$, and \emph{both sides of that bound divide by
the same $\pi$} --- the left through $\|f^{Q}_{X_j}/\pi\|_\infty$ and the right through $\bar c_\pi$
itself; the proof of that remark says so in terms. The pooled density of the
retained sub-population, $\tfrac12(f^{Q}_{X_1}+f^{Q}_{X_2})$, is a different function, and nothing here
claims the two agree --- which is what the closing sentence of \cref{rem:localise}\textup{(2)} already
records when it says that nothing is claimed there about $\pi$. \emph{A localised implementation that
rebuilt $\hat\pi$ on the retained cells alone would be estimating that other function, and this paper
prices no such substitution.}

\textup{(4d)} \emph{The one thing the declaration does not settle: the adjoint.} $\pi$ being unknown,
every display carrying $1/\pi$ is uncomputable as printed. \emph{Five are worth naming and no claim of
exhaustiveness is made}: \textup{(1)} \eqref{eq:hatAadj}, whose $1/\pi(x)$ stands in front of the whole
adjoint; \textup{(2)} \eqref{eq:normconst}, through the entry $\|1/\pi\|_\infty^{1/2}$ of its maximum,
\emph{which is the binding entry only above the threshold
$\varepsilon_\pi=\|1/\pi\|_\infty^{-1}$ and is dominated by the computable $\varepsilon_\pi^{-1/2}$ at
or below it}; \textup{(3)} \eqref{eq:cap}, which carries its square on the same reading; \textup{(4)}
\cref{lem:levels}\textup{(i)}; and \textup{(5)}
the Cauchy--Schwarz step of the measurability verification in the supplement's discussion of the
estimator. \cref{rem:localise}\textup{(2)} lists its own overlapping set
and is likewise not a closure claim. \emph{The one that binds is the first}, because
\cref{def:lfstep}\textup{(b)} runs every Landweber step through $\hat{\mathcal A}^\ast_\varphi$: the
forward operator \eqref{eq:hatAcon} is free of $\pi$ and the iteration is not.
\cref{def:lfstep}'s construction convention makes the substitution of a floored $\hat\pi$ well posed and
its constants computable; \emph{what no result here establishes is that the substitution leaves
\eqref{eq:phirate} in place}, and that is stated as an open item at
\cref{rem:phirscope}\textup{(7)} rather than assumed away --- \emph{in five parts as of 2026-08-27, the
fifth being the relaxation condition \eqref{eq:normconst} is now sized for}.
\emph{And the statement that would otherwise contradict this sentence is scoped rather than left to be
reconciled}: \cref{prp:phirate} states \eqref{eq:phirate} at the infeasible reading $\hat\varphi^{\pi}$
of \cref{def:lfstep} and says in its own conclusion that it does not claim the display at $\hat\varphi$.
\emph{Through 2026-08-27 it claimed exactly that, unhedged}, so this sentence and that proposition were
in printed contradiction; the sentence is unchanged and the proposition moved.
\end{enumerate}
\end{remark}

% ============================================================================
%  END OF PART I.
% ============================================================================
%
% ============================================================================
%  DISPATCH-11 CHANGE LOG (against the Part-I review, 48/100).  FIX-ONLY.
%  NO NEW RESULTS.  ONE NEW LEMMA CLAUSE (lem:sandwich(i')), ONE NEW DEFINITION
%  CLAUSE (def:crossfit(d)), ONE HYPOTHESIS CLAUSE DELETED ((R-b)'s third).
%  Every pre-existing \label still exists and still names the same statement.
%
%  THE ARITHMETIC (MAJOR 1).  (1-kappa)/(4 kappa) = 1/(8 mu), NOT 1/(4 mu):
%  1 - kappa = 1/(2 mu + 1) and 4 kappa = 8 mu/(2 mu + 1).  Illustrative numbers
%  3/20 at mu = 5/6 and 1/16 at mu = 2.  rem:rootnlog(3)(b) now DISPLAYS the
%  identity as eq:rndeltamargin so it can be checked at sight, and records that
%  the constant is EXACT at alpha = 1/(4 kappa) -- hence the infimum over the
%  frontier and not improvable from eq:frontier alone.  THE FINDING SURVIVES: the
%  implication (eq:frontier => r_n delta_n = o(n^{-1/2})) and its direction in mu
%  are unaffected.  Same repair at proofs.tex thm:rootn Step 4(c).
%
%  THE A5(ii) SWEEP (MAJOR 3).  rem:phirscope(2) rewritten: the warning is
%  DISCHARGED, not carried, because A5(ii) now writes the log-inflated form.
%  prp:stopfeas(iv)'s OPENING DISPLAY repaired -- it inflated delta_n a second
%  time, which was the very collision eq:deltalog exists to end, and it made
%  prp:stopfeas(iv) inconsistent WITH ITSELF (log-free display, log-inflated
%  (P-f) paragraph).  rem:rootnlog(1) now records that internal inconsistency,
%  because it STRENGTHENS the ground for route (A).
%
%  cor:ratefloor (MAJOR 4).  "The rate is in primitives" was FALSE on the
%  corollary's own branch.  Now stated by branch: (a) always -- n^{-1/2},
%  delta_n^2, r_n delta_n; (b) never on (R-e.1)+(R-c.2) -- epsilon_n and
%  epsilon^L_n have NO primitive supplier, and s enters no exponent there;
%  (c) entirely, on (R-e.2)+(R-c.1), where both unsupplied symbols vanish.
%  NOTE this is STRONGER than the review asked: scoping to (R-e.2) alone is
%  necessary but NOT sufficient, since epsilon^L_n survives via (R-c.2).
%  A5(vi) with s > 1/2 ADDED to the hypothesis list (lem:epterm needs the
%  embedding H^s -> C(X) and the corollary drops (T-4)).  mu' DELETED: one
%  occurrence corpus-wide, no glossary row; see notation_glossary Sec.8.-1.
%
%  FOLD INDEXING (MODERATE 5).  def:crossfit(d) is new and declares eq:foldrates
%  once; (R-e.1) now carries max_k EXPLICITLY (it is a strengthening in form of
%  A5(iv) and is charged, not read in); lem:epterm and lem:prelimtheta cite the
%  fold reading in their own hypothesis lines, so both are true standalone.
%
%  (R-b) (MODERATE 6a).  Third clause E||b_star||^4 < oo DELETED: b_star enters
%  Part I at second moment only, via A2'(iii).  The fourth-moment need is on
%  b-hat and (R-a) supplies it.
%
%  (T-4) vs STEP 0 (MODERATE 6b).  (T-4) is now ANNOTATED by status: A5(i) and
%  A5(ii) DELIVERED by (T-5); A5(iii) ASSUMED, being a restriction on the
%  primitives (mu, r, d_Z) that no proposition can discharge; A5(vi) with
%  s > 1/2 split into the delivered H^s bound and the primitive restriction.
%  Step 0's "discharges A5(i)-(iii)" corrected to match.  lem:sandwich gains
%  clause (i'), PROVED at Step 6', which is what cor:ratefloor actually needs.
%
%  lem:momzero (MODERATE 6c).  A2'(i.b) MOVED INTO THE LIST: cor:elldiff is
%  stated at A2b, which needs BOTH clauses, and A2'(i)'s closing clause fixes
%  ell_star pointwise from the version satisfying both -- so ell_star is not
%  pointwise well defined from (i.a) alone.  The A2'(ii)-redundancy claim is
%  RE-ATTRIBUTED: implied by A2'(i) IN FULL plus B1(i), (ii), (iv.a), not by
%  (i.a).  Same repair at rem:rootnhyp(1) and (2).  lem:sandwich's list now
%  inherits lem:momzero in full (A1 and the three B1 clauses were missing).
%
%  THE NON-PROMOTION GROUND (MODERATE 6d).  "Every clause of A5 prices a rate in
%  n" is FALSE -- A5(vi) is ||h||_{H^s} = O_P(1) and A5(vi') a w.p.->1 envelope.
%  The DECISION (keep (R-a)-(R-e) on the theorem) STANDS; the REASON is replaced
%  at all three sites (rem:dirmember(a), rem:obligations, and new
%  rem:rootnhyp(5)) by: object existence, rem:obligations' own construction test,
%  misfiling of (R-b), and the (P-a)-(P-g) precedent.  rem:obligations' inventory
%  updated: (O2) and (O3) STILL OPEN -- neither is discharged, since a hypothesis
%  of a theorem is not a verification -- but each now has a named consumer.
%
%  SIX MINORS.  def:crossfit(c) states the three properties def:weight(b) needs
%  of ell-hat, and def:onestep says its identity uses ONE derivative and that the
%  second is hypothesised at thm:rootn Step 1;  rem:rootnhyp(2) re-headed "NOT
%  listed";  lem:epterm's headline no longer claims consistency is free;
%  lem:prelimtheta's "collapse not used" is now TRUE, Step 2 having been restated
%  at the uncollapsed rate;  prp:medianpart's hypothesis is read PER PARTITION
%  over P*F fold complements (and P fixed is used THREE times, not two);
%  lem:biasexp's Psi-vs-R paragraph no longer claims Psi^theta_n is written out.
%  Part I's (P8) paragraph now separates the GENERIC transfer rule from the
%  Part-H CONDITIONING CONVENTION, which def:crossfit re-declares.
% ============================================================================
%
% ============================================================================
%  DISPATCH-15 CHANGE LOG.  TWO FINDINGS AND NOTHING ELSE.  FIX-ONLY.  NO STEP 8,
%  NO NEW MAINTAINED ASSUMPTION, NO CONCLUSION CHANGED.  Every pre-existing
%  \label still exists and still names the same statement.  ONE NEW LABEL.
%
%  NEW LABEL:
%     eq:prodmeas   product measurability of the ESTIMATOR'S OWN POINTWISE
%                   FUNCTION, the one def:crossfit(c) already defines:
%                   (omega, u, z) |-> ell-hat^{(-k)}(omega)(u, z) is
%                   sigma(I_{-k}) tensor Borel(R x R^p) measurable.
%                                                       [def:crossfit(c'.1)]
%
%  THE DEFECT eq:prodmeas REPAIRS.  (c'.1) declared all FIVE nuisances random
%  elements of L^2 -- EQUIVALENCE CLASSES -- and fixed a representative.  (c'.2),
%  fourteen lines below it, argues that the essential sup must be in z and not in
%  (u, z) PRECISELY BECAUSE the shifted argument's law need not be dominated by
%  P_{Delta U, Delta Z}.  If it is not dominated, two representatives of one class
%  DIFFER THERE -- and eq:Jhat evaluates ell-hat^{(-k)} POINTWISE at
%  (Delta U_i(theta, varphi-hat^{(-k)}), Delta Z_i).  Different representative,
%  different J-hat, hence different theta-hat (eq:onestep), Sigma-hat (eq:Sighat)
%  and Psi^{(1)}_k.  THE ESTIMATOR WAS NOT A FUNCTION OF THE OBJECT THE CLAUSE
%  DECLARED.  So (c'.2) is the proof that (c'.1) was wrong.
%
%  WHAT THE REPAIRED CLAUSE SAYS, AND IT IS TWO STATEMENTS AND NOT ONE:
%    varrho-hat, b-hat, varphi-hat KEEP the L^2-random-element reading with a
%      fixed jointly measurable representative.  It is RIGHT for them: each is
%      evaluated only at draws of the law its container is built on (Delta Z;
%      X_{i1}, X_{i2}), so the choice of representative is immaterial.
%    ell-hat and L-hat are governed by eq:prodmeas instead.  NO equivalence class
%      is formed and NO representative is chosen.  L-hat INHERITS it from
%      eq:weightdef by Tonelli; partial_u ell-hat INHERITS it as a pointwise limit
%      of difference quotients (which (c'.2)'s SECOND envelope needs and which a
%      representative supplies nothing of, having no derivative -- a gap the
%      superseded clause also had).
%  TWO FALSE SENTENCES STRUCK.  (1) "The continuity clause of (c) is preserved by
%  the choice" -- FALSE for an arbitrary representative; point evaluation is not
%  continuous on L^2.  Under eq:prodmeas there is no choice, so (c)'s continuity
%  is untouched.  (2) The claim that (c') "re-declares for Part I what lem:balance
%  states for Part F" -- FALSE in both directions: lem:balance states only that
%  ell-hat is sigma(I_a)-measurable with ell-hat in L^2 ALMOST SURELY, which is a
%  MEMBERSHIP statement and strictly weaker.
%  ONE CONTRADICTION CLOSED, in clause (c) itself and because it is the same
%  defect: (c) said (R-a)'s differentiability is "used by no clause of this
%  definition", and (c'.2) -- a clause of this definition -- uses it to make the
%  second envelope's inner supremum measurable.  All five clauses checked; (c'.2)
%  is the only consumer, and (c) now says so and names it.  A reader following (c)
%  alone would have concluded that eq:supnormconv's convention is unconditional.
%  eq:supnormconv ITSELF AND ITS ELEVEN-CONSUMER ENUMERATION ARE UNTOUCHED.
%
%  FINDING 1 of this dispatch is in proofs.tex: the (P8) FREEZING LEMMA is
%  restated at the CONDITIONAL hypothesis its consumers supply, with the
%  conclusion read as an identity of GENERALISED conditional expectations.  What
%  touches THIS file: nothing in any statement.  eq:prodmeas is the hypothesis
%  that lemma consumes at its seven Part-I sites, and (c'.1) now says so.
% ============================================================================
%
% ============================================================================
%  DISPATCH-14 CHANGE LOG (against four narrow audits arriving together).
%  FIX-ONLY.  NO STEP 8, NO NEW MAINTAINED ASSUMPTION, NO CONCLUSION CHANGED.
%  Every pre-existing \label still exists and still names the same statement.
%
%  NEW LABELS, three, each forced by a finding:
%     eq:supnormconv  the reading of ||ell-hat^{(-k)}||_infty: a GENUINE sup in
%                     u (legitimate: continuity in u makes it countable) and an
%                     ESSENTIAL sup in z under P_{Delta Z}.  The z-side reading
%                     is the one the consumers need, because every use evaluates
%                     the envelope at a SHIFTED FIRST ARGUMENT whose law need not
%                     be dominated by P_{Delta U, Delta Z}, while its second
%                     argument is always a draw of Delta Z.        [def:crossfit(c')]
%     eq:sigmapd      Sigma > 0, hence J^{-1} Sigma J^{-1'} > 0.    [lem:momzero(v)]
%
%  NEW CLAUSES:
%     def:crossfit(c'), with (c'.1) and (c'.2).  The five estimated nuisances are
%       sigma(I_{-k})-MEASURABLE RANDOM ELEMENTS of the named L^2 spaces, with a
%       jointly measurable representative fixed once.  THE DEFECT: lem:balance
%       states this for Part F ("ell-hat is sigma(I_a)-measurable with ell-hat in
%       L^2 a.s."), and def:crossfit(c) replaced it with a PER-REALISATION
%       "Borel" -- under which (R-a)'s hypotheses and eq:foldrates are not
%       well-posed probability statements, being statements about the
%       distribution of a supremum and of an integral of a random function.
%     lem:momzero(v): Sigma > 0 under A1b(iv) + C2'(i), PROVED at proofs.tex
%       Step 6.  CONSEQUENCE: thm:rootn's studentised conclusion now quantifies
%       over EVERY unit vartheta; the positivity proviso is IMPLIED by (T-1) and
%       (T-3) and is stated as a consequence.
%     thm:rootn (R-d): DECLARED as a struck placeholder.  There is no (R-d); the
%       list runs (R-a), (R-b), (R-c), (R-e), and the skip was undeclared.
%     rem:rootnscope(6): the theorem is POINTWISE IN P.  No uniformity is claimed
%       or proved, and item (5)'s n* ~ 10^58 is that quantifier's shadow rather
%       than a finite-sample caveat.
%     rem:pasrc(5): the BINDING reading of the glyph mu.  Bare mu is (P-a.3)'s
%       PRIMITIVE index at the compression of the REALISED operator at theta_o;
%       H1's s/(2a) is named mu_{H1} where the two must be told apart.  They are
%       equal only under the proviso rem:pasrc(4) declines.  rem:rootnscope(1)
%       and the two daggered budget-table lines are qualified in place.
%
%  WEAKENED / RETIRED:
%     lem:nusigma(i): "for every h in V" -> "for every h in L^2(X,pi)", which is
%       all its proof establishes.  This REMOVES a hidden A1b(i) need at THREE
%       lemmas at once (lem:prelimtheta Step 4, lem:epterm Step 3, lem:sandwich
%       Steps 3 and 6'), and is strictly better than adding the tag to three
%       hypothesis lists.
%     eq:Sighat: V-hat gains a FALLBACK, V-hat := 0 off {J-hat nonsingular}.  The
%       definition was careful about that event twice and not the third time,
%       while thm:rootn's second conclusion quantifies over it.
%     def:crossfit(d): "the three clauses of A5" -> "the four", A5(vi') being the
%       fourth, with its fold reading and the one consequence (a single R serving
%       all F fold complements).
%
%  NOTATION: prp:medianpart's partition index q -> sans-serif p (COL-1; q is
%     bound to E[Wtilde partial_u ell_star | X_1, X_2] and live BARE in Part I at
%     lem:momzero(iv.3)).  See proofs.tex for COL-2.
%
%  CORRECTED CLAIMS: rem:adjointreconcile(i)'s "A2'(i)--(ii)" -> "A2'(i)", A2'(ii)
%     being retired (same fix in paper/sections/theorems/lem-adjoint.tex);
%     rem:rootnhyp(1) gains C2'(i)'s second, direct consumer; rem:rootnhyp(2)
%     records A2'(ii) as RETIRED rather than merely unlisted; lem:levels' "at
%     index mu" -> "at index mu_{H1} = s/(2a)".
% ============================================================================
%
% ============================================================================
%  ROUND-2 CHANGE LOG (against proofs_review.md, 56/100)
%
%  C-1  lem:bierenscore restated for w = exp; the false cos/sin extension
%       deleted; rem:swintercept added with the counterexample, the
%       intercept-indexed repair, and the correct anchor (S&W 1998 Thm 3.1).
%  C-2  rem:draftrec(ii) rewritten as a KERNEL comparison, with the exact
%       condition p_X = varrho named and rejected; the "same range, same
%       kernel" claim withdrawn.
%  M-1  The pi^{-1} in rem:draftrec(iii) is tied to A2's object (nu/pi);
%       prp:compact gains the clause k(u,t;.) in V that makes this internal.
%  M-2  prp:frechet split into (i) quadratic remainder and (ii) LIPSCHITZ
%       DERIVATIVE MAP, the latter stated with a constant and proved.
%  M-3  rem:R1(ii) routes identification of theta_dagger through A1b(iv), not
%       A3, and says explicitly that neither A2 nor A3 is used.
%  M-5  def:spaces fixes the domain at V once, with the reason.
%  M-6  rem:Tcond restated at the weakened B1(iv); the differentiation-under-
%       the-integral claim names the integral and the dominating function.
%  M-10 A0 and A4(i) removed from prp:wellposed / prp:frechet / prp:compact.
%  m-2  rem:R1(iv) scoped: population statement clean, algorithmic consequence
%       conditional on A1b(iii) and on A6/A7 at the running theta.
%  m-4  lambda (Bierens weight, two arguments) separated from varsigma (generic
%       bounded test function); G freed for the conditional CDF by writing the
%       moment class with g.
%  m-6  Pi_u(R) restored in the quoted form of B1(iv.c); remark-placement note
%       added to the header.
%  m-7  cor:opsol names its parents.
% ============================================================================
%
% ============================================================================
%  DISPATCH-2 CHANGE LOG
%
%  N-10  CLOSED.  lem:adjoint(ii) constructs the induction map W; A2'(0) is
%        deleted from assumptions.tex.  rem:Tcond amended: the t-integral is
%        optional on the ESTIMATION side and forced on the POPULATION side, and
%        the correction measure gamma is identified with Pi_t, which makes the
%        first half of B1(v) a consequence of B1(iv.c) rather than an
%        assumption.
%
%  STEP 3  Part E added: lem:nusigma (L-1), lem:adjfubini (L-4), lem:levels
%        (L-5), lem:adjoint, prp:nophifree with rem:nophifree, cor:udep (O-11),
%        and -- ADDITIVE, flagged as N-12 -- prp:kerchar, cor:a1cprim and
%        rem:ffcompare.
%
%  DRAFT `:172--177'  rem:draftrec(iii) rewritten: what the display MEANS, what
%        it SHOULD have said (eq:adjointdisplay), and THREE defects, not two ---
%        the pi(v) duplication; the vacuous theta subscript SEPARATED from the
%        NON-vacuous varphi subscript, which the draft wrongly suppresses by
%        writing f_U(u); and the missing u-integration, which is also what makes
%        `:124--126' and `:172--177' silently use different range spaces.
%
%  ENVIRONMENT  The formal impossibility statement is prp:nophifree (a
%        proposition, with a proof); rem:nophifree survives as the remark that
%        follows it and carries the interpretation and the Feve--Florens link.
%        Both labels exist, each on the right environment.  See N-11.
% ============================================================================
%
% ============================================================================
%  DISPATCH-3 ROUND-7 CHANGE LOG (fix-only).  NO NEW RESULTS.  Every \label key
%  that existed before this round still exists and still names the same
%  statement; four labels are NEW (def:weight, rem:anchor, def:dirclass,
%  rem:shiftclass, eq:weightdef, eq:shiftnorm) and one has moved file
%  (eq:mtau -> eq:Psitau, in proofs.tex).
%
%  CRITICAL F-2 -- L_star DEFINED.  def:weight supplies the definition the block
%        never had, as the anchored u-antiderivative of a representer;
%        rem:anchor records the two canonical anchors and shows both admissible;
%        A2'(iii) prices the anchor.  prp:orthog(iii)'s perturbation
%        r = eps_varrho sign(E[L_{star,j} | Delta Z]) is now built from a defined,
%        conditionally integrable object.
%
%  CRITICAL F-1 -- THE SHIFTED MOMENT.  def:dirclass adds the shift-uniform
%        seminorm ||.||_{0,h} to the direction class; lem:secondord's Psi is
%        declared finite on [-tau-bar, tau-bar] x R^2 and clause (ii)'s
%        "identically in tau" is scoped to that set.  rem:shiftclass prints the
%        separating example -- a two-character variant of the block's own round-6
%        counterexample, under which Psi is UNDEFINED rather than merely large --
%        and says why the density-ratio alternative was rejected.
%
%  F-4    def:spaces now names the ONE operator read on L^2(pi) rather than V,
%        with the reason (V is not invariant under multiplication) and the
%        codimension-one compression that transfers the Weyl bound.
%  F-8    lem:solset(ii) restricted to ADMISSIBLE members, with the reason.
%  F-9    A0's compact interval consumed at lem:levels.  WITHDRAWN 2026-08-26:
%        the finding was right that A0 was consumed and wrong about which half.
%        lem:levels(i) now integrates over [x_lo, x] cap X -- clause (ii)'s own
%        intersection form -- so the lemma reads COMPACTNESS and not the
%        INTERVAL, and is true on any compact X.  The interval's real consumer
%        is lem:intker(iv), which now carries it.
%  F-10   lem:balance's sample-splitting convention promoted to a hypothesis,
%        with the conditional reading of every nu-object and O_P statement, and
%        with ell-hat in L^2 a.s.
%  M-8 completed  N1 struck from lem:solset, lem:secondord and lem:balance; it
%        was listed at all three and consumed at none.  N1' carries the direction
%        condition.  prp:orthog(i) restated at an ARBITRARY ADMISSIBLE
%        REPRESENTER, which is what makes its second "if and only if"
%        non-vacuous, and which DROPS A2 from that part.
%  NOTATION  Upsilon replaces the cor:rangegap B; Psi(tau) replaces m(tau);
%        sans-serif r replaces the varrho-direction r; vartheta replaces the
%        binding use of varsigma; Xi replaces lem:nusigma(v)'s vartheta;
%        C_0 belongs to cor:udep alone, C3's interval being written out.
% ============================================================================
%
% ============================================================================
%  ROUND-9 CHANGE LOG -- memo Sec.10 STEP 4.  NEW CONTENT, the first since
%  round 6.  Every pre-existing \label still exists and still names the same
%  statement; NO statement in Parts A--F is altered except rem:obligations (O1)
%  and the header's "NOT IN THIS FILE" note.
%
%  NEW LABELS (Part G): lem:reprsmooth, rem:l21route, lem:formstruct,
%        prp:sharpquad, rem:sharpscope; equations eq:embedconst, eq:embed,
%        eq:l21budget, eq:singseq, eq:complower, eq:quadrem, eq:sharpconst,
%        eq:bracket, eq:noweak.  In assumptions.tex: ass:reprsource (A2c),
%        eq:usobolev, eq:lambdamu.
%
%  prp:sharpquad IS STATED AT THE SUP-OVER-BALL QUANTIFIER and not as a necessity
%        claim about the estimator.  rem:sharpscope blocks the four false readings,
%        of which the first is the one that matters: the result bounds the BEST
%        BOUND VALID UNIFORMLY ON THE BALL, and says nothing about the realised
%        remainder along a particular regularisation path.
%
%  THREE THINGS THAT ARE INHERITED VERBATIM AND NOT RE-DERIVED, per the dispatch
%        and glossary Sec.8.4 items 2-3: lem:solset(iii)'s lower-bound chain, which
%        now runs on B_mult; the readability of B = B_mult - 2 vartheta' Ktilde and
%        B_sym = B_mult - vartheta'(Ktilde + Ktilde*); and F-4's codimension-one
%        compression, which is why the sharp constant in eq:sharpconst is the
%        COMPRESSION's norm and not the L^2(pi) norm.
%
%  THREE THINGS THAT MUST NOT SLIP, and each is stated in the file:
%        (a) ||vartheta' sigma_q / pi||_inf is a LOWER bound via the multiplication
%            part, NOT the sharp constant (eq:bracket, rem:sharpscope(2));
%        (b) B is NOT self-adjoint, so (1/2)||B|| is an UPPER bound only
%            (eq:bracket's right-hand member, lem:formstruct(i));
%        (c) sigma_q is R^p-valued, so everything is stated at a fixed direction
%            vartheta with ||vartheta|| = 1 -- vartheta is BINDING and varsigma,
%            being a reserved LOCAL bound variable, cannot be (C25, Sec.8.4 item 2).
%
%  O-16 IS A HYPOTHESIS, IN TWO CLAUSES, AND THEY DO DIFFERENT WORK.  (a) the form
%        does not vanish on V, which is what gives the delta^2 ORDER; (b)
%        vartheta' sigma_q not identically 0, which is what gives the LOWER BOUND,
%        the non-compactness and the weak-norm refutation.  (b) implies (a) and the
%        converse FAILS -- lem:solset(iii)'s own witness is the counterexample.
%        prp:sharpquad DOES NOT inherit memo Sec.5 O-11's "Q(h) != 0 necessarily";
%        that claim does not follow and is escalation N-13, still open and not the
%        theorist's to amend.
%
%  ITALIC Q IS NOT CLAIMED.  The memo's Q(h) is written here as the form
%        <Bh, h>_{L^2(pi)} (eq:quadrem), because Q_1, Q_2, Q_{12} -- (P6b)'s
%        regular conditional distributions -- occupy the letter at six sites in
%        proofs.tex, and Q_{12} is what supplies the jointly measurable version of
%        q, INSIDE the chain prp:sharpquad runs on.  Glossary Sec.1.5b: sweep for
%        the letter you are CLAIMING.  A blocking row is in check_rendering_sync.py.
%
%  L-21 CLOSES (O1), RETROACTIVELY.  rem:obligations(O1) is now marked DISCHARGED,
%        with what "discharged" means spelled out: A2'(i) stays maintained and
%        stops being an unverifiable primitive.  Parts E and F are UNTOUCHED,
%        because each carries A2'(i) or A2b as an explicit hypothesis -- which is
%        why stating every hypothesis explicitly was the right answer to the
%        ordering conflict N-20 raised.
% ============================================================================
